% Merged from nine separately delivered manuscripts on surcomplex analysis,
% all written on 21 September 2026 as independent attempts at the same
% subject.  See MERGE_NOTES.md for the provenance, the notation
% reconciliation, and the six notions that are NOT equivalent across the
% sources.  Every symbol below is fixed by that reconciliation.
\documentclass[11pt,a4paper]{article}
\usepackage[T1]{fontenc}
\usepackage{lmodern}
\usepackage[margin=26mm]{geometry}
\usepackage{amsmath,amssymb,amsthm,mathtools}
\usepackage{booktabs,array,longtable}
\usepackage{microtype}
\usepackage{enumitem}
\usepackage{xcolor}
\usepackage{fancyhdr}
\usepackage{hyperref}
\usepackage{bookmark}

\definecolor{Ink}{HTML}{233240}
\hypersetup{colorlinks=true,linkcolor=Ink,citecolor=Ink,urlcolor=Ink,
 pdftitle={Surcomplex Analysis: germs, coherent Hahn sections, and
 infinitesimal zero geometry},
 pdfauthor={Research notes prepared for Vladimir Reshetnikov},
 pdfsubject={One document merged from nine manuscripts on holomorphic
 function theory over No[i]}}

\pagestyle{fancy}
\fancyhf{}
\fancyhead[L]{\small Surcomplex analysis}
\fancyhead[R]{\small\leftmark}
\fancyfoot[C]{\thepage}
\renewcommand{\sectionmark}[1]{\markboth{#1}{}}
\setlength{\headheight}{14pt}
\setlength{\emergencystretch}{2em}

\numberwithin{equation}{section}
\newtheorem{theorem}{Theorem}[section]
\newtheorem{proposition}[theorem]{Proposition}
\newtheorem{lemma}[theorem]{Lemma}
\newtheorem{corollary}[theorem]{Corollary}
\theoremstyle{definition}
\newtheorem{definition}[theorem]{Definition}
\newtheorem{conjecture}[theorem]{Conjecture}
\newtheorem{example}[theorem]{Example}
\theoremstyle{remark}
\newtheorem{remark}[theorem]{Remark}
\theoremstyle{plain}

% ------------------------------------------------------------- the field ---
% K is the field.  S is NEVER the field: S is a support.
\newcommand{\K}{\mathbb K}
\newcommand{\No}{\mathbf{No}}
\newcommand{\Oz}{\mathbf{Oz}}
\newcommand{\Jpi}{\mathbf J}
\newcommand{\CC}{\mathbb C}
\newcommand{\RR}{\mathbb R}
\newcommand{\ZZ}{\mathbb Z}
\newcommand{\NN}{\mathbb N}

% valuation ring and its maximal ideal.  Hol(U) is ORDINARY holomorphic
% functions: the letter O is reserved for the valuation ring.
\newcommand{\Ov}{\mathcal O}
\newcommand{\mm}{\mathfrak m}
\DeclareMathOperator{\Hol}{Hol}
\DeclareMathOperator{\supp}{supp}
\DeclareMathOperator{\st}{st}
\DeclareMathOperator{\ord}{ord}
\DeclareMathOperator{\Arg}{Arg}
\DeclareMathOperator{\Log}{Log}
\DeclareMathOperator{\Exp}{Exp}

% Conway monomial scale.  NOT Gonshor exponentiation: t^{a+b}=t^a t^b is a
% convention here, not a theorem about surreal exponentiation.
\newcommand{\tm}{t}

% halo over an ordinary domain, and evaluation / canonical lift.  One symbol
% for both, because a lift IS a section supported at exponent zero.
\newcommand{\halo}[1]{#1^{\#}}
\newcommand{\ev}[1]{#1^{\#}}

% the coherent algebra, positive-support part, meromorphic-coefficient class
\newcommand{\Hc}{\mathcal H}
\newcommand{\Mc}{\mathcal M}

% Hahn valuation of a family and its leading coefficient FUNCTION.  Not the
% pointwise valuation of a value.
\newcommand{\vH}{v_{\mathrm H}}

% Neumann monoid generated by a positive well-ordered set
\newcommand{\Neu}[1]{\langle #1\rangle}

% Coefficientwise contour integral.  The superscript is permanent: this is
% not an ordinary integral and not a limit of Riemann sums.
\newcommand{\intH}{{\int^{\mathrm H}}}
\newcommand{\ointH}{{\oint^{\mathrm H}}}

% THE THREE RESIDUES.  Distinct symbols are mandatory.  Res and ResH
% DISAGREE at the same point: for 1/(zeta-eps) with eps infinitesimal,
% ResH_0 = 1 while the local residue at 0 is 0, the actual pole being at eps.
\DeclareMathOperator{\Res}{Res}                 % formal, in K((X))
\newcommand{\ResH}[1]{\Res^{\mathrm H}_{#1}}    % cluster, ORDINARY centre
\newcommand{\Resat}[1]{\Res_{#1}}               % actual, surcomplex point

\newcommand{\abs}[1]{\lvert #1\rvert}
\newcommand{\norm}[1]{\lVert #1\rVert}
\newcommand{\dd}{\mathrm d}

% Inherited from the source manuscripts' bibliographies, which defined it in
% their own preambles.
\newcommand{\arxiv}[1]{\href{https://arxiv.org/abs/#1}{arXiv:#1}}

\title{\textbf{Surcomplex Analysis}\\[4pt]
\large Germs, coherent Hahn sections, and infinitesimal zero geometry\\[2pt]
\normalsize a theory of holomorphic functions on $\K=\No[i]$}
\author{Research notes prepared for Vladimir Reshetnikov}
\date{September 2026}

\begin{document}
\maketitle
\setcounter{tocdepth}{2}
\tableofcontents
\clearpage

%% ==================== p01.tex ====================

% =====================================================================
%  PART 0  --  Foundations: the field, normal forms, and the size
%              discipline; plus the conventions section implementing the
%              notation contract.
%  PART 1  --  Summability, and why it is not convergence.
%  Written for the merged document; sections only, labels prefixed a:.
% =====================================================================

\section{Foundations: the field \texorpdfstring{$\K=\No[i]$}{K = No[i]}}\label{a:sec:field}

All nine source manuscripts begin by fixing the same field, the same
monomial scale, the same valuation and the same class-size discipline.
This section states that material once.  Nothing in it is new, and
nothing in it depends on any analytic notion; every later part is
entitled to use it without further comment.

\subsection{Class-size discipline}\label{a:sec:size}

We work in \textsf{NBG} with global choice. The symbols $\No$ and
$\K$ denote proper classes. Any specified set of coefficient data can
be placed in a set-sized Hahn field, but assertions quantifying over
\emph{all} surreal radii or scales retain their ambient class meaning;
replacing that ambient field by a fixed working field changes those
assertions.  The following standing rule is a hypothesis of the whole
theory and is never relaxed.

\begin{remark}\label{a:rule:size}
In any single construction, every support, every coefficient family,
every polynomial and every formal series is \emph{set}-sized.  A
notation such as $\K[[X]]$ denotes a class of set-indexed coefficient
sequences, not a set of all of them.  No genuinely proper-class family
of nonzero terms is ever summed.  Every integer indexing a power series,
a derivative or a multiplicity is an \emph{ordinary finite} integer, and
every polynomial has ordinary finite degree with a set of coefficients.
\end{remark}

\noindent
The rule is what makes the Hahn machinery of
Section~\ref{a:sec:summability} applicable at all: Neumann~\cite{Neumann}'s lemma is a
statement about well-ordered \emph{sets}, and the valuation constructed
below takes values in a proper class.  Assertions about that valuation
are assertions about explicitly displayed normal forms; they are not
assertions about a class-sized object in the codomain position of a
set-valued function.

\subsection{Algebra, conjugation and the modulus}\label{a:sec:algebra}

The class $\No$ of surreal numbers is a real closed ordered field in the
class-sized sense.  Adjoining a square root of $-1$ gives
\[
 \K=\No[i]=\{x+iy:\ x,y\in\No\},\qquad i^2=-1 ,
\]
and the usual algebraic proof that a real closed field acquires an
algebraically closed quadratic extension applies to each individual
(finite-degree, set-coefficient) polynomial.  Hence $\K$ is
algebraically closed.  Conjugation and the modulus are
\[
 \overline{x+iy}=x-iy,\qquad
 \abs{x+iy}=(x^2+y^2)^{1/2}\in\No_{\ge0},
\]
the square root being supplied by real closedness.

\begin{proposition}[Modulus identities over an ordered field]\label{a:prop:triangle}
For all $z,w\in\K$,
\[
 \abs{zw}=\abs z\,\abs w,\qquad
 \abs{z+w}\le\abs z+\abs w,\qquad
 \abs{\operatorname{Re}z},\ \abs{\operatorname{Im}z}\le\abs z .
\]
Every nonzero $z$ is invertible, with $z^{-1}=\overline z/\abs z^{2}$.
\end{proposition}

\begin{proof}
Multiplicativity follows from $z\bar z\cdot w\bar w=(zw)\overline{(zw)}$.
For the triangle inequality, write $z=x+iy$, $w=u+iv$ and use the
Lagrange identity
\[
 (xu+yv)^2+(xv-yu)^2=(x^2+y^2)(u^2+v^2),
\]
valid in any commutative ring.  Since $(xv-yu)^2\ge0$ in the ordered
field $\No$, this gives the Cauchy--Schwarz inequality
$(xu+yv)^2\le(x^2+y^2)(u^2+v^2)$, hence
$xu+yv\le\abs z\abs w$.  Adding $\abs z^2+\abs w^2$ to twice this
inequality and taking the positive square root of both sides --- legal
because squaring is order preserving on $\No_{\ge0}$ --- yields
$\abs{z+w}\le\abs z+\abs w$.  The coordinate bounds follow from
$x^2\le x^2+y^2$, and the inverse formula from $z\bar z=\abs z^2$.
\end{proof}

No completeness, no limit and no archimedean property is used; only the
order and the existence of positive square roots.  This is why the
estimate survives verbatim over a proper-class field.

\subsection{Complex normal forms and the monomial scale}\label{a:sec:normal}

Combining the Conway~\cite{Conway} normal forms of the real and imaginary parts, and
using that a union of two well-ordered subsets of a linear order is well
ordered, every surcomplex number has a unique representation
\begin{equation}\label{a:eq:normal}
 z=\sum_{\gamma\in S}c_\gamma\,\tm^{\gamma},
 \qquad c_\gamma\in\CC\setminus\{0\},
 \qquad \tm^{\gamma}:=\omega^{-\gamma},
\end{equation}
with $S=\supp(z)\subset\No$ a well-ordered \emph{set}.  Zero has empty
support.  Conversely every such set-supported expression is a surcomplex
number.

\begin{remark}[What $\tm^{\gamma}$ is, and what it is not]\label{a:rem:monomial}
The notation $\tm^{\gamma}=\omega^{-\gamma}$ refers to \emph{Conway's
monomial map} $\gamma\mapsto\omega^{-\gamma}$ from the additive ordered
group $\No$ onto the multiplicative group of monomials.  It is
\emph{not} the Gonshor~\cite{Gonshor} surreal exponential evaluated at
$-\gamma\log\omega$, and the two must never be silently identified for
arbitrary surreal exponents.  In particular the law
\[
 \tm^{\gamma+\delta}=\tm^{\gamma}\tm^{\delta}
\]
is a normal-form \emph{convention} about the monomial map, not a theorem
about surreal exponentiation.  This warning is attached here once and is
in force throughout the article.
\end{remark}

\subsection{The natural valuation, finite elements, standard part}\label{a:sec:valuation}

For $z\ne0$ with normal form \eqref{a:eq:normal} put
\[
 v(z)=\min\supp(z),\qquad
 \operatorname{lc}(z)=c_{v(z)}\in\CC\setminus\{0\},\qquad
 v(0)=+\infty .
\]
Then for all $z,w$
\begin{equation}\label{a:eq:valuation}
 v(zw)=v(z)+v(w),\qquad
 v(z+w)\ge\min\{v(z),v(w)\},\qquad
 v(\abs z)=v(z).
\end{equation}
The last identity is the useful one and deserves its reason: if
$z=\tm^{\gamma}(c+\eta)$ with $c\in\CC^{*}$ and $v(\eta)>0$, then
$z\bar z=\tm^{2\gamma}(\abs c^{2}+\eta_0)$ with $v(\eta_0)>0$, and the
positive square root of a finite element with nonzero standard part has
the same shape; so
\begin{equation}\label{a:eq:modulusleading}
 \abs z=\tm^{\gamma}\bigl(\abs c+\eta_1\bigr),\qquad v(\eta_1)>0 .
\end{equation}
Thus $v$ is an additive valuation with a \emph{proper-class} value
group.  By Rule~\ref{a:rule:size} this is harmless: each displayed
statement quantifies over explicitly given normal forms.

Set
\begin{align*}
 \Ov&=\{z\in\K:\ v(z)\ge0\}
     =\{z\in\K:\abs z<n\ \text{for some } n\in\NN\},\\
 \mm&=\{z\in\K:\ v(z)>0\}
     =\{z\in\K:\abs z<1/n\ \text{for every } n\ge1\}.
\end{align*}
$\Ov$ is the valuation ring of finite elements and $\mm$ its maximal
ideal of infinitesimals.  Extraction of the coefficient at exponent
zero is a surjective ring homomorphism
\begin{equation}\label{a:eq:st}
 \st:\Ov\longrightarrow\CC,\qquad \Ov/\mm\cong\CC ,
\end{equation}
so every finite surcomplex number is uniquely $z=a+\varepsilon$ with
$a=\st(z)\in\CC$ and $\varepsilon\in\mm$.  We call $a+\mm$ the
\emph{monad} of $a$, and more generally $b+r\mm$ a \emph{scaled monad}
for $b\in\K$ and $r\in\K^{*}$.  Note that $c+\mm$ makes sense for every
$c\in\K$, even when $c$ has no standard part.

\begin{lemma}[Crude valuation bound]\label{a:lem:valbound}
If $v(z)\ge\beta$ then $\abs z<\tm^{\beta-1}$.
\end{lemma}

\begin{proof}
$\tm^{-\beta}z$ is finite, hence of absolute value less than the
positive infinite number $\omega=\tm^{-1}$.  Multiply by
$\tm^{\beta}>0$.
\end{proof}

The factor $\omega$ is deliberately generous: it makes the bound
independent of any ordinary real size of a leading coefficient, which is
exactly what is needed when the leading coefficient function is only
known to be holomorphic.  The bound is used again in the Liouville
hierarchy of Section~\ref{e:p8}, where it is the reason a single surreal bound on
the finite halo is vacuous.

\begin{definition}[Halo over an ordinary domain]\label{a:def:halo}
For an ordinary open $U\subseteq\CC$ its \emph{halo} is
\[
 \halo U=\st^{-1}(U)=\{a+\varepsilon:\ a\in U,\ \varepsilon\in\mm\}
 \ \subseteq\ \Ov .
\]
For $c\in\K$ and $r\in\K^{*}$, $c+r\,\halo U$ is a \emph{scale chart},
with normalized coordinate $w=(z-c)/r$.
\end{definition}

Two warnings belong here, at the point of definition, and are not
repeated later.

\begin{remark}[The halo is not a surcomplex ball]\label{a:rem:halonotball}
Writing $\mathbb D$ for the ordinary open unit disk, $\halo{\mathbb D}$
is \emph{not} the surcomplex ball $\{z\in\K:\abs z<1\}$.  The ball
contains $1-\tm$, whose standard part $1$ lies on the ordinary boundary;
the halo omits the entire infinitesimal layer sitting just inside every
standard boundary point.  Every Schwarz-type statement in Section~\ref{e:p8} is a
statement about $\halo{\mathbb D}$ and not about the ball, and the
counterexample $(1+\tm)z$ there turns on precisely the omitted layer.
\end{remark}

\begin{remark}[Puncturing removes a monad]\label{a:rem:puncture}
$\halo{(U\setminus\{0\})}=\halo U\setminus(0+\mm)$ deletes a whole
monad, not a point.  The class $\halo U\setminus\{0\}$ is a different
object.  Removable-singularity statements in Section~\ref{e:p8} concern the former;
keeping the two puncturing operations distinct is what prevents a false
analogy with the classical punctured disk.
\end{remark}

\subsection{Set-sized working fields \texorpdfstring{$\K_\Gamma$}{K Gamma}}\label{a:sec:KGamma}

Occasionally a statement must locate objects inside a set-sized
subfield --- for instance to say that the roots produced by preparation
lie in an algebraically closed \emph{set}, or that a zero set is
countable.  The following device, which costs nothing, is used for that
and only for that.

For a set-sized divisible ordered subgroup $\Gamma\subseteq\No$ put
\[
 \K_\Gamma=\CC((\tm^{\Gamma})),\qquad
 R_\Gamma=\RR((\tm^{\Gamma})),
\]
the Hahn fields of set-supported series with exponents in $\Gamma$.  The
normal-form embeddings place them inside $\K$ and $\No$.  By the
Hahn-field theorem, $\K_\Gamma$ is algebraically closed
\cite[Corollary~4, p.~10 of the author PDF]{Poonen}. Since
$\K_\Gamma=R_\Gamma[i]$ and $R_\Gamma$ is ordered, $R_\Gamma$ is
real closed; in particular every root of a polynomial over
$\K_\Gamma$ already lies in $\K_\Gamma$.

\begin{proposition}[Containment of set-sized data]\label{a:prop:fragment}
Every \emph{set} of surcomplex numbers is contained in $\K_\Gamma$ for
some set-sized divisible ordered subgroup $\Gamma\subseteq\No$.
\end{proposition}

\begin{proof}
Take the union of the supports of the members of the set; it is a set.
Pass to the divisible additive subgroup of $\No$ that it generates;
that is again a set, being a countable union of sets.  Every original
normal form has support inside it.
\end{proof}

Proposition~\ref{a:prop:fragment} contains each specified set of data
and the set of algebraic outputs generated from it. It does not replace
the full ambient class in statements about every fine radius or every
scale. In particular the topology induced from $\K$ on $\K_\Gamma$
must be distinguished from the intrinsic valuation topology of
$\K_\Gamma$ (Remark~\ref{a:rem:notvaluation}).

\section{Conventions: one symbol per concept}\label{a:sec:conventions}

This article is assembled from nine independent developments of the same
subject.  Six of its load-bearing notions are genuinely inequivalent
across those developments, in the strong sense that a theorem proved for
one is \emph{false} for another and the sources exhibit the
counterexamples.  The present section fixes one symbol per concept and,
more importantly, records which concepts must be kept apart.  It proves
nothing; each pointer names the part where the corresponding theory is
developed. The historical reconciliation in \texttt{MERGE\_NOTES.md}
predates the corrections recorded in the maintained \texttt{README.md};
use the statements below for their present mathematical scope.

\subsection{The notation table}\label{a:sec:table}

\begingroup
\small
\begin{longtable}{@{}p{0.30\textwidth}p{0.22\textwidth}p{0.42\textwidth}@{}}
\caption{Fixed notation.  A symbol appearing here has this meaning
throughout and no other.}\label{a:tab:notation}\\
\toprule
\textbf{Concept} & \textbf{Symbol} & \textbf{Remark}\\
\midrule
\endfirsthead
\toprule
\textbf{Concept} & \textbf{Symbol} & \textbf{Remark}\\
\midrule
\endhead
\endfoot
\bottomrule
\endlastfoot
The surcomplex field & $\K=\No[i]$ &
 The letter $S$ never denotes the field; $S$ is a support.\\
Well-ordered set support & $S$; $\supp(z)$, $\supp(F)$ &
 Of a normal form or of a coefficient family.\\
Conway monomial scale & $\tm^{\gamma}=\omega^{-\gamma}$, $\tm=\omega^{-1}$ &
 Monomial map, not Gonshor exponentiation; see
 Remark~\ref{a:rem:monomial}.\\
Valuation ring; its maximal ideal & $\Ov$; $\mm$ &
 $\Ov=\{v\ge0\}$, $\mm=\{v>0\}$.\\
Ordinary holomorphic functions & $\Hol(U)$ &
 Written $\Hol$, never $\Ov$, which is reserved for the valuation
 ring.\\
Halo over an ordinary domain & $\halo U=\st^{-1}(U)$ &
 Definition~\ref{a:def:halo}; not a surcomplex ball.\\
Evaluation of a coherent section; canonical lift & $\ev F$; $\ev f$ &
 One symbol, because a lift \emph{is} a section supported at exponent
 $0$.\\
Coherent algebra; positive-support part; meromorphic coefficients &
 $\Hc(U)$; $\Hc_{>0}(U)$; $\Mc(U)$ &
 Section~\ref{c:p4} and Section~\ref{d:sec:residues}.\\
Family valuation; leading coefficient \emph{function} &
 $\vH(F)$; $f_{\vH(F)}$ &
 Not the pointwise valuation or leading coefficient of a \emph{value}.\\
Derivative; the operator itself & $F'$; $D_z$ &
 See Remark~\ref{a:rem:derivative}.\\
Neumann monoid of a positive well-ordered set & $\Neu S$ &
 Definition~\ref{a:def:neumannmonoid}.\\
Coefficientwise contour integral & $\intH$, $\ointH$ &
 Superscript permanent: not a limit of Riemann sums.\\
Formal residue & $\Res_{X=0}$ & In $\K((X))$; Section~\ref{b:part3} and Section~\ref{d:sec:residues}.\\
Cluster residue at an \emph{ordinary} centre $a$ & $\ResH a$ &
 Section~\ref{d:sec:residues}.\\
Actual local residue at a surcomplex point $b$ & $\Resat b$ &
 Section~\ref{d:sec:residues}.  Never conflate with $\ResH a$.\\
Standard part; finite decomposition & $\st$; $z=a+\varepsilon$ &
 $a=\st(z)\in\CC$, $\varepsilon\in\mm$.\\
Canonical exponential; twisted family & $\Exp$; $E_\lambda$, $E_0=\Exp$ &
 $\lambda$ an ordinary real parameter; Section~\ref{e:p9}.\\
Omnific integers; purely infinite surreals & $\Oz=\Jpi+\ZZ$; $\Jpi$ &
 $\Jpi$ is a real vector space; Section~\ref{e:p9}.\\
Ordinary unit disk; bounded Jordan domain; its boundary; a neighborhood
 of its closure & $\mathbb D$; $\Omega$; $\Gamma=\partial\Omega$; $V$ &
 Fixing $\mathbb D$ and $\Omega$ separately removes the clash with the
 derivative operator $D_z$.\\
Hermite division data & $P_0$, $u_0=f_0/P_0$, $T$, $Q_0$, $d=\deg P_0$ &
 Section~\ref{d:sec:preparation}.\\
The topology & fine topology &
 Definition~\ref{a:def:fine}.\\
\end{longtable}
\endgroup

\subsection{Four function classes and their proved inclusions}\label{a:def:classes}

\begin{definition}[Names of the classes]\label{a:def:classnames}
The following four adjectives are fixed, and are not synonyms.
\begin{enumerate}[label=\textup{(\alph*)},leftmargin=2.2em]
\item \textbf{Germ-analytic} (the \emph{local} class).  A function that
is, near each point, the Hahn sum of a power series
$\sum_{n\ge0}a_n h^{n}$ with \emph{arbitrary} coefficients
$a_n\in\K$, the family $(a_nh^{n})_n$ being required to be Hahn
summable: no ordinary domain, no growth condition, no radius of
convergence.  Developed in Section~\ref{b:part3}, where the germ algebra at any point
is shown to be exactly $\K[[X]]$.  The sources give this class six
names --- microholomorphic, surcomplex analytic germ, Hahn-analytic
germ, germ-analytic, fine-analytic, microanalytic --- all with the same
meaning; \emph{germ-analytic} is the one adopted here.
\item \textbf{Coherent}, the algebra $\Hc(U)$.  A series
$F=\sum_{\gamma\in S}f_\gamma \tm^{\gamma}$ with well-ordered set
support $S$ and every $f_\gamma\in\Hol(U)$ for \emph{one common
connected ordinary domain} $U$, realized as a function $\ev F$ on the
halo $\halo U$.  Developed in Section~\ref{c:p4}.
\item \textbf{Canonical lift} $\ev f$ of an ordinary $f\in\Hol(U)$: the
sub-class of $\Hc(U)$ supported at exponent $0$ alone.
\item \textbf{Differentiable}, reserved for the bare $\varepsilon$--$\delta$
difference-quotient condition over all positive surreal radii, with
\emph{no} power-series requirement.  Section~\ref{b:part2} and Section~\ref{b:part3}.
\end{enumerate}
The superscript $\#$ is used for (b) and (c) alike, because a lift
\emph{is} a section supported at exponent zero; the local class (a)
carries no superscript.  The inclusions are
\[
 \{\ev f\}\ \subsetneq\ \Hc(U)\ \subsetneq\
 \bigl\{\text{germ-analytic on }\halo U\bigr\}\ \subseteq\
 \bigl\{\text{differentiable on }\halo U\bigr\}.
\]
The first two inclusions are strict. Strictness of the last inclusion is
not established here; see Remark~\ref{b:no-goursat}.
\end{definition}

\begin{remark}[The non-transfer rule]\label{a:rem:classtransfer}
\emph{Every theorem in this article names the class it is proved for,
and no theorem is imported across classes.}  Concretely:
\begin{enumerate}[label=\textup{(\roman*)},leftmargin=2.4em]
\item The identity theorem, the maximum principle in global form,
uniqueness of primitives and the open mapping theorem hold in (b)
under the domain hypotheses stated below, and their unrestricted global
forms fail in (a). The coefficientwise Liouville theorem additionally
requires bounded entire coefficient functions.  The witnesses are the
locally constant $L$ of Section~\ref{b:part2} and its sharper variant
$b(z)=1$ exactly when $v(z)$ is an ordinal.
\item The sharp Schwarz and Schwarz--Pick inequalities hold in (c) and
\emph{fail} for general coherent self-maps of the halo, i.e.\ in (b).
The witness is $F=(1+\tm)z$, which fixes $0$, maps $\halo{\mathbb D}$
into itself, and has $\abs{F'(0)}=1+\tm>1$ (Section~\ref{e:p8}).
\item Results in (a) do not descend to (b): $\sum_{n\ge0}n!\,
\varepsilon^{n}$ is a legitimate germ on $\mm$ but is not the
restriction of any element of $\Hc(U)$ on an ordinary neighborhood of
$0$ having those derivatives, since its exponent-zero coefficient
function would need Taylor coefficients $n!$ (Section~\ref{c:p4}).
\item Class (d) supports \emph{only} the calculus rules --- uniqueness
of the derivative and the sum, product and chain rules --- together with
the counterexamples. No theorem deriving a power-series representation
from bare differentiability is proved here.
\end{enumerate}
The one bridge that \emph{is} proved, and that must be cited whenever a
statement about one derivative is used for the other, is the theorem of
Section~\ref{c:p4} that for a coherent section the coefficientwise derivative
$F'=\sum f_\gamma'\tm^{\gamma}$ equals the fine
$\varepsilon$--$\delta$ derivative of $\ev F$.
\end{remark}

\subsection{Three inequivalent residues}\label{a:def:residues}

\begin{definition}[The three residue symbols]\label{a:def:threeres}
\begin{enumerate}[label=\textup{(R\arabic*)},leftmargin=2.8em]
\item $\Res_{X=0}$, the \textbf{formal} residue: the coefficient of
$X^{-1}$ of an element of the Laurent field $\K((X))$, which by
definition has a \emph{finite} principal part.  Purely algebraic, valid
at any surcomplex centre.
\item $\ResH a$, the \textbf{cluster} residue at an \emph{ordinary} centre
$a\in\CC$: for $A=\sum_s a_s\tm^{s}$ with meromorphic coefficient
functions, $\ResH a A=\sum_s\bigl(\Res_a a_s\bigr)\tm^{s}$.
\item $\Resat b$, the \textbf{actual} local residue of a meromorphic germ
at a genuinely surcomplex point $b\in\K$.
\end{enumerate}
\end{definition}

\begin{remark}[Why three symbols]\label{a:rem:resdisagree}
(R1)/(R3) and (R2) \emph{disagree numerically at the same point}.  For
$A(\zeta)=1/(\zeta-\varepsilon)$ with $\varepsilon\in\mm$ nonzero, the
cluster residue at the ordinary centre $0$ is $\ResH0 A=1$, while the
local residue of $A$ at the exact point $0$ is $0$ --- $A$ has no pole
at $0$ at all, its pole sitting at $\varepsilon$.  Writing a single
symbol $\Res_a$ for both would convert a correct statement into a false
one.  The two bridge theorems that connect them, and the standing
example worked out in full, are in Section~\ref{d:sec:residues}.
\end{remark}

\subsection{Three inequivalent Liouville hypotheses}\label{a:rem:liouvillehyp}

Three different boundedness hypotheses circulate under the name
``bounded entire function'', and they give three different theorems.
They are stated as a hierarchy in Section~\ref{e:p8}; the discipline imposed here is
that no conclusion drawn under one may be restated under another.
\begin{enumerate}[label=\textup{(L\arabic*)},leftmargin=2.8em]
\item A single \emph{surreal} bound on the finite halo.  This hypothesis
is \textbf{vacuous}: as Section~\ref{e:p8} proves from
Lemma~\ref{a:lem:valbound}, every nonzero $F\in\Hc(\CC)$ already
satisfies $\abs{\ev F}<\tm^{\vH(F)-1}$ there, so the hypothesis holds
for every such $F$ and selects nothing.
\item A single \emph{ordinary real} bound.  This does \textbf{not} give
constancy; it gives an exact characterization (Section~\ref{e:p8}).  The standing
counterexample is $F=\tm z$, which is bounded by $1$ in ordinary real
terms on the whole finite plane and is not constant.
\item \emph{Coefficientwise} bounds, one ordinary real constant per
exponent.  This is the hypothesis that does force constancy, and whose
degree-$d$ version forces a polynomial of degree $\le d$.
\end{enumerate}

\subsection{Two Cauchy estimates that are not in conflict}\label{a:rem:cauchytwo}

Two estimates appear, and they look like a clash.  They are about
different quantities and both are kept.
\begin{itemize}[leftmargin=1.4em]
\item The \textbf{coefficientwise majorant}
$B_r(F)=\sum_s\bigl(\max\abs{f_s}\bigr)\tm^{s}$, with its Cauchy--Schwarz
majorant lemma.  Its coefficientwise definition is part of its
statement: $B_r(F)$ is \emph{not} in general the least upper bound of
$\abs{\ev F}$ on a surreal circle, and it is not claimed to be a norm.
\item The \textbf{pointwise surreal} estimate
$\abs{(D_z^{n}F)(c)}\le n!\,M/R^{n}$, obtained from a single pointwise
bound $\abs{\ev F}\le M$ on an ordinary circle via the positivity lemma
for the coefficientwise integral.
\end{itemize}
The standing warning is that a single coarse surreal bound does
\emph{not} control every coefficient function uniformly: for
$F=w(1+i\tm)$ on the unit circle one has the pointwise bound
$M=1+\tm^{2}$, while the coefficientwise majorant is $1+\tm>M$.  That
warning forbids a statement about \emph{coefficients}; it does not
forbid the estimate above, which is about the surreal \emph{modulus},
and it must not be used to suppress it.  Both live in Section~\ref{c:p5}.

\subsection{Two logarithms}\label{a:rem:twologs}

Two distinct functions are called ``the'' logarithm in the sources, and
they differ by exactly $2\pi i$ on an explicit class of points.  Both
are kept, with distinct names, in Section~\ref{e:p9}:
\begin{itemize}[leftmargin=1.4em]
\item the \textbf{global} fine-analytic logarithm $\mathcal L$ on all of
$\K^{\times}$, whose branch is chosen from the fine-locally constant
standard part of the unit direction --- this is the primary object, and
it is the one consumed by the contour no-go theorem of Section~\ref{b:part2};
\item the \textbf{principal} logarithm $\Log$ on the cut plane
$\K\setminus\No_{\le0}$, defined by
$\Log z=\log_{\No}\abs z+i\Arg z$ with $-\pi<\Arg z<\pi$.
\end{itemize}
At $z=-1-i\varepsilon$ with $\varepsilon\in\mm_{>0}$ the two values
differ by $2\pi i$.  Neither is a coherent primitive of $1/z$ on
$\halo{(\CC^{\times})}$.

\subsection{The derivative operator}\label{a:rem:derivative}

\begin{remark}\label{a:rule:derivative}
$D_z$ differentiates a \emph{function of a variable} and holds every
surcomplex coefficient fixed.  In particular $D_z(\omega)=0$, $\omega$
being a constant function.  It is \emph{not} the Berarducci~\cite{BMder}--Mantova~\cite{Mantova2026b}
field derivation on surreals viewed as transseries, whose normalization
has $\partial\omega=1$.  Section~\ref{b:part3} upgrades this notational warning into a
theorem: no field derivation on $\K$ fixing $\CC$ with
$\partial\omega=1$ can satisfy $\partial(\Exp z)=\Exp(z)\,\partial z$.
A reader importing surreal differential algebra or transseries results
would silently change the meaning of every derivative in this article.
\end{remark}

\section{Hahn summability over the monomial scales}\label{a:sec:summability}

This is the one load-bearing notion on which all nine sources agree
exactly, clause for clause.  Every infinite operation in the article ---
every power series, every geometric series, every preparation
recursion and every reversion tree --- uses this operation over Hahn
exponents. Ordinary Taylor expansions and contour integrals of the
complex coefficient functions are performed in $\CC$ first; their
values then enter the Hahn sum. These are distinct operations.

\subsection{The definition}\label{a:sec:sumdef}

\begin{definition}[Hahn summable family]\label{a:def:summable}
A \emph{set}-indexed family $(z_j)_{j\in J}$ of normal forms is
\textbf{Hahn summable} if
\begin{enumerate}[label=\textup{(\roman*)},leftmargin=2.4em]
\item \textbf{(well-ordering clause)}
$\bigcup_{j\in J}\supp(z_j)$ is well ordered, and
\item \textbf{(finiteness clause)} each exponent
$\gamma\in\bigcup_j\supp(z_j)$ belongs to $\supp(z_j)$ for only finitely
many $j$.
\end{enumerate}
Its \emph{Hahn sum} is the normal form obtained by performing, at each
exponent, the finite complex addition supplied by clause~(ii).
The same definition, verbatim, applies to families with coefficients in
any commutative ring, in particular in a ring $\Hol(U)$ of ordinary
holomorphic functions.  A Hahn summable family is also called
\emph{strongly summable}; the two terms are synonyms and the first is
used throughout.
\end{definition}

The definition is invariant under reindexing, and a jointly summable
family may be regrouped arbitrarily, each regrouping being a finite
rearrangement at each exponent.  Joint summability of a double family,
rather than the bare existence of two iterated sums, is what licenses an
interchange of summation; that distinction is used at every interchange
below and is never assumed.

\begin{remark}[The finiteness clause is never dropped]\label{a:rule:clauseii}
Clause~(ii) is not a technicality.  Verifying only clause~(i) proves
nothing: a family whose supports are all $\{0\}$ has well-ordered union
of supports and is not summable unless all but finitely many of its
members vanish.  Every summability verification in this article checks
both clauses.
\end{remark}

\subsection{Neumann's support lemma}\label{a:sec:neumann}

\begin{definition}[Neumann monoid]\label{a:def:neumannmonoid}
For a well-ordered set $S$ of strictly positive elements of an ordered
abelian group, write
\[
 \Neu S=\{0\}\ \cup\ \bigcup_{n\ge1}\underbrace{(S+\cdots+S)}_{n}
\]
for the additive monoid it generates.
\end{definition}

\begin{lemma}[Neumann support lemma]\label{a:lem:neumann}
Let $G$ be an ordered abelian group.
\begin{enumerate}[label=\textup{(\alph*)},leftmargin=2.2em]
\item If $A,B\subseteq G$ are well ordered, then $A+B$ is well ordered,
and every element of $A+B$ has only finitely many representations
$a+b$ with $a\in A$, $b\in B$.
\item If $S\subseteq G_{>0}$ is well ordered, then $\Neu S$ is well
ordered, and every element of $G$ is the sum of only finitely many
finite nondecreasing words over the alphabet $S$ --- \emph{counting all
word lengths together}, not merely finitely many for each fixed length.
\end{enumerate}
\end{lemma}

\begin{proof}[Proof sketch]
For (a): an infinite strictly decreasing sequence in $A+B$, or
infinitely many pairs with a fixed sum, can be thinned so that the
$A$-coordinates are nondecreasing, which forces a strictly decreasing
sequence in $B$ and contradicts well-ordering.  For (b), the
combinatorial mechanism is the well-quasi-ordering of finite words over
a well-ordered positive alphabet: a subsequence embedding of one word
into another cannot decrease the sum, since the alphabet is positive.
A decreasing sequence of sums would give a sequence of words with
no earlier word embedding coefficientwise into a later one, contrary
to that well-quasi-ordering. For distinct nondecreasing words, such an
embedding strictly increases the sum unless the words are equal:
either a letter strictly increases or an extra positive letter is
inserted. Thus infinitely many distinct words with a fixed sum are
also impossible.  This is the classical lemma of Neumann; we import it as
a foundational input.  All its applications below take place inside the
set-sized ordered group generated by the displayed supports, even though
the ambient value class is $\No$.
\end{proof}

\begin{remark}[The all-lengths clause is the load-bearing one]\label{a:rule:wordlength}
Part~(a) alone licenses only the product of \emph{two} Hahn series.
What licenses a geometric series $\sum_n(-H)^n$, the homogeneous
preparation recursion of Section~\ref{d:sec:preparation}, Hahn reversion, and the implicit-function
tree of Section~\ref{e:p8} is the all-lengths clause of part~(b): the number of
words of \emph{every} length summing to a given exponent is finite in
total.  Dropping the clause silently invalidates those four arguments
and no others replace it.
\end{remark}

\subsection{Substituting an arbitrary complex sequence}\label{a:sec:substitution}

\begin{corollary}[No growth condition]\label{a:cor:complexsub}
Let $(c_n)_{n\ge0}$ be an \emph{arbitrary} sequence of complex numbers
and $\varepsilon\in\mm$.  Then $(c_n\varepsilon^{n})_{n\ge0}$ is Hahn
summable.  The same holds for a formal series in finitely many
infinitesimal variables.  Evaluation is compatible with addition,
multiplication, formal differentiation, and composition $A(B(X))$ when
$B(0)=0$.
\end{corollary}

\begin{proof}
Let $S=\supp(\varepsilon)\subseteq\No_{>0}$; if $\varepsilon=0$ the
claim is trivial.  Then $\supp(c_n\varepsilon^{n})\subseteq\Neu S$, well
ordered by Lemma~\ref{a:lem:neumann}(b); a complex coefficient has
support $\{0\}$ and does not enlarge it.  A fixed exponent
$\gamma\in\Neu S$ is the sum of only finitely many finite words over
$S$ by the all-lengths clause, hence receives a contribution from only
finitely many $n$.  Both clauses of
Definition~\ref{a:def:summable} hold.  The stated identities are
verified exponent by exponent, where each verification is a finite
computation, and the regroupings they require are exactly the ones
Lemma~\ref{a:lem:neumann} permits.  For finitely many variables use the
union of their positive supports and note that a fixed total degree
contributes only finitely many monomials.
\end{proof}

Thus $\sum_{n\ge0}(n!)^{2}\varepsilon^{n}$ is a perfectly legitimate
element of $\K$ for every surcomplex infinitesimal $\varepsilon$,
although the associated complex series has radius of convergence zero.
No radius-of-convergence formula in terms of the ordinary magnitudes of
the coefficients is asserted anywhere in this article, and none is
available.

\begin{remark}[What this corollary does \emph{not} say]\label{a:rem:nosurcoeff}
Corollary~\ref{a:cor:complexsub} is stated for \emph{complex}
coefficients.  For arbitrary surcomplex coefficients $a_n\in\K$ the
family $(a_n\varepsilon^{n})_n$ need \emph{not} be summable: the union
$\bigcup_n\supp(a_n)$ need not even be well ordered.  What is true is
that after a single monomial rescaling the family becomes summable, by
the separated-degree-block lemma of Section~\ref{b:part3}.  The rescaling is not
cosmetic; it changes the chart, and Section~\ref{c:p4} states the resulting
obstruction beside the positive result.
\end{remark}

\subsection{The separating examples}\label{a:sec:separating}

The following two examples are the ones on which all later summability
discipline rests.  They are stated once, here.

\begin{example}[A convergent complex series that is not a Hahn sum]\label{a:ex:notsummable}
Regard $\sum_{n\ge0}2^{-n}$ as a family of constant surcomplex numbers.
Every member has support $\{0\}$, so the well-ordering clause holds
trivially while the finiteness clause fails at the single exponent $0$:
infinitely many members contribute there.  The family is therefore
\emph{not} Hahn summable.  This is not a defect to be repaired.  The
correct procedure is to compute the ordinary complex sum first, in
$\CC$, and to use the resulting number $2$ as a \emph{single
coefficient}.  Classical summation in the coefficient field and Hahn
summation over scales are different operations, and the article never
substitutes one for the other.
\end{example}

\begin{example}[A Hahn identity whose partial sums do not
converge]\label{a:ex:geometric}
With $\tm=\omega^{-1}$, the family $(\tm^{n})_{n\ge0}$ has supports
$\{n\}$, so both clauses hold and
\begin{equation}\label{a:eq:geom}
 \sum_{n\ge0}\tm^{n}=(1-\tm)^{-1}.
\end{equation}
This is an identity between normal forms, established exponent by
exponent.  It is \emph{not} the convergence of the sequence of finite
partial sums.  Indeed
\[
 (1-\tm)^{-1}-\sum_{n=0}^{N}\tm^{n}=\frac{\tm^{N+1}}{1-\tm},
\]
whose valuation is $N+1$; since $\omega>N+1$ for every ordinary finite
$N$, this remainder exceeds $\tm^{\omega}$.  Hence the ball
$B\bigl((1-\tm)^{-1},\tm^{\omega}\bigr)$ of the fine topology contains
\emph{no} partial sum at all, and \eqref{a:eq:geom} cannot be read as a
limit.
\end{example}

\section{Why summability is not convergence}\label{a:sec:fine}

The previous section supplied an infinite operation.  This section
proves that no topological limit notion can supply one in its place:
the natural topology on $\K$ is so fine that set-indexed limits collapse
entirely.  Everything here is a statement about the field $\K$; no
function class is involved, and nothing proved here is a theorem about
germ-analytic, coherent or lifted functions.

\subsection{The fine topology}\label{a:sec:finedef}

\begin{definition}[Fine topology]\label{a:def:fine}
The \textbf{fine topology} on $\K$ is generated by the balls
\[
 B(a,r)=\{z\in\K:\ \abs{z-a}<r\},\qquad a\in\K,\ r\in\No_{>0},
\]
where the radius ranges over \emph{all} positive surreal numbers.
Equivalently it is generated by the valuation balls
$\{z:v(z-a)>\alpha\}$, $\alpha\in\No$: indeed
$\tm^{v(r)+1}\prec r$ gives one comparison and a smaller monomial ball
the other.
\end{definition}

The sources call this the full surreal topology, the neighborhood
topology and the all-radii topology; all three name this one topology.

\subsection{Uniform discreteness of every set}\label{a:sec:discrete}

\begin{theorem}[Set-sized convergence collapses, with a uniform
radius]\label{a:thm:discrete}
Let $A\subseteq\K$ be a \emph{set}.  Then:
\begin{enumerate}[label=\textup{(\roman*)},leftmargin=2.4em]
\item $A$ is closed and discrete in the fine topology;
\item if $A$ has at least two elements, there is a \emph{single}
$r\in\No_{>0}$ with $\abs{a-b}>r$ for \emph{every} pair of distinct
$a,b\in A$;
\item every net indexed by a set which converges in the fine topology is
eventually equal to its limit.  In particular every convergent ordinary
sequence is eventually constant.
\end{enumerate}
\end{theorem}

\begin{proof}
The mechanism is the surreal cut property (the simplicity theorem), not
an infimum and not an order limit.  For any \emph{set}
$D\subseteq\No_{>0}$, the cut $\{0\mid D\}$ is a surreal number $r$ with
$0<r<d$ for every $d\in D$.

For (ii), let $D=\{\abs{a-b}:a,b\in A,\ a\ne b\}$, a set of positive
surreals because $A$ is a set, and take $r$ as above; then
$\abs{a-b}>r$ for all distinct $a,b\in A$ simultaneously.  Hence
$B(a,r)\cap A=\{a\}$ for each $a\in A$, which is discreteness with one
radius serving all points.

For closedness in (i), fix $x\notin A$ and apply the same construction
to $\{\abs{x-a}:a\in A\}$, obtaining a ball about $x$ missing $A$
entirely.

For (iii), let $(x_j)_{j\in J}$ be a net with $J$ a set and proposed
limit $x$.  The values form a set, so
$D=\{\abs{x_j-x}:j\in J,\ x_j\ne x\}$ is a set of positive surreals;
choose $r$ below all of it.  Convergence forces the net to lie
eventually in $B(x,r)$, which contains no value other than $x$ itself.
The case $D=\emptyset$ is immediate.
\end{proof}

Clause~(ii) is the sharpest available form: the separating radius does
not merely depend on the point, it is uniform over the whole set.  It is
what makes the total disconnectedness of the fine topology quantitative.

\subsection{Clopen cosets, total disconnectedness, no paths}\label{a:sec:nopath}

\begin{proposition}[Clopen infinitesimal cosets]\label{a:prop:clopen}
Every coset $c+\mm$ is both open and closed, as are $\Ov$ and every
affine image $b+r\Ov$, $b+r\mm$ with $r\ne0$.  The fine topology is
totally disconnected.
\end{proposition}

\begin{proof}
A ball of infinitesimal radius about a point of $c+\mm$ stays inside
$c+\mm$, so the coset is open; its complement is a union of cosets and
is open for the same reason, so it is also closed.  For $\Ov$: a point
$z\notin\Ov$ has $\abs z\ge1$, and the ball $B(z,1/2)$ misses $\Ov$;
a point of $\Ov$ has a ball of radius $1$ inside $\Ov$ after translating
by its standard part.  Affine images of clopen classes under
$z\mapsto b+rz$ are clopen because that map is a bijection with
continuous inverse.  Finally, given $z\ne w$, the class
$z+(z-w)\mm$ is clopen, contains $z$, and does not contain $w$, since
$(w-z)/(z-w)=-1\notin\mm$.  So no connected subclass contains both.
\end{proof}

\begin{corollary}[No nonconstant continuous paths]\label{a:cor:nopath}
Every map $\varphi:[0,1]\to\K$ from the ordinary real unit interval
which is continuous for the fine topology is constant.
\end{corollary}

\begin{proof}
$[0,1]$ is a set, so $\varphi([0,1])$ is a set, hence a closed discrete
subspace by Theorem~\ref{a:thm:discrete}(i).  The continuous image of an
ordinary connected space is connected.  A connected discrete space is a
singleton.
\end{proof}

\begin{remark}[Consequence for contours]\label{a:rule:contour}
An ordinary parametrized complex circle, viewed inside $\K$, is
therefore \emph{not} a nonconstant continuous path for the fine
topology.  The coefficientwise contour functional $\intH$ introduced in
Section~\ref{c:p5} is consequently not a limit of Riemann sums along a fine-continuous
path, and its parameterization is nowhere asserted to be fine-continuous.
The superscript $\mathrm H$ is permanent for exactly this reason.
\end{remark}

\subsection{What the discreteness theorem does not say}\label{a:rem:nonconsequences}

Two non-consequences must be recorded, or a later reader will
over-apply Theorem~\ref{a:thm:discrete}.

\begin{remark}[$\K$ itself is not discrete]\label{a:rem:notdiscrete}
Theorem~\ref{a:thm:discrete} concerns \emph{sets}.  It does not make the
class $\K$ discrete: every ball $B(a,r)$ contains points other than $a$
--- for instance $a+r/2$ --- and in fact contains a proper class of
them.  There is no contradiction, because the class of points of a ball
is not a set and the cut construction in the proof does not apply to it.
\end{remark}

\begin{remark}[The induced topology is not a valuation
topology]\label{a:rem:notvaluation}
Theorem~\ref{a:thm:discrete} does not identify the fine topology induced
on a set-sized Hahn subfield $\K_\Gamma\subseteq\K$ with that subfield's
own valuation topology.  The valuation topology of $\K_\Gamma$ has balls
of radii indexed by $\Gamma$; the fine topology supplies, in addition,
radii smaller than $\tm^{\gamma}$ for every $\gamma\in\Gamma$, by
Proposition~\ref{a:prop:fragment} and the cut property.  The induced
topology is discrete; it is strictly finer when $\Gamma\ne\{0\}$.
For $\Gamma=\{0\}$ the intrinsic valuation is trivial and its topology
is discrete as well.
Example~\ref{a:ex:geometric} exhibits the difference concretely: the
partial sums of $\sum_n\tm^{n}$ converge to $(1-\tm)^{-1}$ in the
valuation topology of $\K_{\mathbb Q}$ and do not converge to it in the fine
topology.  A Hahn sum need not give convergence even in an intrinsic valuation
topology: in $\K_{\mathbb Q}$ the summable family
$(\tm^{1-1/(n+1)})_{n\ge1}$ has all its exponents below $1$, so its
partial-sum remainders never enter the valuation ball $\{v>1\}$.
Thus the topology and a separate convergence argument must always be
specified; the existence of a Hahn sum alone does not supply one.
\end{remark}

\subsection{What the fine topology is used for}\label{a:rem:fineuse}

The collapse proved above is not a reason to discard the fine topology.
It is used, genuinely and essentially, in three places, and in none of
them does it appear as a source of limits of sequences.

\begin{enumerate}[label=\textup{(\arabic*)},leftmargin=2.4em]
\item \textbf{The $\varepsilon$--$\delta$ derivative.}  A derivative is
defined by: $F'(a)=L$ iff for every $\eta\in\No_{>0}$ there is
$\delta\in\No_{>0}$ with
\[
 0<\abs h<\delta\ \Longrightarrow\
 \Bigl\lvert\frac{F(a+h)-F(a)}{h}-L\Bigr\rvert<\eta .
\]
Both radii range over \emph{all} positive surreals, and $h$ ranges over
a proper class of increments; no sequence is chosen.  Uniqueness of $L$
follows from Proposition~\ref{a:prop:triangle} alone.
Theorem~\ref{a:thm:discrete} says nothing against this definition: it
forbids nontrivial set-indexed limits, and this quantifier is not one.
\item \textbf{Open classes.}  The open mapping theorem and the strong
maximum modulus principle of Section~\ref{e:p8} are statements about actual
fine-open classes, at arbitrary surcomplex points; they are not
statements about standard parts, and they would be false if weakened to
such.
\item \textbf{Fine-open hypotheses in the identity theorem.}  One of the
four equivalent hypotheses of the identity theorem in Section~\ref{c:p4} is
vanishing on a nonempty fine-open subclass.  That this is a strong
hypothesis, and not a vacuous one, is exactly
Remark~\ref{a:rem:notdiscrete}: a fine neighborhood of a point contains
a proper class of points.
\end{enumerate}

\begin{remark}[The accumulation hypothesis lives in
$\CC$]\label{a:rule:accumulation}
Wherever an accumulation hypothesis appears later --- in the identity
theorem of Section~\ref{c:p4}, in the zero geometry of Section~\ref{d:sec:preparation} --- it is an
accumulation hypothesis on \emph{standard parts, in the ordinary complex
topology of $\CC$}.  By Theorem~\ref{a:thm:discrete} no nontrivial
set-sized sequence accumulates in the fine topology, so no theorem in
this article asserts, or may assert, the existence of one.
\end{remark}

%% ==================== p23.tex ====================

% =====================================================================
% PART 2 and PART 3 of the merged surcomplex document.
% Lead sources: manuscripts 02, 04, 06, 07 (with 01 for the differential
% class, the locally constant witness and the structured witness).
% Every theorem below names its function class.  Nothing proved for the
% coherent algebra $\Hc(U)$ or for canonical lifts $f^{\#}$ is used here.
% =====================================================================

\section{What fails, and the contour no-go theorem}
\label{b:part2}

The constructions of Sections \ref{b:part3} and following carry hypotheses
that look, on a first reading, like caution: one common ordinary domain
for all coefficient functions, a contour integral defined coefficientwise
rather than as a limit, a fundamental theorem restricted to coherent
primitives. None of them is caution. Each is forced by an explicit
counterexample or by a no-go theorem, and this part collects those
obstructions before any of the positive theory is built, so that every
later hypothesis can be read against the thing it excludes.

\subsection{Four function classes, and the rule against transfer}
\label{b:classes}

Four classes of functions occur in this article. They are genuinely
different, and a theorem proved for one of them is false for another.
We fix their names and symbols here and never use one name for another
object.

\begin{definition}[The four classes]
\label{b:def-classes}
Let $U\subseteq\CC$ be a nonempty connected ordinary domain.
\begin{enumerate}[label=\textup{(\alph*)},leftmargin=2.4em]
\item A function on a fine-open class is \emph{germ-analytic} at $a$ if
on some fine ball about $a$ it is the Hahn sum of a formal series
$\sum_{n\ge0}a_n(z-a)^n$ with \emph{arbitrary} coefficients $a_n\in\K$.
There is no growth condition on the coefficients, no ordinary domain, and
no relation between the expansions chosen at different points. The theory
of this class is Section \ref{b:part3}.
\item $\Hc(U)$, the \emph{coherent} algebra, consists of the Hahn series
$F=\sum_{\gamma\in S}f_\gamma\tm^{\gamma}$ with $S\subseteq\No$ a
well-ordered \emph{set} and every $f_\gamma\in\Hol(U)$ an ordinary
holomorphic function on the \emph{one} domain $U$. Its evaluation
$\ev F$ lives on the halo $\halo U=\st^{-1}(U)$.
\item The \emph{canonical lift} $f^{\#}$ of an ordinary $f\in\Hol(U)$ is
the element of $\Hc(U)$ supported at the single exponent $0$; this is why
lifts and evaluations carry the same superscript.
\item A function $F$ is \emph{differentiable} at $a$ (01's bare class) if
some $L\in\K$ satisfies
\[
 \forall\eta\in\No^{>0}\ \exists\delta\in\No^{>0}\ \forall h\in\K:
 \quad 0<\abs h<\delta\ \Longrightarrow\
 \left\lvert\frac{F(a+h)-F(a)}{h}-L\right\rvert<\eta ,
\]
all radii ranging over the positive surreals. No power-series
representation is required. We write $F'(a)=L$; uniqueness of $L$ is the
triangle inequality.
\end{enumerate}
\end{definition}

\begin{proposition}[The inclusions, and what witnesses them]
\label{b:inclusions}
Restricted to a halo $\halo U$,
\[
 \begin{aligned}
 \{f^{\#}:f\in\Hol(U)\}&\subsetneq\Hc(U)\\
 &\subsetneq\{\text{germ-analytic functions on }\halo U\}\\
 &\subseteq\{\text{differentiable functions on }\halo U\}.
 \end{aligned}
\]
The first inclusion is strict: the constant section $\tm\in\Hc(U)$
is not a lift, since its support is $\{1\}$. For the second, choose
$a\in U$ and set $q_a(z)=0$ on $a+\mm$ and $q_a(z)=1$ on the other
monads of $\halo U$. This is locally constant, hence germ-analytic.
If it were coherent, its values at ordinary points would force its
exponent-zero coefficient to be $0$ at $a$ and $1$ everywhere else in
$U$, which is not holomorphic. There is also the stronger local obstruction that a germ-analytic function need not extend to any coherent
section at all: $\sum_{n\ge0}n!\,\varepsilon^{n}$ is a legitimate
germ-analytic function on $\mm$ by
Theorem~\ref{b:realization}, but it is not the restriction of any element
of $\Hc(U)$ on an ordinary neighbourhood $U$ of zero having those same
derivatives, because such a section's exponent-zero coefficient function
would be an ordinary holomorphic function with Taylor coefficients $n!$
at zero, hence of zero radius of convergence. The third inclusion is
Theorem~\ref{b:germalgebra}.
\end{proposition}

\begin{remark}[Two things this article does not prove]
\label{b:no-goursat}
First, strictness of the last inclusion is \emph{not} established here.
No surcomplex Goursat theorem is proved: we do not assert that fine
differentiability at every point of a fine-open class implies a
power-series representation, and we exhibit no function that is
differentiable everywhere and germ-analytic nowhere. The inclusion is
recorded in the direction in which it is proved, and no positive theorem
about the differentiable class beyond the calculus rules --- uniqueness of
the derivative, the sum, product, quotient and chain rules, all of which
follow from finite inequalities in an ordered field with no limiting
sequence --- is used anywhere in this article. Second, nothing below
asserts that an arbitrary assignment of a germ to each point of a halo
arises from a coherent section; Proposition~\ref{b:inclusions} says the
opposite.
\end{remark}

\begin{remark}[The chart changes the function]
\label{b:chart-changes}
Proposition~\ref{b:inclusions} must be read beside the scale-embedding
theorem of Section \ref{c:p4} (Theorem~\ref{c:p4:bridge}), which says
that \emph{every} formal series over $\K$, including
$\sum n!\,X^{n}$, becomes a coherent family on all of $\CC$ with
polynomial coefficient functions after a single monomial rescaling
$z=\rho\zeta$, $\rho=\tm^{\lambda}$. There is no contradiction: the
obstruction above concerns the original chart at zero, and the positive
result concerns a rescaled chart. The rescaling changes which function is
being represented. A germ has \emph{some} sufficiently small coherent
chart; it does not follow, and it is false, that it has a chart of every
larger radius. That is the content of the all-scale rigidity theorem of
Section \ref{f:sec-allscale}.
\end{remark}

\subsection{The locally constant witness}

\begin{example}[Bounded, nonconstant, with zero derivative]
\label{b:L}
Let $\Ov=\{z\in\K:v(z)\ge0\}$ be the class of finite elements and put
\[
 L(z)=\begin{cases}1,&z\in\Ov,\\[2pt] 0,&z\notin\Ov.\end{cases}
\]
The class $\Ov$ is fine-open, since a finite perturbation of a finite
number is finite; its complement is fine-open, since $\abs h<\abs z/2$
leaves an infinite $z$ infinite. Hence $L$ is locally constant. It is
therefore germ-analytic at every point of $\K$, being locally the Hahn sum
of a constant series, and a fortiori differentiable with
$L'\equiv0$ and every positive-order derivative zero.
\end{example}

\begin{proposition}[Five classical statements fail in the germ-analytic class]
\label{b:fivefailures}
The function $L$ of Example~\ref{b:L} is bounded by $1$ on all of $\K$,
is nonconstant, vanishes on a nonempty fine-open class, has local maxima
of its modulus, and is not an open map. Consequently, for the
germ-analytic class \emph{on all of $\K$}, each of the following is false:
\begin{enumerate}[label=\textup{(\roman*)},leftmargin=2.4em]
\item the identity theorem (vanishing on a fine-open class does not force
vanishing);
\item the strong maximum principle;
\item Liouville's theorem (boundedness does not force constancy);
\item uniqueness of primitives (adding $L$ to a primitive changes nothing
about its derivative);
\item the open mapping theorem.
\end{enumerate}
Coherent analogues hold with their stated hypotheses: the identity
theorem is Theorem~\ref{c:p4:identity}; maximum and open mapping
principles use a connected ordinary base; the coefficientwise Liouville
theorem requires $U=\CC$ and bounded coefficient functions; and
coherent primitive uniqueness is in Section~\ref{c:p5}. In particular,
ordinary-real or surreal boundedness alone does not imply constancy
even for coherent sections (Section~\ref{e:p8}).
\end{proposition}

\begin{proof}
All five assertions about $L$ are immediate from local constancy and the
clopenness of $\Ov$. For (iv), note that $L$ is differentiable on $\K$
with derivative $0$, so it may be added to any primitive of any function.
For (v), the image of $\K$ is the two-point set $\{0,1\}$, which contains
no fine ball.
\end{proof}

The witness is not confined to the split between finite and infinite
elements, and in particular it does not need an infinite scale.

\begin{example}[The failure already occurs on one halo]
\label{b:q}
On $\halo{D}$, the halo of the ordinary unit disk, put $q(z)=0$ for
$z\in\mm$ and $q(z)=1$ otherwise. The cosets of $\mm$ are fine-clopen, so
$q$ is locally constant, germ-analytic, bounded, and nonconstant on
$\halo D$ --- although its ordinary base $D$ is connected. Thus coherence,
not connectedness of the base, is what the identity theorem of Section~\ref{c:p4} consumes. Moreover no $\K$-linear functional $\int$ on a
class of germ-analytic functions can satisfy
$\int_a^b H'(z)\,\dd z=H(b)-H(a)$ for all fine-locally analytic $H$ and
arbitrary endpoints: apply it to $q$ with $a$ and $b$ in different cosets
of $\mm$. This is the finite-endpoint shadow of the no-go theorem below.
\end{example}

\subsection{A bounded germ-analytic function with no fine limit}

Example~\ref{b:L} is locally constant in a rather coarse way. The next
witness is sharper: it separates points inside a single monad, at
arbitrarily small scales, and it kills the unrestricted Riemann
removable-singularity theorem.

\begin{example}[No fine limit at zero]
\label{b:ordinalindicator}
For $z\ne0$ define
\[
 b(z)=\begin{cases}
 1,&v(z)\text{ is an ordinal},\\[2pt]
 0,&\text{otherwise}.
 \end{cases}
\]
The valuation is locally constant away from zero --- if
$v(h)>v(z)$ then $v(z+h)=v(z)$ --- so $b$ is locally constant on
$\K^{\times}$, hence germ-analytic there, with $b'\equiv0$, and
$\abs b\le1$. Ordinals are cofinal in $\No$, and for every ordinal
$\alpha$ both $\tm^{\alpha}$ and $\tm^{\alpha+1/2}$ lie in the ball
$\abs z<\tm^{\beta}$ once $\alpha>\beta$, while
$b(\tm^{\alpha})=1$ and $b(\tm^{\alpha+1/2})=0$. Hence $b$ has no fine
limit at zero, although it is bounded and germ-analytic on a full
punctured fine neighbourhood.
\end{example}

\begin{proposition}[No unrestricted removability theorem]
\label{b:noremoval}
There is no theorem asserting that a germ-analytic function bounded on a
punctured fine neighbourhood of a point extends germ-analytically across
that point. The witness is Example~\ref{b:ordinalindicator}.
\end{proposition}

The correct removability statements are coefficientwise and belong to the
coherent class: a coherent section extends across an ordinary puncture
exactly when each ordinary coefficient function is bounded near it
(Section \ref{e:p8}), and a scalar surreal bound on $\ev F$ is not
sufficient there either. Two further points must not be lost. First,
the punctured halo $\halo{(U\setminus\{c\})}$ deletes the \emph{whole
monad} $c+\mm$, not the single point $c$; the two puncturing operations
are different classes and keeping them apart prevents a false analogy.
Second, for a \emph{finite-meromorphic} germ --- an element of $\K((X))$,
with finitely many negative powers --- boundedness on a small punctured
ball \emph{does} force a removable singularity, by
Proposition~\ref{b:formalres}. Example~\ref{b:ordinalindicator} is not
finite-meromorphic at zero, and it is exactly the gap between ``germ-analytic
on a punctured neighbourhood'' and ``finite principal part'' that it
occupies.

\subsection{A structured witness: unit modulus, all derivatives zero}

The locally constant examples can be dismissed as pathologies of the
indicator function. They cannot: the same phenomenon occurs for a
quotient of two everywhere-germ-analytic exponential homomorphisms.

\begin{proposition}[The twisted-exponential witness]
\label{b:twistwitness}
Let $\Exp=E_0$ be the canonical surcomplex exponential and $E_\lambda$
($\lambda\in\RR$) the twisted family of Section \ref{e:p9}
(Theorem~\ref{e:thm-exp}). For $\lambda\notin2\pi\ZZ$ the quotient
\[
 \Theta_\lambda:=E_\lambda/E_0:\K\longrightarrow\K^{\times}
\]
is a nonconstant function of unit modulus on all of $\K$, germ-analytic
everywhere, with every positive-order derivative identically zero. Indeed
$\Theta_\lambda(i\omega)=e^{i\lambda}\ne1=\Theta_\lambda(0)$.
\end{proposition}

\begin{proof}
Both $E_\lambda$ and $E_0$ are homomorphisms $(\K,+)\to(\K^{\times},\cdot)$
with the same modulus law $\abs{E_\lambda(x+iy)}=\exp_{\No}(x)$, so the
quotient is a homomorphism into the unit circle of $\K$. It differs from
$E_0$ only through the real coefficient of the monomial $\omega$ in the
imaginary part; for an \emph{infinitesimal} increment $h$ that
coefficient is zero, so $\Theta_\lambda(a+h)=\Theta_\lambda(a)$ for all
$h\in\mm$. Hence $\Theta_\lambda$ is constant on every monad --- in
particular germ-analytic with constant germ, so all positive-order
derivatives vanish --- while the displayed values show it is globally
nonconstant.
\end{proof}

This is the same obstruction as Example~\ref{b:L}, but now carried by a
function with a complete differential and algebraic description. It shows
that $E'=E$, $E(0)=1$, agreement with $\exp_{\CC}$ and with the real
surreal exponential, and the correct modulus law do not determine a
function on $\K$: local analytic data impose no continuation between
parts of $\K$ that are infinitely far apart.

\subsection{The contour no-go theorem}

Section \ref{c:p5} defines the contour functional coefficientwise,
$\intH_\Gamma F\,\dd z:=\sum_\gamma\tm^{\gamma}\int_\Gamma f_\gamma\,\dd z$
along an \emph{ordinary} piecewise-$C^1$ path $\Gamma$, and restricts its
fundamental theorem to \emph{coherent} primitives. The following theorem
is the reason both decisions are forced. Its input is the global
fine-analytic logarithm $\mathcal L$ of Section \ref{e:p9}
(Theorem~\ref{e:thm-globallog}): a function defined on \emph{all} of
$\K^{\times}$, germ-analytic at every point, with
$\mathcal L'(z)=1/z$, whose branch is chosen from the fine-locally
constant standard part of the unit direction $z/\abs z$.

\begin{theorem}[Contour no-go]
\label{b:nogo}
Let $\mathcal C$ be any class of functions that are germ-analytic on a
fine neighbourhood of the ordinary unit circle
$\Gamma_1=\{\abs z=1\}\subset\CC$ and that contains $z\mapsto1/z$
together with $G'$ for every $G$ germ-analytic on all of $\K^{\times}$.
There is no functional
\[
 I:\mathcal C\longrightarrow\K
\]
satisfying both of the following:
\begin{enumerate}[label=\textup{(\roman*)},leftmargin=2.4em]
\item \textup{(classical normalization)}
$I(1/z)=2\pi i$;
\item \textup{(fundamental theorem for germ-analytic primitives)}
$I(G')=0$ for every $G$ that is germ-analytic on all of $\K^{\times}$.
\end{enumerate}
\end{theorem}

\begin{proof}
Take $G=\mathcal L$, the global fine-analytic logarithm of
Theorem~\ref{e:thm-globallog}. It is germ-analytic at every point of
$\K^{\times}$ and $\mathcal L'(z)=1/z$ identically there. Hypothesis
(ii) gives $I(1/z)=I(\mathcal L')=0$; hypothesis (i) gives
$I(1/z)=2\pi i\ne0$.
\end{proof}

The theorem is a genuine obstruction, not a defect of one construction.
Three consequences fix the architecture of the rest of the article.

\begin{corollary}[The fundamental theorem cannot be widened]
\label{b:nogo-fundamental}
The identity $\intH_\Gamma F'\,\dd z=\ev F(\Gamma(1))-\ev F(\Gamma(0))$ of
Section \ref{c:p5} holds for $F\in\Hc(U)$ and cannot be extended to
germ-analytic primitives, even ones that are germ-analytic on all of
$\K^{\times}$. The restriction of the fundamental theorem to the coherent
category is therefore forced.
\end{corollary}

\begin{corollary}[$\mathcal L$ is not a coherent primitive of $1/z$]
\label{b:log-not-coherent}
$\mathcal L$ is not the evaluation of any element of
$\Hc(\CC^{\times})$ on $\halo{(\CC^{\times})}$.
\end{corollary}

\begin{proof}
If it were, Corollary~\ref{b:nogo-fundamental} applied to the closed
ordinary unit circle would give
$\ointH_{\Gamma_1}\dd z/z=0$, whereas the coefficientwise integral of
$1/z$ along $\Gamma_1$ is its ordinary value $2\pi i$.
\end{proof}

\begin{corollary}[Uniqueness of primitives is a coherent statement]
\label{b:primitive-uniqueness-scope}
Within $\Hc(U)$ a primitive is unique up to an additive constant of $\K$.
No such uniqueness holds in the germ-analytic or the differentiable class:
by Example~\ref{b:L} a primitive there may be modified by any locally
constant function.
\end{corollary}

\begin{remark}[What the no-go does \emph{not} say]
\label{b:nogo-scope}
Theorem~\ref{b:nogo} does not say that no useful surreal integration
theory exists, and the coefficientwise functional $\intH$ is not proposed
as the unique notion of surreal integration; it addresses a different
input class from, for instance, integration programs for surreal
extensions of real functions. What it says is that on the ordinary unit
circle the classical normalization and an unrestricted fundamental theorem
are incompatible, so any functional deserving the name must give up one of
them, and this article gives up the second in the explicit and minimal way
recorded in Corollary~\ref{b:nogo-fundamental}. Two further points
survive with it. The local Laurent residue obstruction of
Proposition~\ref{b:formalprimitive} is untouched: $\mathcal L$ is
germ-analytic at every nonzero point but is \emph{not} a
finite-meromorphic germ at zero, so it is not a counterexample to the
formal statement that a Laurent differential has a primitive in
$\K((X))$ iff its residue vanishes. And the coefficientwise definition of
$\intH$ is independently forced by Section~\ref{a:sec:fine}: an ordinary
parametrized circle is not a fine-continuous path, since every
fine-continuous map $[0,1]\to\K$ is constant, so no Riemann-sum reading of
the symbol is available in the first place.
\end{remark}

\subsection{Summary of what Section \ref{b:part2} forbids}

\begin{enumerate}[label=\textup{(\arabic*)},leftmargin=2.4em]
\item No unrestricted identity theorem, strong maximum principle,
Liouville theorem, primitive-uniqueness theorem or open mapping theorem
for the germ-analytic class (Proposition~\ref{b:fivefailures}); each
requires coherence, and each is restated with that hypothesis in Sections~\ref{c:p4} and \ref{e:p8}.
\item No unrestricted removability theorem for germ-analytic functions
bounded near a point (Proposition~\ref{b:noremoval}).
\item No determination of a global exponential, or of a global function on
$\K$, by local analytic data alone (Proposition~\ref{b:twistwitness}).
\item No contour functional on the ordinary unit circle with both the
classical value for $1/z$ and an unrestricted fundamental theorem
(Theorem~\ref{b:nogo}).
\end{enumerate}

\section{The local theory: germs over \texorpdfstring{$\K$}{K}}
\label{b:part3}

We now develop the local power-series class: arbitrary surcomplex
coefficients, no growth condition, no ordinary domain. Everything in this
part is a statement about germ-analytic functions in the sense of
Definition~\ref{b:def-classes}(a). No result of this part is used for
$\Hc(U)$ or for canonical lifts except through the evaluation theorem of
Section \ref{c:p4}, which is proved there; conversely, nothing proved for
$\Hc(U)$ is available here, as Section \ref{b:part2} has just shown.

The one surprise of the local theory is how large it is: the germ algebra
at any point is exactly $\K[[X]]$, with no radius-of-convergence condition
whatsoever. Classical growth restrictions disappear because the sole
infinite operation in use is Hahn summability (Section \ref{a:sec:summability}) --- a
support condition, not a limit.

\subsection{Separated degree blocks}

Everything rests on one piece of interval arithmetic in the surreal value
class. The point is to bound \emph{supports}, not merely leading
valuations.

\begin{lemma}[Separated degree blocks]
\label{b:blocks}
Let $(a_\nu)_{\nu\in\NN^{d}}$ be a family in $\K$ indexed by the
multi-indices of $d<\infty$ variables. Put
$S=\bigcup_\nu\supp(a_\nu)$; this is a set, although it need not be well
ordered. Choose $\Delta\in\No^{>0}$ with $\abs{s}<\Delta$ for every
$s\in S$, which is possible by surreal interpolation, and choose
$\lambda>2\Delta$. Put $\rho=\tm^{\lambda}$. Then:
\begin{enumerate}[label=\textup{(\roman*)},leftmargin=2.4em]
\item for each total degree $n=\abs\nu$ the support of $a_\nu\rho^{\,n}$
lies in the open interval $(n\lambda-\Delta,\;n\lambda+\Delta)$, and these
intervals are pairwise disjoint and increasingly ordered in $n$;
\item the family $(a_\nu\rho^{\,\abs\nu})_{\nu}$ is Hahn summable, and an
exponent occurring in it determines the total degree $n$;
\item for every $u\in\Ov^{d}$ the fully expanded family
$(a_\nu\rho^{\,\abs\nu}u^{\nu})_\nu$ is Hahn summable.
\end{enumerate}
\end{lemma}

\begin{proof}
(i) is the choice of $\Delta$ and $\lambda$: the upper endpoint
$n\lambda+\Delta$ of the $n$th interval is below the lower endpoint
$(n+1)\lambda-\Delta$ of the next precisely because $\lambda>2\Delta$.
For (ii), within a fixed total degree there are only finitely many $\nu$,
so the union of their shifted well-ordered supports is well ordered and
each exponent occurs finitely often there; the ordered union of the
disjoint increasing blocks is then well ordered, and by (i) an exponent
lies in exactly one block, which fixes $n$.

For (iii) write $u_j=c_j+\eta_j$ with $c_j\in\CC$ and $\eta_j\in\mm$, let
$B$ be the union of the supports of the $\eta_j$ --- a well-ordered subset
of $\No^{>0}$ --- and expand each finite product $u^{\nu}$ by the binomial
theorem. All resulting supports lie in $A+\Neu B$, where $A$ is the
well-ordered union from (ii). By Neumann~\cite{Neumann}'s lemma (Section \ref{a:sec:summability})
$A+\Neu B$ is well ordered, and a fixed exponent admits only finitely many
decompositions into a member of $A$ and a member of $\Neu B$, only
finitely many finite nondecreasing words from $B$ for the second factor,
only finitely many $\nu$ of the total degree pinned by the first, and only
finitely many binomial terms for each such $\nu$. This is joint
summability of the full expansion, which is what licenses regrouping.
\end{proof}

Two examples show that the hypothesis really is a bound on supports.

\begin{example}[A valuation bound is not enough]
\label{b:badcoeff}
For $n\ge1$ put $a_n=\tm^{\,\omega-n+1/n}$. Every $a_n$ is infinitesimal,
so every coefficient is small in valuation. Nevertheless at $h=\tm$ the
terms are $a_nh^{n}=\tm^{\,\omega+1/n}$, whose exponents strictly decrease
and have no least element: the family is not Hahn summable. An assertion
that ``all infinitesimal coefficients may be evaluated at every
infinitesimal argument'' is therefore false. It is the common bound
$\Delta$ on the whole union of supports in Lemma~\ref{b:blocks}, not a
bound on each $v(a_n)$, that does the work.
\end{example}

\begin{example}[A radius unavailable at rank one]
\label{b:rank}
The series $\sum_{n\ge0}\tm^{-n^{2}}X^{n}$ cannot be evaluated at any
nonzero argument of positive \emph{real} valuation $\lambda$: the
exponents $n\lambda-n^{2}$ are eventually strictly decreasing. At
$X=\tm^{\omega}$ the exponents are $n\omega-n^{2}$, with consecutive
differences $\omega-(2n+1)>0$, so the support is well ordered and the sum
exists. The usable radius can be smaller than $\tm^{\,n}$ for every
ordinary integer $n$; the value class of $\No$ has rank far above one and
Lemma~\ref{b:blocks} uses all of it.
\end{example}

\subsection{Universal realization}

\begin{theorem}[Realization of every formal series]
\label{b:realization}
Every $P\in\K[[X_1,\ldots,X_d]]$ has a Hahn-summed evaluation on a fine
neighbourhood of zero: with $\rho$ as in Lemma~\ref{b:blocks}, the family
defining $P(\rho u_1,\ldots,\rho u_d)$ is summable for every
$u\in\Ov^{d}$. On that scaled finite polydisk $\rho\Ov^d$ the values are bounded in
surcomplex modulus by one fixed surreal, namely
$\abs{P(\rho u)}<\tm^{-\Delta-1}$. Formal sums, products, derivatives and
compositions with zero constant inner term agree with the corresponding
operations on evaluations, after shrinking the neighbourhood if necessary.
For each nonzero one-variable series $P$, its value $P(h)$ is nonzero
for every sufficiently small $h\ne0$, with the scale allowed to depend
on $P$. In any finite number of variables, evaluation as a
\emph{function germ} is injective. There is no injective evaluation of
all of $\K[[X]]$ at a fixed point $h$: the polynomial $X-h$ already
precludes that claim, and no one radius realizes every series.
\end{theorem}

\begin{proof}
Summability on the scaled finite polydisk is Lemma~\ref{b:blocks}: every
exponent occurring is at least $-\Delta$, so the modulus of the value is
below $\tm^{-\Delta-1}$, a factor of $\omega$ dominating every real
leading coefficient; cancellation only raises the valuation.

Sums and products are expansions of jointly summable families. For a
composition with zero inner constant term, each coefficient of the formal
composite is a finite algebraic expression in the input coefficients;
include the supports of all those individual products, taken \emph{before}
cancellation, in the set used to choose $\Delta$. Then the separated-block
argument applies to the fully expanded family, there being only finitely
many algebraic terms at each fixed total degree, and grouping by total
degree reproduces the formal composite exactly. A further shrinking puts
the inner values inside a realization neighbourhood of the outer series.

For one-variable nonvanishing, let $m$ be the least index with
$a_m\ne0$ and write $\alpha_n=v(a_n)$, $\beta=v(h)\ge\lambda$. For $n>m$,
\[
 (\alpha_n+n\beta)-(\alpha_m+m\beta)
 =(\alpha_n-\alpha_m)+(n-m)\beta>0 ,
\]
because $\abs{\alpha_n-\alpha_m}<2\Delta<\lambda\le\beta$. So
$a_mh^{m}$ contributes the unique least exponent and the value is nonzero.

For injectivity on multivariable germs, let $P_m$ be the first nonzero
homogeneous part of a nonzero formal series. A nonzero polynomial over
$\K$ cannot vanish on all of $\mathbb Q^d$: induction on the number of
variables reduces this to the finite root bound in one variable. Choose
$u\in\mathbb Q^d$ with $P_m(u)\ne0$. The one-variable series $P(uX)$
is nonzero, so its evaluations are nonzero at all sufficiently small
nonzero $X$. These points lie in any prescribed fine neighborhood of
zero. Thus $P$ cannot represent the zero germ, even when evaluation is
restricted to surreal-real coordinate tuples.

Derivative compatibility is proved in Theorem~\ref{b:germalgebra} using
Proposition~\ref{b:remainder}. Those proofs use only the summability,
algebra and faithfulness clauses established here, not that compatibility
as an input.
\end{proof}

\begin{proposition}[Uniform Taylor remainder]
\label{b:remainder}
For a fixed $A=\sum_{n\ge0}a_nX^{n}\in\K[[X]]$ there are
$r,M\in\No^{>0}$ such that for every $N\ge0$ and every $h$ with
$\abs h\le r$,
\[
 \left\lvert A(h)-\sum_{n<N}a_nh^{n}\right\rvert
 \ \le\ \frac{M\abs h^{N}}{1-\abs h}\ \le\ 2M\abs h^{N}.
\]
\end{proposition}

\begin{proof}
Choose $\Delta>0$ bounding the absolute values of all exponents in
\emph{both} families $(a_n)$ and $(\abs{a_n})$, and choose
$\lambda>2\Delta$ as in Lemma~\ref{b:blocks}. Set
$r=\tm^\lambda$ and $M=\tm^{-\Delta-1}$.
Lemma~\ref{a:lem:valbound} gives $\abs{a_n}<M$ for every $n$.
The separated-block lemma applied to $(\abs{a_n})$ makes
$(\abs{a_n}\abs h^n)_n$ summable for $\abs h\le r$; the family
$(M\abs h^n)_n$ is summable by the positive-support geometric series.
Their termwise differences are therefore summable and nonnegative.
A Hahn sum of a summable family of nonnegative surreals is nonnegative:
at the least exponent in the union of nonempty supports, every
contributing leading coefficient is positive.

For completeness, the triangle inequality for this tail follows without
any limiting argument. Put $T=\sum_{n\ge N}a_nh^n$. If $T\ne0$, multiply
by $\theta=\bar T/\abs T$ and sum
$\operatorname{Re}(\theta a_nh^n)\le\abs{a_n}\abs h^n$; multiplication
by the fixed $\theta$ preserves Hahn summability. The preceding positivity
argument gives $\abs T\le\sum_{n\ge N}\abs{a_n}\abs h^n$; for $T=0$
this is immediate. Summing the bounds by $M\abs h^n$ and using the
geometric identity gives the first estimate. Since $r$ is infinitesimal,
$1/(1-\abs h)<2$, giving the second. Here $\abs h^0=1$, including at
$h=0$.
\end{proof}

The next two results separate infinitesimal realization with ordinary
coefficients from the much stronger requirement of one series being
summable at every surcomplex argument.

\begin{corollary}[No growth condition]
\label{b:nogrowth}
For every $\varepsilon\in\mm$ and every complex sequence $(c_n)$, the
family $(c_n\varepsilon^{n})_{n\ge0}$ is Hahn summable, with no growth
hypothesis on $(c_n)$. In particular $\sum_{n\ge0}(n!)^{2}\varepsilon^{n}$
and $\sum_{n\ge0}n!\,\varepsilon^{n}$ are germ-analytic functions on
$\mm$, although the corresponding classical power series have radius of
convergence zero.
\end{corollary}

\begin{proof}
Complex coefficients have support $\{0\}$, so Neumann's lemma applied to
$\Neu{\supp\varepsilon}$ suffices: at each exponent only finitely many
finite words from $\supp\varepsilon$ contribute, and the word length fixes
$n$.
\end{proof}

\begin{proposition}[Everywhere-summable means polynomial]
\label{b:globalseries}
If $\sum_{n\ge0}a_nX^{n}$ is Hahn summable at \emph{every} $X\in\K$, it is
a polynomial.
\end{proposition}

\begin{proof}
Suppose infinitely many $a_n$ are nonzero, and choose $\Delta$ bounding
all $\abs{\supp a_n}$ as in Lemma~\ref{b:blocks}, and take
$\lambda>2\Delta$. At $X=\tm^{-\lambda}$ the exponents
$v(a_n)-n\lambda$ strictly decrease along the nonzero indices: for
$m>n$ their difference is
$v(a_m)-v(a_n)-(m-n)\lambda<2\Delta-\lambda<0$.
They all lie in the union of the term supports, which is therefore
not well ordered.
\end{proof}

Proposition~\ref{b:globalseries} is a statement about one formal series,
not about functions: it emphatically does not say that a globally
germ-analytic function on $\K$ is a polynomial, and Section \ref{b:part2}
supplies the counterexamples. The classification of the functions that are
\emph{coherent at every scale} is the rigidity theorem of Section~\ref{f:sec-allscale}.

\subsection{The germ algebra}

\begin{definition}[Germ-analytic germs]
\label{b:def-germ}
For $a\in\K$, two functions germ-analytic at $a$ define the same germ if
they agree on some fine ball about $a$. Write $\mathcal G_a$ for the
$\K$-algebra of germs at $a$.
\end{definition}

\begin{theorem}[{The germ algebra is $\K[[X]]$}]
\label{b:germalgebra}
Hahn evaluation is an isomorphism of differential $\K$-algebras
\[
 \K[[X]]\ \xrightarrow{\ \sim\ }\ \mathcal G_a,
 \qquad X=z-a .
\]
Every germ-analytic function is infinitely differentiable near its centre
in the sense of Definition~\ref{b:def-classes}(d), its derivatives are
germ-analytic, differentiation is termwise, and
\[
 a_n=\frac{F^{(n)}(a)}{n!}.
\]
It admits a Hahn-summed Taylor expansion about every sufficiently nearby
point:
\[
 A(h_0+k)=\sum_{j\ge0}c_jk^{j},
 \qquad c_j=\sum_{n\ge j}\binom nj a_nh_0^{\,n-j},
\]
all displayed families being summable. Sums, products, quotients by germs
with nonzero constant term, and compositions of compatible germs are
germ-analytic, with the ordinary sum, product, quotient and chain rules.
Every germ has a primitive $C+\sum_{n\ge0}\frac{a_n}{n+1}X^{n+1}$, unique
up to an additive constant \emph{as a germ}.
\end{theorem}

\begin{proof}
Surjectivity is Definition~\ref{b:def-germ} together with
Theorem~\ref{b:realization}; injectivity is the series-dependent nonvanishing clause of that
theorem, since two series representing the same germ have a difference
vanishing at every sufficiently small nonzero $h$.

For the derivative at the centre write $A(h)=a_0+a_1h+h^{2}B(h)$; by
Proposition~\ref{b:remainder} with $N=2$ the difference quotient differs
from $a_1$ by at most $2M\abs h$, so for a prescribed target tolerance
$\eta$ it suffices to take $\abs h<\min\{r,\eta/(2M+1)\}$. Note that this
is the full $\varepsilon$--$\delta$ derivative of
Definition~\ref{b:def-classes}(d), with both radii ranging over all
positive surreals and in particular over infinitesimal tolerances; the
quantification is over a proper class of increments, so the discreteness
theorem of Section~\ref{a:sec:fine} does not make it vacuous.

The recentring formula follows by expanding $a_n(h_0+k)^{n}$ with the
finite binomial theorem --- $n+1$ terms at degree $n$ --- and regrouping,
which Lemma~\ref{b:blocks}(iii) licenses. Applying the first argument at
$h_0$ identifies $F'(h_0)=c_1=\sum_{n\ge1}na_nh_0^{\,n-1}$; iteration gives
all higher derivatives and the coefficient formula. The algebraic
operations and the chain rule are the corresponding formal identities,
transported by the composition clause of Theorem~\ref{b:realization}.
Formal integration divides $a_n$ by $n+1$ and shifts the degree, which
changes no support; a formal series with zero derivative has no
positive-degree coefficient in characteristic zero, which is the
uniqueness clause.
\end{proof}

\begin{remark}[Where uniqueness of the primitive stops]
\label{b:primitive-germ-only}
The uniqueness in Theorem~\ref{b:germalgebra} is uniqueness \emph{of a
germ} at one point. It does not propagate: by Example~\ref{b:L} and
Corollary~\ref{b:primitive-uniqueness-scope}, a germ-analytic primitive on
a disconnected-in-the-fine-sense domain may be modified by any locally
constant function without changing its derivative anywhere.
\end{remark}

\begin{corollary}[Isolated zeros of a germ]
\label{b:isolated}
A nonzero germ at $a$ has a unique finite order $m\ge0$, factors as
$A(h)=a_mh^{m}U(h)$ with $U(0)=1$, and after shrinking the ball
$U(h)-1$ is infinitesimal, so there are no zeros there if $m=0$, and otherwise the centre is
the only zero there.
Consequently, if a germ-analytic function has zeros meeting every
sufficiently small punctured fine ball about $a$, its germ at $a$ is zero.
\end{corollary}

\begin{proof}
Factor formally at the first nonzero coefficient and evaluate; choosing
the scale small enough makes every exponent in the evaluated support of
$U-1$ strictly positive, by Lemma~\ref{b:blocks}. A nonzero germ then has
nonzero values at all sufficiently small nonzero increments.
\end{proof}

The last sentence of Corollary~\ref{b:isolated} permits a proper class of
zeros. It is \emph{not} a sequential accumulation statement: by the
discreteness theorem of Section~\ref{a:sec:fine}, no nontrivial \emph{set} of
distinct points accumulates in the fine topology, so no argument in this
article may be phrased as ``the zeros have an accumulation point''. The
global identity theorem of Section \ref{c:p4} instead uses accumulation of
\emph{standard parts} in $\CC$, which is a statement about the ordinary
plane.

\subsection{Which derivative? The derivation obstruction}

\begin{remark}[$D_z$ differentiates the variable]
\label{b:Dz}
The operator $D_z$ used throughout this article differentiates a function
of a variable and holds every surcomplex \emph{coefficient} fixed. In
particular $D_z(\omega)=0$, $\omega$ being a constant function. It is not
the Berarducci~\cite{BMder}--Mantova~\cite{Mantova2026b} field derivation on surreal numbers viewed as
transseries, whose normalization has $\partial\omega=1$. Confusing the two
roles of a surreal number would invalidate every derivative in this
article. Nothing here imports surreal differential algebra, and nothing
here may be combined with it without first checking which operator is
meant.
\end{remark}

That warning is not a matter of taste. It is a theorem that the two
operators cannot be reconciled through the exponential.

\begin{theorem}[Derivation obstruction]
\label{b:derivation-obstruction}
Let $\partial$ be any field derivation on $\K$ fixing $\CC$ pointwise with
$\partial\omega=1$; the complexification of the normalized
Berarducci--Mantova derivation is one. Then the canonical surcomplex
exponential $\Exp$ of Section \ref{e:p9} cannot satisfy
\begin{equation}
\label{b:eq-derivobstruction}
 \partial\,\Exp(z)=\Exp(z)\,\partial z\qquad\text{for all }z\in\K .
\end{equation}
More generally, if $E:(\K,+)\to(\K^{\times},\cdot)$ is a homomorphism
satisfying \eqref{b:eq-derivobstruction} in place of $\Exp$, then every
element of $\ker E$ is a differential constant of $\partial$. Since
$\ker\Exp=2\pi i\,\Oz=i(\Jpi+2\pi\ZZ)$ contains $i\omega$, and
$\partial(i\omega)=i\ne0$, no such $\partial$ exists for $\Exp$.
\end{theorem}

\begin{proof}
For the general assertion, let $k\in\ker E$, so $E(k)=1$. Applying
\eqref{b:eq-derivobstruction} at $z=k$ gives
$0=\partial(1)=\partial E(k)=E(k)\,\partial k=\partial k$, so $k$ is a
differential constant. For the specific assertion, $\Exp(i\omega)=1$, so
the left side of \eqref{b:eq-derivobstruction} at $z=i\omega$ is
$\partial(1)=0$ while the right side is
$\Exp(i\omega)\,\partial(i\omega)=i\ne0$.
\end{proof}

\begin{remark}
\label{b:no-conflict-Dz}
There is no conflict with $D_z\Exp(z)=\Exp(z)$, which is proved in Section~\ref{e:p9}. That identity differentiates a variable while holding every
surcomplex coefficient fixed; \eqref{b:eq-derivobstruction} differentiates
the field's own elements as transseries. The two are different
mathematical tasks, and Theorem~\ref{b:derivation-obstruction} says
precisely that the omnific-period kernel is incompatible with the
normalization $\partial\omega=1$. No composition law compatible with the
Berarducci--Mantova derivation is claimed anywhere in this article.
\end{remark}

\subsection{Cauchy--Riemann equations and harmonic germs}

Call a function \emph{real germ-analytic} in the surreal variables
$(x,y)$ if it is the Hahn evaluation of a formal series in two variables
over $\K$, as in Theorem~\ref{b:realization} with $d=2$. Such a function
has a Fr\'echet differential in the fine sense: subtracting its
degree-zero and degree-one parts leaves a remainder bounded by
$M\norm h^{2}$ on a small polydisk, by the support estimates of
Lemma~\ref{b:blocks} and a finite sum over degree-two factors. Here
``$\No$-linear'' means linear on the two surreal-real coordinates, not
merely $\RR$-linear on $\No$ viewed as an infinite-dimensional real vector
space.

\begin{theorem}[Cauchy--Riemann criterion for germs]
\label{b:CRgerm}
Let $F=u+iv$ be real germ-analytic near a point of $\No^{2}$. Then $F$ is
germ-analytic in $z=x+iy$ near that point if and only if, on a
neighbourhood,
\begin{equation}
\label{b:eq-CR}
 u_x=v_y,\qquad u_y=-v_x .
\end{equation}
At a single point, \eqref{b:eq-CR} is exactly the condition that the real
Fr\'echet differential be $\K$-linear, i.e.\ multiplication by one element
of $\K$. The real and imaginary parts of a germ-analytic function satisfy
$u_{xx}+u_{yy}=v_{xx}+v_{yy}=0$.
\end{theorem}

\begin{proof}
A real $2\times2$ matrix commutes with multiplication by $i$, represented
by $\left(\begin{smallmatrix}0&-1\\1&0\end{smallmatrix}\right)$, exactly
when it has the form
$\left(\begin{smallmatrix}p&-q\\q&p\end{smallmatrix}\right)$; that is
\eqref{b:eq-CR} for the differential, and the remainder condition converts
it into the complex difference quotient.

For the germ statement introduce the formal coordinates $Z=X+iY$ and
$W=X-iY$, invertible over $\K$ because $2i\ne0$. Equations
\eqref{b:eq-CR} say $\partial_WF=\tfrac12(\partial_X+i\partial_Y)F=0$. If
they hold as functions on a neighbourhood, injectivity of multivariable
evaluation as a function germ (Theorem~\ref{b:realization}) makes this a formal identity, and
in characteristic zero it forces every term of positive $W$-degree to
vanish; the remaining formal series depends on $Z$ alone and realizes a
germ-analytic function. Conversely $\partial_W$ annihilates a series in
$Z$. Differentiating \eqref{b:eq-CR} once more, with formal mixed partials
commuting, gives the two Laplace equations.
\end{proof}

\begin{theorem}[Harmonic conjugate germs]
\label{b:harmonicgerm}
Every real germ-analytic harmonic germ $u(x,y)$ over $\No$ is the real part
of a germ-analytic germ. Its harmonic conjugate is unique up to a constant in $\No$ as a germ
(the constant need not belong to the ordinary real field $\RR$).
\end{theorem}

\begin{proof}
Translate the centre to $(0,0)$ and define, with coefficientwise formal
primitives,
\[
 v(X,Y)=-\int_0^{X}u_Y(s,0)\,\dd s+\int_0^{Y}u_X(X,\tau)\,\dd\tau .
\]
Then $v_Y=u_X$, and since $u_{XX}=-u_{YY}$,
\[
 v_X=-u_Y(X,0)+\int_0^{Y}u_{XX}(X,\tau)\,\dd\tau=-u_Y(X,Y).
\]
So $u+iv$ satisfies \eqref{b:eq-CR}; realize it and apply
Theorem~\ref{b:CRgerm}. A difference of two conjugates has both formal
partial derivatives zero, hence is a constant formal series.
\end{proof}

\begin{remark}[Two scope limits]
\label{b:CRscope}
Theorem~\ref{b:CRgerm} does not assert that the existence of two partial
derivatives at a point suffices for complex differentiability, and it does
not assert that an arbitrary pointwise differentiable function has a
normal-form expansion: the Fr\'echet or real germ-analytic hypothesis is
load-bearing (Remark~\ref{b:no-goursat}). Theorem~\ref{b:harmonicgerm} is
a differential identity at the level of germs; it asserts nothing about a
global Dirichlet problem. The coefficientwise harmonic-conjugate theorem
for $\Hc(U)$ on simply connected ordinary domains, the Poisson formula and
the coherent Dirichlet problem are in Section \ref{e:p8} and have
different hypotheses --- ordinary coefficientwise boundary data on an
ordinary disk.
\end{remark}

\subsection{Inverse, implicit function, and finite ramification}

\begin{theorem}[Local inverse and implicit function theorems for germs]
\label{b:inverse}
Let $F=(F_1,\ldots,F_d)$ be a germ-analytic germ at $a\in\K^{d}$ with
invertible Jacobian $J_F(a)$. Then $F$ has a germ-analytic inverse between
fine-open neighbourhoods of $a$ and $F(a)$.

If $\Phi(X,Y)$ is a tuple of $q$ germ-analytic germs with $\Phi(0,0)=0$
and $\partial\Phi/\partial Y(0,0)$ an invertible $q\times q$ matrix, there
is a unique germ $Y=\varphi(X)$ with $\varphi(0)=0$ and
$\Phi(X,\varphi(X))=0$; after shrinking neighbourhoods it gives all
sufficiently small solutions.
\end{theorem}

\begin{proof}
Translate centres and write $F(X)=AX+\sum_{n\ge2}F_n(X)$ with $A$
invertible and $F_n$ homogeneous of degree $n$. Seek
$G(Y)=A^{-1}Y+\sum_{n\ge2}G_n(Y)$. Once $G_2,\ldots,G_{n-1}$ are known,
the degree-$n$ part of $F(G(Y))-Y$ is $AG_n$ plus a known homogeneous
polynomial, so exactly one $G_n$ cancels it. The same recursion, or
substitution, makes the formal inverse two-sided. Realize the two
composition identities by Theorem~\ref{b:realization}. Germ-analytic
functions are continuous by Proposition~\ref{b:remainder}; choose a fine
ball $B$ on which $G\circ F=\mathrm{id}$, then a smaller ball $V$ about $F(a)$ with
$F\circ G=\mathrm{id}$ on $V$ and $G(V)\subseteq B$. Then $W=B\cap F^{-1}(V)$ is
fine-open and $F:W\to V$, $G:V\to W$ are inverse.

For the implicit statement apply the inverse theorem to
$(X,Y)\mapsto(X,\Phi(X,Y))$, whose Jacobian is block triangular with
invertible diagonal blocks, and restrict the inverse to pairs $(X,0)$;
uniqueness also follows degree by degree from the invertible
$Y$-Jacobian.

Every recursion here is a set-sized list of coefficient equations, solved
degree by degree. There is no appeal to completeness, to a contraction
principle, or to convergence of a sequence of Newton iterates --- none of
which is available in the fine topology.
\end{proof}

\begin{theorem}[Finite ramification normal form]
\label{b:ramification}
Let $F$ be a nonconstant germ-analytic germ at $a$ and let $m\ge1$ be the
least index with $F^{(m)}(a)\ne0$, so that
$F(a+X)-F(a)=bX^{m}+O(X^{m+1})$ with $b\ne0$. There is a germ-analytic
local coordinate $\phi$ with $\phi(0)=0$ and $\phi'(0)=1$ such that
\begin{equation}
\label{b:eq-ramification}
 F(a+X)=F(a)+b\,\phi(X)^{m}.
\end{equation}
The zero of $F-F(a)$ at $a$ has finite order $m$ and is isolated; $F$ is an
open map at $a$; on a suitable neighbourhood every target
$F(a)+w$ with $w\ne0$ sufficiently small has exactly $m$ distinct
preimages, all simple.
\end{theorem}

\begin{proof}
Factor $F(a+X)-F(a)=bX^{m}(1+Q(X))$ with $Q(0)=0$ and put
$\phi(X)=X(1+Q(X))^{1/m}$ using the formal binomial series; $\phi'(0)=1$,
so $\phi$ is an invertible local coordinate by Theorem~\ref{b:inverse},
and realization preserves the $m$th-power identity. In the coordinate
$w=\phi(X)$ the equation for the target $F(a)+y$ is $bw^{m}=y$. Since $\K$
is algebraically closed, for $y\ne0$ there are exactly $m$ distinct roots,
each of modulus $(\abs y/\abs b)^{1/m}$, so all lie in a prescribed small
coordinate ball once $\abs y$ is small enough; the derivative
$bmw^{m-1}\phi'(X)$ is nonzero at each. For $y=0$ the only root is $w=0$,
of multiplicity $m$. The same root argument shows that the image contains
a fine ball about $F(a)$.
\end{proof}

\begin{corollary}[Local openness and no local modulus maximum]
\label{b:localopen}
At a point where its germ is nonconstant, a germ-analytic function maps
every sufficiently small fine ball onto a class containing a fine ball
about its value. Its modulus has no local maximum there, and its real part
has no local maximum or minimum.
\end{corollary}

\begin{proof}
Openness is Theorem~\ref{b:ramification}. If $F(a)\ne0$, every fine ball
about $F(a)$ contains $(1+\delta)F(a)$ with $\delta>0$ infinitesimal,
of strictly larger modulus; if $F(a)=0$, it contains nonzero points. For
the real part, a ball about $b$ contains $b\pm\delta$ with $\delta>0$
infinitesimal real.
\end{proof}

\begin{remark}
\label{b:modulus-not-valuation}
The modulus in Corollary~\ref{b:localopen} is $\abs z=\sqrt{z\bar z}$ in
the ordered field $\No$, not an ultrametric absolute value read off the
valuation $v$. The valuation is locally constant on nonzero values, as
Example~\ref{b:ordinalindicator} exploits, and satisfies no such maximum
principle. Note also that ``nonconstant germ'' cannot be relaxed to
``globally nonconstant function'': $L$ of Example~\ref{b:L} is globally
nonconstant and is not an open map.
\end{remark}

\begin{remark}[Scope of the coherent inverse theorem]
\label{b:coherent-implicit-elsewhere}
Theorem~\ref{b:inverse} is a statement about germs; it produces no section
over a common ordinary base and says nothing about supports. The coherent result proved in Section~\ref{e:p8} is the
\emph{one-variable} local inverse theorem
(Theorem~\ref{e:thm-localinverse}), using positive-support Hahn reversion
after a scale change. The multivariable germ theorem above does not by
itself produce a coherent section on a common ordinary parameter domain.
An earlier version of this remark incorrectly referred to itself as a
proved multivariable coherent implicit-function theorem; no such theorem
is established in this article.
\end{remark}

\subsection{Local identity and maximum modulus, as germ statements}

\begin{corollary}[Local identity and maximum-modulus principles]
\label{b:localidentity}
A germ-analytic germ all of whose derivatives vanish at its centre is the
zero germ. A germ-analytic function whose modulus has a local maximum at a
point is constant \emph{as a germ} there.
\end{corollary}

\begin{proof}
The first is uniqueness of Taylor coefficients
(Theorem~\ref{b:germalgebra}). The second is
Corollary~\ref{b:localopen}: a nonconstant germ is locally open, which
contradicts a local maximum.
\end{proof}

These are germ statements and nothing more. Three things must travel with
them. They do not propagate to a whole fine-open class: $L$ and $q$
(Examples~\ref{b:L}, \ref{b:q}) have constant germs at every point and are
globally nonconstant. They are not sequential accumulation theorems, since
by Section~\ref{a:sec:fine} no nontrivial set of distinct points accumulates in
the fine topology; in particular no proof in this article may construct a
convergent sequence of zeros. And the passage from ``constant germ'' to
``one global constant'' is exactly what coherence buys, in
Theorem~\ref{c:p4:identity} and the global maximum principle of
Section \ref{e:p8}, and it is available for no other reason.

\subsection{Formal residues in \texorpdfstring{$\K((X))$}{K((X))}}

The field of fractions of $\K[[X]]$ is the Laurent field $\K((X))$,
whose elements have only finitely many negative powers; multiplication by
a large enough power of $X$ makes such an element germ-analytic, and it
evaluates on a sufficiently small punctured fine ball. We call these
\emph{finite-meromorphic germs} and set
\[
 \Res_{X=0}\Bigl(\sum_{n\ge n_0}a_nX^{n}\,\dd X\Bigr)=a_{-1}\in\K .
\]

\begin{remark}[Which residue this is]
\label{b:residue-discipline}
$\Res$ is the \emph{formal} residue, the coefficient of $X^{-1}$ in
$\K((X))$, taken at any surcomplex centre in an analytic local coordinate
$X$; the letter $X$ is a local coordinate and never the Hahn monomial
parameter $\tm^{\gamma}$. The actual residue $\Resat b$ of Section~\ref{d:sec:residues} is this same coefficient operation on the Laurent
germ at $b$, with $X=z-b$; the formal laws apply after that identification.
The different operation is the cluster residue $\ResH a$ at an ordinary
centre. For $1/(\zeta-\varepsilon)$ with
$\varepsilon\in\mm\setminus\{0\}$ one has $\ResH0=1$ while the actual
local residue at the exact point $0$ is zero, the pole being at
$\varepsilon$. The two bridge theorems of Section~\ref{d:sec:residues}
relate this cluster operation to actual local residues; no bridge is
needed between a Laurent germ's formal residue and its actual residue.

\end{remark}

\begin{proposition}[The only local primitive obstruction]
\label{b:formalprimitive}
A Laurent differential $\alpha\in\K((X))\,\dd X$ has a primitive in
$\K((X))$ if and only if $\Res_{X=0}\alpha=0$. Equivalently
\[
 \K((X))\,\dd X\ /\ \dd\,\K((X))\ \simeq\ \K\,\frac{\dd X}{X}.
\]
\end{proposition}

\begin{proof}
A formal derivative has no $X^{-1}$ coefficient, since the only possible
source is the degree-zero term, which differentiates to zero. Conversely
each $a_nX^{n}\dd X$ with $n\ne-1$ integrates to $a_nX^{n+1}/(n+1)$, and
these coefficients again form a Laurent series with finite negative part.
The only omitted term is $a_{-1}\dd X/X$.
\end{proof}

\begin{theorem}[Change of variable, with the ramification factor]
\label{b:reschange}
Let $g(X)=cX^{m}(1+O(X))$ with $c\ne0$ and $m\ge1$. For every
$A(Y)\in\K((Y))$,
\begin{equation}
\label{b:eq-reschange}
 \Res_{X=0}A(g(X))\,g'(X)\,\dd X=m\,\Res_{Y=0}A(Y)\,\dd Y .
\end{equation}
In particular residues are invariant under an invertible local coordinate
change, the case $m=1$.
\end{theorem}

\begin{proof}
Only the finitely many negative powers of $Y$ can affect the residue, so
it suffices to treat monomials $A(Y)=Y^{n}$. For $n\ne-1$ the pullback is
$\dd\bigl(g^{n+1}/(n+1)\bigr)$, of residue zero by
Proposition~\ref{b:formalprimitive}. For $n=-1$ write $g=cX^{m}u$ with
$u(0)=1$; then $g'/g=m/X+u'/u$ and $u'/u$ is a power series, whose residue
is zero, leaving the factor $m$. Positive powers of $Y$ pull back to
series with no negative powers and contribute nothing at exponent $-1$.
\end{proof}

\begin{corollary}[Local argument principle]
\label{b:localargument}
For a nonzero finite-meromorphic germ $F=X^{m}U$ with $U(0)\ne0$,
\[
 \Res_{X=0}\frac{F'}{F}\,\dd X=m\in\ZZ ,
\]
a positive $m$ being a zero order and a negative $m$ minus a pole order.
\end{corollary}

\begin{proposition}[Formal residue facts and removability]
\label{b:formalres}
For finite-meromorphic germs: the residue of a derivative is zero; a germ
bounded on some sufficiently small punctured fine ball by one fixed surreal
constant has zero principal part, hence a removable singularity.
\end{proposition}

\begin{proof}
The first is Proposition~\ref{b:formalprimitive}. For the second, a germ
with no principal part is germ-analytic and bounded on a small ball by
Theorem~\ref{b:realization}. If instead a pole of order $m>0$ is present,
write $A(X)=cX^{-m}(1+V(X))$ with $V(0)=0$, so $\abs{1+V}>1/2$ on a small
ball; given any proposed bound $M>0$, choose a nonzero $X$ small enough
that $\abs c/(2\abs X^{m})>M$.
\end{proof}

Proposition~\ref{b:formalres} is the exact scope of local removability,
and Proposition~\ref{b:noremoval} is its boundary: the hypothesis
``finite-meromorphic'' cannot be weakened to ``germ-analytic on a punctured
fine neighbourhood''.

There is also a formal Cauchy identity. In the \emph{explicit} ring
$\K((X))[[Y]]$, interpret
\[
 \frac1{X-Y}=\sum_{n\ge0}Y^{n}X^{-n-1};
\]
then for $P\in\K[[X]]$, coefficientwise residue in $X$ gives
\begin{equation}
\label{b:eq-formalCauchy}
 \Res_{X=0}\frac{P(X)}{X-Y}\,\dd X=P(Y).
\end{equation}
The ambient ring matters and is part of the statement: an infinite
negative Laurent tail in $X$ has \emph{not} been silently declared a
member of $\K((X))$.

\begin{theorem}[Lagrange--B\"urmann inversion]
\label{b:lagrange}
Let $F(X)=a_1X+a_2X^{2}+\cdots$ with $a_1\ne0$ and let $G$ be its
compositional inverse germ. For $n\ge1$ and every $H\in\K[[X]]$,
\[
 [Y^{n}]\,H(G(Y))
 =\frac1n\,[X^{n-1}]\,H'(X)\left(\frac{X}{F(X)}\right)^{n}.
\]
In particular,
\[
 [Y^{n}]\,G(Y)=\frac1n\,[X^{n-1}]\left(\frac{X}{F(X)}\right)^{n}.
\]
These are identities of realized germ-analytic germs as well as of formal
series.
\end{theorem}

\begin{proof}
By Theorem~\ref{b:reschange} with $Y=F(X)$ and $m=1$,
\[
 [Y^{n}]H(G(Y))=\Res_{X=0}\frac{H(X)F'(X)}{F(X)^{n+1}}\,\dd X .
\]
The residue of $\dd\bigl(H(X)F(X)^{-n}\bigr)$ vanishes; expanding that
derivative turns the right side into
$\frac1n\Res_{X=0}H'(X)F(X)^{-n}\dd X$, which is the stated coefficient.
Taking $H(X)=X$ gives the second formula. Theorems~\ref{b:realization} and
\ref{b:inverse} identify the formal inverse with the actual inverse germ
and realize both sides on a sufficiently small ball.
\end{proof}

\subsection{What the local theory does and does not give}

The germ algebra is $\K[[X]]$ exactly: arbitrary coefficients, no growth
condition, no radius of ordinary convergence, and the full
characteristic-zero formal calculus with actual function values and actual
fine derivatives. What it does not give is any global statement whatever.
Specifically, and to be read together with Section \ref{b:part2}:

\begin{enumerate}[label=\textup{(\arabic*)},leftmargin=2.4em]
\item No radius-of-convergence formula is claimed from the ordinary
magnitudes of the coefficients; the usable scale is produced by
Lemma~\ref{b:blocks} from a bound on the \emph{union of supports}, and
Example~\ref{b:badcoeff} shows a valuation bound will not do.
\item A Hahn sum is not a fine limit of its partial sums, and no theorem
here may be reproved by exhibiting a convergent sequence.
\item Theorem~\ref{b:realization} does not assert summability at every
infinitesimal for arbitrary surcomplex coefficients: the input must be
small relative to a field containing the coefficient data. For purely
complex coefficients every infinitesimal works
(Corollary~\ref{b:nogrowth}).
\item Agreement at ordinary complex points does not determine a
germ-analytic function; $\st$ itself is locally constant on $\Ov$, agrees
with the identity at every ordinary point, and differs at infinitesimal
ones.
\item The local identity, maximum modulus and open mapping statements are
germ statements. Their global forms require coherence and are proved in
Sections \ref{c:p4} and \ref{e:p8}, for $\Hc(U)$ only.
\item No classification of all functions on a punctured fine neighbourhood
is offered: finite-meromorphic germs and coherent quotients are specified
subclasses, not an exhaustive list, and no general theory of essential
singularities is developed here. The one value-distribution theorem in
this article concerns a single specified function and is in Section~\ref{e:p9}.
\end{enumerate}

%% ==================== p45.tex ====================

% =====================================================================
% PART 4.  The coherent theory: H(U).
% PART 5.  Contours, and the order-valued estimate.
% Lead sources: 01, 07, 02.  All labels carry the c: prefix.
% CLASS DISCIPLINE.  Every theorem below is stated for exactly one of
%   (a) germ-analytic  (b) coherent, H(U)  (c) the canonical lift f^#.
% Nothing is imported across those three classes.
% =====================================================================

\section{The coherent theory: \texorpdfstring{$\Hc(U)$}{H(U)}}
\label{c:p4}

Section~\ref{b:part3} developed the local power-series layer: germ-analytic functions,
with arbitrary surcomplex coefficients, no ordinary domain and no growth
condition. That layer is too wide to carry a global theorem. The
locally constant witness $L$ of Section~\ref{b:part2}
(the locally constant germ-analytic witness of Section~\ref{b:part2}) and its germ-analytic refinement
$b$ (the germ-analytic witness of Section~\ref{b:part2}) are germ-analytic on $\K$ and
$\K^\times$, respectively, and
destroy, at one stroke, the identity theorem, the global maximum
principle, Liouville's theorem, uniqueness of primitives and the open
mapping theorem.

This part introduces the coherent layer: one \emph{ordinary} connected
complex domain carries all coefficient functions at once. This common-domain
condition ties different monads to one another; additional hypotheses
for each global principle, such as entire coefficientwise boundedness
for Liouville, are stated with that theorem, and every global theorem in
Sections~\ref{c:p5}--\ref{e:p8} is a theorem about this class and not about
the germ class.

\subsection{The halo}
\label{c:p4:halo}

\begin{definition}[Halo of an ordinary domain]\label{c:p4:def-halo}
For an ordinary open $U\subseteq\CC$ put
\begin{equation}\label{c:p4:eq-halo}
 \halo U=\st^{-1}(U)=\{a+\varepsilon: a\in U,\ \varepsilon\in\mm\}
 \subseteq\Ov .
\end{equation}
This is an open class in the fine topology, being a union of cosets of
$\mm$.
\end{definition}

Two warnings belong to the definition and are made once here.

\begin{remark}[The halo is not a surcomplex ball]\label{c:p4:rk-notball}
$\halo{\mathbb D}$ consists of the \emph{finite} points whose standard
part has modulus strictly less than one. It is strictly smaller than
$\{z\in\K:\abs z<1\}$: for positive infinitesimal $\varepsilon$ the
point $1-\varepsilon$ lies in the surcomplex unit ball but not in
$\halo{\mathbb D}$, because its standard part lies on the unit circle.
The whole boundary monad is missing. No statement below about
$\halo{\mathbb D}$ is a statement about the surcomplex unit ball; in
particular the Schwarz--Pick theorem of Section~\ref{e:p8} is not, and
the failure witness $(1+\tm)z$ there turns on exactly the omitted
infinitesimal boundary layer.
\end{remark}

\begin{remark}[Puncturing deletes a monad]\label{c:p4:rk-punctured}
$\halo{(\mathbb D\setminus\{0\})}$ is the halo of the punctured ordinary
disk, and it omits the \emph{entire} monad $\mm$, not the single point
$0$. It is therefore different from $\halo{\mathbb D}\setminus\{0\}$.
The removable-singularity statements of Section~\ref{e:p8} concern the
former. Keeping the two puncturing operations apart is what prevents a
false analogy with the classical punctured disk.
\end{remark}

\subsection{Coherent Hahn sections}
\label{c:p4:def}

\begin{definition}[The coherent algebra]\label{c:p4:def-HU}
Let $U\subseteq\CC$ be a nonempty connected ordinary open set. A
\emph{coherent Hahn section} on $U$ is a formal expression
\begin{equation}\label{c:p4:eq-HU}
 F=\sum_{\gamma\in S}f_\gamma(Z)\,\tm^\gamma,
 \qquad f_\gamma\in\Hol(U),
 \qquad S=\{\gamma:f_\gamma\not\equiv0\}
 \ \text{a well-ordered \emph{set}}.
\end{equation}
The class of these is written $\Hc(U)$; it is a $\K$-algebra under
coefficientwise addition and Hahn multiplication, the latter being
well defined because at each exponent only finitely many products
contribute (Section~\ref{a:sec:summability}, \ref{a:sec:neumann}). Put
\[
 \Hc_{>0}(U)=\{F\in\Hc(U): \vH(F)>0\}\cup\{0\},
 \qquad
 F'=\sum_{\gamma\in S}f_\gamma'\,\tm^\gamma,
\]
and write $D_z$ for the operator $F\mapsto F'$ when the operator itself
is the subject. For $F\ne0$ define the \emph{family valuation} and the
\emph{leading coefficient function}
\[
 \vH(F)=\min S,\qquad f_{\vH(F)}\in\Hol(U).
\]
As for the scalar valuation, put $\vH(0)=\infty$, an added symbol.
\end{definition}

\begin{remark}[What ``coherent'' means, and what it does not]
\label{c:p4:rk-coherent}
The entire content of the word is: \emph{one} ordinary holomorphic
function per exponent, all of them on \emph{one} common connected
ordinary domain, with well-ordered set support. Three disclaimers
travel with it.
\begin{enumerate}[label=(\roman*)]
\item This is not the technical assertion that some sheaf is coherent as
a module in the sense of complex geometry. The word is used in the
sense just given and in no other.
\item The hypothesis is not free. An arbitrary family of ordinary
holomorphic germs, one at each point, whose radii of convergence shrink
to zero, is not an element of $\Hc(U)$ for any ordinary neighborhood
$U$. It is precisely this hypothesis that links monads with different
standard parts, and it is why the identity theorem
(\ref{c:p4:identity}), the global maximum principle
(Section~\ref{e:p8}) and the zero-cluster theorem (Section~\ref{d:sec:preparation}) hold at
all.
\item Coherence is a genuine restriction on functions, not a
reparameterization. The function $z\mapsto\st(z)$ on $\halo U$ is not
in $\Hc(U)$ when $U$ has more than one point: its values at ordinary
points would force it to be the canonical lift of the identity, which it
is not.
\end{enumerate}
The synonym ``holomorphic profile'' is recorded once and then dropped.
\end{remark}

\begin{remark}[$\vH$ is not a pointwise valuation]\label{c:p4:rk-vH}
$\vH(F)$ and $f_{\vH(F)}$ are attached to the \emph{family}. They are
not the valuation $v(F^{\#}(z))$ and the leading complex coefficient of
the surcomplex \emph{value} $F^{\#}(z)$, and the two disagree exactly
where the theory is interesting: at a point $z$ of the halo lying over a
zero of $f_{\vH(F)}$, the value has strictly larger valuation than
$\vH(F)$. The inequality in one direction is Theorem~\ref{c:p4:eval}.
\end{remark}

Because $U$ is connected, $\Hol(U)$ is an integral domain, so
$\vH(FG)=\vH(F)+\vH(G)$ and $\vH$ is a valuation on the fraction field
of $\Hc(U)$.

\subsection{Faithful evaluation on the halo}
\label{c:p4:evalsec}

\begin{definition}[Evaluation]\label{c:p4:def-eval}
For $F\in\Hc(U)$ and $z=a+\varepsilon\in\halo U$ set
\begin{equation}\label{c:p4:eq-eval}
 \ev F(a+\varepsilon)
 =\sum_{\gamma\in S}\sum_{n\ge0}
   \frac{f_\gamma^{(n)}(a)}{n!}\,\varepsilon^n\,\tm^\gamma .
\end{equation}
\end{definition}

\begin{theorem}[Faithful coherent evaluation]\label{c:p4:eval}
\emph{(Class (b).)} Let $U$ be connected and $F\in\Hc(U)$.
\begin{enumerate}[label=(\roman*)]
\item The double family in \eqref{c:p4:eq-eval} is a single jointly
Hahn summable family, so $\ev F$ is a well-defined function
$\halo U\to\K$.
\item $F\mapsto\ev F$ is an injective homomorphism of $\K$-algebras onto
its image, and it commutes with differentiation,
$\ev{(F')}=(\ev F)'$, with restriction to a nonempty connected ordinary
open subset, and with conjugation: if
$U^{*}=\{\bar z:z\in U\}$ and
$F^{*}=\sum_\gamma\overline{f_\gamma(\bar Z)}\,\tm^\gamma\in\Hc(U^{*})$,
then $\ev{(F^{*})}(z)=\overline{\ev F(\bar z)}$.
\item If $F\ne0$ and $\lambda=\vH(F)$, then for every $z=a+\varepsilon
\in\halo U$
\begin{equation}\label{c:p4:eq-reduction}
 v\bigl(\ev F(z)\bigr)\ge\lambda,
 \qquad
 \st\bigl(\tm^{-\lambda}\ev F(a+\varepsilon)\bigr)=f_\lambda(a).
\end{equation}
\end{enumerate}
\end{theorem}

\begin{proof}
Let $E=\supp(\varepsilon)\subseteq\No_{>0}$. Every term of
\eqref{c:p4:eq-eval} has support inside $S+\Neu E$. By the Neumann~\cite{Neumann}
support lemma (\ref{a:sec:neumann}) that class is a well-ordered set and
each of its elements admits only finitely many decompositions
$\gamma+\mu$ with $\mu$ a finite nondecreasing word in $E$; the word
length is exactly the Taylor order $n$, so only finitely many pairs
$(\gamma,n)$ contribute at a fixed exponent. That is precisely the
two-clause summability criterion of Section~\ref{a:sec:summability}, and it also
licenses the regroupings used below.

Sums and products: apply the same argument to two families and use the
ordinary Leibniz and Taylor product identities at $a$; each Hahn
exponent involves a finite convolution only. Constants of $\K$ are
represented by constant coefficient functions, so the map is $\K$-linear
and not merely $\CC$-linear.

Injectivity: at an ordinary $a\in U$ the value is
$\ev F(a)=\sum_\gamma f_\gamma(a)\tm^\gamma$, already in normal form.
If $\ev F\equiv0$ then uniqueness of normal forms gives $f_\gamma(a)=0$
for every $\gamma$ and every $a\in U$, so every coefficient function
vanishes identically.

The identification with the fine derivative is proved below in
Theorem~\ref{c:p4:derivbridge}, using the jointly summable Taylor
expansion; no derivative identity is needed for the present algebraic
arguments. Restriction and conjugation are coefficientwise. For (iii), factor $\tm^{\lambda}$ out of
\eqref{c:p4:eq-eval}. Every remaining term except the $n=0$, $\gamma=
\lambda$ term has strictly positive valuation: either $\gamma>\lambda$,
or $n>0$ and the factor $\varepsilon^n$ has positive support. Hence the
exponent-zero coefficient of $\tm^{-\lambda}\ev F(z)$ is $f_\lambda(a)$.
\end{proof}

\begin{corollary}[Leading coefficients control zero shadows]
\label{c:p4:shadow}
\emph{(Class (b).)} If $F\in\Hc(U)$ is nonzero and $\ev F(z)=0$ for some
$z\in\halo U$, then $f_{\vH(F)}(\st z)=0$. Equivalently, no zero of
$\ev F$ lies over a non-zero of the leading coefficient function.
\end{corollary}

\begin{proof}
Immediate from the second identity in \eqref{c:p4:eq-reduction}.
\end{proof}

\begin{theorem}[Exact monad-wide Taylor expansion]\label{c:p4:monadtaylor}
\emph{(Class (b).)} Let $F\in\Hc(U)$, let $a\in U$ and let
$z_0,z$ lie in the same monad $a+\mm$. Then
\begin{equation}\label{c:p4:eq-monadtaylor}
 \ev F(z)=\sum_{n\ge0}\frac{\ev{(F^{(n)})}(z_0)}{n!}\,(z-z_0)^n,
\end{equation}
the family being Hahn summable. The expansion is exact
\emph{throughout the whole monad}, not merely on some smaller ball
around $z_0$.
\end{theorem}

\begin{proof}
Write $z_0=a+\varepsilon_0$ and $z=z_0+h$ with
$\varepsilon_0,h\in\mm$. Expand each coefficient $f_\gamma$ in its
ordinary Taylor series at $a$ and substitute the two infinitesimals.
The total support lies in $S+\Neu{\supp\varepsilon_0\cup\supp h}$, which
by \ref{a:sec:neumann} is well ordered with finite multiplicity at each
exponent. The binomial regrouping
$(\varepsilon_0+h)^n=\sum_j\binom nj\varepsilon_0^{\,n-j}h^{j}$ is
therefore legitimate, and collecting powers of $h$ produces exactly the
coefficients $\ev{(F^{(n)})}(z_0)/n!$. No hypothesis on $\abs h$ beyond
$h\in\mm$ was used, which is the assertion of exactness on the whole
monad.
\end{proof}

The next theorem is the bridge without which nothing proved with one
derivative transfers to the other. It is used silently throughout
Sections~\ref{c:p5}--\ref{e:p8}.

\begin{theorem}[The section derivative is the fine derivative]
\label{c:p4:derivbridge}
\emph{(Class (b).)} Let $F\in\Hc(U)$ be nonzero with
$\lambda=\vH(F)$, and let $b\in\halo U$. For every finite $N\ge0$,
\begin{equation}\label{c:p4:eq-remainder}
 \ev F(b+h)-\sum_{n=0}^{N}\frac{\ev{(F^{(n)})}(b)}{n!}h^n
 =h^{N+1}Q_N(b,h),
 \qquad \abs{Q_N(b,h)}<\tm^{\lambda-1}
 \quad(h\in\mm).
\end{equation}
Consequently $\ev{(F')}$ is the derivative of $\ev F$ in the
$\varepsilon$--$\delta$ sense of the fine topology: for every
$\eta\in\No^{>0}$ there is $\delta\in\No^{>0}$ with
\[
 \left|\frac{\ev F(b+h)-\ev F(b)}{h}-\ev{(F')}(b)\right|<\eta
 \qquad(0<\abs h<\delta).
\]
The quantifiers range over all sufficiently small field elements, not
over a chosen sequence.
\end{theorem}

\begin{proof}
Apply Theorem~\ref{c:p4:monadtaylor} at $z_0=b$ and factor $h^{N+1}$
out of the tail. Every exponent remaining in $Q_N$ is at least
$\lambda$, so $v(Q_N)\ge\lambda$, and the valuation bound of
Section~\ref{a:sec:field} (\ref{a:lem:valbound}) gives
$\abs{Q_N}<\tm^{\lambda-1}$; the deliberately generous factor
$\omega=\tm^{-1}$ makes the bound independent of the ordinary real size
of any leading coefficient. For the derivative take $N=1$ and
\[
 0<\delta<\min\{\omega^{-1},\ \eta\,\tm^{1-\lambda}\},
\]
so that $\abs{hQ_1(b,h)}<\delta\tm^{\lambda-1}<\eta$ whenever
$0<\abs h<\delta$.
\end{proof}

\begin{remark}[Which derivative this is]\label{c:p4:rk-deriv}
$D_z$ differentiates a function of the variable $z$ and holds every
surreal coefficient fixed; in particular $D_z(\omega)=0$. It is not a
field derivation on $\K$ in the sense of surreal differential algebra,
whose normalization has $d(\omega)=1$, and the two are not merely
differently normalized but genuinely incompatible --- see
Section~\ref{b:part3}, \ref{b:derivation-obstruction}. A reader importing
transseries results under the same symbol would silently change the
meaning of every derivative in this article.
\end{remark}

\subsection{Units and the infinitesimal functional calculus}
\label{c:p4:units}

\begin{proposition}[Units]\label{c:p4:prop-units}
\emph{(Class (b).)} Let $F\in\Hc(U)$ be nonzero and put $\lambda=\vH(F)$.
Then $F$ is a unit of $\Hc(U)$ if and only if
$f_\lambda$ is nowhere zero on $U$, in which case $E=\tm^{-\lambda}F/f_\lambda-1$ belongs to
$\Hc_{>0}(U)$ and
\begin{equation}\label{c:p4:eq-unit}
 F^{-1}=\tm^{-\lambda}f_\lambda^{-1}\sum_{n\ge0}(-E)^n
\end{equation}
and $\ev F$ is nowhere zero on $\halo U$.
\end{proposition}

\begin{proof}
If $f_\lambda$ is zero free, $f_\lambda^{-1}\in\Hol(U)$ and the
geometric series is Hahn summable, since $\supp(E^n)\subseteq\Neu{\supp
E}$ with only finitely many words of each exponent
(\ref{a:sec:neumann}); multiplying by $1+E$ verifies
\eqref{c:p4:eq-unit}. Nonvanishing of $\ev F$ follows from
\eqref{c:p4:eq-reduction}. Conversely an inverse forces the product of
the two leading coefficient functions to be $1$, so $f_\lambda$ cannot
vanish anywhere.
\end{proof}

\begin{proposition}[Infinitesimal functional calculus]
\label{c:p4:prop-calculus}
\emph{(Class (b).)} For $H\in\Hc_{>0}(U)$ and any $\varphi\in\CC[[X]]$,
the substitution $\varphi(H)=\sum_n\varphi_nH^n$ lies in $\Hc(U)$, with
support inside $\Neu{\supp H}$. In particular
\begin{equation}\label{c:p4:eq-smallcalc}
 (1+H)^{-1}=\sum_{n\ge0}(-H)^n,\qquad
 \log(1+H)=\sum_{n\ge1}\frac{(-1)^{n+1}}{n}H^n,\qquad
 e^{H}=\sum_{n\ge0}\frac{H^n}{n!}
\end{equation}
are coherent sections, and $\log$ and $\exp$ are mutually inverse
isomorphisms between $(\Hc_{>0}(U),+)$ and
$(1+\Hc_{>0}(U),\cdot)$.
\end{proposition}

\begin{proof}
Positivity of $\supp H$ and the word-length clause of
\ref{a:sec:neumann} give summability; the coefficient of each exponent is
a finite sum of products of the $f_\gamma$, hence holomorphic on $U$.
No growth condition on $(\varphi_n)$ is needed, and none is imposed.
The inverse-isomorphism claim is the formal identity
$\exp\log(1+X)=1+X$ evaluated inside a summable family.
\end{proof}

\subsection{Support-controlled composition}
\label{c:p4:comp}

\begin{theorem}[Admissible composition]\label{c:p4:composition}
\emph{(Class (b).)} Let $V,U\subseteq\CC$ be connected ordinary domains,
$F\in\Hc(V)$, and let
\[
 G=g_0+E\in\Hc(U),\qquad g_0\in\Hol(U)\ \text{with}\ g_0(U)\subseteq V,
 \qquad E\in\Hc_{>0}(U).
\]
Then
\begin{equation}\label{c:p4:eq-composition}
 F\circ G:=\sum_{\gamma\in\supp F}\ \sum_{n\ge0}
   \frac{f_\gamma^{(n)}\circ g_0}{n!}\,E^{\,n}\,\tm^\gamma
\end{equation}
belongs to $\Hc(U)$, satisfies
$\ev{(F\circ G)}=\ev F\circ\ev G$ on $\halo U$, and obeys the chain rule
$(F\circ G)'=\bigl((F'\circ G)\bigr)\cdot G'$.
\end{theorem}

\begin{proof}
The support of \eqref{c:p4:eq-composition} lies in
$\supp F+\Neu{\supp E}$, well ordered with finite multiplicity at each
exponent by \ref{a:sec:neumann}; each coefficient is a finite sum of
products of ordinary holomorphic functions on $U$, hence lies in
$\Hol(U)$. At $z\in\halo U$, $\ev G(z)$ has standard part
$g_0(\st z)\in V$, so $\ev G(z)\in\halo V$ and the right-hand side is
defined; expanding both sides at $\st z$ and using the ordinary formal
Taylor composition identity inside one summable family gives equality,
and injectivity (Theorem~\ref{c:p4:eval}) transports the chain rule from
the evaluated identity back to the sections.
\end{proof}

The shape of the hypothesis is the theorem's content. The inner section
must be an \emph{ordinary holomorphic map plus a positive-support
perturbation}. Composition with an inner section having a negative
exponent, or having an exponent-zero part that is not an ordinary
holomorphic map into $V$, is not defined here and is not claimed: the
support monoid $\Neu{\supp E}$, which supplies the finiteness clause,
exists only because $\supp E>0$.

\subsection{Affine chart changes}
\label{c:p4:chart}

Coherence is a property relative to a chart. The following criterion
says exactly which affine changes of chart preserve it.

\begin{proposition}[Chart-change criterion]\label{c:p4:chartchange}
\emph{(Class (b).)} Let $F\in\Hc(U)$ and let $z=\alpha+\beta w$ with
$\alpha,\beta\in\Ov$ \emph{finite}. Write
$\alpha=\alpha_0+\delta_\alpha$ and $\beta=\beta_0+\delta_\beta$ with
$\alpha_0,\beta_0\in\CC$ and $\delta_\alpha,\delta_\beta\in\mm$. If $V$
is a connected ordinary domain with $\alpha_0+\beta_0V\subseteq U$, then
\begin{equation}\label{c:p4:eq-chartchange}
 F(\alpha+\beta w)
 =\sum_{s\in\supp F}\ \sum_{n\ge0}
   \frac{f_s^{(n)}(\alpha_0+\beta_0w)}{n!}
   \bigl(\delta_\alpha+\delta_\beta w\bigr)^{n}\tm^{s}
\end{equation}
defines a section of $\Hc(V)$ whose evaluation is
$w\mapsto\ev F(\alpha+\beta w)$. The degenerate case $\beta_0=0$ is
allowed, provided $\alpha_0\in U$.
\end{proposition}

\begin{proof}
All support growth comes from the fixed positive supports of
$\delta_\alpha$ and $\delta_\beta$, so the total support lies in
$\supp F+\Neu{\supp\delta_\alpha\cup\supp\delta_\beta}$ and
\ref{a:sec:neumann} applies; the coefficients are finite sums of
$f_s^{(n)}(\alpha_0+\beta_0w)$ times polynomials in $w$, hence
holomorphic on $V$. When $\beta_0=0$ the inner argument is the constant
$\alpha_0\in U$ and the coefficients are polynomials in $w$; the series
is still summable for the same reason.
\end{proof}

\begin{remark}[Where the criterion fails, and why it matters]
\label{c:p4:rk-infinitecoeff}
An \emph{infinite} normalized coefficient $\beta$ does not satisfy the
criterion: the decomposition $\beta=\beta_0+\delta_\beta$ with
$\beta_0\in\CC$ and $\delta_\beta\in\mm$ does not exist, and the powers
$\beta^n$ have decreasing valuations. Thus substitution into a general
infinite Taylor series is not justified by this criterion; finite
polynomial substitution remains valid. This is the exact mechanism by
which a function can be coherent in every finite translated chart and
still fail to be coherent at all scales, and it is why the canonical
exponential $\Exp$ is excluded from the all-scale class of
Section~\ref{f:sec-allscale} (\ref{f:thm-rigidity}) although it is locally coherent.
The rigidity theorem states the exclusion; the failure of
\eqref{c:p4:eq-chartchange} for infinite $\beta$ explains it.
\end{remark}

\subsection{The identity theorem}
\label{c:p4:identitysec}

\begin{theorem}[Identity theorem for coherent sections]
\label{c:p4:identity}
\emph{(Class (b).)} Let $U\subseteq\CC$ be connected and
$F\in\Hc(U)$. Each of the following four conditions implies $F=0$.
\begin{enumerate}[label=(\roman*)]
\item $\ev F$ vanishes at points of $\halo U$ whose distinct
\emph{standard parts} have an accumulation point in $U$ in the ordinary
complex topology.
\item All fine derivatives $\ev{(F^{(n)})}(z_0)$, $n\ge0$, vanish at one
single point $z_0\in\halo U$.
\item $\ev F$ vanishes on a nonempty fine-open subset of $\halo U$.
\item $\ev F$ vanishes as a germ at a single point of $\halo U$.
\end{enumerate}
In particular two coherent sections agreeing under any one of these
conditions are equal, and their evaluations agree on all of $\halo U$.
\end{theorem}

\begin{proof}
(i) By Corollary~\ref{c:p4:shadow} the leading coefficient function
$f_\lambda$, $\lambda=\vH(F)$, vanishes at each of the indicated
standard parts. If these accumulate in $U$, the ordinary identity
theorem gives $f_\lambda\equiv0$, contradicting $\lambda=\vH(F)$ unless
$F=0$.

(ii) Put $a=\st z_0$. Theorem~\ref{c:p4:monadtaylor} applied to every
$F^{(j)}$ with increment $-(z_0-a)$ gives
$\ev{(F^{(j)})}(a)=\sum_{n\ge0}\ev{(F^{(j+n)})}(z_0)(-(z_0-a))^n/n!=0$
for all $j$. At the ordinary point $a$ these values are already in
normal form, so uniqueness of normal forms yields
$f_\gamma^{(j)}(a)=0$ for all $\gamma$ and all $j$. Each coefficient has
identically zero Taylor series at $a$ and hence vanishes on the
connected $U$.

(iii) A fine-open set contains a point together with a ball of some
positive surreal radius; repeated differentiation of the zero function
there gives (ii) at that point.

(iv) The zero-germ hypothesis at $z_0$ supplies a fine neighborhood
on which $\ev F$ vanishes. Thus (iv) is the same hypothesis
as (iii) with a specified point, and (iii) already gives $F=0$.
No preparation theorem is needed for this implication. Apply the result to $F-G$ for the last assertion.
\end{proof}

\begin{remark}[Two things the identity theorem does not say]
\label{c:p4:rk-identity}
The accumulation in (i) is accumulation of \emph{standard parts in
$\CC$}. Nothing here asserts, and nothing in this article may assert,
the existence of a nontrivial set-sized convergent sequence in the fine
topology: by the discreteness theorem of Section~\ref{a:sec:fine}
(\ref{a:thm:discrete}) every set is closed and discrete and every
set-indexed convergent net is eventually equal to its limit.

Second, the theorem is false in the germ class (a). The function $L$ of
the locally constant germ-analytic witness of Section~\ref{b:part2} is germ-analytic on all of $\K$ and vanishes on a
nonempty fine-open class without being zero. Conditions (ii)--(iv) above
do not force global vanishing in the germ class. In the coherent class,
monad-wide Taylor expansion and the ordinary identity theorem for each
coefficient function supply that propagation.
\end{remark}

\begin{theorem}[Gluing over the ordinary base]\label{c:p4:gluing}
\emph{(Class (b).)} Let a connected ordinary domain $U$ be covered by
connected ordinary open sets $U_j$, and let $F_j\in\Hc(U_j)$ have
$\ev{F_j}=\ev{F_k}$ on $\halo{(U_j\cap U_k)}$ whenever that is nonempty.
Then there is a unique $F\in\Hc(U)$, with a single well-ordered support,
restricting to each $F_j$. Restriction to a nonempty connected ordinary
open subset is injective. On a disconnected base one defines sections
component by component, allowing \emph{different supports on different
components}; this yields a class-valued sheaf over the ordinary complex
base.
\end{theorem}

\begin{proof}
On a nonempty overlap, evaluation at ordinary points and uniqueness of
normal forms force $f_{j,\gamma}=f_{k,\gamma}$ for every $\gamma$. A
holomorphic function that is nonzero on the connected $U_j$ cannot
vanish identically on a nonempty open overlap, so the \emph{exact}
supports of $F_j$ and $F_k$ coincide there. The intersection graph of a
connected open cover is connected, so one common well-ordered support
$S$ propagates through it. At each $\gamma\in S$ the coefficient
functions glue to an element of $\Hol(U)$, and coefficientwise
uniqueness gives uniqueness of $F$. Injectivity of restriction is the
ordinary identity theorem applied to each coefficient.
\end{proof}

This is a sheaf over the \emph{ordinary} base. It is not an assertion
that arbitrary fine-local analytic data glue to a coherent section; that
assertion is false, and the locally constant germ-analytic witness of Section~\ref{b:part2} is the counterexample.
Requiring one well-ordered support across arbitrarily many components of
a disconnected base is a strictly stronger and unmotivated condition,
and it is not imposed.

\subsection{Canonical classical lifts}
\label{c:p4:lifts}

\begin{definition}[Canonical lift]\label{c:p4:def-lift}
For $f\in\Hol(U)$, its \emph{canonical lift} is the evaluation of the
coherent section supported at exponent $0$ alone:
\begin{equation}\label{c:p4:eq-lift}
 \ev f(a+\varepsilon)=\sum_{n\ge0}\frac{f^{(n)}(a)}{n!}\varepsilon^{n},
 \qquad a\in U,\ \varepsilon\in\mm .
\end{equation}
The same superscript is used for lifts and for evaluations of coherent
sections because a lift \emph{is} a section, namely one supported at
exponent zero. Class (c) denotes the lifts, class (b) the coherent
sections.
\end{definition}

\begin{theorem}[Functoriality of the lift]\label{c:p4:lift}
\emph{(Class (c).)} $f\mapsto\ev f$ is an injective
$\CC$-algebra homomorphism $\Hol(U)\to\{\halo U\to\K\}$, and
\[
 \st(\ev f(z))=f(\st z),\qquad (\ev f)'=\ev{(f')},
\]
and
\[
 \ev{(g\circ f)}=\ev g\circ\ev f
 \quad\text{if }f(U)\subseteq V,\ g\in\Hol(V),
\]
together with the reflection identity
$\ev{(f^{*})}(z)=\overline{\ev f(\bar z)}$ for
$f^{*}(Z)=\overline{f(\bar Z)}\in\Hol(U^{*})$. Lifting commutes with
restriction and with ordinary analytic continuation, because one
continues the ordinary coefficient function first and then applies
\eqref{c:p4:eq-lift}.
\end{theorem}

\begin{proof}
This is Theorem~\ref{c:p4:eval} for the support $\{0\}$, together with
one extra point: at $z=c+\varepsilon$ the difference $\ev f(z)-f(c)$ is
infinitesimal, so substituting its Hahn-summed Taylor series into the
Taylor series of $g$ is licensed by the positive-support clause of
\ref{a:sec:neumann}; collection by total degree is exactly the ordinary
formal chain rule. Taking the exponent-zero coefficient of
\eqref{c:p4:eq-lift} gives $\st(\ev f(z))=f(\st z)$. Conjugation is
coefficientwise.
\end{proof}

\begin{proposition}[A lift creates no displaced zeros]
\label{c:p4:liftzeros}
\emph{(Class (c).)} Let $f\in\Hol(U)$ be nonzero and $a\in U$ a zero of
$f$ of order $m$. Then the zeros of $\ev f$ in the monad $a+\mm$ are
exactly the single point $a$, with multiplicity $m$; and $\ev f$ has no
zeros over a non-zero of $f$.
\end{proposition}

\begin{proof}
Write $f(Z)=(Z-a)^m u(Z)$ with $u$ ordinary holomorphic and
$u(a)\ne0$ on a small disk about $a$. Then
$\ev f(a+\varepsilon)=\varepsilon^{m}\,\ev u(a+\varepsilon)$, and
$\st(\ev u(a+\varepsilon))=u(a)\ne0$, so the second factor is a unit;
hence $\varepsilon=0$ is the only solution in the monad. The second
assertion is Corollary~\ref{c:p4:shadow} for the support $\{0\}$.
\end{proof}

\begin{remark}[The point of the previous proposition]
\label{c:p4:rk-liftzeros}
A canonical lift cannot split an ordinary multiple zero into
infinitesimally separated roots. The splitting phenomenon of
Section~\ref{d:sec:preparation} therefore requires genuinely surcomplex coefficients at
positive exponents: class (b) strictly contains class (c), and the
containment is visible already in the zero geometry. The three classes
satisfy
\[
 \begin{aligned}
 \text{(c) lifts}&\subsetneq\text{(b) }\Hc(U)\\
 &\subsetneq\{\text{(a) germ-analytic functions on }\halo U\}\\
 &\subseteq\{\text{bare differentiable functions on }\halo U\}.
 \end{aligned}
\]
The first two inclusions are proper; the factorial series gives the
second obstruction (Proposition~\ref{c:p4:obstruction}), while a
locally constant nonconstant function on the halo gives a direct global
witness. Strictness of the last inclusion is not proved here
(Remark~\ref{b:no-goursat}). In particular $L$ cannot witness it, since
$L$ is itself germ-analytic. Beyond the calculus rules, no regularity theorem for the
bare differentiable class is proved here.
\end{remark}

\subsection{The scale-embedding bridge, and its limit}
\label{c:p4:bridge-sec}

The germ layer and the coherent layer meet in the following theorem.
Coherence is \emph{not} restricted to classically convergent germs.

\begin{theorem}[Scale embedding of arbitrary formal series]
\label{c:p4:bridge}
\emph{(From class (a) to class (b).)} For every
$P(X)=\sum_{n\ge0}a_nX^{n}\in\K[[X]]$ --- including series whose
coefficients are infinite and whose ordinary radius of convergence is
zero --- there are a positive monomial $\rho=\tm^{\lambda}$ and a
section $\mathcal P\in\Hc(\CC)$ with
\begin{equation}\label{c:p4:eq-bridge}
 \ev{\mathcal P}(\zeta)=P(\rho\zeta)\qquad(\zeta\in\Ov).
\end{equation}
Every coefficient function of $\mathcal P$ may be taken to be a
polynomial; in the separated construction each nonzero coefficient is a
complex constant times a single power $Z^{n}$.
\end{theorem}

\begin{proof}
Choose $\lambda$ by the separated-block lemma of Section~\ref{b:part3}
(\ref{b:blocks}): a single $\Delta$ bounds the --- possibly not well
ordered --- union $\bigcup_n\supp(a_n)$, by the surreal cut property,
and $\lambda>2\Delta$ places the support of $a_n\rho^{n}$ inside the
band $(n\lambda-\Delta,\,n\lambda+\Delta)$. These bands are pairwise
disjoint and increasingly ordered, so writing
$a_n\rho^{n}=\sum_{\gamma\in S_n}c_{n,\gamma}\tm^{\gamma}$ the union
$S=\bigcup_nS_n$ is well ordered and every $\gamma\in S$ lies in exactly
one block, determining $n=n(\gamma)$ uniquely. Set the coefficient
function at $\gamma$ to be $c_{n,\gamma}Z^{n}$, an entire polynomial.
Evaluating at finite $\zeta$ reproduces the Hahn summable family
$a_n\rho^{n}\zeta^{n}$, including its finite binomial expansion about
$\st\zeta$, which is \eqref{c:p4:eq-bridge}.
\end{proof}

A germ centred at an arbitrary $a\in\K$ is obtained by substituting
$z=a+\rho\zeta$. The theorem must be read next to the following
obstruction, which concerns a \emph{different chart}.

\begin{proposition}[An obstruction in the original chart]
\label{c:p4:obstruction}
\emph{(Classes (a) and (b), contrasted.)} The series
$\sum_{n\ge0}n!\,\varepsilon^{n}$ is a legitimate germ-analytic function
on $\mm$, by the substitution corollary of Section~\ref{a:sec:summability}
(\ref{a:sec:substitution}). It is \emph{not} the restriction to a monad
of any $F\in\Hc(U)$, for any ordinary neighborhood $U$ of $0$, having
those same derivatives at $0$.
\end{proposition}

\begin{proof}
Such an $F$ would have $\ev{(F^{(n)})}(0)=n!\cdot n!$ for all $n$.
Evaluating at the ordinary point $0$ and using uniqueness of normal
forms, the exponent-zero coefficient function $f_0\in\Hol(U)$ would have
$f_0^{(n)}(0)/n!=n!$, i.e.\ Taylor coefficients $n!$ at $0$. No ordinary
holomorphic function on a neighborhood of $0$ has Taylor coefficients of
that size, since a positive complex radius of convergence bounds them
geometrically.
\end{proof}

\begin{remark}[The two statements are compatible]
\label{c:p4:rk-reconcile}
Theorem~\ref{c:p4:bridge} and Proposition~\ref{c:p4:obstruction} look
like a contradiction and are not. The obstruction concerns the
\emph{original} chart: a section on an ordinary neighborhood of zero with
the prescribed derivatives at zero. The bridge concerns a
\emph{rescaled} chart: after $z=\rho\zeta$ the degree-$n$ term lands in
its own Hahn exponent band, each Hahn coefficient becomes a finite sum
and hence a polynomial, and what is being represented is a different
function of a different variable. The rescaling changes the chart; it
does not manufacture a holomorphic function with factorial Taylor
coefficients.

The positive result has an exact limit, which must travel with it:
every germ has \emph{some} sufficiently small coherent chart, but there
is no implication that it has a chart of every larger radius. That
distinction is the content of the all-scale rigidity theorem of
Section~\ref{f:sec-allscale} (\ref{f:thm-rigidity}), and the mechanism of failure is
Remark~\ref{c:p4:rk-infinitecoeff}.
\end{remark}

% =====================================================================

\section{Contours, and the order-valued estimate}
\label{c:p5}

Section~\ref{a:sec:fine} showed that an ordinary circle is not a path in $\K$:
every fine-continuous map from a real interval into $\K$ is constant
(\ref{a:sec:nopath}). A contour integral over $\K$ therefore cannot be a
limit of Riemann sums along a fine-continuous parameterization. It must
be defined some other way, and Section~\ref{b:part2} showed that there is
essentially no choice: the contour no-go theorem (\ref{b:nogo})
proves that no functional on the ordinary unit circle can simultaneously
give the classical value $2\pi i$ for $1/z$ and satisfy the fundamental
theorem of calculus for every globally fine-analytic primitive on
$\K^{\times}$. The definition below, and the restriction of the
fundamental theorem to \emph{coherent} primitives in
\S\ref{c:p5:primsec}, are forced by that theorem and are not
conservatism.

\subsection{The coefficientwise contour functional}
\label{c:p5:defsec}

\begin{definition}[Contour and coefficientwise integral]
\label{c:p5:def-integral}
A \emph{contour} in $U$ is an ordinary piecewise-$C^{1}$ path
$\gamma:[a,b]\to U$ in the \emph{ordinary} complex domain. For
$F=\sum_{s\in S}f_s\tm^{s}\in\Hc(U)$ set
\begin{equation}\label{c:p5:eq-integral}
 \intH_{\gamma}F(\zeta)\,\dd\zeta
 :=\sum_{s\in S}\left(\int_{\gamma}f_s(\zeta)\,\dd\zeta\right)\tm^{s}.
\end{equation}
The integrals on the right are ordinary complex contour integrals; the
outer sum is a Hahn sum whose support can only shrink, so the result is
a well-defined surcomplex number. The superscript $\mathrm H$ is
permanent.
\end{definition}

\begin{proposition}[Elementary calculus of the functional]
\label{c:p5:calculus}
\emph{(Class (b).)} The map $F\mapsto\intH_{\gamma}F\,\dd\zeta$ is
$\K$-linear --- not merely $\CC$-linear --- additive under concatenation
of contours, sign-reversing under reversal of orientation, and invariant
under ordinary piecewise-$C^{1}$ reparameterization. The same definition
applies verbatim when the coefficient functions are merely continuous
along $\gamma$, and when they are meromorphic on a common ordinary
domain with no pole on $\gamma$. It commutes with any Hahn summable
family of sections whose coefficientwise integrals are defined.
\end{proposition}

\begin{proof}
Every clause except $\K$-linearity is the corresponding ordinary
property applied at each exponent. $\K$-linearity holds because
multiplication by a fixed Hahn series has finite convolution at each
exponent (\ref{a:sec:neumann}), so
$\intH_\gamma(cF)=c\intH_\gamma F$ for $c\in\K$ by comparing
coefficients. The last clause holds because each exponent of the sum
involves only finitely many members of the family.
\end{proof}

\begin{definition}[Affine transport]\label{c:p5:affine}
Fix $b\in\K$ and $r\in\K^{\times}$, and let $F\in\Hc(U)$ define a
function on $b+r\halo U$ by $f(b+rw)=\ev F(w)$. For an ordinary contour
$\gamma$ in $U$ set
\begin{equation}\label{c:p5:eq-affine}
 \intH_{\,b+r\gamma}f(z)\,\dd z:=r\intH_{\gamma}F(w)\,\dd w.
\end{equation}
The centre $b$ may be infinite and $\abs r$ may be infinitesimal or
infinite.
\end{definition}

Differentiation in the transported chart introduces $r^{-n}$ and the
differential introduces $r$; the two cancel on a simple-pole residue, so
Cauchy's formula, root counts and the residue theorem transfer to the
chart with the usual affine change-of-variable factors. Two admissible
descriptions giving the same pointwise pulled-back integrand have equal
integrals, by uniqueness of normal forms applied to their coefficient
functions.

\begin{remark}[What the symbol is not]\label{c:p5:rk-notriemann}
$\intH$ is a coefficientwise functional. It is not a fine-topological
Riemann integral, and $b+r\gamma$ is a chart-defined algebraic contour,
not a claim that its parameterization is fine-continuous into $\K$. The
superscript is kept everywhere for exactly this reason. Nothing below
proposes $\intH$ as \emph{the} unique notion of surreal integration;
what is proved is that a functional with the two properties in the no-go
theorem does not exist, so some restriction of this kind is unavoidable.
\end{remark}

\subsection{Positivity and the surcomplex ML inequality}
\label{c:p5:MLsec}

A coefficientwise definition does not by itself yield an inequality
about the surreal \emph{modulus}. The next lemma supplies one, and it is
the load-bearing step of this part.

\begin{lemma}[Positivity of the coefficientwise integral]
\label{c:p5:positivity}
Let $A(s)=\sum_{\lambda\in S}a_\lambda(s)\tm^{\lambda}$ be a family with
one common well-ordered support $S$ and continuous \emph{real}
coefficient functions $a_\lambda$ on a compact interval $[\alpha,\beta]$.
If $A(s)\ge0$ in $\No$ for every ordinary $s\in[\alpha,\beta]$, then
\[
 \intH_{\alpha}^{\beta}A(s)\,\dd s\ \ge\ 0 .
\]
\end{lemma}

\begin{proof}
If $A\equiv0$ or the interval is degenerate there is nothing to prove.
Otherwise discard identically zero coefficient functions and let
$\lambda_0$ be the least remaining exponent. For each fixed $s$, the sign of the
surreal number $A(s)$ is the sign of its leading nonzero coefficient;
since $A(s)\ge0$ for all $s$, we get $a_{\lambda_0}(s)\ge0$ for every
$s$ --- were it negative at some $s$, the value $A(s)$ would be negative
there. Now $a_{\lambda_0}$ is continuous, nonnegative and not identically
zero, so $\int_\alpha^\beta a_{\lambda_0}>0$. Hence the coefficient of
$\tm^{\lambda_0}$ in $\intH A$ is positive and all smaller exponents are
absent, so the integral is positive. Piecewise versions follow by
splitting into finitely many intervals and adding.
\end{proof}

\begin{theorem}[Surcomplex ML inequality]\label{c:p5:ML}
\emph{(Class (b), and more generally any Hahn family with continuous
coefficient functions along the path.)} Let $\gamma$ be a
piecewise-$C^{1}$ ordinary path of ordinary length $L$, let $F$ have
continuous complex coefficient functions with common well-ordered
support along $\gamma$, and suppose
\[
 \abs{\ev F(\gamma(s))}\le M\qquad\text{for every ordinary }s,
\]
for one fixed $M\in\No^{\ge0}$. Then
\begin{equation}\label{c:p5:eq-ML}
 \left|\ \intH_{\gamma}F(z)\,\dd z\ \right|\ \le\ LM .
\end{equation}
\end{theorem}

\begin{proof}
Write $I=\intH_\gamma F\,\dd z$. If $I=0$ the bound is trivial.
Otherwise put $\alpha=\overline I/\abs I$, so that $\abs\alpha=1$ and
$\operatorname{Re}(\alpha I)=\abs I$. Pointwise along the path,
\[
 M\abs{\gamma'(s)}-\operatorname{Re}
   \bigl(\alpha\,\ev F(\gamma(s))\,\gamma'(s)\bigr)\ \ge\ 0,
\]
because $\abs{\operatorname{Re}(\alpha F\gamma')}\le
\abs{F}\abs{\gamma'}\le M\abs{\gamma'}$ in the ordered field $\No$.
The left-hand side is a Hahn family of continuous real coefficient
functions on the parameter interval, with one common support (the
union of $\supp M$ and $\supp\alpha+\supp F$). Apply
Lemma~\ref{c:p5:positivity} and then $\K$-linearity of $\intH$
(Proposition~\ref{c:p5:calculus}) to obtain $LM-\abs I\ge0$.
\end{proof}

The separating scalar $\alpha$ is what converts a coefficientwise
statement into an ordered-field positivity statement; without it the
theorem is not available.

\begin{remark}[Two Cauchy-type statements that are not in conflict]
\label{c:p5:rk-noclash}
Theorem~\ref{c:p5:ML} and the coefficientwise estimates of
\S\ref{c:p5:majorantsec} look like a clash and are not. They bound
different quantities.
\begin{itemize}
\item It is \emph{false} that a single coarse surreal bound on $\ev F$
controls every coefficient function uniformly. Witness:
$F(w)=w(1+i\tm)$ on the ordinary unit circle has
$\abs{\ev F}=\sqrt{1+\tm^{2}}<1+\tm^{2}$ pointwise, so $M=1+\tm^{2}$ is
a legitimate pointwise bound, whereas the coefficientwise majorant
$B_1(F)=1+\tm$ is strictly larger than $M$. No estimate below asserts
the contrary.
\item It is \emph{true} that a single surreal bound controls the surreal
modulus of the integral, and hence of the derivatives; that is
\eqref{c:p5:eq-ML} and Corollary~\ref{c:p5:cauchyest}. The conclusion is
weaker than the false one and is proved above.
\end{itemize}
The two are consistent on the standard test case: for $F=\varepsilon Z$
with $\varepsilon\in\mm\cap\No_{>0}$, the bound on $\abs z=R$ is $M=\varepsilon R$,
and Corollary~\ref{c:p5:cauchyest} gives
$\abs{\ev{(F')}(0)}\le\varepsilon$, which is exact. Reading the
coefficientwise caution as a prohibition would delete the strongest
estimate available, and that reading is rejected here.
\end{remark}

\subsection{Cauchy's theorem and formula at surcomplex points}
\label{c:p5:cauchysec}

\begin{theorem}[Cauchy's theorem and integral formula]
\label{c:p5:cauchy}
\emph{(Class (b).)} Let $U$ be a connected ordinary domain and
$F\in\Hc(U)$.
\begin{enumerate}[label=(\roman*)]
\item If the ordinary closed path $\gamma$ is null-homologous in $U$,
then $\ointH_{\gamma}F(\zeta)\,\dd\zeta=0$.
\item Let $\Omega$ be a bounded ordinary Jordan domain with
piecewise-$C^{1}$ positively oriented boundary
$\Gamma=\partial\Omega$ and $\overline\Omega\subset U$. Then for every
genuinely surcomplex point $z=c+\varepsilon\in\halo\Omega$ and every
ordinary finite $k\ge0$,
\begin{equation}\label{c:p5:eq-cauchy}
 \frac{\ev{(F^{(k)})}(z)}{k!}
 =\frac1{2\pi i}\ointH_{\Gamma}
    \frac{F(\zeta)}{(\zeta-z)^{k+1}}\,\dd\zeta,
\end{equation}
the kernel being understood through its uniform Hahn expansion on an
ordinary neighborhood of $\Gamma$.
\end{enumerate}
\end{theorem}

\begin{proof}
(i) is the ordinary Cauchy theorem applied at each exponent; this is
exactly what is inherited coefficientwise from complex analysis, and
nothing more.

(ii) Since $c\in\Omega$ is an ordinary interior point, $\zeta-c$ is
bounded away from zero on $\Gamma$, and on an ordinary neighborhood of
$\Gamma$
\[
 \frac1{(\zeta-c-\varepsilon)^{k+1}}
 =\sum_{n\ge0}\binom{n+k}{k}
   \frac{\varepsilon^{n}}{(\zeta-c)^{\,n+k+1}} .
\]
The support of $\varepsilon$ is positive, so this is a Hahn summable
family with coefficient functions holomorphic near $\Gamma$; its product
with $F$ is jointly summable inside
$\supp F+\Neu{\supp\varepsilon}$. Integrating coefficientwise and
applying the ordinary Cauchy derivative formula to each $f_\gamma$,
\[
 \frac{k!}{2\pi i}\ointH_{\Gamma}\frac{F(\zeta)}{(\zeta-z)^{k+1}}
   \dd\zeta
 =\sum_{\gamma,n}\tm^{\gamma}\varepsilon^{n}\,
    k!\binom{n+k}{k}\frac{f_\gamma^{(n+k)}(c)}{(n+k)!}
 =\sum_{\gamma,n}\tm^{\gamma}
    \frac{f_\gamma^{(n+k)}(c)}{n!}\varepsilon^{n},
\]
which is $\ev{(F^{(k)})}(z)$ by Theorem~\ref{c:p4:monadtaylor}. Only
finitely many terms are combined at any fixed Hahn exponent, so every
interchange is justified by summability and not by an unproved
convergence theorem; no real limit is exchanged with a surreal limit.
\end{proof}

\begin{corollary}[Cauchy estimates from a pointwise surreal bound]
\label{c:p5:cauchyest}
\emph{(Class (b).)} Let $c\in\CC$, let $R>0$ be ordinary with
$\overline{B(c,R)}\subset U$, and suppose $\abs{\ev F(\zeta)}\le M$ for
one fixed $M\in\No^{\ge0}$ at every point $\zeta$ of the \emph{ordinary}
circle $\abs{\zeta-c}=R$. Then for every ordinary finite $n\ge0$
\begin{equation}\label{c:p5:eq-cauchyest}
 \bigl|\ev{(F^{(n)})}(c)\bigr|\ \le\ \frac{n!\,M}{R^{n}},
\end{equation}
and for $z=c+\varepsilon$ with $\varepsilon\in\mm$,
\begin{equation}\label{c:p5:eq-cauchyest2}
 \bigl|\ev{(F^{(n)})}(z)\bigr|
 \ \le\ \frac{n!\,R\,M}{(R-\abs\varepsilon)^{\,n+1}} .
\end{equation}
\end{corollary}

\begin{proof}
Apply Theorem~\ref{c:p5:ML} to \eqref{c:p5:eq-cauchy} with $L=2\pi R$
and, for \eqref{c:p5:eq-cauchyest2}, with the lower bound
$\abs{\zeta-z}\ge R-\abs\varepsilon$ on the kernel. The bound $M$ is
allowed to be an arbitrary nonnegative surreal, infinite or
infinitesimal.
\end{proof}

\begin{corollary}[Cauchy estimates in an affine chart]
\label{c:p5:scaledCauchy}
\emph{(Class (b), in a transported chart.)} Suppose $f(a+rw)=\ev F(w)$
with $a\in\K$, $r\in\No^{>0}$ and $F\in\Hc(B(0,2))$. If
$\abs{f(a+re^{i\theta})}\le M$ for every ordinary real $\theta$, then
\begin{equation}\label{c:p5:eq-scaledCauchy}
 \abs{f^{(n)}(a)}\ \le\ \frac{n!\,M}{r^{n}}\qquad(n\ge0).
\end{equation}
Cauchy's formula, root counts and the residue theorem transfer to this
chart with the usual affine factors.
\end{corollary}

\begin{proof}
The chain rule gives $\ev{(F^{(n)})}(0)=r^{n}f^{(n)}(a)$. Apply
\eqref{c:p5:eq-cauchyest} on the unit circle, where the bound $M$ may be
surreal. The contour assertions follow by substitution in
\eqref{c:p5:eq-affine}, the differential $\dd z=r\,\dd w$ cancelling the
affine scaling of a simple-pole residue.
\end{proof}

The qualification ``coherent chart'' is essential and is not decorative:
a germ-analytic function has \emph{some} sufficiently small coherent
chart (Theorem~\ref{c:p4:bridge}), but there is no implication that it
has a chart of every larger radius. See Section~\ref{f:sec-allscale}.

\subsection{The coefficientwise majorant}
\label{c:p5:majorantsec}

The second estimate has a different shape. It bounds each coefficient
separately and assembles the bounds into a surreal number, and it is
retained as a distinct tool.

\begin{lemma}[Cauchy--Schwarz majorant lemma]\label{c:p5:csmajorant}
Let $Z=\sum_{s}c_s\tm^{s}\in\K$ with $c_s\in\CC$ and put
$B=\sum_{s}\abs{c_s}\tm^{s}\in\No^{\ge0}$. Then $\abs Z\le B$.
\end{lemma}

\begin{proof}
\[
 B^{2}-\abs Z^{2}
 =\sum_{s,u}\bigl(\abs{c_s}\abs{c_u}
   -\operatorname{Re}(c_s\overline{c_u})\bigr)\tm^{\,s+u}\ \ge\ 0,
\]
since every coefficient of this Hahn summable family is a nonnegative
real number and the family is summable by \ref{a:sec:neumann}. Both $B$
and $\abs Z$ are nonnegative, so $\abs Z\le B$.
\end{proof}

\begin{proposition}[Coefficientwise Cauchy estimate]
\label{c:p5:majorant}
\emph{(Class (b).)} Let $\overline{B(a,r)}\subset U$ with $a\in\CC$ and
$r>0$ ordinary real. For $F=\sum_sf_s\tm^{s}\in\Hc(U)$ put
\[
 M_s=\max_{\abs{\zeta-a}=r}\abs{f_s(\zeta)},
 \qquad
 B_r(F)=\sum_sM_s\tm^{s}\in\No^{\ge0},
\]
ordinary compact maxima followed by a Hahn sum. Then for every ordinary
finite $n\ge0$
\begin{equation}\label{c:p5:eq-majorant}
 \bigl|\ev{(F^{(n)})}(a)\bigr|\ \le\ \frac{n!}{r^{n}}\,B_r(F).
\end{equation}
\end{proposition}

\begin{proof}
The ordinary Cauchy estimates give
$\abs{f_s^{(n)}(a)}\le n!M_s/r^{n}$ for every $s$. By
Lemma~\ref{c:p5:csmajorant}, the surreal modulus of
$\ev{(F^{(n)})}(a)=\sum_sf_s^{(n)}(a)\tm^{s}$ is at most
$\sum_s\abs{f_s^{(n)}(a)}\tm^{s}$, and increasing nonnegative
coefficients preserves the surreal order on a common support.
\end{proof}

\begin{remark}[The scope of the majorant]\label{c:p5:rk-majorant}
$B_r(F)$ is not in general the least upper bound of $\abs{\ev F}$ on a
surreal circle; its coefficientwise definition is part of the statement
of \eqref{c:p5:eq-majorant}. It is not claimed to be a norm. The
coefficientwise order used here is much stronger than the surreal order,
and the two estimates \eqref{c:p5:eq-cauchyest} and
\eqref{c:p5:eq-majorant} use different input data. They are not
logically incomparable as modulus bounds: the majorant lemma gives
$\abs{\ev F(\zeta)}\le B_r(F)$ at every ordinary point of the circle,
so the ML estimate with $M=B_r(F)$ implies
\eqref{c:p5:eq-majorant}. For
$F(w)=w(1+i\tm)$ on $r=1$ we have $B_1(F)=1+\tm$ while the pointwise
bound is $M=1+\tm^{2}<B_1(F)$, so the ML estimate is the sharper one
there. The coefficientwise bounds still retain information that a
single surreal inequality loses: the Liouville and growth theorem
\ref{e:thm-coeffliouville} applies ordinary estimates to each coefficient
separately. A pointwise surreal bound always exists on a fixed circle
and, by Proposition~\ref{e:prop-autobounded}, even on the finite halo
for macroscopically entire sections; such a bound alone does not recover
those coefficientwise growth hypotheses.
\end{remark}

\subsection{Morera's criterion}
\label{c:p5:morerasec}

\begin{theorem}[Coefficientwise Morera]\label{c:p5:morera}
Let $S$ be a well-ordered set and let $A=\sum_{s\in S}a_s\tm^{s}$ have
\emph{ordinary continuous} coefficient functions $a_s$ on an ordinary
domain $U$, with contour integrals defined coefficientwise by
\eqref{c:p5:eq-integral}. The following are equivalent.
\begin{enumerate}[label=(\roman*)]
\item Every $a_s$ is holomorphic on $U$, i.e.\ $A\in\Hc(U)$.
\item $\ointH_{\partial T}A(\zeta)\,\dd\zeta=0$ for every ordinary
triangle $T$ whose closed interior lies in $U$.
\end{enumerate}
If the $a_s$ are $C^{1}$, both are further equivalent to the
coefficientwise Cauchy--Riemann equations
$\partial_{\bar Z}a_s=0$ for every $s$. When they hold, $A$ has a unique
coherent evaluation on $\halo U$ by \eqref{c:p4:eq-eval}.
\end{theorem}

\begin{proof}
Uniqueness of normal forms makes the vanishing of the Hahn-valued
triangle integral equivalent to the vanishing of every ordinary
coefficient integral $\oint_{\partial T}a_s$. Ordinary Morera's theorem
applied to each continuous $a_s$ gives (ii)$\Rightarrow$(i), and
ordinary Cauchy (i)$\Rightarrow$(ii). For $C^{1}$ coefficients use the
ordinary Cauchy--Riemann criterion at each exponent. Existence and
uniqueness of the evaluation are Theorem~\ref{c:p4:eval}.
\end{proof}

The regularity hypothesis is part of the theorem. It is continuity of
the \emph{ordinary coefficient functions} on the ordinary domain, and
the theorem asserts nothing whatever about a function that is merely
continuous, or even differentiable, for the fine topology. Before the
conclusion is reached, the object $A$ is a Hahn-valued function on an
ordinary domain; no Taylor extension of a merely continuous coefficient
to the halo is presumed at that stage.

\subsection{Coherent primitives, periods and endpoint corrections}
\label{c:p5:primsec}

\begin{theorem}[Coherent primitives]\label{c:p5:primitive}
\emph{(Class (b).)} Let $U$ be a connected ordinary domain and
$F\in\Hc(U)$.
\begin{enumerate}[label=(\roman*)]
\item If $U$ is simply connected and $a_0\in U$, there is a unique
$P\in\Hc(U)$ with $P'=F$ and $\ev P(a_0)=0$, and
$\supp P\subseteq\supp F$. Every other coherent primitive differs from
$P$ by a constant of $\K$, and for an ordinary path $\gamma$ in $U$ from
$a_1$ to $a_2$,
\[
 \intH_{\gamma}F(\zeta)\,\dd\zeta=\ev P(a_2)-\ev P(a_1).
\]
\item For general connected $U$, a coherent primitive exists if and only
if all ordinary closed-path coefficientwise integrals vanish, i.e.\
$\ointH_{\sigma}F\,\dd\zeta=0$ for every ordinary closed path $\sigma$
in $U$.
\end{enumerate}
\end{theorem}

\begin{proof}
(i) On a simply connected $U$ each $f_s$ has a unique ordinary primitive
$p_s$ with $p_s(a_0)=0$, by path independence; set $P=\sum_sp_s\tm^{s}$.
Its support is contained in $\supp F$ --- a primitive can only lose
exponents, when $f_s$ has an identically zero primitive, which happens
only for $f_s\equiv0$ --- so $P\in\Hc(U)$ and $P'=F$. If $H'=0$ then
every $h_s'=0$, and connectedness of $U$ makes each $h_s$ a complex
constant, so $H$ is a constant of $\K$; this is uniqueness. The endpoint
identity holds at each exponent.

(ii) Vanishing of the Hahn period $\ointH_\sigma F$ is, by uniqueness of
normal forms, equivalent to vanishing of every ordinary period
$\oint_\sigma f_s$, which is the classical criterion for each $f_s$ to
have an ordinary primitive on $U$. Run the construction of (i) with
those primitives.
\end{proof}

For endpoints that are genuinely surcomplex, $z_j=a_j+\varepsilon_j\in
\halo U$ with $a_j\in U$ and $\varepsilon_j\in\mm$, define the
microscopic endpoint correction
\begin{equation}\label{c:p5:eq-endpointcorr}
 T_F(a,\varepsilon)
 =\sum_{n\ge0}\frac{\ev{(F^{(n)})}(a)}{(n+1)!}\,\varepsilon^{\,n+1},
\end{equation}
Hahn summable by the support calculation of Theorem~\ref{c:p4:eval}, and
set, for an ordinary path $\gamma$ from $a_1$ to $a_2$,
\begin{equation}\label{c:p5:eq-endpointintegral}
 \intH_{z_1}^{z_2}F\,\dd z
 :=-T_F(a_1,\varepsilon_1)+\intH_{\gamma}F\,\dd z
   +T_F(a_2,\varepsilon_2).
\end{equation}
For a coherent primitive $P$ this equals $\ev P(z_2)-\ev P(z_1)$, by
Theorem~\ref{c:p4:monadtaylor}; on a simply connected base it is
independent of the ordinary path chosen. This gives a concrete endpoint
integral between surcomplex endpoints without pretending that the
connecting ordinary path is fine-continuous.

\begin{remark}[Uniqueness belongs to the coherent category only]
\label{c:p5:rk-uniqueness}
The uniqueness clause in Theorem~\ref{c:p5:primitive}(i) is a statement
about class (b) and is false outside it. An arbitrary differentially
holomorphic primitive may still be modified by a locally constant
function: $P+L$ has the same derivative as $P$ everywhere, with $L$ as in
the locally constant germ-analytic witness of Section~\ref{b:part2}. The restriction cannot be lifted, and the theorem
that shows it cannot is the contour no-go theorem
(\ref{b:nogo}).
\end{remark}

\begin{remark}[The unit circle, and the global logarithm]
\label{c:p5:rk-logarithm}
For $F(z)=1/z$ on $\CC^{\times}$ the coefficientwise unit-circle period
is $\ointH_{\abs z=1}\dd z/z=2\pi i\ne0$. By
Theorem~\ref{c:p5:primitive}(ii) there is therefore \emph{no coherent
primitive of $1/z$} on $\halo{(\CC^{\times})}$ --- no coherent global
logarithm there.

This is consistent with, and is the exact complement of, the global
fine-analytic logarithm $\mathcal L$ constructed in Section~\ref{e:p9}
(\ref{e:thm-globallog}) on all of $\K^{\times}$, which satisfies
$\mathcal L'(z)=1/z$ everywhere. That function is germ-analytic, class
(a); its restriction to $\halo{(\CC^\times)}$ is not the evaluation
of an element of $\Hc(\CC^\times)$, as the computation above proves.
On an ordinary domain avoiding the argument cut, it does agree with
the canonical lift of an ordinary logarithm branch. If it were coherent, the fundamental theorem
(Theorem~\ref{c:p5:primitive}(i)) would force the unit-circle integral to
be $0$ rather than $2\pi i$. The monodromy obstruction survives in the
coherent category and disappears in the more permissive local one; that
is the whole content of the no-go theorem of Section~\ref{b:part2}, and it is
why the fundamental theorem above is stated only for coherent
primitives.
\end{remark}

%% ==================== p67.tex ====================

% =====================================================================
% Part 6 and Part 7 of the merged surcomplex document.
% Lead sources: 02, 08, 01 (with 03, 04, 05, 06, 07, 09 for the
% additions recorded in the merge contract).
% CLASS DISCIPLINE.  Every theorem of Part 6 is a theorem about the
% COHERENT class H(U); none of it is asserted for germ-analytic
% functions, and no canonical-lift theorem is used.  Part 7 works with
% three separate classes and says which one each statement belongs to.
% =====================================================================

\section{Preparation, division, and infinitesimal zero geometry}
\label{d:sec:preparation}

Everything in this part is a statement about the \emph{coherent} class
$\Hc(U)$ of Section~\ref{c:p4}: one common connected ordinary domain $U$, one
ordinary holomorphic coefficient function per exponent, well-ordered
set support.  Nothing here is claimed for the germ-analytic class of
Section~\ref{b:part3}, and nothing here is a statement about canonical lifts.  The
reason is structural rather than cautious: the whole mechanism below is
ordinary Hermite division applied to the coefficient functions
\emph{simultaneously on one domain}, and it is exactly that common
domain which ties the monads over different standard points together.
A family of germs whose radii shrink to zero admits no such division,
and the zero-cluster theorem is false for it.

Throughout, $U\subseteq\CC$ is a connected ordinary domain,
$\Omega\subset\CC$ is a bounded Jordan domain with piecewise $C^1$
boundary $\Gamma=\partial\Omega$ and $\overline\Omega\subset U$, and $V$
denotes an ordinary open neighbourhood of $\overline\Omega$ with
$\overline\Omega\subset V\subseteq U$.  We write a nonzero
$F\in\Hc(U)$ in the normalized form
\begin{equation}\label{d:eq:normalform}
 F=\tm^{\beta}\bigl(f_0+E\bigr),
 \qquad \beta=\vH(F),\quad f_0=f_{\vH(F)}\in\Hol(U)\setminus\{0\},
 \quad \supp E\subset\No_{>0},
\end{equation}
so that $f_0$ is the leading coefficient \emph{function} of $F$ --- not
the leading coefficient of any one of its values --- and $E\in
\Hc_{>0}(U)$.

\subsection{Ordinary Hermite division data}
\label{d:sub:hermite}

\begin{definition}[The division data attached to $(f_0,\Omega)$]
\label{d:def:divisiondata}
Assume $f_0$ has no zero on $\Gamma$.  Then $f_0$ has finitely many
zeros $c_1,\dots,c_r$ in $\Omega$, of multiplicities
$m_1,\dots,m_r$, because $f_0$ is holomorphic near the compact set
$\overline\Omega$ and its zeros are isolated.  Put
\[
 d=\sum_{j=1}^{r}m_j,\qquad
 P_0(Z)=\prod_{j=1}^{r}(Z-c_j)^{m_j}\in\CC[Z],
\]
choose $V$ containing no further zeros of $f_0$, and set
\[
 u_0=f_0/P_0\in\Hol(V)^{\times}.
\]
If $f_0$ has no zero in $\Omega$, take $d=0$, $P_0=1$, $u_0=f_0$.
\end{definition}

The compactness used here is compactness of an ordinary complex set.
No compactness of a surcomplex ball is invoked anywhere in this part,
and none is available.

\begin{lemma}[Ordinary division, extended coefficientwise]
\label{d:lem:ordinarydivision}
There are $\CC$-linear maps
$T\colon\Hol(V)\to\CC[Z]_{<d}$ and $Q_0\colon\Hol(V)\to\Hol(V)$ with
\begin{equation}\label{d:eq:ordinarydivision}
 h=T h+P_0\,Q_0h\qquad(h\in\Hol(V)),
\end{equation}
and the decomposition \eqref{d:eq:ordinarydivision} is unique among
pairs (polynomial of degree $<d$, element of $\Hol(V)$).  Applied
exponent by exponent, $T$ and $Q_0$ extend to $\K$-linear maps on
$\Hc(V)$ that \emph{never enlarge a support}: $\supp Th$ and
$\supp Q_0h$ are contained in $\supp h$.  For $d=0$ one takes $T=0$ and
$Q_0=\mathrm{id}$.
\end{lemma}

\begin{proof}
$Th$ is the Hermite interpolation polynomial: the unique polynomial of
degree $<d$ whose derivatives through order $m_j-1$ at $c_j$ agree with
those of $h$, for every $j$.  Then $h-Th$ vanishes to order $m_j$ at
each $c_j$, so $(h-Th)/P_0$ has removable singularities and defines
$Q_0h\in\Hol(V)$.  If $Th+P_0Q_0h=\widetilde T h+P_0\widetilde Q_0h$
then the difference of the polynomial parts has degree $<d$ and is
divisible by $P_0$, hence is zero.  The coefficientwise extension is
immediate, and it cannot create exponents because it acts inside each
exponent separately.  That last clause is not bookkeeping: it is what
makes the recursion of Theorem~\ref{d:thm:preparation} live inside the
Neumann~\cite{Neumann} monoid of a fixed positive set.
\end{proof}

\subsection{Hahn--Weierstrass preparation: the global form}
\label{d:sub:preparation}

The theorem is stated on the whole Jordan domain at once.  The local
single-zero form is recovered as Corollary~\ref{d:cor:localprep}.  The
order matters: the exact weighted argument principle
(Theorem~\ref{d:thm:weightedarg}) and the contour-moment reconstruction
(Theorem~\ref{d:thm:moments}) both require one monic polynomial valid
over the entire domain, and neither is derivable from local
factorizations at the separate zeros, since assembling those would
require a common factorization that the local theorem does not supply.

\begin{theorem}[Hahn--Weierstrass preparation, global form]
\label{d:thm:preparation}
Let $F\in\Hc(U)$ be nonzero, normalized as in
\eqref{d:eq:normalform}, and suppose $f_0$ has no zero on $\Gamma$.
Let $P_0,u_0,d,V$ be the division data of
Definition~\ref{d:def:divisiondata}, and put
\[
 H=E/u_0\in\Hc_{>0}(V),\qquad S=\supp H\subset\No_{>0}.
\]
Then there exist unique
\[
 p\in\K[Z],\quad \deg p<d,\qquad q\in\Hc_{>0}(V)
\]
such that
\begin{equation}\label{d:eq:preparation}
 F=\tm^{\beta}\,u_0\,(1+q)\,W,\qquad W=P_0+p .
\end{equation}
Every coefficient of $p$ lies in $\mm$, the factor $u_0(1+q)$ is a unit
of $\Hc(V)$, and
\[
 \supp p\cup\supp q\subset\Neu{S}\setminus\{0\},
\]
where $\Neu{S}$ is the Neumann monoid generated by $S$.  The degree
$d$ is the \emph{total} ordinary zero count of $f_0$ in $\Omega$.  When
$d=0$ one has $W=1$ and $1+q=1+H$.
\end{theorem}

\begin{proof}
Dividing \eqref{d:eq:normalform} by $\tm^{\beta}u_0$ turns the assertion
into the equation
\begin{equation}\label{d:eq:factoreq}
 H=p+P_0q+pq
\end{equation}
for $p$ of degree $<d$ and $q$ of positive support.  Introduce a formal
bookkeeping parameter $s$ and look for homogeneous corrections
$p=\sum_{n\ge1}p_n$, $q=\sum_{n\ge1}q_n$, where $p_n,q_n$ carry formal
degree $n$ in $H$.  Define
\begin{equation}\label{d:eq:preprecursion}
 p_1=T H,\quad q_1=Q_0H,
 \qquad
 C_n=\sum_{j=1}^{n-1}q_jp_{n-j},\quad
 p_n=-T C_n,\quad q_n=-Q_0C_n\quad(n\ge2).
\end{equation}
All operations take place in $\Hc(V)$, and each $p_n$ is a polynomial
of degree $<d$ by construction.  By induction, using
Lemma~\ref{d:lem:ordinarydivision} (the operators do not enlarge
supports) and the fact that a product adds supports,
\[
 \supp p_n\cup\supp q_n\subset S^{(n)}
 :=\{\sigma_1+\dots+\sigma_n:\sigma_i\in S\}.
\]
Since $S$ is a well-ordered set of positive surreals, Neumann's lemma
in its finite-word form (Sections~\ref{a:sec:summability}--\ref{a:sec:fine}) says that
$\Neu{S}=\bigcup_{n\ge0}S^{(n)}$ is well ordered and that each of its
elements is $S^{(n)}$-representable for only finitely many $n$, with
finitely many words in each case.  Hence $(p_n)_{n\ge1}$ and
$(q_n)_{n\ge1}$ are Hahn summable, and
\[
 p=\sum_{n\ge1}p_n,\qquad q=\sum_{n\ge1}q_n
\]
exist, with supports in $\Neu{S}\setminus\{0\}$.  Because the
polynomial degrees are uniformly bounded by $d$, the coefficients of
$p$ are single elements of $\mm$ and $p$ is one polynomial over $\K$.

Equation \eqref{d:eq:ordinarydivision} makes the degree-one part of
\eqref{d:eq:factoreq} read $H=p_1+P_0q_1$, and for $n\ge2$ the
degree-$n$ part reads $0=p_n+P_0q_n+C_n$, which is
\eqref{d:eq:preprecursion}.  Summing the identities --- a licensed Hahn
sum of a summable family, evaluated at the bookkeeping value $s=1$, and
not a claim that the partial sums converge for any topology --- gives
\eqref{d:eq:factoreq}, hence \eqref{d:eq:preparation}.  The leading
coefficient function of $u_0(1+q)$ is $u_0$, nowhere zero on $V$, so
that factor is a unit of $\Hc(V)$ by the unit criterion of Section~\ref{c:p4}.

Uniqueness.  Suppose $(p,q)$ and $(\widetilde p,\widetilde q)$ both
solve \eqref{d:eq:factoreq} with the stated support conditions, and let
$\alpha$ be the least exponent occurring in
$\supp(p-\widetilde p)\cup\supp(q-\widetilde q)$, assuming this set is
nonempty.  The difference of the two copies of the nonlinear term $pq$
is $(p-\widetilde p)q+\widetilde p(q-\widetilde q)$; each summand is a
difference of valuation $\ge\alpha$ times a factor of strictly positive
support, so it has valuation $>\alpha$ and contributes nothing at
exponent $\alpha$.  Comparing coefficients of $\tm^{\alpha}$ in
\eqref{d:eq:factoreq} therefore gives
\[
 [\tm^{\alpha}](p-\widetilde p)+P_0\,[\tm^{\alpha}](q-\widetilde q)=0
\]
in $\Hol(V)$, with the first summand a polynomial of degree $<d$.
Uniqueness in \eqref{d:eq:ordinarydivision} forces both to vanish,
contradicting the choice of $\alpha$.
\end{proof}

\begin{remark}[A second proof, by recursion on the exponent]
\label{d:rem:exponentrecursion}
The same factorization may be built by transfinite recursion along the
well-ordered set $M=\Neu{S}_{>0}$ instead of along the formal degree.
Writing $H=\sum_{\delta}\tm^{\delta}h_\delta$,
$p=\sum_{\delta}\tm^{\delta}p_\delta$,
$q=\sum_{\delta}\tm^{\delta}q_\delta$, equation
\eqref{d:eq:factoreq} at exponent $\delta$ reads
\[
 r_\delta=h_\delta-\!\!\sum_{\substack{\alpha,\beta\in M\\ \alpha+\beta=\delta}}\!\!p_\alpha q_\beta,
 \qquad p_\delta=T(r_\delta),\quad q_\delta=Q_0(r_\delta),
\]
the sum being finite with all indices strictly below $\delta$.  This is
ordinary transfinite recursion on a well-ordered set; limit stages
require no analytic limit.  The two constructions produce the same
$(p,q)$ by the uniqueness clause.
\end{remark}

\begin{remark}[What the proof does not use]
\label{d:rem:nonorm}
No norm on $\Hc(V)$ is introduced, no contraction estimate is made, no
completeness is invoked, and it is \emph{not} assumed that
$n\min S$ is cofinal in the surreal value class.  The last point is the
one that would be needed for a metric fixed-point argument and is
exactly what a high-rank Hahn field does not supply: corrections of
valuation $n\delta$ need not eventually be smaller than every
prescribed surreal tolerance.  The construction avoids the issue
because coefficients are assigned on a well-ordered support with
finitely many earlier contributions at each exponent.  The evaluation at
$s=1$ in the first proof is a licensed normal sum, not convergence of a
sequence of partial sums.
\end{remark}

\begin{corollary}[Local preparation at a single zero]
\label{d:cor:localprep}
Let $F\in\Hc(U)$ be nonzero with leading coefficient function $f_0$,
and let $a\in U$ satisfy $\ord_a(f_0)=m\ge1$.  On a sufficiently small
ordinary disk $a+D_\varrho\subset U$ on which $f_0$ has no other zero,
\begin{equation}\label{d:eq:localprep}
 F(a+w)=\tm^{\beta}\Bigl(w^{m}+\sum_{j=0}^{m-1}p_jw^{j}\Bigr)V(w),
 \qquad p_j\in\mm,
\end{equation}
where $V\in\Hc(D_\varrho)$ has leading coefficient function
$u_0(w)=f_0(a+w)/w^{m}$, nowhere zero on $D_\varrho$; hence $V$ is a
unit and $\ev V$ vanishes nowhere on $\halo{D_\varrho}$.
\end{corollary}

\begin{proof}
This is Theorem~\ref{d:thm:preparation} with $\Omega=a+D_\varrho$, for
which $P_0(w)=w^m$ and $r=1$.  Equivalently, one may run the recursion
\eqref{d:eq:preprecursion} with the single-zero division operators
\[
 T_mg=\sum_{j=0}^{m-1}\frac{g^{(j)}(0)}{j!}w^{j},
 \qquad R_mg=\frac{g-T_mg}{w^{m}},
\]
which are the specialization of $T,Q_0$ to $P_0=w^m$; the apparent
singularity of $R_mg$ at $0$ is removable.  The most restrictive input
form --- leading part exactly $w^m$, i.e.\ $A=w^m+H$ with
$H\in\Hc_{>0}$ --- is obtained by first dividing by $u_0$, and the
general statement follows by multiplying the unit back in.
\end{proof}

No uniform complex bound on the infinitely many perturbation
coefficients is required anywhere above.  The common holomorphic domain
together with the positive well-ordered support replaces every
hypothesis of that kind; in particular the coefficient functions of $E$
may grow arbitrarily fast with the exponent.

\subsection{Division by the prepared polynomial}
\label{d:sub:division}

\begin{theorem}[Hahn--Weierstrass division]
\label{d:thm:division}
Let $W=P_0+p$ be the monic polynomial of
Theorem~\ref{d:thm:preparation}.  Every $A\in\Hc(V)$ has a unique
decomposition
\begin{equation}\label{d:eq:division}
 A=WQ+R,\qquad Q\in\Hc(V),\quad R\in\K[Z],\quad \deg R<d .
\end{equation}
\end{theorem}

\begin{proof}
Let $\mathcal L(Q)=Q_0(pQ)$, a $\K$-linear operator on $\Hc(V)$.  Each
application of $\mathcal L$ adds an element of the positive set
$\supp p$ to the support, while $Q_0$ alters only coefficient
functions.  Hence the family $\bigl((-\mathcal L)^{n}Q_0A\bigr)_{n\ge0}$
is supported in $\supp A+\Neu{\supp p}$, with finitely many
contributions at each exponent by Neumann's lemma, and
\[
 Q=\sum_{n\ge0}(-\mathcal L)^{n}Q_0A\in\Hc(V)
\]
is a Hahn sum satisfying $Q=Q_0A-Q_0(pQ)=Q_0(A-pQ)$.  Put
$R=T(A-pQ)$; since the degrees are uniformly bounded by $d$, $R$ is a
single polynomial over $\K$.  Applying
\eqref{d:eq:ordinarydivision} to $A-pQ$ gives $A-pQ=P_0Q+R$, i.e.\
\eqref{d:eq:division}.

If $WQ+R=0$ with $(Q,R)\ne(0,0)$, let $\alpha$ be the least exponent
occurring in $Q$ or $R$.  The term $pQ$ has support in
$\supp p+\supp Q$, all of whose elements exceed $\alpha$, so at
exponent $\alpha$ the identity reduces to $P_0\,q+r=0$ with
$q\in\Hol(V)$ and $\deg r<d$.  Uniqueness in
\eqref{d:eq:ordinarydivision} forces $q=r=0$, a contradiction.  For
$d=0$ the statement is the identity $A=1\cdot A+0$.
\end{proof}

\subsection{Roots of a monic polynomial with infinitesimal corrections}
\label{d:sub:roots}

\begin{lemma}[Finite roots and infinitesimal corrections]
\label{d:lem:smallroots}
Let $W(Z)=Z^{d}+\sum_{j<d}b_jZ^{j}$ with every $b_j\in\Ov$ and $d\ge1$.
Then $W$ has exactly $d$ roots in $\K$ counted with algebraic
multiplicity, and every root is finite. If all the $b_j$ are
infinitesimal, every root lies in $\mm$.  Moreover, if all the $b_j$
lie in a set-sized subfield $\K_G=\CC((\tm^{G}))$ for a divisible
set-sized ordered subgroup $G\subset\No$, then all the roots lie in
$\K_G$.
\end{lemma}

\begin{proof}
$\K$ is algebraically closed, which gives the $d$ roots.  Let $\xi$ be
a root.  If $v(\xi)<0$ then
$v(b_j\xi^{j})-v(\xi^{d})=v(b_j)+(j-d)v(\xi)>0$ for $j<d$, so the monic
term $\xi^{d}$ is the unique term of least valuation and cannot be
cancelled --- this is the exact place where a root of negative
valuation is excluded, and it uses only the ordered-field valuation, no
analysis. Thus every root is finite. If all the $b_j$ are
infinitesimal and $v(\xi)=0$, reducing $W(\xi)=0$ modulo $\mm$ gives
$\st(\xi)^d=0$, impossible for $\st(\xi)\ne0$. In this special case
every root is infinitesimal (including the possible root $0$).  For the last clause, $\K_G$ is itself
algebraically closed (Poonen), so it already splits $W$; passing to the
full class field adds no further roots.
\end{proof}

\subsection{The zero-cluster theorem}
\label{d:sub:zerocluster}

\begin{theorem}[Exact zero clustering; specialization of the zero divisor]
\label{d:thm:zerocluster}
Let $F\in\Hc(U)$ be nonzero, normalized as in
\eqref{d:eq:normalform}, with $f_0$ having no zero on $\Gamma$, and let
$d$, $W$ be as in Theorem~\ref{d:thm:preparation}.  Then:
\begin{enumerate}[label=(\roman*)]
\item $\ev F$ has exactly $d$ zeros in the halo $\halo\Omega$, counted
with analytic multiplicity, and they are precisely the roots in $\K$ of
$W$;
\item for every $a\in\Omega$, the monad $a+\mm$ contains exactly
$\ord_a(f_0)$ of them; in particular there are no zeros over points
where $f_0$ does not vanish, a simple zero of $f_0$ lifts to exactly one
simple surcomplex zero, and a zero of order $m$ splits into exactly $m$
roots (repetitions allowed);
\item equivalently, with zero divisors carrying multiplicities,
\begin{equation}\label{d:eq:divisor}
 \st_{*}\bigl(\operatorname{div}_0F\bigr)=\operatorname{div}_0(f_0);
\end{equation}
\item the zeros lie in a set-sized subfield $\K_G$, and the full zero
set of $\ev F$ in $\halo U$ is at most countable.
\end{enumerate}
\end{theorem}

\begin{proof}
(i) By \eqref{d:eq:preparation}, $\ev F=\tm^{\beta}\,\ev{u_0}\,
\ev{(1+q)}\,\ev W$ on $\halo V$.  The first three factors are
units: $u_0$ is nowhere zero on $V$, and $1+q$ has inverse
$\sum_{n\ge0}(-q)^{n}$, a Hahn sum by Neumann's lemma, so $\ev{(1+q)}$
takes values in $1+\mm$ at every halo point.  Hence the zeros of
$\ev F$, with their local analytic multiplicities, are exactly the
roots of $W$ with their polynomial multiplicities: multiplying by a
germ-invertible factor changes neither the location nor the least
nonzero Taylor degree at a root.

The lower coefficients of $W=P_0+p$ are finite, so by
Lemma~\ref{d:lem:smallroots} every root of $W$ is finite, so it has
a standard part.  Applying the standard-part ring homomorphism to the
factorization $W=\prod_{k=1}^{d}(Z-\xi_k)$ gives
\[
 P_0(Z)=\prod_{k=1}^{d}\bigl(Z-\st\xi_k\bigr)\quad\text{in }\CC[Z],
\]
so the $\st\xi_k$ are exactly the zeros $c_1,\dots,c_r$ of $f_0$ in
$\Omega$, each occurring $m_j$ times.  This proves (ii) and
\eqref{d:eq:divisor}, and also proves that all $d$ roots lie in
$\halo\Omega$.  Over a point where $f_0(a)\ne0$ the faithful-evaluation
estimate of Section~\ref{c:p4}, $\st\bigl(\tm^{-\beta}\ev F(a+\epsilon)\bigr)=f_0(a)$,
already excludes a zero; so does the standard-part identity just
derived.

(iv) The coefficients of $W$ form a set, so they lie in some
$\K_G=\CC((\tm^{G}))$ with $G$ a divisible set-sized ordered subgroup
of $\No$; by Lemma~\ref{d:lem:smallroots} the roots lie in $\K_G$
too.  For countability, apply
Corollary~\ref{d:cor:localprep} at each zero of $f_0$ in $U$: the zeros
of the nonzero ordinary holomorphic function $f_0$ on the planar domain
$U$ form a discrete, hence at most countable, set, and each
standard-part fibre contains at most $\ord_a(f_0)$ distinct roots.  The
union of countably many finite sets is at most countable.
\end{proof}

\begin{remark}[Three things the theorem does not say]
\label{d:rem:clusterscope}
First, the count is exact but it is \emph{not} a reduction modulo
infinitesimals: the root positions may carry surreal terms at
arbitrarily deep scales, indeed at scales below every finite power of a
chosen infinitesimal, while the integer $d$ is exact.  Second,
$\halo\Omega=\st^{-1}(\Omega)$ is not the surcomplex ball over
$\Omega$: points infinitesimally inside a standard boundary point are
absent from the halo, so a root sitting in the monad of a boundary
point is counted by neither side.  Third, neither this theorem nor the
Rouch\'e statements below count zeros whose standard-scale coordinate
is infinite; those live outside every halo of a bounded ordinary domain
and require a different affine scale chart (Section~\ref{f:sec-allscale}).
\end{remark}

\begin{corollary}[Simple-root lifting and its displacement]
\label{d:cor:displacement}
Suppose $f_0(c)=0$ and $a:=f_0'(c)\ne0$.  Then $\ev F$ has exactly one
zero $\xi=c+\delta$ in $c+\mm$, it is simple, and $\delta$ has support
inside $\Neu{\supp E}$.  Writing $H=E$ and
$\sigma=\min\supp H>0$ when $H\ne0$,
\begin{equation}\label{d:eq:displacement}
 \delta=-\frac{\ev H(c)}{a}
       +\frac{\ev H(c)\,\ev{(H')}(c)}{a^{2}}
       -\frac{f_0''(c)\,\ev H(c)^{2}}{2a^{3}}+R,
 \qquad v(R)\ge3\sigma .
\end{equation}
In particular, if $E=\tm^{\rho}g+(\text{exponents}>\rho)$ with
$\rho>0$, the leading correction is
\begin{equation}\label{d:eq:leadingcorrection}
 \xi=c-\tm^{\rho}\frac{g(c)}{f_0'(c)}+o(\tm^{\rho}),
\end{equation}
and if $g(c)=0$ the displayed correction vanishes and the actual
displacement is smaller.
\end{corollary}

\begin{proof}
Existence, uniqueness and simplicity are
Theorem~\ref{d:thm:zerocluster}(ii) with $m=1$; there $W=Z-c+p_0$ and
$\xi=c-p_0$, whose support lies in $\Neu{\supp E}$ by
Theorem~\ref{d:thm:preparation}.  For \eqref{d:eq:displacement},
substitute $\delta=\delta_1+\delta_2+\cdots$ into the Taylor identity
\[
 0=a\delta+\tfrac12f_0''(c)\delta^{2}+\cdots
   +\ev H(c)+\ev{(H')}(c)\,\delta+\cdots
\]
assigning formal degree one to $H$ and to each of its derivatives.  The
degree-one and degree-two equations give the two displayed corrections;
at every higher degree, division by $a$ solves the next term uniquely.
The degree-$n$ term has support inside $(\supp H)^{(n)}$, so Neumann's
lemma sums the expansion and bounds the remaining valuation below by
$3\sigma$.  Formal substitution identifies the resulting element with
the root already produced by preparation.  This is a homogeneous
support argument; it is not an appeal to convergence of Newton
iterates.
\end{proof}

\subsection{Three Rouch\'e theorems, and the inequality that is not enough}
\label{d:sub:rouche}

The three statements below differ in where the strict inequality is
placed, and --- a distinction which must not be flattened --- in whether
they preserve the count monad by monad or only in total.

\begin{theorem}[Valuation Rouch\'e: fibrewise stability]
\label{d:thm:rouchevaluation}
Let $F\in\Hc(U)$ be normalized as in \eqref{d:eq:normalform} with $f_0$
nonzero on $\Gamma$, and let $G\in\Hc(U)$ satisfy $\vH(G)>\vH(F)$.
Then $\ev F$ and $\ev{(F+G)}$ have the same number of zeros in
$\halo\Omega$ counted with multiplicity, and, more precisely, the same
number in \emph{every individual monad} $a+\mm$ with $a\in\Omega$.
\end{theorem}

\begin{proof}
$F+G$ has the same leading exponent and the same leading coefficient
function $f_0$ as $F$, and its perturbation
$E+\tm^{-\beta}G$ still has strictly positive support.  Apply
Theorem~\ref{d:thm:zerocluster}(ii) to each of $F$ and $F+G$: both
counts in $a+\mm$ equal $\ord_a(f_0)$.
\end{proof}

\begin{theorem}[Ordinary-margin Rouch\'e: stability of the total]
\label{d:thm:rouchemargin}
Let $F,G\in\Hc(U)$ and $\beta\in\No$ be such that
\[
 \tm^{-\beta}F=f+E_F,\qquad \tm^{-\beta}G=g+E_G,
 \qquad E_F,E_G\in\Hc_{>0}(U),\quad f,g\in\Hol(U),
\]
and suppose the \emph{ordinary} strict inequality
\[
 \abs{g(z)}<\abs{f(z)}\qquad(z\in\Gamma)
\]
holds on the boundary.  Then $\ev F$ and $\ev{(F+G)}$ have the same
total number of zeros in $\halo\Omega$, counted with multiplicity.
\end{theorem}

\begin{proof}
The boundary inequality makes $f$ and $f+g$ nonzero on $\Gamma$, and
ordinary Rouch\'e gives them equal total zero counts in $\Omega$.
Theorem~\ref{d:thm:zerocluster}(i) transfers each of those ordinary
counts to the corresponding surcomplex count in $\halo\Omega$.
\end{proof}

\begin{remark}[Total, not fibrewise]
\label{d:rem:totalonly}
Theorem~\ref{d:thm:rouchemargin} generally preserves only the
\emph{total}: $f$ and $f+g$ may have zeros at different ordinary
points, so the distribution among monads can change completely.
Theorem~\ref{d:thm:rouchevaluation} is the form that preserves the
count in each fibre, because it does not move the leading coefficient
function at all.
\end{remark}

\begin{corollary}[Real-margin Rouch\'e: a surcomplex-valued hypothesis]
\label{d:cor:rouchereal}
Write $\tm^{-\beta}F=f_0+E_F$ and
$\tm^{-\beta}G=g_0+E_G$ with $f_0,g_0\in\Hol(U)$ and
$E_F,E_G\in\Hc_{>0}(U)$. Thus $f_0,g_0$ are the exponent-zero
coefficient functions (a zero coefficient is allowed). Suppose $f_0$
has no zero on $\Gamma$.  If for \emph{one ordinary real} $q$ with
$0\le q<1$
\begin{equation}\label{d:eq:realmargin}
 \bigl\lvert\ev G(z)-\ev F(z)\bigr\rvert\le q\,\bigl\lvert\ev F(z)\bigr\rvert
 \qquad(z\in\Gamma),
\end{equation}
then $\ev F$ and $\ev G$ have the same number of zeros in
$\halo\Omega$, counted with multiplicity.
\end{corollary}

\begin{proof}
Divide \eqref{d:eq:realmargin} by the positive surreal $\tm^{\beta}$
and take standard parts.  Because $q$ is an ordinary real strictly
below $1$, the reduced inequality is
$\abs{g_0-f_0}\le q\abs{f_0}<\abs{f_0}$ on $\Gamma$, which is strict.
Ordinary Rouch\'e now applies and also shows $g_0$ has no zero on
$\Gamma$; Theorem~\ref{d:thm:zerocluster}(i) lifts both counts.
\end{proof}

\begin{example}[Why a bare strict surreal inequality is not enough]
\label{d:ex:nomargin}
Replace \eqref{d:eq:realmargin} by the bare strict surreal inequality
$\abs{\ev G-\ev F}<\abs{\ev F}$ on $\Gamma$.  Strictness can disappear
under the standard part, and then nothing is left to apply.  Take
$\Omega=D$ the unit disk, $\Gamma$ the unit circle,
\[
 F(z)=z-2,\qquad \ev G-\ev F\equiv 1-\tm .
\]
On $\Gamma$ one has $\abs{\ev F}=\abs{z-2}\ge1$ with equality exactly at
$z=1$, while $\abs{\ev G-\ev F}=1-\tm<1$, so the strict surreal
inequality holds at every boundary point.  Its standard-part reduction
is $1\le\abs{z-2}$, which is not strict at $z=1$; and indeed
$G=z-1-\tm$ has leading coefficient function $z-1$, which vanishes at
the boundary point $1$.  So $G$ does not satisfy the hypothesis of the
zero-cluster theorem at all, and its unique zero $1+\tm$ sits in the
monad of a boundary point, outside $\halo\Omega$. Since the halo is fine-clopen, this
point does lie in the fine interior of its complement. This example
exhibits loss of the boundary-reduction hypothesis, but both halo zero
counts happen to be zero.

For an actual failure of the zero-count conclusion, instead take
$F(Z)=Z$ and $G(Z)=Z-1+\tm$ on the same unit circle. Then
$\abs{\ev G-\ev F}=1-\tm<1=\abs{\ev F}$ everywhere on the circle.
Yet $F$ has one zero in $\halo\Omega$ and $G$ has none: its unique
zero $1-\tm$ has standard part $1$ and is excluded from that halo.
No manuscript in the collection asserts the
version of Rouch\'e with a bare surreal inequality, and this merged
document does not either: the ordinary real margin $q<1$ of
Corollary~\ref{d:cor:rouchereal}, or the ordinary boundary inequality of
Theorem~\ref{d:thm:rouchemargin}, is exactly the hypothesis the proofs
use.
\end{example}

\subsection{Contour moments and Newton's identities}
\label{d:sub:moments}

The prepared polynomial $W$ has a second, completely independent
constructive description: it is determined by finitely many Hahn-valued
contour integrals.  The proof consumes the exact weighted argument
principle, Theorem~\ref{d:thm:weightedarg} of Section~\ref{d:sec:residues}, which is proved
there from Theorem~\ref{d:thm:preparation} and is independent of this
subsection.

\begin{theorem}[Contour moments of a zero cluster]
\label{d:thm:moments}
In the situation of Theorem~\ref{d:thm:preparation}, let
$\xi_1,\dots,\xi_d\in\halo\Omega$ be the zeros of $\ev F$ listed with
multiplicity, and define the moments
\begin{equation}\label{d:eq:moments}
 s_k=\frac1{2\pi i}\ointH_{\Gamma}w^{k}\,
      \frac{F'(w)}{F(w)}\,\dd w\qquad(k\ge0).
\end{equation}
Then $s_0=d$ and $s_k=\sum_{j=1}^{d}\xi_j^{k}$ for all $k\ge0$.  The
moments $s_1,\dots,s_d$ determine $W$ exactly, by Newton's identities:
\begin{equation}\label{d:eq:newton}
 e_0=1,\qquad k\,e_k=\sum_{j=1}^{k}(-1)^{j-1}e_{k-j}s_j
 \quad(1\le k\le d),
\end{equation}
\[
 W(Z)=Z^{d}-e_1Z^{d-1}+e_2Z^{d-2}-\cdots+(-1)^{d}e_d .
\]
\end{theorem}

\begin{proof}
The moment identity is Theorem~\ref{d:thm:weightedarg} applied to the
weight $A(Z)=Z^{k}$, whose evaluation at $\xi_j$ is $\xi_j^{k}$.
Newton's identities \eqref{d:eq:newton} are an identity of symmetric
functions over any field of characteristic zero --- obtained by
comparing coefficients in the logarithmic derivative of
$\prod_j(1+\xi_jT)$ --- and $\K$ has characteristic zero, so the
divisions by $k$ are legitimate.  The recursion determines
$e_1,\dots,e_d$, hence the monic polynomial with those elementary
symmetric functions, which by Theorem~\ref{d:thm:zerocluster}(i) is
$W$.
\end{proof}

\begin{remark}
The moments $s_k$ are genuinely surcomplex numbers, and neither
description of $W$ --- the recursion \eqref{d:eq:preprecursion} by
holomorphic division, or the moments \eqref{d:eq:moments} by symmetric
functions --- discards infinitesimal corrections to the roots.  Both
require the \emph{global} preparation: the moment integral is taken
around the whole of $\Gamma$ and is matched against one polynomial with
all $d$ roots.
\end{remark}

% =====================================================================

\section{Residues: three notions, and the theorems that join them}
\label{d:sec:residues}

Three different objects in this subject are called ``the residue'',
they are attached to different classes of function, and two of them
take \emph{different numerical values at the same point}.  They are
kept apart here by symbol:
\[
 \begin{aligned}
 \Res&\quad\text{formal, in }\K((X)),\\
 \ResH a&\quad\text{cluster, at an ordinary centre }a,\\
 \Resat b&\quad\text{actual, at a surcomplex point }b.
 \end{aligned}
\]
Using one symbol for the first two turns a correct statement into a
false one; the standing example of
\S\ref{d:sub:standingexample} shows exactly how.  Two bridge theorems
say when and how they agree, and no manuscript in the collection
contains both.

\subsection{R1: the formal residue in \texorpdfstring{$\K((X))$}{K((X))}}
\label{d:sub:R1}

\begin{definition}[Formal residue]
\label{d:def:formalres}
For $A=\sum_{n\ge-N}a_nX^{n}\in\K((X))$ --- a Laurent series with
coefficients in $\K$ and only \emph{finitely many} negative powers ---
set
\[
 \Res_{X=0}\bigl(A(X)\,\dd X\bigr)=a_{-1}.
\]
For a meromorphic germ at a point $b\in\K$ one takes $X=z-b$ and writes
$\Resat{b}$ for the resulting coefficient; this is the notion R3 of
\S\ref{d:sub:R3}.
\end{definition}

This is a purely algebraic notion, valid at \emph{any} surcomplex
centre, and it belongs to the germ layer of Section~\ref{b:part3}, not to the coherent
class.  Its two basic laws are proved there and are used freely below:
the change-of-variable law
$\Res_{Y=0}A(\psi(Y))\psi'(Y)\dd Y=m\,\Res_{X=0}A(X)\dd X$ for
$\psi$ of order $m$ (with $m=1$ the invertible case), and the local
argument principle $\Res_{X=0}(A'/A)\dd X=m$ for $A=X^{m}U$ with
$U(0)\ne0$.  The primitive obstruction is likewise a statement about
$\K((X))$: a formal meromorphic differential has a formal meromorphic
primitive if and only if its residue vanishes, every monomial except
$X^{-1}$ being integrable.

\subsection{R2: the cluster residue at an ordinary centre}
\label{d:sub:R2}

\begin{definition}[The class $\Mc(U)$ and the cluster residue]
\label{d:def:clusterres}
Let $\Mc(U)$ consist of the Hahn series $A=\sum_{s\in S}a_s\tm^{s}$
with $S\subset\No$ a well-ordered \emph{set} and every $a_s$
meromorphic on the connected ordinary domain $U$.  The poles are free to
vary with the exponent $s$, and their orders are not assumed bounded in
$s$.  For $a\in U$ define the \emph{cluster residue} at the
\emph{ordinary} centre $a$ by
\begin{equation}\label{d:eq:clusterres}
 \ResH{a}\bigl(A(\zeta)\dd\zeta\bigr)
  =\sum_{s\in S}\Bigl(\Res_{\zeta=a}a_s(\zeta)\,\dd\zeta\Bigr)\tm^{s}.
\end{equation}
The support of the right-hand side is contained in $S$, so it is a
surcomplex number.
\end{definition}

The coefficient ring of $\Mc(U)$ is a field, so a nonzero element of
$\Mc(U)$ is invertible by leading-term factorization together with the
infinitesimal functional calculus of Section~\ref{c:p4}.  The adjective ``cluster''
is the whole point: $\ResH{a}$ is attached to an ordinary centre and
records a whole infinitesimal cluster sitting above it.  It is
\emph{not} the residue of a meromorphic germ at the exact point $a$.

\subsection{R3: the actual residue at a surcomplex point}
\label{d:sub:R3}

\begin{definition}[Actual local residue]
\label{d:def:actualres}
Let $A,B\in\Hc(U)$ with $B\ne0$, and let $b\in\halo U$ be a point at
which the germ of $A/B$ is a formal Laurent series with finite
principal part,
\[
 \frac{A(b+h)}{B(b+h)}=\sum_{n\ge-N}C_nh^{n}\in\K((h)),
\]
which holds at every $b\in\halo U$: the coherent identity theorem
excludes the zero germ of $B$, so its order there is finite.
Its \emph{actual local residue} is
$\Resat{b}\bigl((A/B)(z)\dd z\bigr)=C_{-1}$.  It is centred at $b$, not
at $\st b$.
\end{definition}

R3 is defined on quotients of coherent families.  It is \emph{not}
available on all of $\Mc(U)$: an arbitrary element of $\Mc(U)$ need not
be such a quotient, and need not have any actual surcomplex pole in a
given monad at all.

\subsection{The standing example: \texorpdfstring{$1/(\zeta-\epsilon)$}{1/(zeta-eps)}}
\label{d:sub:standingexample}

\begin{example}[The three residues disagree, and both disagreements are correct]
\label{d:ex:movingpole}
Let $\epsilon\in\mm$, $\epsilon\ne0$, and
$R(\zeta)=1/(\zeta-\epsilon)$.  On an ordinary annulus about $0$,
\begin{equation}\label{d:eq:movingpole}
 \frac1{\zeta-\epsilon}=\sum_{n\ge0}\epsilon^{n}\zeta^{-n-1},
\end{equation}
which, after expanding each $\epsilon^{n}$ in the monomial scale, is an
element of $\Mc(\CC)$ whose coefficient functions all have their poles
at the single ordinary point $0$.  Then
\[
 \ResH{0}\bigl(R(\zeta)\dd\zeta\bigr)=1,
 \qquad
 \Resat{\epsilon}\bigl(R(\zeta)\dd\zeta\bigr)=1,
 \qquad
 \Resat{0}\bigl(R(\zeta)\dd\zeta\bigr)=0 .
\]
The first value is read off \eqref{d:eq:movingpole}: only the $n=0$
coefficient $\zeta^{-1}$ contributes a residue at $0$, with value $1$.
The second is immediate, the actual pole being the simple pole at
$\epsilon$.  The third is also immediate and is not a competing
computation: on the fine neighborhood $\epsilon\mm$, equivalently when
$\zeta/\epsilon$ is infinitesimal, the function expands as $-\epsilon^{-1}\sum_{n\ge0}(\zeta/\epsilon)^{n}$, a power
series with no negative powers, so the meromorphic germ at the exact
point $0$ is regular there. The weaker inequality
$\abs\zeta<\abs\epsilon$ does not ensure Hahn summability: at
$\zeta=\epsilon/2$ all nonzero geometric terms have valuation zero
before multiplication by $-\epsilon^{-1}$.  The cluster residue at the ordinary centre
$0$ correctly accounts for the actual pole at $\epsilon$; the local
residue at $0$ correctly reports that there is no pole at $0$.

Two scope statements travel with the example.  The expansion
\eqref{d:eq:movingpole} is \emph{not} asserted to evaluate throughout
the monad containing the pole; it is a coefficientwise identity valid
away from that monad.  And the term $\epsilon^{n}\zeta^{-n-1}$
has a pole of order $n+1$ in its ordinary coefficient. To check that
these orders survive Hahn regrouping, put $\alpha=v(\epsilon)>0$ and
$c=\operatorname{lc}(\epsilon)$. At exponent $n\alpha$, the $n$th term
contributes $c^n\zeta^{-n-1}$; terms with larger index have valuation
$>n\alpha$, and smaller indices have pole order at most $n$. Hence the
coefficient at $n\alpha$ has pole order exactly $n+1$. These unbounded
orders are excluded by the hypothesis of Theorem~\ref{d:thm:bridge2}.
\end{example}

\subsection{The coefficientwise residue theorem, in maximal generality}
\label{d:sub:Hresidue}

\begin{theorem}[Cluster residue theorem on $\Mc(U)$]
\label{d:thm:Hresidue}
Let $A\in\Mc(U)$, let $\Omega$ be a bounded Jordan domain with
piecewise $C^1$ boundary $\Gamma$ and $\overline\Omega\subset U$, and
suppose no coefficient $a_s$ has a pole on $\Gamma$.  Then
\begin{equation}\label{d:eq:Hresidue}
 \frac1{2\pi i}\ointH_{\Gamma}A(\zeta)\,\dd\zeta
  =\sum_{a\in\Omega}\ResH{a}\bigl(A(\zeta)\dd\zeta\bigr).
\end{equation}
Only a set of ordinary centres contributes.  The family on the right is
Hahn summable even when infinitely many centres contribute in total.
\end{theorem}

\begin{proof}
Fix an exponent $s\in S$.  The coefficient $a_s$ is meromorphic on $U$
and $\overline\Omega\subset U$ is compact, so $a_s$ has finitely many
poles in $\overline\Omega$, none on $\Gamma$.  Hence only finitely many
centres contribute at the fixed exponent $s$, and the ordinary residue
theorem gives
$\frac1{2\pi i}\int_{\Gamma}a_s=\sum_{a\in\Omega}\Res_{a}a_s$, a finite
sum.  The union of the supports of the family
$\bigl(\ResH{a}A\bigr)_{a\in\Omega}$ is contained in the single fixed
well-ordered set $S$, and each exponent $s$ occurs for only finitely
many $a$ by the previous sentence: these are exactly the two clauses of
Hahn summability.  Summing the ordinary identity coefficientwise,
which is the definition \eqref{d:eq:clusterres} of $\ResH{a}$ and the
definition of $\ointH$, gives \eqref{d:eq:Hresidue}.
\end{proof}

\begin{remark}[The hypotheses that are not needed, and the conclusion that does not transfer]
\label{d:rem:Hresiduescope}
Two stronger hypotheses appear in the sources and are unnecessary: a
\emph{common} closed discrete pole set for all coefficients, and a
single \emph{fixed finite} pole set.  The proof above uses neither;
meromorphy of each coefficient on the connected domain $U$ is enough,
because compactness of $\overline\Omega$ localizes the count at each
fixed exponent.

The price of that generality must be stated.  On $\Mc(U)$ there need be
no factorization into a quotient of coherent families, so the
conclusion of Theorem~\ref{d:thm:bridge1} --- finitely many actual
poles, and a \emph{finite} sum of actual residues --- is \emph{not}
available and must not be imported.  For a concrete separation, take
$U=D$, $\Omega=\{\abs z<\tfrac12\}$, distinct points
$b_n=1/(2n+4)\in\Omega$, and
\[
 A=\sum_{n\ge1}\frac{\tm^{n}}{\zeta-b_n}\in\Mc(D).
\]
Its support $\{1,2,3,\dots\}$ is a well-ordered set and each coefficient
function is meromorphic on $D$ with one pole, so $A\in\Mc(D)$; but
infinitely many ordinary centres contribute, and
$\sum_{a\in\Omega}\ResH{a}A=\sum_{n\ge1}\tm^{n}=\tm/(1-\tm)$, a Hahn sum
with infinitely many nonzero terms.  The poles accumulate at
$0\in\Omega$, so the coefficient family is not that of any quotient of
coherent sections, whose coefficient poles lie in the finite zero set of
one leading coefficient function inside $\overline\Omega$.  Summability,
not finiteness, is what Theorem~\ref{d:thm:Hresidue} provides.
\end{remark}

\begin{proposition}[Quotients of coherent families lie in $\Mc(U)$]
\label{d:prop:quotientinM}
Let $A,B\in\Hc(U)$ with $B\ne0$, and write
$B=\tm^{\beta}f_0(1+E_B)$ with $f_0$ the leading coefficient function of
$B$ and $E_B$ of strictly positive support (with meromorphic
coefficients, poles only at the zeros of $f_0$).  Then
$A/B\in\Mc(U)$: each of its coefficient functions is a finite sum of
products of meromorphic functions on $U$, hence meromorphic, and its
poles lie over the zeros of $f_0$.  The inclusion is proper.
\end{proposition}

\begin{proof}
$1/B=\tm^{-\beta}f_0^{-1}\sum_{n\ge0}(-E_B)^{n}$, and the family
$\bigl((-E_B)^{n}\bigr)_{n\ge0}$ is Hahn summable by Neumann's lemma
since $\supp E_B$ is a well-ordered set of positive surreals.  Hence at
each exponent only finitely many terms contribute, and each coefficient
of $1/B$ is a finite sum of finite products of the meromorphic
functions $f_0^{-1}$ and the coefficients of $E_B$.  Multiplying by $A$
preserves this.  Properness is the example of
Remark~\ref{d:rem:Hresiduescope}.
\end{proof}

This containment is what makes the two residue theorems compatible
rather than competing: on the quotient class both apply and agree,
which is the content of Bridge~1.  What it does not give is the
converse, and the converse is false.

\subsection{Bridge 1: the cluster residue is a finite sum of actual residues}
\label{d:sub:bridge1}

\begin{theorem}[Residue-cluster theorem]
\label{d:thm:bridge1}
Let $A,B\in\Hc(U)$ with $B\ne0$, let $a\in U$, and suppose the leading
coefficient function $f_0$ of $B$ has a zero of order $m\ge1$ at $a$ and
no other zero on a closed ordinary disk $\overline{D_\varrho(a)}\subset U$.
Let $\gamma_a=\partial D_\varrho(a)$, positively oriented.  Then $A/B$
has at most $m$ poles in the monad $a+\mm$, all of them actual points of
$\K$, and
\begin{equation}\label{d:eq:bridge1}
 \ResH{a}\Bigl(\frac{A}{B}(z)\,\dd z\Bigr)
 =\frac1{2\pi i}\ointH_{\gamma_a}\frac{A(z)}{B(z)}\,\dd z
 =\sum_{\substack{b\in a+\mm\\ b\ \mathrm{a\ pole\ of}\ A/B}}
    \Resat{b}\Bigl(\frac{A}{B}(z)\,\dd z\Bigr),
\end{equation}
a \emph{finite} sum of actual surcomplex residues.  Cancelled zeros of
the denominator contribute no pole.
\end{theorem}

\begin{proof}
The first equality is Theorem~\ref{d:thm:Hresidue} applied on
$D_\varrho(a)$, legitimate by Proposition~\ref{d:prop:quotientinM}: the
only ordinary centre inside is $a$.

For the second, translate $a$ to $0$.  Prepare
$B=\tm^{\beta}u_0(1+q)W$ by Theorem~\ref{d:thm:preparation} on a
neighbourhood of $\overline{D_\varrho}$, with $W$ monic of degree $m$
and infinitesimal lower coefficients.  The section
$C=\tm^{-\beta}A/\bigl(u_0(1+q)\bigr)$ lies in $\Hc(V)$ because
$u_0(1+q)$ is a unit.  Theorem~\ref{d:thm:division} gives
$C=WQ+R$ with $Q\in\Hc(V)$ and $R\in\K[Z]$, $\deg R<m$, so
\begin{equation}\label{d:eq:meromdecomp}
 \frac AB=Q+\frac RW .
\end{equation}
The first summand is coherent holomorphic on $V$, so its coefficientwise
contour integral vanishes by Cauchy's theorem (Section~\ref{c:p5}).  Over the
algebraically closed field $\K$, the proper fraction $R/W$ has a finite
partial-fraction expansion
\[
 \frac{R(Z)}{W(Z)}=\sum_{j}\sum_{k=1}^{m_j}
   \frac{c_{j,k}}{(Z-b_j)^{k}},
\]
where $b_1,\dots$ are the distinct roots of $W$, all infinitesimal by
Lemma~\ref{d:lem:smallroots}, and absent or cancelled principal parts
are recorded by zero coefficients.  At $b_j$ the other partial
fractions and $Q$ are analytic on a small fine disk, so the actual local
residue of $A/B$ at $b_j$ is exactly $c_{j,1}$; this identifies the
right-hand side of \eqref{d:eq:bridge1} with $\sum_jc_{j,1}$, and shows
that the poles are finite in number and lie in $\mm$.

Finally, on the translated ordinary circle $\gamma_a$, where $\abs z$ is a fixed
positive real and $b_j$ is infinitesimal,
\begin{equation}\label{d:eq:kernelexp}
 \frac1{(z-b_j)^{k}}
 =\sum_{n\ge0}\binom{k+n-1}{n}\frac{b_j^{n}}{z^{k+n}} ,
\end{equation}
a Hahn summable family by the all-lengths clause of Neumann's lemma
applied to the positive support of $b_j$. Increasing leading valuations
alone would not justify the finiteness clause in higher rank. The case
$b_j=0$ leaves only the $n=0$ term.  Its coefficientwise circle
integral divided by $2\pi i$ is the ordinary residue of each displayed
monomial at $0$: only the term with $k+n=1$, that is $k=1$ and $n=0$,
survives, giving $1$; all terms with $k\ge2$ give $0$.  Multiplying by
the fixed constants $c_{j,k}$ and adding finitely many terms gives
$\sum_jc_{j,1}$.  This is the second equality.

Note that the argument is finite partial fractions plus one summable
kernel expansion.  It is not a limit, and at no point are distinct
nearby surcomplex poles replaced by their common standard part:
standard parts organize the cluster, the residues themselves stay in
$\K$.
\end{proof}

\begin{theorem}[Residue theorem for coherent quotients, global form]
\label{d:thm:residuequotient}
Let $A,B\in\Hc(U)$ with $B\ne0$ and the leading coefficient function of
$B$ nonzero on $\Gamma=\partial\Omega$.  Then $M=A/B$ has finitely many
poles in $\halo\Omega$, all of them actual points of $\K$, and
\begin{equation}\label{d:eq:residuequotient}
 \ointH_{\Gamma}M(w)\,\dd w
  =2\pi i\sum_{b\in\halo\Omega}\Resat{b}\bigl(M(z)\dd z\bigr).
\end{equation}
The sum is finite; its individual terms may be infinite surcomplex
numbers.
\end{theorem}

\begin{proof}
By Theorem~\ref{d:thm:zerocluster} the leading coefficient function of
$B$ has finitely many zeros in $\Omega$ and every pole of $M$ lies in
one of their monads.  By Proposition~\ref{d:prop:quotientinM},
$M\in\Mc(U)$ with coefficient poles only at those ordinary points, so
Theorem~\ref{d:thm:Hresidue} replaces $\Gamma$ by finitely many small
disjoint ordinary circles around them.  On each circle apply
Theorem~\ref{d:thm:bridge1}.  Alternatively, prepare $B$ once on the
whole of $V$ and use \eqref{d:eq:meromdecomp} globally: the coherent
term integrates to zero and the proper rational term is handled by
\eqref{d:eq:kernelexp} exactly as above, with all roots of $W$ lying in
$\halo\Omega$.
\end{proof}

\begin{remark}
The decomposition \eqref{d:eq:meromdecomp} is considerably more than a
coefficientwise residue computation at fixed ordinary poles: it
separates the \emph{actual} poles of $\K$, including infinitesimally
split ones, from a coherent holomorphic remainder.  That separation is
supplied by preparation and division, and by nothing else in this
article.
\end{remark}

\subsection{Bridge 2: when the cluster residue is the residue at the exact point}
\label{d:sub:bridge2}

\begin{theorem}[Uniformly bounded pole order]
\label{d:thm:bridge2}
Let $F=\sum_{s\in S}f_s\tm^{s}$ have common well-ordered support $S$,
with each $f_s$ holomorphic on $D_\varrho(a)\setminus\{a\}$, and suppose
there is \emph{one} finite integer $d\ge1$ such that
\[
 (z-a)^{d}f_s(z)\ \text{extends holomorphically across }a
 \quad\text{for every }s\in S .
\]
Then $G:=(z-a)^{d}F$ is a coherent section in $\Hc\bigl(D_\varrho(a)\bigr)$
with the same support. The quotient $\ev G/(z-a)^d$ extends the
original evaluation from $\halo{(D_\varrho(a)\setminus\{a\})}$ to
$\halo{D_\varrho(a)}\setminus\{a\}$ and has at most a pole of
order $d$ at
the \emph{exact} point $a$, and
\begin{equation}\label{d:eq:bridge2}
 \ResH{a}\bigl(F(z)\dd z\bigr)
  =\Resat{a}\bigl(\ev F(z)\dd z\bigr)
  =\frac{G^{(d-1)}(a)}{(d-1)!}.
\end{equation}
\end{theorem}

\begin{proof}
The extensions have the same index set $S$, which is well ordered, so
$G\in\Hc(D_\varrho(a))$.  At the exact point $a$, writing $h=z-a$, the
germ of $\ev F$ is
$\sum_{n\ge0}\frac{G^{(n)}(a)}{n!}h^{\,n-d}\in\K((h))$, a Laurent series
with finite principal part; its coefficient of $h^{-1}$ is
$G^{(d-1)}(a)/(d-1)!$, which is the middle and right members of
\eqref{d:eq:bridge2} by Definition~\ref{d:def:formalres}.  For the left
member, the ordinary residue of $f_s$ at $a$ is
$g_s^{(d-1)}(a)/(d-1)!$ with $g_s=(z-a)^{d}f_s$, so by
\eqref{d:eq:clusterres}
\[
 \ResH{a}\bigl(F\dd z\bigr)
  =\sum_{s\in S}\frac{g_s^{(d-1)}(a)}{(d-1)!}\tm^{s}
  =\frac{G^{(d-1)}(a)}{(d-1)!},
\]
the last step because differentiation of a coherent section and
evaluation at an ordinary point are both coefficientwise (Section~\ref{c:p4}).
\end{proof}

\begin{remark}[Exactly what the hypothesis excludes, and what is not claimed]
\label{d:rem:bridge2scope}
The uniform bound is what excludes the displaced-pole phenomenon.  In
the standing Example~\ref{d:ex:movingpole} the coefficient of
$\epsilon^{n}$ has a pole of order $n+1$ at $0$, so no single finite $d$
works, and indeed $\ResH{0}=1$ while $\Resat{0}=0$.  Conversely, under
the hypothesis of Theorem~\ref{d:thm:bridge2} the whole array is one
meromorphic object at the exact point and the two residues coincide.
What is \emph{not} asserted is that uniform boundedness is necessary for
numerical equality in every individual instance: the theorem is a
sufficient criterion, and the example shows it cannot simply be dropped.
Note also that R2 and R3 may both be defined and still differ in
location: Theorem~\ref{d:thm:bridge1} equates $\ResH{a}$ with a sum over
the whole cluster $a+\mm$, while \eqref{d:eq:bridge2} equates it with the
single value at the centre. The uniform pole-order hypothesis makes
all possible poles lie at that centre and thus ensures agreement;
numerical equality by cancellation can also occur without that hypothesis.
\end{remark}

\begin{remark}[Two meanings of a Laurent expansion]
\label{d:rem:twolaurents}
The same distinction appears one level up.  For $F\in\Hc(\mathcal A)$ on
an ordinary annulus $\mathcal A=\{r<\abs{z-a}<R\}$, each coefficient has
its ordinary Laurent expansion
$f_s(z)=\sum_{n\in\ZZ}c_{s,n}(z-a)^{n}$, and the resulting object is a
\emph{coefficientwise} Laurent expansion: the inner series is summed in
the ordinary complex sense first and the results are then lifted.  There
is no general justification for reading the whole array $(s,n)$ as one
Hahn sum at a nonzero ordinary argument, since infinitely many ordinary
scalar terms would land at the single exponent $s$.  By contrast a
\emph{meromorphic germ} at an arbitrary point $b\in\K$ is an element of
$\K((z-b))$ with finitely many negative powers.  R2 belongs to the first
picture and R1/R3 to the second.
\end{remark}

\subsection{The exact argument principle}
\label{d:sub:argument}

\begin{theorem}[Exact weighted argument principle]
\label{d:thm:weightedarg}
Let $F\in\Hc(U)$ be nonzero with leading coefficient function $f_0$
having no zero on $\Gamma$, and let $\xi_1,\dots,\xi_d\in\halo\Omega$ be
the zeros of $\ev F$ repeated with multiplicity, as supplied by
Theorem~\ref{d:thm:zerocluster}.  Then for every weight $A\in\Hc(U)$,
\begin{equation}\label{d:eq:weightedarg}
 \frac1{2\pi i}\ointH_{\Gamma}A(w)\,\frac{F'(w)}{F(w)}\,\dd w
  =\sum_{j=1}^{d}\ev A(\xi_j).
\end{equation}
In particular, with $A=1$,
\begin{equation}\label{d:eq:argument}
 \frac1{2\pi i}\ointH_{\Gamma}\frac{F'(w)}{F(w)}\,\dd w
  =\sum_{z\in\halo\Omega}\ord_z\ev F
  =\sum_{a\in\Omega}\ord_a f_0=d\in\NN,
\end{equation}
an ordinary integer with \emph{no} infinitesimal error term, and
locally $\ResH{a}\bigl((F'/F)\dd\zeta\bigr)=\ord_a f_0$.
\end{theorem}

\begin{proof}
Global preparation (Theorem~\ref{d:thm:preparation}) writes
$F=\mathcal U W$ on $V$ with $\mathcal U=\tm^{\beta}u_0(1+q)$ a unit of
$\Hc(V)$ and $W$ the single monic cluster polynomial of degree $d$.
Then
\[
 \frac{F'}{F}=\frac{\mathcal U'}{\mathcal U}+\frac{W'}{W}.
\]
The first summand lies in $\Hc(V)$, so $A\,\mathcal U'/\mathcal U$ is
coherent holomorphic and its coefficientwise integral around the closed
contour $\Gamma$ vanishes by Cauchy's theorem (Section~\ref{c:p5}).  Algebraically
$W'/W=\sum_{j=1}^{d}(w-\xi_j)^{-1}$ in $\K(w)$.  Cauchy's integral
formula at the genuinely surcomplex point $\xi_j\in\halo\Omega$
(Section~\ref{c:p5}) evaluates
$\frac1{2\pi i}\ointH_\Gamma A(w)(w-\xi_j)^{-1}\dd w=\ev A(\xi_j)$, the
kernel being handled there by its uniform Hahn expansion near the
ordinary contour.  Adding the finitely many terms gives
\eqref{d:eq:weightedarg}.  The identification of $d$ with both zero
counts in \eqref{d:eq:argument} is
Theorem~\ref{d:thm:zerocluster}.
\end{proof}

\begin{proof}[Second proof of the unweighted identity \eqref{d:eq:argument}]
Near $\Gamma$ write $F=\tm^{\beta}f_0(1+E)$ with $\supp E>0$, which is
legitimate in $\Mc$ because $f_0$ has no zero there.  Then
\[
 \frac{F'}{F}=\frac{f_0'}{f_0}+\bigl(\log(1+E)\bigr)',
\]
where $\log(1+E)=\sum_{n\ge1}(-1)^{n-1}E^{n}/n$ is a Hahn sum by
Neumann's lemma and is a \emph{single-valued} coherent section along the
whole contour.  Each coefficient of the second summand is the derivative
of a function holomorphic near $\Gamma$, so it has zero residue
everywhere and zero integral around the closed contour.  The remaining
term contributes the ordinary winding integer
$\sum_{a\in\Omega}\ord_af_0$ by the classical argument principle.  This
proof isolates where the integrality comes from; preparation is what
identifies that integer with the actual surcomplex roots, split
multiplicities included.
\end{proof}

\begin{corollary}[Zeros minus poles]
\label{d:cor:zerosminuspoles}
Let $A,B\in\Hc(U)$ be nonzero with both leading coefficient functions
nonzero on $\Gamma$, and let $M=A/B$.  Then
\[
 \frac1{2\pi i}\ointH_{\Gamma}\frac{M'(w)}{M(w)}\,\dd w
  =\sum_{b\in\halo\Omega}\ord_b\ev M
  =N_A(\halo\Omega)-N_B(\halo\Omega)\in\ZZ,
\]
where $N$ counts zeros with multiplicity.  Only finitely many terms in
the middle are nonzero, and cancellations are counted with their signs.
\end{corollary}

\begin{proof}
$M'/M=A'/A-B'/B$ and the integral is $\K$-linear; apply
\eqref{d:eq:argument} to each.  Equivalently, $M'/M$ is a coherent
quotient whose only poles in $\halo\Omega$ are at the zeros of $A$ and
$B$, with actual local residue $\ord_b\ev M$ there by the local
argument principle of \S\ref{d:sub:R1}; then apply
Theorem~\ref{d:thm:residuequotient}.
\end{proof}

\begin{remark}
The normalized integral is an exact ordinary integer even though the
zeros, the derivatives and the residues used to compute it live in
$\K$.  It is an analytic index for the specified contour pairing.  It is
\emph{not} identified with the winding number of a continuous loop in
the fine topology: every continuous map $[0,1]\to\K$ is constant
(Sections~\ref{a:sec:summability}--\ref{a:sec:fine}), so no such nonconstant loop exists.
\end{remark}

\subsection{Deformation of the contour}
\label{d:sub:deformation}

\begin{theorem}[Deformation invariance under admissible infinitesimal deformations]
\label{d:thm:deformation}
Let $\gamma_0$ be a closed ordinary piecewise-$C^1$ path in $U$ and let
\[
 \widetilde\gamma(\tau)=\gamma_0(\tau)+\eta(\tau),
 \qquad \eta(\tau)=\sum_{\lambda\in S}\eta_\lambda(\tau)\tm^{\lambda},
 \qquad S\subset\No_{>0}\ \text{a well-ordered set},
\]
with every $\eta_\lambda$ piecewise $C^1$ and with matching endpoints,
the integral over $\widetilde\gamma$ being defined by Taylor
substitution in the integrand followed by coefficientwise integration of
$R(\widetilde\gamma(\tau))\widetilde\gamma'(\tau)$.  Let $R=A/B$ be a
quotient of coherent families whose denominator's leading coefficient
function has no zero on $\gamma_0$.  Then
\[
 \ointH_{\widetilde\gamma}R(z)\,\dd z=\ointH_{\gamma_0}R(z)\,\dd z .
\]
\end{theorem}

\begin{proof}
Use the ordinary real homotopy $\gamma_u=\gamma_0+u\eta$, $0\le u\le1$.
The support calculus makes every substitution admissible: $\eta$ has
positive well-ordered support, so the Taylor expansion of
$R(\gamma_u)$ about $\gamma_0$ is Hahn summable, and at each exponent
the expression is a \emph{finite polynomial in $u$} with piecewise
continuous coefficient functions of $\tau$.  The chain rule gives
\[
 \frac{\partial}{\partial u}\bigl(R(\gamma_u)\gamma_u'\bigr)
  =\frac{\dd}{\dd \tau}\bigl(R(\gamma_u)\,\eta\bigr),
\]
and integrating coefficientwise in $\tau$ kills the right-hand side,
the endpoint terms cancelling because the parametrizations are closed.
Hence each coefficient is independent of $u$; evaluate at $u=0$ and
$u=1$.
\end{proof}

This is invariance under a specified algebraic-analytic deformation of a
classical contour.  It is \emph{not} an unrestricted homotopy theorem
for maps into the fine-topological class field, and none is proved
anywhere in this article.

\subsection{The global rational residue theorem on
\texorpdfstring{$\mathbb P^{1}(\K)$}{P1(K)}}
\label{d:sub:rationalresidue}

The last theorem of this part uses no contour at all.  It is a theorem
about rational functions over $\K$, hence about the formal/actual
residue R1--R3 and not about $\ResH{}$, and it holds over the entire
proper-class field.

\begin{definition}
For $R\in\K(z)$ put
$\Resat{\infty}\bigl(R(z)\dd z\bigr)
 =\Res_{w=0}\bigl(-w^{-2}R(1/w)\bigr)\dd w$, using $w=1/z$.
\end{definition}

\begin{theorem}[Global rational residue theorem]
\label{d:thm:rationalresidue}
A rational differential $R(z)\dd z$ over $\K$ has finitely many poles in
$\mathbb P^{1}(\K)=\K\cup\{\infty\}$ and
\begin{equation}\label{d:eq:rationalresidue}
 \sum_{b\in\mathbb P^{1}(\K)}\Resat{b}\bigl(R(z)\dd z\bigr)=0 .
\end{equation}
Moreover $R(z)\dd z$ has a rational primitive if and only if all its
affine residues vanish.
\end{theorem}

\begin{proof}
Algebraic closedness of $\K$ gives the finite partial-fraction
decomposition
\[
 R(z)=P(z)+\sum_{b}\sum_{j=1}^{m_b}\frac{c_{b,j}}{(z-b)^{j}},
 \qquad P\in\K[z],
\]
with finitely many $b$.  The affine residues are the $c_{b,1}$.  In the
expansion at infinity the coefficient of $z^{-1}$ is $\sum_bc_{b,1}$, so
the residue at infinity is its negative, proving
\eqref{d:eq:rationalresidue}: a simple-pole term contributes its
coefficient at its pole and the negative at infinity, while polynomial
terms and terms with $j\ge2$ contribute zero in total.  For the second
assertion, every summand with $j\ge2$ and every polynomial term has a
rational primitive (characteristic zero), while a derivative of a
rational function has zero residue at every point, so the $j=1$ terms
are the whole obstruction.
\end{proof}

No contour, no completeness, no topological connectedness, and no
compactness enters the proof.  In the same spirit,
Definition~\ref{d:def:formalres} gives the local statement behind it:
modulo derivatives, $\K((X))\dd X/\dd\K((X))\cong\K\,\dd X/X$.

\subsection{What this part does not assert}
\label{d:sub:residuescope}

\begin{enumerate}[label=(\alph*)]
\item $\Res$, $\ResH{a}$ and $\Resat{b}$ are three symbols for three
notions, and at $1/(\zeta-\epsilon)$ the first and second take the
values $0$ and $1$ at the same ordinary point.  Nothing in this article
licenses writing one symbol for two of them.
\item Theorem~\ref{d:thm:Hresidue} produces a Hahn-summable sum of
cluster residues.  It does not produce actual surcomplex poles; on
$\Mc(U)$ there need be none, and the finite-sum conclusion of
Theorem~\ref{d:thm:bridge1} is licensed only for quotients of coherent
families.
\item Theorem~\ref{d:thm:bridge2} is a sufficient criterion.  Uniform
boundedness of the pole orders is not claimed to be necessary for
numerical coincidence in an individual case.
\item The expansions \eqref{d:eq:movingpole} and
\eqref{d:eq:kernelexp} are coefficientwise identities valid away from
the monad containing the relevant pole.  They are not asserted to
evaluate inside it.
\item Theorem~\ref{d:thm:deformation} moves a classical contour by an
admissible infinitesimal deformation.  No unrestricted homotopy
invariance is claimed, and the coefficientwise contour functional is not
offered as the unique possible notion of surcomplex integration.
\end{enumerate}

%% ==================== p89.tex ====================

% =====================================================================
% PART 8.  Global principles on the halo
% Lead sources: 09, 07, 03, 08, 06, 02, 01, 05.
% CLASS DISCIPLINE.  Every theorem in this part names its class.
%   * coherent, i.e. H(U): normalization, open mapping, maximum modulus,
%     reversion, the whole Liouville hierarchy, removability, halo
%     conformal equivalence, Montel, harmonic/Poisson/reflection.
%   * canonical lift f^#: the SHARP Schwarz and Schwarz-Pick inequalities
%     ONLY.  They are false in H(U); the witness (1+t)z is kept.
%   * germ-analytic: used only as the ambient local language (the
%     ramification normal form of Part 3), never as a hypothesis under
%     which a global principle of this part is asserted.
% =====================================================================

\section{Global principles on the halo}\label{e:p8}

This part collects the classical global theorems that survive on a halo
$\halo U$, each with the exact hypothesis that makes it true and, where
one exists, with the counterexample that shows the hypothesis cannot be
weakened.  Every statement below is a statement about the \emph{coherent}
class $\Hc(U)$ of Section~\ref{c:p4}, with the sole exception of
\S\ref{e:sec-schwarz}, whose sharp inequalities hold only for canonical
classical lifts $\ev f$ and are \emph{false} for general coherent
self-maps of the halo.  Nothing in this part is asserted for the
germ-analytic class of Section~\ref{b:part3}: the identity theorem, the global maximum
principle, Liouville, uniqueness of primitives and the open mapping
theorem all fail there, as the locally constant witnesses of Section~\ref{b:part2} show.

\subsection{The normalization lemma and local finite degree}
\label{e:sec-openmapping}

The engine of the local mapping theorems is a rescaling whose radius is an
explicitly computed \emph{monomial}.  Its motivating case is a point $b$
at which $\ev F{}'(b)$ is nonzero but itself infinitesimal; such a point is
invisible to any argument that looks only at the reduction to $\CC$.

\begin{lemma}[Infinitesimal rescaling; monomial radius]\label{e:lem-normalize}
Let $U\subseteq\CC$ be a connected ordinary domain, let $F\in\Hc(U)$ be
nonconstant, and let $b\in\halo U$.  There is a least integer $m\ge1$ with
\[
 a_m:=\frac{(D_z^mF)^{\#}(b)}{m!}\ne0 .
\]
Write $b=a+\varepsilon$ with $a=\st(b)\in U$ and $\varepsilon\in\mm$, and
let $\sigma$ be any lower bound for the exponents occurring in the common
support
$\supp(F)+\Neu{\supp(\varepsilon)}$
of the Taylor coefficients $a_n=(D_z^nF)^{\#}(b)/n!$; one may take
$\sigma=\vH(F)$.  Choose a surreal
\begin{equation}\label{e:eq-rho}
 \rho>\max\bigl\{0,\;v(a_m)-\sigma\bigr\},
 \qquad r:=\tm^{\rho}>0 .
\end{equation}
Then the normalized function
\begin{equation}\label{e:eq-normalize}
 G(w)\;=\;\frac{\ev F(b+rw)-\ev F(b)}{a_m\,r^{m}}
\end{equation}
is the evaluation of a section of $\Hc(\CC)$ of the form $w^m+H(w)$ with
$\vH(H)>0$.  Increasing $\rho$ makes $r$ smaller than any prescribed
positive surreal radius, so $b+r\mm$ may be placed inside any prescribed
fine-open neighbourhood of $b$.
\end{lemma}

\begin{proof}
That some $m$ exists is the identity theorem of Section~\ref{c:p4} applied to
$F-\ev F(b)$: were every positive-order derivative to vanish at $b$, that
section would vanish identically.  Concretely, its first nonzero leading
coefficient function has finite vanishing order at $a$, so some
positive-order derivative is nonzero at $b$.

The exact Taylor expansion of Section~\ref{c:p4}, valid throughout the monad and here
applied with the infinitesimal increment $rw$ for finite $w$, gives
\[
 G(w)\;=\;w^m+\sum_{n>m}\frac{a_n}{a_m}\,r^{\,n-m}w^n .
\]
The support of the tail lies in
$\bigl(\supp(F)+\Neu{\supp(\varepsilon)}+\supp(a_m^{-1})\bigr)
 +\{\rho,2\rho,3\rho,\dots\}$.
The reciprocal support must be included in full; its minimum is
$-v(a_m)$, but it need not be a singleton. The displayed set is a sum
of a fixed well-ordered set with the positive multiples of $\rho$, so it is well ordered, and each exponent
receives a contribution from only finitely many $n$; hence the coefficient
of each exponent is an ordinary \emph{polynomial} in $w$, in particular an
entire function of $w$.  By \eqref{e:eq-rho} the least exponent occurring
in the tail is at least $\sigma-v(a_m)+\rho>0$, so $\vH(H)>0$.  The last
assertion is immediate, since $\rho\mapsto\tm^{\rho}$ is strictly
decreasing.
\end{proof}

The choice of a \emph{monomial} radius rather than a real one is what makes
the lemma usable at arbitrary surcomplex points: $\rho$ is allowed to
absorb the valuation $v(a_m)$ of an infinitesimal derivative, which no
real radius can do.

\begin{theorem}[Local finite degree and the open mapping theorem]
\label{e:thm-openmapping}
With the notation of Lemma~\ref{e:lem-normalize},
\begin{equation}\label{e:eq-localdegree}
 \ev F\bigl(b+r\mm\bigr)\;=\;\ev F(b)+a_mr^m\,\mm ,
\end{equation}
and every value on the right has exactly $m$ preimages in $b+r\mm$,
counted with multiplicity.  Consequently every nonconstant $F\in\Hc(U)$
evaluates to a map $\ev F:\halo U\to\K$ that carries fine-open subclasses
to fine-open subclasses.
\end{theorem}

\begin{proof}
For $\eta\in\mm$ the section $G-\eta$ of $\Hc(\CC)$ has leading
coefficient function $w^m$, whose only zero in $\CC$ is a zero of order
$m$ at the origin.  The zero-cluster theorem of Section~\ref{d:sec:preparation} therefore produces
exactly $m$ solutions of $G(w)=\eta$ in $\mm$, counted with multiplicity,
and no solutions in any other monad.  Conversely $G(w)\in\mm$ for $w\in\mm$,
because $G-w^m$ has positive support and $w^m\in\mm$.  Undoing the two
affine substitutions of \eqref{e:eq-normalize} gives
\eqref{e:eq-localdegree} and the count.  For openness, given a fine-open
neighbourhood of $b$, choose $\rho$ by the last clause of
Lemma~\ref{e:lem-normalize} so that $b+r\mm$ lies inside it; its image
contains the fine-open neighbourhood $\ev F(b)+a_mr^m\mm$ of $\ev F(b)$.
\end{proof}

This is an assertion about actual open classes and actual surcomplex
preimages, not merely about the reduction to $\CC$.

\begin{corollary}[Strong local maximum modulus]\label{e:cor-maxmod}
Let $U$ be connected and $F\in\Hc(U)$ nonconstant.  Then $\abs{\ev F}$ has
no local maximum at any point of $\halo U$ in the fine topology.
Equivalently: if $\abs{\ev F}$ has a fine-local maximum at some
$b\in\halo U$, then $F$ is a constant element of $\K$ and $\ev F$ is that
constant on all of $\halo U$.
\end{corollary}

\begin{proof}
By Theorem~\ref{e:thm-openmapping} every fine-neighbourhood of $\ev F(b)$
contains values of $\ev F$; and every fine-neighbourhood of a point of
$\K$ contains points of strictly larger modulus --- take $(1+\eta)c$ with
$\eta>0$ a sufficiently small surreal if $c=\ev F(b)\ne0$, and any
sufficiently small nonzero value if $c=0$.  This contradicts a local
maximum.  For the second formulation, a local maximum forces the germ of
$\ev F$ at $b$ to be constant; subtracting that constant, which is itself
an element of $\Hc(U)$ with constant coefficient functions, and applying the
identity theorem of Section~\ref{c:p4} in its vanishing-germ form, gives $F$ constant
on the whole halo over the connected ordinary base.
\end{proof}

\begin{remark}[Exactly what the maximum principle does and does not say]
\label{e:rem-maxscope}
Two scope statements must travel with
Corollary~\ref{e:cor-maxmod}.  First, this is a genuine
\emph{neighbourhood} maximum principle in the fine topology, not merely a
statement about standard parts; coherence enters only in the last step,
the passage from a constant germ to one global constant.  Second, nothing
here asserts that $\abs{\ev F}$ \emph{attains} a maximum on a closed
surcomplex disk.  No closed disk in $\K$ is compact, and it is precisely
this missing compactness that prevents the classical route from the
maximum principle to a boundary-maximum or Liouville statement.  The
Liouville question is settled instead by \S\ref{e:sec-liouville}, and by
hypotheses of an entirely different shape.
\end{remark}

\subsection{Hahn reversion and the coherent local inverse}
\label{e:sec-reversion}

\begin{lemma}[Hahn reversion]\label{e:lem-reversion}
Let $G(w)=w+H(w)$ with $H\in\Hc_{>0}(\CC)$.  There is a unique
$K\in\Hc_{>0}(\CC)$ such that $J(y)=y+K(y)$ is a two-sided inverse of $G$
under infinitesimal Hahn composition, and $\ev J$ is a bijection
$\Ov\to\Ov$ inverse to $\ev G$.  Explicitly,
\begin{equation}\label{e:eq-reversion}
 K(y)\;=\;\sum_{n\ge1}\frac{(-1)^n}{n!}\,
   D_y^{\,n-1}\bigl(H(y)^n\bigr),
\end{equation}
a Hahn sum whose $n$-th term has support inside $n\cdot\supp(H)$.
\end{lemma}

\begin{proof}
Expanding $H(y+K)$ as the formal Taylor series in the positive-support
section $K$, the equation $K=-H(y+K)$ determines the part of $K$ of degree
$n$ in occurrences of $H$ from the parts of smaller degree.  Derivatives
of entire coefficient functions are entire, so every term lies in
$\Hc_{>0}(\CC)$, and the support of a degree-$n$ term is a sum of $n$
elements of $\supp(H)$.  Neumann~\cite{Neumann}'s lemma in its finite-word form (Sections~\ref{a:sec:summability}--\ref{a:sec:fine})
makes the family summable, since each exponent of
$\Neu{\supp(H)}$ admits only finitely many representations as a word.
Uniqueness follows by comparing least exponents: altering $K$ on the right
multiplies the difference by a section of strictly positive valuation.
For \eqref{e:eq-reversion}, introduce an auxiliary ordinary formal
parameter $q$ and solve $u=-qH(y+u)$; formal Lagrange inversion gives
$[q^n]u=\frac{(-1)^n}{n!}D_y^{n-1}H(y)^n$, and summability licenses the
graded evaluation at $q=1$ --- a normal sum, not the convergence of
partial sums.  Finally, for a finite target $y=c+\eta$ the zero-cluster
theorem applied to $G-y$ gives exactly one solution in $c+\mm$ and none
elsewhere, so $\ev G$ is a bijection of $\Ov$; faithful evaluation
identifies its inverse with $\ev J$.
\end{proof}

\begin{theorem}[Coherent local inverse]\label{e:thm-localinverse}
Let $F\in\Hc(U)$ and $b\in\halo U$ with $\ev F{}'(b)\ne0$.  Then $\ev F$
has a local inverse $\Phi$ between fine-open neighbourhoods of $b$ and of
$\ev F(b)$; after the affine rescaling of Lemma~\ref{e:lem-normalize} with
$m=1$ the inverse is itself coherent, and
\[
 \Phi'\bigl(\ev F(z)\bigr)=\frac1{\ev F{}'(z)}
\]
on a sufficiently small such neighbourhood.
\end{theorem}

\begin{proof}
Apply Lemma~\ref{e:lem-normalize} with $m=1$, invert by
Lemma~\ref{e:lem-reversion}, and undo the affine normalizations; the
derivative formula is the chain rule.  Shrinking once more makes
$\ev F{}'$ nonvanishing throughout, by the normalized leading-derivative
estimate.  The hypothesis $\ev F{}'(b)\ne0$ is the right one: it does
\emph{not} require $\ev F{}'(b)$ to have nonzero standard part, and the
monomial radius of \eqref{e:eq-rho} is exactly what accommodates an
infinitesimal derivative.
\end{proof}

\subsection{The Liouville hierarchy: three inequivalent hypotheses}
\label{e:sec-liouville}

Call $F\in\Hc(\CC)$ \emph{macroscopically entire}; its evaluation lives on
$\halo{\CC}=\Ov$, the finite elements, not on all of $\K$.  Three
different boundedness hypotheses are in circulation for such $F$, and they
give three genuinely different theorems.  They are stated here together,
with their separate conclusions. The ordinary-real and coefficientwise
bounds are incomparable; neither should be substituted for the other.

\paragraph{(a) A single surreal bound is vacuous.}

\begin{proposition}[Automatic surreal boundedness]\label{e:prop-autobounded}
Every nonzero $F\in\Hc(\CC)$ satisfies
\begin{equation}\label{e:eq-automaticbound}
 \abs{\ev F(z)}\;<\;\tm^{\,\vH(F)-1}
 \qquad\bigl(z\in\halo{\CC}=\Ov\bigr),
\end{equation}
with a bound independent of $z$.  There are even nonconstant
$F\in\Hc(\CC)$ bounded on $\Ov$ by the ordinary constant $1$.
\end{proposition}

\begin{proof}
Evaluation at a finite point adds only positive exponents to the
coefficient supports, so $v(\ev F(z))\ge\vH(F)$ whenever the value is
nonzero.  Since $v(\abs w)=v(w)$ and $\tm^{\gamma-1}$ has valuation
$\gamma-1<\gamma$, any $w$ with $v(w)\ge\gamma$ satisfies
$\abs w<\tm^{\gamma-1}$; take $\gamma=\vH(F)$.  For the second assertion
take $F(Z)=\tm Z$: its value at every finite point is infinitesimal, hence
of modulus less than $1$, yet $F'=\tm\ne0$.
\end{proof}

Thus ``bounded by some surreal constant on the finite plane'' is a
property of \emph{every} macroscopically entire function and cannot be the
hypothesis of any Liouville theorem.  Note also why the order-valued
Cauchy estimate of Section~\ref{c:p5} does not rescue the argument: for $F=\tm Z$ and
the real bound $M=1$ on the ordinary circle $\abs z=R$, that estimate
gives $\tm\le1/R$ for every ordinary real $R>0$, which is true for every
infinitesimal $\tm$.  The classical step is an archimedean limit in $R$,
and there is no such limit through radii cofinal in $\No$.

\paragraph{(b) A single ordinary real bound: an exact characterization,
not constancy.}

\begin{theorem}[Real-bounded macroscopically entire functions]
\label{e:thm-realliouville}
For $F\in\Hc(\CC)$ the following are equivalent.
\begin{enumerate}[label=\textup{(\roman*)}]
\item There is an \emph{ordinary real} $M>0$ with
      $\abs{\ev F(z)}\le M$ for every $z\in\Ov$.
\item There are $c\in\CC$ and $R\in\Hc_{>0}(\CC)$ with $F=c+R$.
\end{enumerate}
Hence the algebra of real-bounded macroscopically entire functions is
exactly
\begin{equation}\label{e:eq-liouvillealgebra}
 \mathcal B\;=\;\CC\oplus\Hc_{>0}(\CC),
 \qquad \mathcal B/\Hc_{>0}(\CC)\;\cong\;\CC .
\end{equation}
More generally, a bound $\abs{\ev F(z)}\le M\tm^{\beta}$ with $M$ ordinary
real holds if and only if $F=\tm^{\beta}(c+R)$ with $c\in\CC$ and $R$ of
positive support.
\end{theorem}

\begin{proof}
Assume (i).  If $\vH(F)<0$, pick an ordinary point $c$ at which the
leading coefficient function $f_{\vH(F)}$ is nonzero; the value
$\ev F(c)$ then has negative valuation and is infinite, contradicting a
real bound.  So $F\in\Hc_{\ge0}(\CC)$.  Let $f_0$ be its exponent-zero
coefficient function (possibly zero).  Taking standard parts of the bound
at ordinary points gives $\abs{f_0(c)}\le M$ for all $c\in\CC$, so the
ordinary Liouville theorem makes $f_0$ a constant $c\in\CC$.  All remaining
support is strictly positive, which is (ii).

Conversely, if $F=c+R$ with $\supp(R)>0$ then $\ev R(z)\in\mm$ for every
finite $z$, so $\abs{\ev F(z)}<\abs c+1$: a single ordinary real bound,
despite the fact that the individual coefficient functions of $R$ may grow
arbitrarily fast.  The quotient statement is immediate, and multiplying by
$\tm^{-\beta}$ gives the graded version.
\end{proof}

This is the only sharp statement in the collection of what a scalar bound
buys, and it is what reconciles the other two hypotheses:
$\tm Z=0+\tm Z$ lies in $\CC\oplus\Hc_{>0}(\CC)$, exactly as
\eqref{e:eq-liouvillealgebra} predicts.

\paragraph{(c) Coefficientwise growth: the classical hypothesis.}

\begin{theorem}[Coefficientwise growth and polynomial degree]
\label{e:thm-coeffliouville}
Let $d\ge0$ be an ordinary integer and $F=\sum_{\gamma\in S}f_\gamma\tm^\gamma
\in\Hc(\CC)$.  The following are equivalent.
\begin{enumerate}[label=\textup{(\roman*)}]
\item For every $\gamma\in S$ there is an ordinary real $M_\gamma>0$ with
      $\abs{f_\gamma(c)}\le M_\gamma(1+\abs c)^d$ for all $c\in\CC$.
\item There are $A_0,\dots,A_d\in\K$ with
      $\ev F(z)=\sum_{j=0}^{d}A_jz^{\,j}$ for all $z\in\Ov$.
\end{enumerate}
The case $d=0$ is the coefficientwise Liouville theorem: if every
$f_\gamma$ is bounded on $\CC$ by its own real constant, then $F$ is a
constant element of $\K$.
\end{theorem}

\begin{proof}
Assume (i) and apply the ordinary Cauchy estimates to $f_\gamma$ on the
circle of ordinary radius $R$ about $0$: for $n>d$,
\[
 \abs{f_\gamma^{(n)}(0)}\;\le\;n!\,M_\gamma\frac{(1+R)^d}{R^{\,n}} .
\]
Letting the \emph{ordinary real} $R$ tend to infinity --- a limit inside
$\CC$, taken separately for each $\gamma$, not a limit through surreal
radii --- makes this derivative vanish.  So each $f_\gamma$ is a
polynomial of degree at most $d$.  Put
$A_j=\sum_{\gamma\in S}\frac{f_\gamma^{(j)}(0)}{j!}\tm^{\gamma}$; these
sums exist with support inside $S$, and regrouping the finitely many
powers of $z$ gives (ii).  Conversely, expand the finitely many $A_j$ into
normal form: each ordinary coefficient function is a polynomial of degree
at most $d$ and obeys a bound of the required shape, and injectivity of
coherent evaluation identifies those coefficients with the original
$f_\gamma$.
\end{proof}

\begin{remark}[The hierarchy, and the two counterexamples]
\label{e:rem-liouvillehierarchy}
For the boundedness case $d=0$, hypotheses (b) and (c) both imply (a),
but are incomparable. The constant section $F=\omega$ has bounded
ordinary coefficient functions and no ordinary-real bound on its values,
so (c) does not imply (b). Conversely $F=\tm Z$ satisfies (b) and
violates (c). The earlier assertion of a linear hierarchy
(c)$\Rightarrow$(b)$\Rightarrow$(a) was therefore false.  Importing the conclusion of
Theorem~\ref{e:thm-coeffliouville} under hypothesis (b) is false:
$F=\tm Z$ is real-bounded by $1$ on the whole finite plane and is not
constant.  Importing it under hypothesis (a) is worse, since (a) holds for
every macroscopically entire $F$ by
Proposition~\ref{e:prop-autobounded}.  The two standing counterexamples
are recorded once here: $\tm Z$, and $\varepsilon Z$ for a positive real
infinitesimal $\varepsilon$, whose nonzero coefficient functions are
nonzero scalar multiples of $Z$ and are therefore unbounded, so that (c)
is genuinely violated.

These are a \emph{different} obstruction from the locally constant
function $L$ of Section~\ref{b:part2}.  The functions $\tm Z$ and $\varepsilon Z$ are
coherent, and their domain is $\halo{\CC}$, the halo of the ordinary
plane; $L$ is defined on all of $\K$ and is germ-analytic. Its restriction
to every finite halo is the coherent constant $1$; it is its global
behavior on $\K$ that gives the different counterexample.  Different analytic class, different domain; the
two must not be presented as instances of one phenomenon.
\end{remark}

\subsection{Removable singularities}\label{e:sec-removable}

\begin{theorem}[Coefficientwise removal is necessary and sufficient]
\label{e:thm-removable}
Let $U\subseteq\CC$ be a domain, $c\in U$, and
$F=\sum_{\gamma\in S}f_\gamma\tm^\gamma\in\Hc(U\setminus\{c\})$.  Then $F$
extends to an element of $\Hc(U)$ if and only if every $f_\gamma$ has a
removable singularity at $c$.  In particular the extension exists whenever
each $f_\gamma$ is bounded on some ordinary punctured neighbourhood of
$c$, with a bound allowed to depend on $\gamma$.  The extension is unique.
\end{theorem}

\begin{proof}
Necessity: any coherent extension restricts coefficientwise to the given
functions, by uniqueness of normal forms at ordinary points, so each
$f_\gamma$ extends holomorphically.  Sufficiency: the ordinary removable
singularity theorem extends each $f_\gamma$; the extended family has the
same well-ordered support $S$, hence is an element of $\Hc(U)$.
Uniqueness is the ordinary identity theorem applied coefficient by
coefficient.
\end{proof}

\begin{remark}[A scalar bound is not enough, and the puncture is a monad]
\label{e:rem-punctured}
A bound on the surcomplex value $\ev F$ does not suffice.  On the halo of
the ordinary punctured disk $0<\abs Z<1$ the family $\tm/Z$ is
infinitesimal at every point, hence bounded by the ordinary constant $1$;
but its coefficient at exponent $1$ is $1/Z$, which has a pole, so by
Theorem~\ref{e:thm-removable} it has no coherent extension across the
origin.  The family $\tm\,e^{1/Z}$ is bounded by $1$ in the same way and
its coefficient has an essential singularity.  This is the same
phenomenon as in Remark~\ref{e:rem-liouvillehierarchy}: a bound on the
leading scale of the value says nothing at all about the smaller scales.

Note also which domain is in play.  The halo
$\halo{(\mathbb D\setminus\{0\})}$ excludes the \emph{whole monad} of the
origin; it is not the class obtained by deleting the single point $0$ from
$\halo{\mathbb D}$.  Keeping these two puncturing operations distinct
prevents a false analogy with the ordinary punctured disk.
\end{remark}

\subsection{Schwarz--Pick for canonical classical lifts}\label{e:sec-schwarz}

This subsection, alone in Section~\ref{e:p8}, is about the class of canonical
classical lifts $\ev f$ of ordinary holomorphic $f$.  The sharp
inequalities below are \emph{false} for general coherent self-maps of
$\halo{\mathbb D}$; see Example~\ref{e:ex-schwarzfail}.  Let
$\mathbb D=\{Z\in\CC:\abs Z<1\}$, let $f:\mathbb D\to\mathbb D$ be
ordinary holomorphic, and put, for $z,w\in\halo{\mathbb D}$,
\[
 d(z,w)=\left\lvert\frac{z-w}{1-\bar wz}\right\rvert .
\]
The denominator is invertible because its standard part $1-\bar{w_0}z_0$
is a nonzero complex number, $z_0,w_0$ being the standard parts.  Since
standard parts commute with evaluation, $\ev f$ maps $\halo{\mathbb D}$
into $\halo{\mathbb D}$.

\begin{theorem}[Schwarz--Pick for canonical lifts]\label{e:thm-schwarzpick}
For all $z,w\in\halo{\mathbb D}$,
\begin{equation}\label{e:eq-schwarzpick}
 d\bigl(\ev f(z),\ev f(w)\bigr)\;\le\;d(z,w),
 \qquad
 \frac{\abs{(\ev f)'(z)}}{1-\abs{\ev f(z)}^{2}}
 \;\le\;\frac1{1-\abs z^{2}} .
\end{equation}
For $z\ne w$, equality in the first inequality holds if and only if $f$ is
an ordinary disk automorphism; equality in the second at any point has the
same characterization.
\end{theorem}

\begin{proof}
Suppose first that $f$ is not an ordinary disk automorphism.  Then the
ordinary Schwarz--Pick inequality is strict at distinct ordinary points,
and its infinitesimal form is strict at every ordinary point.

If $\st z\ne\st w$, the quotient of the two sides of the first inequality
in \eqref{e:eq-schwarzpick} has standard part strictly less than one, so
the surreal inequality is strict.

If $\st z=\st w=c$ and $z\ne w$, this argument is unavailable: both
pseudohyperbolic distances have standard part zero, and taking standard
parts of the classical inequality yields only $0\le0$.  One must divide
out $z-w$ \emph{before} reducing.  Consider the divided difference
\[
 G(z,w)=\frac{\ev f(z)-\ev f(w)}{z-w}.
\]
The coherent Taylor expansion at $c$ gives a strongly summable extension
of $G$ across the diagonal, with $\st G(z,w)=f'(c)$.  Hence
\[
 \st\!\left(\frac{d(\ev f(z),\ev f(w))}{d(z,w)}\right)
 =\abs{f'(c)}\,\frac{1-\abs c^{2}}{1-\abs{f(c)}^{2}}<1 ,
\]
by the strict ordinary infinitesimal Schwarz--Pick inequality at $c$.  So
the strict inequality holds even when $z-w$ is arbitrarily small.  The
derivative inequality follows from the same strict ordinary inequality by
taking standard parts.

If $f$ is an automorphism, write $f(Z)=e^{i\theta}(Z-a)/(1-\bar aZ)$ with
$a\in\mathbb D$ and $\theta$ ordinary real.  The usual algebraic
identities for this expression are polynomial identities valid over the
real closed field $\No$ and its complexification $\K$, so both inequalities
in \eqref{e:eq-schwarzpick} become equalities.  This proves the equality
characterizations.
\end{proof}

\begin{corollary}[Schwarz lemma for canonical lifts]\label{e:cor-schwarz}
If in addition $f(0)=0$ then
\[
 \abs{\ev f(z)}\le\abs z\quad(z\in\halo{\mathbb D}),
 \qquad \abs{(\ev f)'(0)}\le1,
\]
with equality at a nonzero $z$, or in the derivative bound, if and only if
$f(Z)=uZ$ for an ordinary complex $u$ with $\abs u=1$.
\end{corollary}

\begin{proof}
Put $w=0$ in Theorem~\ref{e:thm-schwarzpick}; an automorphism fixing the
origin is an ordinary rotation.
\end{proof}

\begin{corollary}[Stability away from the equality case]
\label{e:cor-spstability}
Let $F=f_0+E\in\Hc(\mathbb D)$ with $E=0$ or $\vH(E)>0$, and suppose
$f_0:\mathbb D\to\mathbb D$ is ordinary holomorphic and is \emph{not} a
disk automorphism.  Then $\ev F$ satisfies both inequalities of
\eqref{e:eq-schwarzpick}, strictly, the first one at distinct points.
\end{corollary}

\begin{proof}
The standard parts of $\ev F$ and of $\ev F{}'$ are $f_0$ and $f_0'$.  At
distinct standard parts use the strict ordinary inequality for $f_0$; in a
single monad use the coherent Taylor expansion, whose divided difference
has standard part $f_0'(c)$.  The proof of
Theorem~\ref{e:thm-schwarzpick} then applies verbatim.
\end{proof}

Corollary~\ref{e:cor-spstability} is a statement about the coherent class,
but it is strictly weaker than
Theorem~\ref{e:thm-schwarzpick}: it asserts only the strict inequalities,
and it requires the ordinary part $f_0$ to be non-extremal.  The sharp
inequality, with its equality case, is not available in $\Hc(\mathbb D)$
at all.

\begin{example}[Why the restriction to lifts is substantive]
\label{e:ex-schwarzfail}
For the positive infinitesimal $\tm=\omega^{-1}$ the coherent function
$F(Z)=(1+\tm)Z$ maps $\halo{\mathbb D}$ bijectively onto itself, since it
does not change a point's standard part, and it fixes the origin.
Nevertheless
\[
 \ev F{}'(0)=1+\tm>1,\qquad \abs{\ev F(z)}>\abs z\quad(z\ne0).
\]
So Schwarz's lemma is false for arbitrary coherent self-maps of the halo.
Two features of the setting are exposed by this example.  First, the
ordinary part of $F$ here \emph{is} an automorphism, so
Corollary~\ref{e:cor-spstability} does not apply --- the equality case of
the ordinary theorem is exactly what is unstable under infinitesimal
perturbation.  Second, $\halo{\mathbb D}$ is not the surcomplex unit ball
$\{z\in\K:\abs z<1\}$: the latter contains points such as $1-\varepsilon$
with $\varepsilon$ a positive infinitesimal, whose standard part lies on
the ordinary unit circle, and no point of that boundary monad belongs to
$\halo{\mathbb D}$.  It is precisely this omitted infinitesimal boundary
layer that leaves room for $(1+\tm)Z$.
\end{example}

\subsection{Conformal equivalence of halos}\label{e:sec-conformal}

\begin{theorem}[Lifting of infinitesimally deformed biholomorphisms]
\label{e:thm-biholo}
Let $f_0:U\to V$ be an ordinary biholomorphism of ordinary complex domains
and let $F=f_0+H\in\Hc(U)$ with $H=0$ or $\vH(H)>0$.  Then there is
$G=g_0+J\in\Hc(V)$ with $g_0=f_0^{-1}$ and $J=0$ or $\vH(J)>0$, such that
\[
 F\circ G=\mathrm{id}_V,\qquad G\circ F=\mathrm{id}_U
\]
as coherent sections.  Consequently $\ev F:\halo U\to\halo V$ is a
bijection with germ-analytic inverse $\ev G$, and $\ev F{}'$ vanishes
nowhere.
\end{theorem}

\begin{proof}
Let $M=\Neu{\supp(H)}$ and seek
$J=\sum_{\gamma\in M,\ \gamma>0}j_\gamma\tm^\gamma$.  In the coefficient
equation at exponent $\gamma$ of $F(g_0+J)=\mathrm{id}$, the only term
containing the unknown $j_\gamma$ is $f_0'(g_0)\,j_\gamma$; every other
term involves strictly smaller exponents of $J$ or known coefficients of
$H$, and by Neumann's lemma each such equation has finitely many terms.
Since $f_0'(g_0)$ is nowhere zero on $V$, the equation determines a single
ordinary holomorphic $j_\gamma$ on $V$.  Well-ordered recursion over
$M_{>0}$ constructs $J$, and the support-controlled composition law of
Section~\ref{c:p4} (inner ordinary holomorphic map plus positive-support perturbation)
justifies the composition identities.

For the evaluated statement: $\st(\ev F(a+\eta))=f_0(a)$, so the image lies
in $\halo V$, and the evaluated identity $\ev F\circ\ev G=\mathrm{id}$
gives surjectivity.  For injectivity take $w=b+\varepsilon\in\halo V$; the
section $F-w$ has leading coefficient function $f_0-b$, which has a unique
and simple zero at $a=g_0(b)$, so the zero-cluster theorem of Section~\ref{d:sec:preparation} gives
exactly one solution of $\ev F(z)=w$ in $a+\mm$ and none in any other
monad.  Faithfulness of evaluation gives the remaining composition
identity, and
$\st(\ev F{}'(a+\eta))=f_0'(a)\ne0$ gives the last assertion.
\end{proof}

\begin{corollary}[Lifted conformal equivalence]\label{e:cor-liftbiholo}
The case $H=0$: if $f:U\to V$ is an ordinary biholomorphism, then
$\ev f:\halo U\to\halo V$ is a bijection, both it and its inverse
$\ev{(f^{-1})}$ are germ-analytic, and $(\ev f)'$ is nowhere zero.
\end{corollary}

\begin{corollary}[Standard-halo Riemann mapping theorem]
\label{e:cor-riemannmap}
If $U\subsetneq\CC$ is nonempty, open and simply connected, then $\halo U$
is analytically equivalent to $\halo{\mathbb D}$ by a canonical classical
lift.
\end{corollary}

\begin{proof}
Choose an ordinary Riemann map and apply
Corollary~\ref{e:cor-liftbiholo}.
\end{proof}

The existence input in Corollary~\ref{e:cor-riemannmap} is precisely the
\emph{ordinary} Riemann mapping theorem.  No uniformization of arbitrary
fine-open subclasses of $\K$ is claimed, and none is proved anywhere in
this article.  Theorem~\ref{e:thm-biholo} is strictly stronger than
Corollary~\ref{e:cor-liftbiholo}, and the extra strength is real: for
example $G(Z)=Z+\tm Z^2$ is a non-identity biholomorphism of
$\halo{\mathbb D}$ fixing the origin with $\ev G{}'(0)=1$, so even the
normalization that would characterize an ordinary disk automorphism leaves
infinitesimal coherent freedom.  There is no conflict with
Corollary~\ref{e:cor-schwarz}, whose hypothesis is a canonical lift.

\subsection{A coefficientwise Montel theorem}\label{e:sec-montel}

\begin{definition}[Coefficientwise compact-open convergence]
\label{e:def-coefconv}
Fix a well-ordered support set $S$ and let
$F_n=\sum_{\gamma\in S}f_{n,\gamma}\tm^\gamma$ and
$F=\sum_{\gamma\in S}f_\gamma\tm^\gamma$ lie in $\Hc(U)$.  Say
$F_n\to F$ \emph{coefficientwise compact-open} if, for each fixed
$\gamma\in S$, $f_{n,\gamma}\to f_\gamma$ uniformly on every ordinary
compact subset of $U$.
\end{definition}

\begin{theorem}[Countable-support Montel theorem]\label{e:thm-montel}
Let $(F_n)_{n\ge1}$ in $\Hc(U)$ have a common \emph{countable}
well-ordered support $S$, and suppose that for each $\gamma\in S$ and each
ordinary compact $K\Subset U$,
\[
 \sup_n\;\sup_{z\in K}\abs{f_{n,\gamma}(z)}<\infty .
\]
Then some subsequence converges coefficientwise compact-open to an element
of $\Hc(U)$, and every fixed coefficientwise derivative $D_z^kF_n$
converges in the same sense along that subsequence.
\end{theorem}

\begin{proof}
Enumerate $S$ as $\gamma_1,\gamma_2,\dots$.  The ordinary Montel theorem
applied to $(f_{n,\gamma_1})_n$ gives a locally uniformly convergent
subsequence; refine for $\gamma_2$, and so on.  The diagonal subsequence
converges for every $\gamma\in S$.  Each limit function is ordinary
holomorphic and the limits are supported in the same well-ordered set $S$,
so they assemble to an element of $\Hc(U)$.  Local uniform convergence of
the ordinary derivatives, obtained from the ordinary Cauchy formula, gives
the statement about $D_z^k$.
\end{proof}

\begin{remark}[What this theorem is a theorem about]\label{e:rem-montelscope}
Definition~\ref{e:def-coefconv} describes a topology on \emph{coefficient
data}.  It is not, and must not be read as, pointwise convergence of
evaluated values in the fine topology: by the discreteness theorem of
Sections~\ref{a:sec:summability}--\ref{a:sec:fine} every set-indexed convergent net in $\K$ is eventually equal to its
limit, so a sequential Montel theorem in the fine topology would be the
wrong statement.  Countability of $S$ is also not cosmetic.  The diagonal
argument extracts one subsequence controlling countably many independent
coefficient sequences; it does not automatically produce a subsequence
controlling uncountably many.  Product-compactness arguments for suitable
coefficient spaces and nets can be substituted, but they too concern the
coefficient topology and not the topology of values.
\end{remark}

\subsection{Harmonic components, Poisson's formula, and reflection}
\label{e:sec-harmonic}

If $F\in\Hc(U)$ and $\ev F=u+iv$ with $z=x+iy$, then the Cauchy--Riemann
equations of Section~\ref{b:part3}, together with the equality of mixed second
derivatives, give
\[
 u_{xx}+u_{yy}=0,\qquad v_{xx}+v_{yy}=0,
\]
where the derivatives are genuine fine partial derivatives, by the bridge
theorem of Section~\ref{c:p4} identifying the coefficientwise derivative with the fine
difference-quotient derivative.  These are differential identities; no
global Dirichlet problem is asserted by them.  Theorem~\ref{e:thm-poisson}
is what fills that gap, and only on ordinary disks with coefficientwise
boundary data.

A \emph{coherent harmonic datum} on an ordinary domain $U$ is a family
$u=\sum_{\gamma\in S}u_\gamma\tm^\gamma$ with common well-ordered support
and ordinary real-valued harmonic coefficient functions $u_\gamma$.  Its
evaluation at $z=c+\varepsilon$ is defined by the two-real-variable Taylor
expansion of each $u_\gamma$ at $c$; harmonic functions are real analytic,
and the support lemma applies to $\operatorname{Re}\varepsilon$ and
$\operatorname{Im}\varepsilon$, so $\ev u$ is well defined and
real-surreal valued.

\begin{theorem}[Coefficientwise harmonic conjugates]\label{e:thm-harmonic}
On a simply connected ordinary domain $U$, every coherent harmonic datum
is the real part of an element of $\Hc(U)$.  A harmonic conjugate is
unique modulo one real surreal constant, $\ev u$ satisfies the Laplace
equation in the fine differentiable sense, and a nonconstant coherent
harmonic datum on a connected ordinary base has no fine-local maximum or
minimum.
\end{theorem}

\begin{proof}
Fix $a\in U$ and let $v_\gamma$ be the ordinary harmonic conjugate of
$u_\gamma$ normalized by $v_\gamma(a)=0$.  The functions $u_\gamma+iv_\gamma$
are ordinary holomorphic, and their Hahn sum with weights $\tm^\gamma$ is
an element of $\Hc(U)$ whose real part evaluates to $\ev u$, by the
ordinary Cauchy--Riemann equations and uniqueness of Taylor expansions.
Two conjugates differ coefficientwise by real constants, hence by one real
surreal constant.  The Laplace equation holds coefficientwise and is
preserved by Taylor evaluation.  For the extremum assertion, work on a
small simply connected ordinary disk about the standard part of a proposed
extremum; the local holomorphic conjugate is nonconstant unless the datum
is constant on that disk, in which case ordinary unique continuation
propagates constancy through the connected base.  In the nonconstant case
Theorem~\ref{e:thm-openmapping} supplies values with both larger and
smaller real part in every fine-neighbourhood, excluding an extremum.
\end{proof}

\begin{theorem}[Poisson's formula and a coherent Dirichlet problem]
\label{e:thm-poisson}
Let the coefficients of a coherent harmonic datum $u$ be harmonic on a
neighbourhood of the closed ordinary disk $\overline{B(c,R)}$, with $c\in\CC$
and $R>0$ ordinary.  Then for every $z\in\halo{B(c,R)}$,
\begin{equation}\label{e:eq-poisson}
 \ev u(z)=\frac1{2\pi}\intH_{[0,2\pi]}
 \frac{R^{2}-\abs{z-c}^{2}}{\abs{c+Re^{i\theta}-z}^{2}}\,
 u\bigl(c+Re^{i\theta}\bigr)\,\dd\theta .
\end{equation}
More generally, a well-ordered Hahn family of ordinary continuous real
boundary data on the ordinary circle has a unique coherent harmonic
extension whose coefficient functions extend continuously to the closed
disk with the prescribed boundary values.
\end{theorem}

\begin{proof}
At an ordinary interior point apply the ordinary Poisson formula to each
coefficient.  For $z=a+\varepsilon$ with $a$ ordinary and $\varepsilon$
infinitesimal, the kernel's denominator has a positive ordinary lower
bound on the circle before perturbation; expand its inverse in the
infinitesimal coordinates of $\varepsilon$.  The resulting support lies in
the sum of the original support with a positive-support Neumann monoid, so
the family is Hahn summable and each coefficient is an ordinary continuous
function of $\theta$, integrable in the coefficientwise sense of Section~\ref{c:p5}.
Ordinary differentiation under the compact integral identifies each Taylor
coefficient of the right-hand side with the corresponding Taylor
coefficient of $u$ at $a$, proving \eqref{e:eq-poisson}.  For general
continuous boundary coefficients take their ordinary Poisson extensions
and reassemble with the same support; uniqueness is the uniqueness of the
ordinary Dirichlet problem, coefficient by coefficient.
\end{proof}

The boundary condition in Theorem~\ref{e:thm-poisson} is an ordinary
coefficientwise trace.  It is not a claim about compactness, nor about
convergence to the boundary in the fine topology, and it says nothing
about boundary data that are not presented as a common-support Hahn family
of ordinary continuous functions.
Boundary data given as coefficientwise measures or distributions, and the
Herglotz positivity problem, are treated in the companion report held in
\path{docs/surcomplex/hahn-herglotz-positivity/}
(\texttt{herg:thm:hierarchy}).

\begin{theorem}[Coherent Schwarz reflection]\label{e:thm-reflection}
Let $F=\sum_{\gamma\in S}f_\gamma\tm^\gamma$ be a coherent datum on an
upper half-disk, and suppose each $f_\gamma$ extends continuously to the
real diameter and is real valued there.  Then the coefficientwise
reflection $\widetilde f_\gamma(Z)=\overline{f_\gamma(\overline Z)}$ on
the lower half-disk joins each $f_\gamma$ to a function holomorphic on the
full disk, the unchanged support gives a coherent extension
$\widetilde F$ of $F$ to the full disk, and
$\ev{\widetilde F}(\bar z)=\overline{\ev{\widetilde F}(z)}$.  The
extension is unique.
\end{theorem}

\begin{proof}
Ordinary Schwarz reflection applies to each $f_\gamma$.  The support is
unchanged, so the reflected family is coherent, and evaluation commutes
with conjugation of coefficients and of the variable.  Uniqueness is the
coherent identity theorem of Section~\ref{c:p4}.
\end{proof}

The reality hypothesis in Theorem~\ref{e:thm-reflection} concerns
\emph{all} Hahn coefficients on the ordinary diameter.  It cannot be
replaced by a hypothesis about an arbitrarily selected fine-open component
of the halo, since the fine topology supplies no connection between such
components.

% =====================================================================
% PART 9.  Exponentials and logarithms: local rigidity, global freedom
% Lead sources: 01, 02, 04, 05, 06, 07, 03, 09.
% CLASS DISCIPLINE.  Everything in this part lives in the GERM-ANALYTIC
% class of Part 3 -- Exp, the twisted family E_lambda, the global
% logarithm L, the principal branch Log, the local charts Log_lambda and
% the function Exp(1/z).  NONE of them is coherent on the relevant
% domain; the global logarithm is proved NOT to be a coherent primitive
% of 1/z on the halo of the punctured plane.  Consequently NO theorem of
% Part 8 is applied to any function in this part.
% =====================================================================

\section{Exponentials and logarithms: local rigidity, global freedom}
\label{e:p9}

Every function constructed in this part is germ-analytic in the sense of
Section~\ref{b:part3} --- it has a Hahn-summable power-series expansion on a fine
neighbourhood of every point of its domain. Their coherent restrictions
and their class-wide domains must be distinguished: the exponential
and its twists are coherent on finite translated charts, while the
global logarithm is not coherent on the halo of the whole punctured
ordinary plane. That specific obstruction, proved in
\S\ref{e:sec-globallog}, is what the contour no-go theorem consumes.

\subsection{The infinitesimal and finite exponentials, and polar form}
\label{e:sec-infexp}

\begin{proposition}[Infinitesimal exponential and logarithm]
\label{e:prop-infexp}
For $h\in\mm$ and $\eta\in\mm$ the families
\[
 \exp_{\mm}(h)=\sum_{n\ge0}\frac{h^n}{n!},
 \qquad
 \log_{\mm}(1+\eta)=\sum_{n\ge1}\frac{(-1)^{n+1}}{n}\eta^{\,n}
\]
are Hahn summable, and they are mutually inverse isomorphisms
\[
 \exp_{\mm}:(\mm,+)\;\xrightarrow{\ \sim\ }\;(1+\mm,\cdot),
 \qquad
 \log_{\mm}:(1+\mm,\cdot)\;\xrightarrow{\ \sim\ }\;(\mm,+).
\]
Both commute with conjugation.  In particular, if $v\in1+\mm$ has
$\abs v=1$, then $\log_{\mm}(v)$ is purely imaginary.
\end{proposition}

\begin{proof}
Summability is Neumann's lemma applied to the positive well-ordered set
$\supp(h)$, or $\supp(\eta)$: complex scalar coefficients add no exponents
and each exponent of $\Neu{\supp(h)}$ has finitely many word
representations.  The identities
$\exp_{\mm}(h+k)=\exp_{\mm}(h)\exp_{\mm}(k)$,
$\log_{\mm}\bigl((1+\eta)(1+\xi)\bigr)=\log_{\mm}(1+\eta)+\log_{\mm}(1+\xi)$,
$\log_{\mm}(\exp_{\mm}(h))=h$ and
$\exp_{\mm}(\log_{\mm}(1+\eta))=1+\eta$ are formal identities in finitely
many infinitesimal variables and hold after evaluation by the same support
argument.  Conjugation acts on the complex coefficients only.  For the
last assertion, $\abs v=1$ means $v\bar v=1$, so
$\log_{\mm}(v)+\overline{\log_{\mm}(v)}=\log_{\mm}(v\bar v)=0$.
\end{proof}

\begin{proposition}[The finite exponential and the polar form]
\label{e:prop-polar}
For $z=c+h\in\Ov$ with $c=\st(z)\in\CC$ and $h\in\mm$, set
\[
 E_{\mathrm{fin}}(z)=e^{c}\exp_{\mm}(h).
\]
Then $E_{\mathrm{fin}}$ is a surjective group homomorphism from
$(\Ov,+)$ onto the group of finite units of $\K$ with nonzero standard
part, and $\ker E_{\mathrm{fin}}=2\pi i\ZZ$.  Moreover every $z\in\K^{\times}$
has a polar representation
\begin{equation}\label{e:eq-polar}
 z=\abs z\,E_{\mathrm{fin}}(i\theta),
 \qquad \theta\in\No\ \text{finite},
\end{equation}
and $\theta$ is unique modulo $2\pi\ZZ$ among finite angles.
\end{proposition}

\begin{proof}
The homomorphism property combines $e^{c+c'}=e^ce^{c'}$ with
Proposition~\ref{e:prop-infexp}.  A finite unit is $c(1+h)$ with $c\ne0$
complex and $h\in\mm$; choosing an ordinary logarithm of $c$ and applying
$\log_{\mm}$ to $1+h$ gives surjectivity.  For the kernel, take standard
parts first, obtaining $e^{c}=1$ and hence $c\in2\pi i\ZZ$, then use the
injectivity of $\exp_{\mm}$ on $\mm$.

For \eqref{e:eq-polar}, put $u=z/\abs z$, so $\abs u=1$ and $u$ is finite.
Its standard part $u_0\in\CC$ has $\abs{u_0}=1$; write $u=u_0(1+\delta)$
with $\delta\in\mm$.  Since $\abs{1+\delta}=1$,
Proposition~\ref{e:prop-infexp} makes $\log_{\mm}(1+\delta)=i\eta$ with
$\eta$ a real infinitesimal.  Choose an ordinary real $\theta_0$ with
$u_0=e^{i\theta_0}$ and set $\theta=\theta_0+\eta$, a finite real surreal.
Uniqueness modulo $2\pi\ZZ$ is the kernel computation.  Shifting by an
ordinary integer multiple of $2\pi$ normalizes $\theta$ to lie in
$(-\pi,\pi]$; for a negative real $z$ this normalized angle is $\pi$.
\end{proof}

Note what \eqref{e:eq-polar} does \emph{not} require: no value of a sine or
cosine at an infinite real argument is used, because every angle occurring
is finite.

\subsection{The canonical surcomplex exponential}\label{e:sec-canonicalexp}

For a real surreal $y$ split the normal form uniquely as
\begin{equation}\label{e:eq-ysplit}
 y=y_{\mathrm{PI}}+y_0+y_{\mathrm{inf}},
 \qquad y_{\mathrm{PI}}\in\Jpi,\ y_0\in\RR,\ y_{\mathrm{inf}}\in\mm\cap\No,
\end{equation}
where $y_{\mathrm{PI}}$ collects the monomials $\tm^{\gamma}$ with
$\gamma<0$, and put
$\operatorname{Fin}(y)=y_0+y_{\mathrm{inf}}$.  Here $\Jpi$ is the additive
group of purely infinite real surreals; it is a real \emph{vector space},
a fact used in Remark~\ref{e:rem-kernel}.  All three projections in
\eqref{e:eq-ysplit} are additive, and $\operatorname{Fin}$ retains all
terms at nonnegative exponents, not merely the standard part.

\begin{definition}[Canonical exponential]\label{e:def-canonicalexp}
For $x,y\in\No$ put
\begin{equation}\label{e:eq-canonicalexp}
 \Exp(x+iy)=\exp_{\No}(x)\,E_{\mathrm{fin}}\bigl(i\operatorname{Fin}(y)\bigr),
\end{equation}
where $\exp_{\No}$ is the Gonshor~\cite{Gonshor} real surreal exponential.
\end{definition}

\begin{theorem}[Properties of the canonical exponential]\label{e:thm-exp}
The map $\Exp:(\K,+)\to(\K^{\times},\cdot)$ is a surjective group
homomorphism extending both $\exp_{\No}$ and the ordinary complex
exponential, with
\begin{equation}\label{e:eq-expprops}
 \abs{\Exp(x+iy)}=\exp_{\No}(x),
 \qquad
 \ker\Exp=2\pi i\,\Oz,\qquad \Oz=\Jpi+\ZZ .
\end{equation}
It is germ-analytic at every point of $\K$: for every $b\in\K$ and every
$h\in\mm$,
\begin{equation}\label{e:eq-explocal}
 \Exp(b+h)=\Exp(b)\sum_{n\ge0}\frac{h^n}{n!}=\Exp(b)\exp_{\mm}(h),
\end{equation}
and consequently $\Exp'=\Exp$ in the fine sense.
\end{theorem}

\begin{proof}
Additivity of $\operatorname{Fin}$ together with the group laws for
$\exp_{\No}$ and $E_{\mathrm{fin}}$ gives the homomorphism property.  The
modulus identity holds because $E_{\mathrm{fin}}$ of a purely imaginary
finite argument has modulus one, by
Proposition~\ref{e:prop-infexp}.

Kernel: the modulus identity forces $x=0$; then
$E_{\mathrm{fin}}(i\operatorname{Fin}(y))=1$, so
$\operatorname{Fin}(y)\in2\pi\ZZ$ by
Proposition~\ref{e:prop-polar}, which says exactly
$y\in\Jpi+2\pi\ZZ$.  Hence $\ker\Exp=i(\Jpi+2\pi\ZZ)$; see
Remark~\ref{e:rem-kernel} for the identification with $2\pi i\Oz$.

Surjectivity is \eqref{e:eq-polar}: for $z\ne0$ write
$z=\abs z\,E_{\mathrm{fin}}(i\theta)$ with $\theta$ finite, and then
$z=\Exp(\log_{\No}\abs z+i\theta)$, since
$\operatorname{Fin}(\theta)=\theta$ for finite $\theta$.

For \eqref{e:eq-explocal}, let $h=p+iq\in\mm$.  Then $q$ is infinitesimal,
so $y_{\mathrm{PI}}(q)=0$ and $\operatorname{Fin}(y+q)=\operatorname{Fin}(y)+q$;
the real infinitesimal Taylor identity for $\exp_{\No}$ and
Proposition~\ref{e:prop-infexp} give $\Exp(h)=\exp_{\mm}(h)$, and the
group law gives \eqref{e:eq-explocal}.  Factoring $h$ out of the linear
remainder bounds the difference-quotient error by
$\abs{\Exp(b)}\,\omega\,\abs h$, which is fine-small, so $\Exp'=\Exp$.
\end{proof}

The existence of a canonical surcomplex exponential with kernel
$2\pi i\Oz$ is established in Ehrlich and Kaplan~\cite[Proposition~11.2 and
Theorem~11.1, pp.~33--34 of arXiv:2002.07739v3]{EK}; what is
proved here is its local analytic calculus.  Note also that $\Exp$ does
\emph{not} have a Hahn-summable Maclaurin series at every argument: at an
infinite positive real argument it is defined by $\exp_{\No}$, not by
declaring $\sum z^n/n!$ summable.

\begin{remark}[The kernel written two ways]\label{e:rem-kernel}
The two standard presentations of the kernel,
\[
 \ker\Exp=2\pi i\,\Oz\quad\text{with }\Oz=\Jpi+\ZZ,
 \qquad\text{and}\qquad
 \ker\Exp=i(\Jpi+2\pi\ZZ),
\]
are equal, and the identity that makes them equal must be stated, because
the first form invites a false reading.  Indeed
$2\pi(\Jpi+\ZZ)=2\pi\Jpi+2\pi\ZZ$, and $\Jpi$ is a real vector space, so
$2\pi\Jpi=\Jpi$.  Written as $2\pi i\Oz$ the kernel looks like a uniform
scaling of a discrete group, suggesting that multiplication by $2\pi$ changes the purely infinite subgroup.
It changes individual elements but preserves that subgroup as a whole.
\end{remark}

\begin{remark}[$\Exp(i\omega)=1$, and how it is \emph{not} obtained]
\label{e:rem-iomega}
Since $\operatorname{Fin}(\omega)=0$, formula
\eqref{e:eq-canonicalexp} gives $\Exp(i\omega)=1$.  This value is
\emph{not} inferred by summing $\sum_{n\ge0}(i\omega)^n/n!$: that family is
not Hahn summable, its term valuations $-n$ forming a strictly decreasing
sequence with no well-ordered union of supports.  The value records the
chosen treatment of infinite imaginary phases in
\eqref{e:eq-canonicalexp}, and \S\ref{e:sec-twisted} shows that the choice
is not forced by any local analytic axiom.
\end{remark}

\subsection{The twisted family: local axioms do not fix the infinite phase}
\label{e:sec-twisted}

For a real surreal $y$ with normal form $\sum_\gamma y_\gamma\tm^{\gamma}$,
let $c_\omega(y)=y_{-1}\in\RR$ be the real coefficient of the monomial
$\omega=\tm^{-1}$.  This is an additive real-linear functional vanishing on
every finite $y$.

\begin{theorem}[The twisted global exponentials]\label{e:thm-twisted}
For each ordinary real $\lambda$ define
\begin{equation}\label{e:eq-twisted}
 E_\lambda(x+iy)=\exp_{\No}(x)\,
 E_{\mathrm{fin}}\bigl(i(\operatorname{Fin}(y)+\lambda\,c_\omega(y))\bigr),
 \qquad E_0=\Exp .
\end{equation}
Each $E_\lambda:(\K,+)\to(\K^{\times},\cdot)$ is a surjective group
homomorphism which
\begin{enumerate}[label=\textup{(\alph*)}]
\item extends $\exp_{\No}$ and the ordinary complex exponential, and
      agrees with $\Exp$ on \emph{all finite} surcomplex arguments;
\item satisfies $E_\lambda(0)=1$, $\abs{E_\lambda(x+iy)}=\exp_{\No}(x)$
      and, for $h\in\mm$,
      $E_\lambda(b+h)=E_\lambda(b)\exp_{\mm}(h)$, hence
      $E_\lambda'=E_\lambda$ and germ-analyticity at every point of $\K$;
\item is locally coherent in finite translated charts:
      $E_\lambda(b+w)=E_\lambda(b)\,\ev{(e^{w})}$ for $w\in\Ov$, a
      surcomplex constant times a canonical classical lift;
\item satisfies $E_\lambda(i\omega)=e^{i\lambda}$, so that
      $E_0(i\omega)=1$ while $E_\pi(i\omega)=-1$, and $E_\lambda\ne E_\mu$
      whenever $\lambda\ne\mu$.
\end{enumerate}
\end{theorem}

\begin{proof}
The map $y\mapsto\operatorname{Fin}(y)+\lambda c_\omega(y)$ is additive and
is the identity on finite $y$, so (a) and the group law follow from
Theorem~\ref{e:thm-exp}.  The correction factor
$E_{\mathrm{fin}}(i\lambda c_\omega(y))$ lies on the ordinary unit circle,
giving the modulus identity.  For an infinitesimal increment $h=p+iq$ one
has $c_\omega(q)=0$ and $\operatorname{Fin}(q)=q$, so the correction factor
is unchanged and the local expansion of Theorem~\ref{e:thm-exp} applies
verbatim; this is (b).  For (c), on the translated finite tube $b+\Ov$ the
same computation gives $E_\lambda(b+w)=E_\lambda(b)\ev{(e^w)}$, a product
of a constant of $\K$ with the lift of the ordinary exponential, hence a
coherent section in that normalized chart.  For (d), $c_\omega(\omega)=1$
and $\operatorname{Fin}(\omega)=0$; taking $y=r\omega$ with
$r=\pi/(\lambda-\mu)$ separates $E_\lambda$ from $E_\mu$.  Surjectivity
follows from the preimages of
Proposition~\ref{e:prop-polar}, whose imaginary parts are finite and hence
untouched by the twist.
\end{proof}

Theorem~\ref{e:thm-twisted} is a diagnostic statement.  It says that
\[
 E'=E,\quad E(0)=1,\quad
 \text{agreement with }\exp_{\No}\text{ and }\exp_{\CC},\quad
 \abs{E(x+iy)}=\exp_{\No}(x)
\]
do \emph{not} determine an exponential's values at infinite imaginary
arguments.  The additional requirement that the phase annihilate the
purely infinite imaginary summand --- that is,
$c_\omega$ and all lower monomials contribute nothing --- is what selects
$\Exp$ from the family.  This is a normalization, not a theorem.

\begin{corollary}[A structured unit-modulus witness]\label{e:cor-twistwitness}
For $\lambda\notin2\pi\ZZ$ the quotient $E_\lambda/\Exp$ is a nonconstant
function on all of $\K$ with values in the ordinary unit circle; it is
germ-analytic everywhere and every one of its positive-order derivatives
vanishes identically.
\end{corollary}

\begin{proof}
By \eqref{e:eq-twisted} the quotient is
$z\mapsto E_{\mathrm{fin}}(i\lambda c_\omega(\operatorname{Im}z))$, which
is unchanged by any finite increment of $z$, hence locally constant on
every translated tube $b+\Ov$ and therefore germ-analytic with vanishing
derivatives.  It is nonconstant by
Theorem~\ref{e:thm-twisted}(d).
\end{proof}

This is a more structured version of the locally constant counterexample of
Section~\ref{b:part2}: local analytic Taylor data alone impose no continuation between
infinitely separated regions of $\K$.  In particular
Corollary~\ref{e:cor-twistwitness} is exactly why no theorem of Section~\ref{e:p8} may
be applied to a function that is merely germ-analytic on $\K$.

\begin{theorem}[The period obstruction]\label{e:thm-periodobstruction}
Let $E:(\K,+)\to(\K^{\times},\cdot)$ be a group homomorphism which
extends the strip exponential --- that is, which agrees with
$\exp_{\No}$ on $\No$ and with $E_{\mathrm{fin}}$ on finite imaginary
arguments --- and which satisfies $E(\bar z)=\overline{E(z)}$.  Then for
every \emph{infinite} real surreal $y$ there is a \emph{finite} real
surreal $v$ with
\begin{equation}\label{e:eq-periodobstruction}
 i(y-v)\in\ker E .
\end{equation}
Consequently no such global extension has kernel exactly $2\pi i\ZZ$.
\end{theorem}

\begin{proof}
The homomorphism property and conjugation symmetry give
$E(iy)\overline{E(iy)}=E(iy)E(-iy)=E(0)=1$, so $E(iy)$ lies on the
surcomplex unit circle.  By the polar form
\eqref{e:eq-polar}, the whole surcomplex unit circle is parametrized by
\emph{finite} real angles: choose a finite $v$ with
$E(iy)=E_{\mathrm{fin}}(iv)=E(iv)$, the last equality by the hypothesis
that $E$ agrees with $E_{\mathrm{fin}}$ on finite imaginary arguments.
Then $E(i(y-v))=1$, and $y-v$ is infinite, so this period is not an
ordinary integer multiple of $2\pi$.
\end{proof}

Theorem~\ref{e:thm-periodobstruction} is the general statement behind the
kernel computations: the other constructions compute the kernel of one
particular exponential, whereas this argument uses only the homomorphism
property, conjugation symmetry and the finite-angle parametrization.  The
obstruction is algebraic, not topological.  Extending the ordinary complex
exponential to all of $\K$ while keeping conjugation and the finite-angle
values necessarily changes the period structure.

\subsection{The global fine-analytic logarithm}\label{e:sec-globallog}

Fix once and for all the ordinary principal argument
$\Arg:\{u\in\CC:\abs u=1\}\to(-\pi,\pi]$.  This is a choice of values on
ordinary unit complex numbers only; it is not a claim that the ordinary
principal argument is continuous across its cut.

\begin{definition}[The global logarithm]\label{e:def-globallog}
For $z\in\K^{\times}$ put $r=\abs z$, $u=z/r$ and $u_0=\st(u)\in\CC$, so
$\abs{u_0}=1$ and $u/u_0\in1+\mm$.  Define
\begin{equation}\label{e:eq-globallog}
 \mathcal L(z)=\log_{\No}(r)+i\Arg(u_0)+\log_{\mm}\!\left(\frac{u}{u_0}\right).
\end{equation}
\end{definition}

The branch is chosen by $u_0=\st(z/\abs z)$, the standard part of the unit
direction --- a quantity that is constant on a sufficiently small
fine-neighbourhood of each nonzero point, since
an infinitesimal relative perturbation does not move a standard part.  That
is what makes \eqref{e:eq-globallog} analytic rather than merely defined.

\begin{theorem}[A logarithm on the whole punctured class plane]
\label{e:thm-globallog}
The function $\mathcal L:\K^{\times}\to\K$ is germ-analytic at every
point, has finite imaginary part, and satisfies
\[
 \mathcal L'(z)=\frac1z,
 \qquad E_\lambda\bigl(\mathcal L(z)\bigr)=z
 \quad\text{for every ordinary real }\lambda .
\]
More precisely, whenever $h/z\in\mm$,
\begin{equation}\label{e:eq-loglocal}
 \mathcal L(z+h)-\mathcal L(z)=\log_{\mm}\!\left(1+\frac hz\right)
 =\sum_{n\ge1}\frac{(-1)^{n+1}}{n\,z^{n}}h^{\,n}.
\end{equation}
Its restriction to ordinary nonzero complex points has the values of the
ordinary principal logarithm, including the selected value $i\pi$ of the
argument on the negative real axis.
\end{theorem}

\begin{proof}
Set $\eta=h/z\in\mm$.  The unit direction of $z(1+\eta)$ has the same
standard part $u_0$ as that of $z$, so the argument term in
\eqref{e:eq-globallog} is unchanged.  The modulus becomes $r\abs{1+\eta}$
and the unit direction becomes $u(1+\eta)/\abs{1+\eta}$, so by the group
laws for $\log_{\No}$ and $\log_{\mm}$,
\[
 \mathcal L(z+h)-\mathcal L(z)
 =\log_{\No}\abs{1+\eta}
  +\log_{\mm}\!\left(\frac{1+\eta}{\abs{1+\eta}}\right)
 =\log_{\mm}(1+\eta),
\]
which is \eqref{e:eq-loglocal}; the displayed series is Hahn summable and
has linear coefficient $1/z$, giving germ-analyticity and the derivative.
The last term of \eqref{e:eq-globallog} is purely imaginary (by
Proposition~\ref{e:prop-infexp}, since $\abs{u/u_0}=1$ exactly) and infinitesimal, and $\Arg(u_0)$ is an ordinary real, so
the imaginary part of $\mathcal L(z)$ is finite.  The right-inverse
identity is then the surjectivity computation of
Theorem~\ref{e:thm-twisted}, because $\mathcal L(z)$ has finite imaginary
part and every $E_\lambda$ agrees with $\Exp$ whenever the imaginary part
is finite, even when the real part is infinite.  At an
ordinary nonzero complex $z$ one has $u=u_0$ and $r$ ordinary, so
\eqref{e:eq-globallog} reduces to $\log\abs z+i\Arg(z/\abs z)$.
\end{proof}

\begin{corollary}[$\mathcal L$ is not a coherent primitive of $1/z$]
\label{e:cor-lognotcoherent}
$\mathcal L$ is not the evaluation of any element of
$\Hc(\CC^{\times})$ on $\halo{(\CC^{\times})}$.
\end{corollary}

\begin{proof}
If it were, then $\mathcal L$ would be a coherent primitive of $1/z$ on
$\halo{(\CC^{\times})}$, and the coherent fundamental theorem of Section~\ref{c:p5}
would force
$\ointH_{\abs z=1}\dd z/z=0$.  But that coefficientwise integral is
$2\pi i$.  Equivalently, the ordinary coefficient data of $\mathcal L$
have a jump across the selected argument cut and do not glue
holomorphically around the circle.
\end{proof}

This is the property that must travel with $\mathcal L$ wherever it goes.
It is also what makes the architecture of Sections~\ref{c:p5} and~\ref{d:sec:preparation} forced rather than
merely cautious: $\mathcal L$ is a globally germ-analytic $G$ on
$\K^{\times}$ with $G'=1/z$, so by the contour no-go theorem of Section~\ref{b:part2} no
integral along the ordinary unit circle can simultaneously return the
classical value $2\pi i$ for $1/z$ and annihilate $G'$ for every such $G$.
The coefficientwise definition of the contour, and the restriction of the
fundamental theorem to \emph{coherent} primitives, are consequences of
that, not stylistic conservatism.

\begin{remark}[No paradox at the negative real axis]\label{e:rem-noparadox}
It is tempting to object that $\mathcal L$ cannot be analytic at a
negative real point, where the ordinary principal logarithm is
discontinuous.  The objection confuses two topologies.  An infinitesimal
perturbation of a fixed negative real point leaves the standard unit
direction $u_0=-1$ unchanged, so it does not cross the chosen
standard-part cut; the branch is fine-locally constant.  Ordinary points
approaching the axis from below do \emph{not} form a fine-convergent
family --- by the discreteness theorem of Sections~\ref{a:sec:summability}--\ref{a:sec:fine} no nontrivial set-indexed
family converges in the fine topology at all --- so no fine-continuity
requirement is violated.  The classical topological obstruction to a
global logarithm and the fine-local differential calculation concern
different structures.  The obstruction survives in the coherent category,
where it is Corollary~\ref{e:cor-lognotcoherent}; it disappears in the more
permissive germ-analytic category, which is Theorem~\ref{e:thm-globallog}.
\end{remark}

\subsection{The principal branch on the cut plane, and the
\texorpdfstring{$2\pi i$}{2 pi i} discrepancy}\label{e:sec-principallog}

A second, genuinely different function is also called ``the'' logarithm in
the sources.  It is retained here under its own name.

\begin{theorem}[Principal logarithm on the cut plane]\label{e:thm-principallog}
On $\K\setminus\No_{\le0}$ define
\[
 \Log z=\log_{\No}\abs z+i\Arg_{\mathrm c}(z),
 \qquad -\pi<\Arg_{\mathrm c}(z)<\pi,
\]
where $\Arg_{\mathrm c}(z)$ is the unique finite angle in $(-\pi,\pi)$
supplied by the polar form \eqref{e:eq-polar}.  Then $\Log$ is
germ-analytic, $\Exp(\Log z)=z$, $\Log'(z)=1/z$, and $\Exp$ restricts to a
bijection of the strip $\{x+iy:x\in\No,\,-\pi<y<\pi\}$ onto the cut plane,
with inverse $\Log$.
\end{theorem}

\begin{proof}
Finite-angle uniqueness in Proposition~\ref{e:prop-polar} gives the
bijection and the inverse identities.  For fixed $z$ off the cut and
sufficiently small $h$, the candidate value
$\Log z+\log_{\mm}(1+h/z)$ still has imaginary part inside $(-\pi,\pi)$,
and its image under $\Exp$ is $z+h$; uniqueness in the strip therefore
gives
$\Log(z+h)=\Log z+\sum_{n\ge1}\frac{(-1)^{n+1}}{n}(h/z)^n$, which proves
germ-analyticity and the derivative.  The required shrinking of $h$ is
legitimate even when the angle is only infinitesimally far from an
endpoint of the strip, since a smaller fine radius is always available.
\end{proof}

\begin{proposition}[$\mathcal L$ and $\Log$ are different functions]
\label{e:prop-logdiscrepancy}
$\mathcal L$ and $\Log$ agree wherever
$\Arg_{\mathrm c}(z)$ and $\Arg(\st(z/\abs z))$ agree, and differ by
exactly $2\pi i$ on the lower half of the monad of every negative ordinary real.
Concretely, let $\varepsilon>0$ be a real infinitesimal and
$z=-1-i\varepsilon$.  Then
\[
 \mathcal L(z)=i\bigl(\pi+\varepsilon\bigr)+O_{\mathrm{formal}}(\varepsilon^{2}),
 \qquad
 \Log(z)=i\bigl(-\pi+\varepsilon\bigr)+O_{\mathrm{formal}}(\varepsilon^{2}),
 \qquad
 \mathcal L(z)-\Log(z)=2\pi i .
\]
\end{proposition}

\begin{proof}
Here $\abs z=1+\varepsilon^{2}/2+\cdots$ is infinitesimally near $1$, so
the real parts of both expressions are infinitesimal and equal.  For
$\mathcal L$: the unit direction is
$u=(-1-i\varepsilon)/\abs z$ with $\st(u)=u_0=-1$, and
$\Arg(-1)=+\pi$ under the convention $(-\pi,\pi]$; the correction term is
$\log_{\mm}(u/u_0)=i\varepsilon+O_{\mathrm{formal}}(\varepsilon^{2})$.
For $\Log$: the finite angle of $z$ in $(-\pi,\pi)$ is
$-\pi+\varepsilon+O_{\mathrm{formal}}(\varepsilon^{2})$, since $z$ lies
infinitesimally below the negative real axis.  Subtracting gives
$2\pi i$ exactly, the infinitesimal corrections cancelling.  On the upper
half-monad, $z=-1+i\varepsilon$, both give
$i(\pi-\varepsilon)+\cdots$ and the two agree.
\end{proof}

Both branches are legitimate, and neither is wrong; they are different
functions and must not be identified.  This document adopts $\mathcal L$
as the primary logarithm, for three reasons.  Its domain is all of
$\K^{\times}$, not $\K$ minus a class-sized ray.  It is germ-analytic
everywhere, because its branch is selected by the fine-locally constant
quantity $\st(z/\abs z)$, and $\Log$ is likewise fine-analytic at every point of its cut-plane
domain. The discontinuity concerns extending the chosen principal
values across a deleted negative-axis point: the values from its two
sides differ by $2\pi i$. Proposition~\ref{e:prop-logdiscrepancy}
does not show a discontinuity at any point where $\Log$ is defined.  And the contour no-go theorem of
Section~\ref{b:part2} requires a global primitive of $1/z$ on all of $\K^{\times}$, which
$\Log$ does not supply.

\subsection{Local logarithm charts}\label{e:sec-logcharts}

\begin{theorem}[Logarithm charts on multiplicative monads]
\label{e:thm-logcharts}
Let $b\in\K^{\times}$ and choose $\lambda\in\K$ with $\Exp(\lambda)=b$.
For $w\in b(1+\mm)$ define
\[
 \Log_\lambda(w)=\lambda+\sum_{n\ge1}\frac{(-1)^{n+1}}{n}
   \left(\frac{w-b}{b}\right)^{n}
 =\lambda+\log_{\mm}\!\left(\frac wb\right).
\]
This is a bijection $b(1+\mm)\to\lambda+\mm$ inverse to the restriction of
$\Exp$, and $\Log_\lambda'(w)=1/w$.  Any two admissible choices of
$\lambda$ differ by an element of $2\pi i\Oz$, and the corresponding charts
differ by that constant.
\end{theorem}

\begin{proof}
The formal identities $\log(e^h)=h$ and $e^{\log(1+h)}=1+h$ evaluate at
infinitesimals by Proposition~\ref{e:prop-infexp}; the derivative is the
evaluated derivative of the formal logarithm followed by the affine chain
rule.  The last assertion is the kernel computation of
Theorem~\ref{e:thm-exp}.
\end{proof}

These charts describe branch choices without requiring a
fine-continuous principal argument on all of $\K^\times$ or a
path-connected covering-space structure. A principal argument can be
chosen as a function, but its conventional values do not extend
continuously across the negative real axis; independently, Sections~\ref{a:sec:summability}--\ref{a:sec:fine}
excludes nonconstant ordinary-parameter fine-continuous paths.  Note that the ambiguity is by an element of
$2\pi i\Oz$, not merely by an ordinary integer multiple of $2\pi i$.

\subsection{A value-distribution theorem at an explicit essential
singularity}\label{e:sec-picard}

Put $P(z)=\Exp(1/z)$ for $z\in\K^{\times}$.  By
Theorem~\ref{e:thm-exp} and the chain rule $P$ is germ-analytic on
$\K^{\times}$ with $P'(z)=-z^{-2}P(z)$, and it never takes the value zero.

\begin{theorem}[Every nonzero value at every surreal scale]
\label{e:thm-picard}
For every $w\in\K^{\times}$ and every positive surreal $\delta$ there is a
\emph{proper class} of distinct $z$ with
\[
 0<\abs z<\delta,\qquad \Exp(1/z)=w .
\]
Consequently $P$ has at zero neither a removable singularity nor a pole in
the fine topology.
\end{theorem}

\begin{proof}
Choose $\lambda$ with $\Exp(\lambda)=w$, possible by
Theorem~\ref{e:thm-exp}.  For $q\in\Oz$ with $\lambda+2\pi iq\ne0$ put
\[
 z_q=\frac1{\lambda+2\pi iq} .
\]
These are distinct, and $P(z_q)=\Exp(\lambda+2\pi iq)=w$ because
$2\pi i\Oz=\ker\Exp$.  The omnific integers are cofinal in $\No^{>0}$ ---
indeed $\Oz$ contains every ordinal --- so for all sufficiently large
positive $q\in\Oz$,
\[
 \abs{\lambda+2\pi iq}\ \ge\ 2\pi q-\abs\lambda\ >\ \delta^{-1},
\]
whence $0<\abs{z_q}<\delta$.  A proper-class tail of such $q$ supplies the
required solutions.  Taking $w=1$ and $w=2$ shows that $P$ has no fine
limit at zero, since every punctured fine-neighbourhood of zero contains
both values; taking $w=1$ also excludes a pole, which would require
$\abs P$ to exceed every fixed surreal bound near zero.
\end{proof}

\begin{remark}[Scope, and the phase dependence of the twisted variants]
\label{e:rem-picardscope}
Theorem~\ref{e:thm-picard} is an exact value-distribution theorem for one
specified function.  It is not a general Picard theorem for germ-analytic
functions, and no such theorem is proved or claimed anywhere in this
article: a general statement of that breadth would be incompatible with
the locally constant witnesses of Section~\ref{b:part2} and with
Corollary~\ref{e:cor-twistwitness}.  Nor does the phrase ``essential
singularity'' here assert a global Laurent classification; it records
exactly the two proved failures, of removability and of a pole.

The value-distribution conclusion also holds for every twist
$E_\lambda$. The canonical periods $2\pi i\Oz$ need not remain
periods of a twist, but a cofinal family of periods does remain:
\[
 E_\lambda(i\omega^\alpha)=1\qquad(\alpha\ge2\text{ an ordinal}),
\]
because $\omega^\alpha$ is purely infinite and has zero coefficient at
$\omega$. These positive monomials are cofinal in $\No$, and every tail
is a proper class. Given $w\ne0$, choose $\ell$ with $E_\lambda(\ell)=w$.
The distinct points $z_\alpha=(\ell+i\omega^\alpha)^{-1}$ then have
$E_\lambda(1/z_\alpha)=w$, and all sufficiently large $\alpha$ give
$0<\abs{z_\alpha}<\delta$, by the same triangle estimate as above.
Thus the earlier claim that the canonical phase was necessary for this
conclusion was incorrect.

The kernels are not generally nested. For example $i\omega$ belongs to
$\ker\Exp$ but not to $\ker E_\pi$, while $i(\omega-\pi)$ belongs to
$\ker E_\pi$ but not to $\ker\Exp$. The functions still differ at
specified points: at $z=-i/\omega$ their values are respectively $1$ and
$-1$. Agreement on ordinary arguments and the ordinary Laurent series
does not fix those values, even though all the twists share the stated
value-distribution property.

Finally, a summability boundary that belongs beside the theorem.  The
canonical lift $\ev{(e^{1/Z})}$ is a coherent section on
$\halo{(\CC^{\times})}$, but that halo excludes the entire monad of zero.
At a nonzero infinitesimal $z$ the formal expression
$\sum_{n\ge0}z^{-n}/n!$ is \emph{not} Hahn summable: its term valuations
$-n\,v(z)$ form a strictly decreasing sequence.  So the ordinary Laurent
series does not itself define any extension into that monad, and
$P=\Exp(1/z)$ is an extension chosen by the exponential, not derived from
the series.  On ordinary annuli avoiding zero, $P$ does agree with
$\ev{(e^{1/Z})}$, so its coefficientwise contour integral on an ordinary
positively oriented circle is $2\pi i$, the ordinary Laurent coefficient of
$Z^{-1}$ being one.  That is consistent with, but not determined by, the
behaviour of $P$ on the infinitesimal monads near the singularity.
\end{remark}

%% ==================== p1012.tex ====================

% =====================================================================
% PART 10 -- Affine charts at all scales, and the rigidity theorems
% PART 11 -- Worked examples with exact arithmetic
% PART 12 -- Boundaries: what is and is not established
% Sources: 07, 01, 02, 08 (lead); 03, 04, 05, 06, 09 (supporting).
% Every theorem below names the function class it is proved for.
% =====================================================================

\section{Affine charts at all scales, and the rigidity theorems}
\label{f:sec-allscale}

The field $\K$ carries a proper class of multiplicative scales, and
nothing in the coherent theory of $\Hc(U)$ privileges the scale $1$.
This part first transports that theory to an arbitrary centre and an
arbitrary nonzero scale, then asks what survives if one insists on a
coherent description at \emph{every} scale at once.  The answer is a
pair of classification theorems: all-scale coherence admits exactly the
polynomials, and its meromorphic counterpart admits exactly the rational
functions.  Every statement in this part is a statement about the
\emph{coherent} class $\Hc(U)$ and about class functions assembled from
coherent charts; none of it is asserted for germ-analytic functions,
whose class contains the locally constant counterexample of
the locally constant witness of Section~\ref{b:part2} and therefore admits no identity theorem to
run the argument with.

\subsection{Coherent affine charts at an arbitrary centre and scale}
\label{f:sub-charts}

\begin{definition}[Scaled coherent chart]\label{f:def-chart}
Let $b\in\K$, let $r\in\K^{\times}$, and let $U\subseteq\CC$ be an
ordinary domain.  The \emph{scaled chart} with centre $b$ and scale $r$
is the fine-open class
\[
 \mathcal U(b,r;U)\;=\;b+r\,\halo U\;=\;\{\,b+rw:\;w\in\halo U\,\}.
\]
A function $f$ on $\mathcal U(b,r;U)$ is \emph{coherent in the
coordinate $(z-b)/r$} if there is $F\in\Hc(U)$ with
\[
 f(z)=\ev F\!\left(\frac{z-b}{r}\right)
 \qquad\bigl(z\in\mathcal U(b,r;U)\bigr).
\]
The affine coordinate is part of the specified structure, not a
normalization that may afterwards be suppressed.  The centre $b$ may
have infinite real or imaginary part, and $\abs r$ may be infinitesimal,
finite or infinite.
\end{definition}

\begin{theorem}[Affine transport of the coherent theory]\label{f:thm-affine}
Let $f$ be coherent in the coordinate $(z-b)/r$, with $F\in\Hc(U)$ as in
Definition~\textup{\ref{f:def-chart}}.  Then:
\begin{enumerate}[label=\textup{(\roman*)},leftmargin=2.4em]
\item Derivatives rescale:
 $f^{(n)}(b+rw)=r^{-n}\,\ev{(F^{(n)})}(w)$ for every ordinary
 $n\ge0$ and every $w\in\halo U$.
\item For an ordinary piecewise-$C^{1}$ path $\gamma$ in $U$, the chart
 contour functional is \emph{defined} by
 \begin{equation}\label{f:eq-chartintegral}
  \intH_{\,b+r\gamma}f(z)\,\dd z\;:=\;r\intH_{\gamma}F(w)\,\dd w .
 \end{equation}
 It is $\K$-linear, additive under concatenation, orientation
 reversing, and invariant under ordinary reparameterization of
 $\gamma$.
\item The identity, open-mapping, local maximum-modulus, preparation,
 division and zero-counting theorems transport verbatim from
 $\halo U$ to $\mathcal U(b,r;U)$, and Cauchy's theorem, the Cauchy
 integral formula and the residue theorem hold for
 \eqref{f:eq-chartintegral}.
\item If $f$ is a coherent meromorphic quotient in the chart and
 $p=b+r\lambda$ with $\lambda\in\halo U$, then
 \begin{equation}\label{f:eq-scaledresidue}
  \Resat{p}\bigl(f(z)\,\dd z\bigr)
  \;=\;r\,\Resat{\lambda}\bigl(\ev F(w)\,\dd w\bigr).
 \end{equation}
 Zero and pole multiplicities and the normalized logarithmic-derivative
 index are unchanged by the affine coordinate.
\end{enumerate}
\end{theorem}

\begin{proof}
(i) is the finite chain rule applied inside the coefficientwise
derivative, which introduces no new Hahn exponents.  (ii) is a
definition together with the corresponding properties of the
coefficientwise contour functional, \ref{a:rule:contour}.  For (iii), $w\mapsto b+rw$ is a
fine homeomorphism with analytic inverse and nonzero linear term, so it
preserves local multiplicities; pulling a preparation factorization
back through it gives the factorization in the chart coordinate with the
same finite number of roots, and the Cauchy and residue formulas follow
from \eqref{f:eq-chartintegral} by substitution.

For (iv), write $F(\lambda+X)=\sum_{n}a_{n}X^{n}$ as a formal Laurent
germ with finite principal part.  Then
$f(p+Y)=\sum_{n}a_{n}r^{-n}Y^{n}$, whose $Y^{-1}$ coefficient is
$r\,a_{-1}$.  In the logarithmic derivative the factor $r^{-1}$ produced
by differentiation cancels the factor $r$ carried by the differential
$\dd z=r\,\dd w$, so the index computed by
$\ointH F'/F$ is exactly the ordinary integer of the argument
principle, \ref{d:eq:argument}.
\end{proof}

\begin{remark}[What a scaled contour is and is not]\label{f:rem-chartpath}
The symbol $b+r\gamma$ in \eqref{f:eq-chartintegral} records an affine
contour \emph{datum}: an ordinary path $\gamma$ together with a centre
and a scale.  It is not an assertion that $b+r\gamma$ is a nonconstant
fine-continuous image of a real interval, and by the discreteness
theorem \ref{a:thm:discrete} no such image exists.  This caveat does not
weaken when the centre or the scale changes.
\end{remark}

\begin{remark}[No unrestricted chart independence]\label{f:rem-nogluing}
Ordinary complex affine coordinate changes preserve the coherent class
by coefficientwise composition.  Changes involving surreal centres and
scales may relate very different ordinary base domains, and no
unrestricted gluing or chart-independence theorem is asserted here.
Compatibility of two charts means equality of their \emph{evaluated}
functions on the overlap; every contour computation below names the
chart in which it is performed.  When two admissible descriptions yield
the same pointwise pulled-back integrand, uniqueness of normal forms
identifies their coefficient functions and their integrals therefore
agree.
\end{remark}

\begin{example}[Charts far from the ordinary plane]\label{f:ex-farchart}
Take $b=\omega+i\omega^{2}$ and $r=\omega^{-\omega}$.  The resulting
chart is an exceedingly small neighbourhood of an infinite point, and
its zero clusters, preparation polynomials and residues are as
accessible as those in the halo of the ordinary unit disk
$\halo{\mathbb D}$.  Taking $r=\omega^{\omega}$ instead gives a chart at
an infinitely large scale.  Nothing in Theorem~\ref{f:thm-affine}
distinguishes these cases from $b=0$, $r=1$.
\end{example}

\subsection{The chart-change criterion, and the source of the
obstruction}\label{f:sub-chartchange}

The transport of Theorem~\ref{f:thm-affine} is along a \emph{fixed}
affine map.  A finer question is when a normalized change of chart
carries a coherent section to a coherent section.  The criterion is
explicit, and its failure is exactly what separates local coherence from
coherence at all scales.

\begin{proposition}[Chart-change criterion]\label{f:prop-chartchange}
Let $F=\sum_{s\in S}f_{s}t^{s}\in\Hc(U)$, and consider the normalized
coordinate change $z=\alpha+\beta w$ with $\alpha,\beta\in\Ov$
\emph{finite}.  Write
$\alpha=\alpha_{0}+\delta_{\alpha}$ and
$\beta=\beta_{0}+\delta_{\beta}$ with
$\alpha_{0}=\st(\alpha)$, $\beta_{0}=\st(\beta)$ in $\CC$ and
$\delta_{\alpha},\delta_{\beta}\in\mm$.  Let $V\subseteq\CC$ be an
ordinary domain with $\alpha_{0}+\beta_{0}V\subseteq U$.  Then
\begin{equation}\label{f:eq-chartchange}
 F(\alpha+\beta w)\;=\;
 \sum_{s\in S}\;\sum_{n\ge0}
 \frac{f_{s}^{(n)}(\alpha_{0}+\beta_{0}w)}{n!}
 \bigl(\delta_{\alpha}+\delta_{\beta}w\bigr)^{n}\,t^{s}
\end{equation}
defines a section of $\Hc(V)$, whose evaluation on $\halo V$ agrees with
$w\mapsto \ev F(\alpha+\beta w)$.  The degenerate case
$\beta_{0}=0$ is permitted, provided $\alpha_{0}\in U$.
\end{proposition}

\begin{proof}
For fixed $s$ and $n$ the function
$w\mapsto f_{s}^{(n)}(\alpha_{0}+\beta_{0}w)$ is ordinary holomorphic on
$V$, and in the degenerate case $\beta_{0}=0$ it is the constant
$f_{s}^{(n)}(\alpha_{0})$, which is legitimate because
$\alpha_{0}\in U$.  The support of the coefficient of $t^{s}$ in
\eqref{f:eq-chartchange} is contained in
$s+\Neu{\supp\delta_{\alpha}\cup\supp\delta_{\beta}}$, a translate of
the Neumann~\cite{Neumann} monoid of one fixed well-ordered set of positive exponents.
By Neumann's support lemma \ref{a:sec:neumann} that union is well
ordered and each exponent receives only finitely many contributions, so
the double family is Hahn summable and its sum is a section over the one
ordinary domain $V$.  Agreement of the evaluations is the exact
monad-wide Taylor expansion of \ref{c:p4:eval}.
\end{proof}

\begin{remark}[An infinite normalized coefficient breaks the
criterion]\label{f:rem-infinitebeta}
Proposition~\ref{f:prop-chartchange} requires $\alpha$ and $\beta$
\emph{finite}.  If the normalized scale coefficient $\beta$ is infinite
there is no standard part $\beta_{0}\in\CC$ to serve as the ordinary
inner map, the remainder $\delta_{\beta}$ is not infinitesimal, and the
substitution is no longer justified by \eqref{f:eq-chartchange}.
For example, substituting an infinite $\beta$ into the exponential
Taylor series gives decreasing valuations $n v(\beta)$. Polynomial
coefficient functions can still be substituted at infinite scales;
failure of the criterion is not a claim that every such substitution fails.  This is the
precise mechanism by which a function can be coherent in every finite
translated chart and yet fail to be coherent at all scales.  It is the
structural reason behind Theorem~\ref{f:thm-rigidity}, and it is why the
canonical exponential $\Exp$ of \ref{e:thm-exp}, which is locally
coherent, is excluded by that theorem.
\end{remark}

\subsection{Cauchy estimates at an arbitrary coherent scale}
\label{f:sub-scaledcauchy}

\begin{corollary}[Cauchy estimates at an arbitrary coherent
scale]\label{f:cor-scaledCauchy}
Let $a\in\K$, $r\in\No_{>0}$, and suppose $f(a+rw)=\ev F(w)$ for some
$F\in\Hc(B(0,2))$.  If there is a single $M\in\No$ with
\[
 \abs{f(a+re^{i\theta})}\le M
 \qquad\text{for every \emph{ordinary} real }\theta,
\]
then
\begin{equation}\label{f:eq-scaledCauchy}
 \abs{f^{(n)}(a)}\;\le\;\frac{n!\,M}{r^{n}}
 \qquad(n\ge0).
\end{equation}
\end{corollary}

\begin{proof}
By Theorem~\ref{f:thm-affine}(i), $F^{(n)}(0)=r^{n}f^{(n)}(a)$.  Apply
the pointwise surreal Cauchy estimate \ref{c:p5:eq-cauchyest} on the
ordinary unit circle to $F$; its hypothesis is a single \emph{pointwise}
surreal bound on the path, and its bound $M$ is allowed to be surreal
because the underlying $\mathrm{ML}$ inequality \ref{c:p5:eq-ML} is proved
from the positivity of the coefficientwise integral rather than from a
coefficientwise majorant.
\end{proof}

\begin{remark}[Which estimate this is]\label{f:rem-whichestimate}
Corollary~\ref{f:cor-scaledCauchy} uses the \emph{pointwise} estimate
\ref{c:p5:eq-cauchyest}, resting on \ref{c:p5:eq-ML}, which bounds
the surreal \emph{modulus} of a derivative.  It is not the
coefficientwise majorant $B_{r}(F)=\sum_{s}(\max\abs{f_{s}})t^{s}$ of
\ref{c:p5:majorant}, and it does not assert that a single surreal
bound controls every coefficient function: that assertion is false, as
$F(w)=w(1+it)$ shows, since $\abs{\ev F}=\sqrt{1+t^{2}}$ on the unit
circle admits the pointwise bound $1+t^{2}$ while its coefficientwise
majorant is $1+t>1+t^{2}$.  The two estimates are about different
quantities and both are retained.
\end{remark}

\begin{remark}[Some small chart, but not every large one]
\label{f:rem-somechart}
Every germ-analytic germ admits \emph{at least one} coherent affine
chart, by the scale-embedding bridge \ref{c:p4:bridge}: after a
single monomial rescaling $\rho=t^{\lambda}$ chosen by block separation,
each Hahn coefficient of the rescaled series is a finite sum and hence an
ordinary polynomial, so the rescaled datum is coherent on all of $\CC$.
There is \emph{no} implication that the same germ has a chart of every
larger radius.  Making that distinction precise is exactly the content
of the next subsection.
\end{remark}

\subsection{All-scale coherence forces polynomiality}
\label{f:sub-rigidity}

\begin{definition}[All-scale Hahn-entire]\label{f:def-allscale}
A class function $f:\K\to\K$ is \emph{all-scale Hahn-entire} if for
every $r\in\No_{>0}$ there exists $F_{r}\in\Hc(B(0,2))$ with
\begin{equation}\label{f:eq-allscale}
 f(rw)=\ev{F_{r}}(w)\qquad\bigl(w\in\halo{B(0,2)}\bigr).
\end{equation}
Only charts centred at zero are required, and the well-ordered supports
$S_{r}=\supp F_{r}$ may depend on $r$; \emph{no} common support across
scales is imposed.
\end{definition}

An equivalent formulation replaces $B(0,2)$ by the ordinary unit disk
$\mathbb D$ and quantifies over $R\in\No^{>0}$; the two conditions
define the same class of functions, and the proof below is the same
argument with different constants.  These charts cover $\K$ and make $f$
fine-locally analytic; but although each individual chart is
set-supported, the requirement ranges over a proper class of radii, and
that is what makes it restrictive.

\begin{theorem}[All-scale polynomial rigidity]\label{f:thm-rigidity}
A class function $f:\K\to\K$ is all-scale Hahn-entire if and only if it
is a polynomial in $z$ with coefficients in $\K$ and ordinary finite
degree.
\end{theorem}

\begin{proof}
($\Leftarrow$)  Let $P(z)=\sum_{n=0}^{N}a_{n}z^{n}$ with $a_{n}\in\K$.
For a fixed $r$, expanding each $a_{n}r^{n}$ into normal form uses
finitely many well-ordered supports, whose union is well ordered, and
the coefficient function attached to each exponent is an ordinary
polynomial in $w$ of degree at most $N$.  Hence $P(rw)$ is the
evaluation of a member of $\Hc(B(0,2))$ for every $r$.

($\Rightarrow$)  Put $a_{n}=f^{(n)}(0)/n!\in\K$ for $n\ge0$; the chart
at $r=1$ supplies these derivatives, and $(a_{n})_{n\in\NN}$ is a
\emph{set}.  Fix $r$ and let $S_{r}$ be the support of $F_{r}$.  Since
coefficientwise differentiation of an ordinary holomorphic coefficient
function introduces no new Hahn exponents, the $n$th Taylor coefficient
of $F_{r}$ at $0$, which by the chain rule equals $a_{n}r^{n}$, has
normal-form support contained in $S_{r}$; in particular
$v(a_{n}r^{n})\in S_{r}$ whenever $a_{n}\ne0$.

Suppose infinitely many $a_{n}$ are nonzero.  The valuations
$\{v(a_{n}):a_{n}\ne0\}$ form a set of surreals, so by the surreal cut
property there is $\Delta>0$ with
\[
 \abs{v(a_{n})}<\Delta/4\qquad\text{whenever }a_{n}\ne0 .
\]
Take the scale $r=t^{-\Delta}$, so $v(r)=-\Delta$.  For two indices
$m>n$ with $a_{m},a_{n}\ne0$,
\[
 v(a_{m}r^{m})-v(a_{n}r^{n})
 =v(a_{m})-v(a_{n})-(m-n)\Delta
 <\tfrac{\Delta}{2}-\Delta<0 .
\]
Thus the valuations of the nonzero $a_{n}r^{n}$ form an infinite
strictly decreasing sequence of elements of the single set $S_{r}$,
contradicting its well ordering.

Hence $a_{n}=0$ for all $n>N$, for some ordinary finite $N$.  Put
$P(z)=\sum_{n=0}^{N}a_{n}z^{n}$.  For each $r$ the coherent section
$F_{r}-P(r\,\cdot)$ lies in $\Hc(B(0,2))$ and has all derivatives zero
at $0$, so by the coherent identity theorem \ref{c:p4:identity} it is
the zero section and $f(rw)=P(rw)$ on $\halo{B(0,2)}$.  Given any
$z\in\K$, choose $r\in\No_{>0}$ with $r>2\abs z$; then $z/r$ is finite
with standard part of modulus at most $1/2$, hence lies in
$\halo{B(0,2)}$, and therefore $f(z)=P(z)$.
\end{proof}

\begin{remark}[The mechanism is support-theoretic, not
Archimedean]\label{f:rem-rigiditymechanism}
Theorem~\ref{f:thm-rigidity} is not Liouville's theorem with a larger
word for infinity, and its proof contains no growth estimate.  Two
proper-class facts do the work: the surreal cut property supplies
\emph{one} exponent $\Delta$ dominating a whole set of coefficient
valuations, and $\K$ then supplies a scale $r=t^{-\Delta}$ large enough
to reverse the order of infinitely many nonzero Taylor coefficients.  A
single fixed set-sized Hahn field $K_{\Gamma}$ need contain no such
$\Delta$ or $r$, so the theorem is a statement about $\K$ and must not
be asserted inside a chosen set-sized workspace.
\end{remark}

\begin{corollary}[Global classical consequences in the rigid
category]\label{f:cor-globalclassical}
Let $f$ be all-scale Hahn-entire.
\begin{enumerate}[label=\textup{(\roman*)},leftmargin=2.4em]
\item If $\abs{f(z)}\le M$ for all $z\in\K$, for some $M\in\No$, then
 $f$ is constant.
\item If $\abs{f(z)}\le C(1+\abs z)^{N}$ for all $z\in\K$, with
 $C\in\No_{>0}$ and $N$ an ordinary nonnegative integer, then
 $\deg f\le N$.
\item If $f$ is nonconstant it is surjective onto $\K$.
\item If $f$ is injective it is affine with nonzero slope.
\item There is no all-scale Hahn-entire solution of $f'=f$, $f(0)=1$.
\end{enumerate}
\end{corollary}

\begin{proof}
By Theorem~\ref{f:thm-rigidity}, $f$ is a polynomial
$\sum_{n\le d}a_{n}z^{n}$ with $a_{d}\ne0$.  For (ii), if $d>N$ choose a
positive real-coordinate surreal $x$ large enough that the leading term
dominates the sum of the lower terms and $\abs{a_{d}}x^{d-N}>2^{N+1}C$;
then $\abs{f(x)}\ge\abs{a_{d}}x^{d}/2>C(1+x)^{N}$, a contradiction.
(i) is the case $N=0$ after enlarging $M$ to a nonnegative bound.
(iii) is algebraic closedness of $\K$ applied to $f(z)-c$.  For (iv), a
polynomial of degree at least two has a critical point, since its
nonconstant derivative has a root in the algebraically closed $\K$; the
finite-ramification normal form \ref{b:ramification} at that point
gives more than one preimage of sufficiently close noncritical values,
so $f$ is not injective, and degree $\le1$ with nonzero leading
coefficient is exactly the affine case.  (v) is a comparison of degrees:
no nonzero polynomial equals its own derivative.
\end{proof}

\begin{remark}[This is not the Liouville hierarchy]\label{f:rem-notliouville}
Corollary~\ref{f:cor-globalclassical}(i) must not be confused with any
of the three Liouville hypotheses tabulated in
\S\ref{f:sub-threehierarchies}.  Hypothesis (a) there --- a single surreal bound on a coherent
section over the finite halo $\halo\CC$ --- is \emph{vacuous}, by
\ref{e:prop-autobounded}.  The hypothesis here is different in two ways:
the bound is required on all of the class $\K$, not on the finite halo,
and the function is required to be coherent at \emph{every} scale, not
on one.  It is the all-scale hypothesis, through
Theorem~\ref{f:thm-rigidity}, that produces constancy, not the bound by
itself.
\end{remark}

\begin{remark}[All-scale coherence is a rigidity condition, not a
definition]\label{f:rem-notadefinition}
Condition \eqref{f:eq-allscale} should not be offered as the definition
of ``entire surcomplex analyticity.''  It excludes the canonical
exponential $\Exp$, whose Taylor coefficients at zero are all nonzero,
and by Remark~\ref{f:rem-infinitebeta} the exclusion is caused by an
infinite normalized affine coefficient rather than by any defect of
$\Exp$ as an analytic object: $\Exp$ and each twist $E_{\lambda}$ of
\ref{e:thm-twisted} is germ-analytic everywhere and coherent in every
finite translated chart.  What \eqref{f:eq-allscale} does is classify
exactly what is possible when one demands the coherence of this article
on every scale simultaneously.
\end{remark}

\subsection{Pole-free fractions and all-scale rational rigidity}
\label{f:sub-rational}

\begin{lemma}[A pole-free coherent fraction is coherent
holomorphic]\label{f:lem-polefree}
Let $U\subseteq\CC$ be a connected ordinary domain, let
$F,G\in\Hc(U)$ with $G\ne0$, and suppose the evaluated fraction
$\ev F/\ev G$ has no actual poles in $\halo U$, its removable values
being filled in.  Then it is the evaluation of a member of $\Hc(U)$.
\end{lemma}

\begin{proof}
Let $c\in U$.  If the leading coefficient function $g_{\vH(G)}$ of $G$
does not vanish at $c$, then $G$ is a unit of $\Hc(U')$ on a small
ordinary disk $U'\ni c$ by the coherent unit criterion, and
$F/G\in\Hc(U')$ directly.

Otherwise $g_{\vH(G)}$ vanishes at $c$, and Hahn--Weierstrass
preparation \ref{d:thm:preparation} on a small ordinary disk about $c$
gives, in the local coordinate $u=z-c$,
\[
 G(c+u)=t^{\lambda}\,q_{0}(u)\,P(u)\,Q(u),
\]
where $q_{0}$ and $Q$ are coherent units, $\lambda=\vH(G)$, and $P$ is
monic of degree $m=\ord_{c}(g_{\lambda})$ with infinitesimal lower
coefficients.  Absorb $t^{\lambda}q_{0}Q$ into the numerator, obtaining
a coherent $N$, and apply the division theorem
\ref{d:thm:division}:
\[
 \frac{N}{P}=J+\frac{R}{P},\qquad J\in\Hc(U'),\quad
 \deg R<m .
\]
Every root of $P$ is infinitesimal, since a monic polynomial with
infinitesimal lower coefficients has only infinitesimal roots: a root of
valuation $\le0$ would leave the monic term uncancellable.  As $J$ is
coherent holomorphic and the original fraction has no pole in $\halo{U'}$,
the rational function $R/P\in\K(u)$ has no pole at any root of $P$.
Hence $P\mid R$ in $\K[u]$, and $\deg R<m=\deg P$ forces $R=0$.  So
$\ev F/\ev G=\ev J$ on $\halo{U'}$.

These local representations agree as evaluated functions on overlaps, so
the coherent gluing theorem \ref{c:p4:gluing} joins them into a single
member of $\Hc(U)$.  The gluing step is what supplies a \emph{common
global} support; it is not being assumed from an arbitrary union of
local supports, which need not be well ordered.
\end{proof}

\begin{definition}[All-scale Hahn-meromorphic]\label{f:def-allscalemero}
A function $f:\K\to\mathbb P^{1}(\K)$, not identically $\infty$, is
\emph{all-scale Hahn-meromorphic} if for every $r\in\No_{>0}$ the
pullback $w\mapsto f(rw)$ is represented on $\halo{B(0,2)}$ by a
fraction $\ev{F_{r}}/\ev{G_{r}}$ with $F_{r},G_{r}\in\Hc(B(0,2))$ and
$G_{r}\ne0$.  Fractions are understood as local meromorphic germs,
including their uniquely filled removable values and their pole values
in $\mathbb P^{1}(\K)$.
\end{definition}

\begin{theorem}[All-scale rational rigidity]\label{f:thm-rationalrigidity}
The all-scale Hahn-meromorphic functions on $\K$ are precisely the
rational functions with surcomplex coefficients.
\end{theorem}

\begin{proof}
First, $f$ has only finitely many poles in the whole class plane.
Otherwise choose a countable set of distinct poles $a_{1},a_{2},\dots$
and, by the surreal cut property, a positive surreal $r$ exceeding
$n\abs{a_{j}}$ for all ordinary positive integers $n$ and all $j$; then
every $a_{j}/r$ is infinitesimal, hence lies in the monad of $0$ inside
$\halo{B(0,2)}$.  The radius-$r$ representation writes $f(rw)$ as a
coherent fraction $\ev{F_{r}}/\ev{G_{r}}$ there.  By the zero-cluster
theorem \ref{d:thm:zerocluster} applied to the leading coefficient
function of $G_{r}$, the denominator has exactly
$\ord_{0}(g_{\vH(G_{r})})$ zeros in that monad, a finite number.  So
$f(r\,\cdot)$ cannot have infinitely many distinct poles there, a
contradiction.

Let the actual poles be $a_{1},\dots,a_{k}$ with orders
$m_{1},\dots,m_{k}$, and set
$Q(z)=\prod_{j=1}^{k}(z-a_{j})^{m_{j}}\in\K[z]$.  Fill the removable
values of $g=Qf$.  In every scale chart, $g$ is a coherent fraction with
no actual poles, hence by Lemma~\ref{f:lem-polefree} is the evaluation
of a member of $\Hc(B(0,2))$.  So $g$ is all-scale Hahn-entire, and
Theorem~\ref{f:thm-rigidity} makes $g=P$ a polynomial.  Therefore
$f=P/Q$ is rational.

Conversely, the numerator and denominator of a rational function over
$\K$ become polynomials with surcomplex coefficients after every affine
scaling, and by the converse half of Theorem~\ref{f:thm-rigidity} these
supply coherent fractions at every scale.
\end{proof}

\begin{remark}\label{f:rem-divisor}
The proof has a global divisor reading: all-scale coherence forces the
pole divisor of $f$ on the class plane to be finite, and multiplying by
the corresponding polynomial removes it.  The conclusion resembles the
classical rationality of meromorphic functions on the Riemann sphere,
but the hypothesis is scale coherence on a proper-class plane, not
compactness of a sphere; no compactness, completeness or normal-family
argument is used anywhere in
Theorems~\ref{f:thm-rigidity} and~\ref{f:thm-rationalrigidity}.
\end{remark}

% =====================================================================

\section{Worked examples with exact arithmetic}\label{f:sec-examples}

Each example below is chosen to separate a hypothesis proved earlier: it
exhibits a function satisfying all but one clause of some theorem, or it
computes exactly a quantity that an approximate argument would blur.
Every displayed series is a Hahn sum in the sense of
\ref{a:def:summable}, never a limit of partial sums.  The finite
verifications are collected in \S\ref{f:sub-verify}.

\subsection{Classically divergent series: what is germ-analytic and what
is coherent}\label{f:sub-divergent}

\begin{example}[A germ-analytic function with zero classical
radius]\label{f:ex-factorialsq}
Let $\eta\in\mm$ and put
\[
 A(\eta)=\sum_{n\ge0}(n!)^{2}\eta^{n}.
\]
The family $\bigl((n!)^{2}\eta^{n}\bigr)_{n\ge0}$ is Hahn summable: the
supports lie in $\Neu{\supp\eta}$, which is well ordered, and
Neumann's all-lengths clause gives only finitely many contributing
word lengths at each exponent. This is stronger than merely observing
that the leading valuations increase. The case $\eta=0$ is immediate.  Hence $A$ is a
\emph{germ-analytic} function on $\mm$, and its formally differentiated
series is its actual fine $\varepsilon$--$\delta$ derivative, by the
germ-algebra isomorphism \ref{b:germalgebra}.  At every nonzero
ordinary complex argument $h$ the displayed classical power series
diverges, the ratio of successive absolute values being
$(n+1)^{2}\abs h$.  The germ-analytic class is therefore strictly larger
than the class of Taylor lifts of ordinary analytic germs.  No growth
condition on the coefficients was used, and none is available to be
used.
\end{example}

\begin{example}[A germ with no coherent representative in the original
chart]\label{f:ex-euler}
Let $E(X)=\sum_{n\ge0}n!\,X^{n}$.  Its germ at $0$ is germ-analytic, by
the same computation.  But there is \emph{no} connected ordinary
neighbourhood $U$ of $0$ and \emph{no} $F=\sum_{\gamma\in S}f_{\gamma}
t^{\gamma}\in\Hc(U)$ whose evaluated germ at $0$ is this germ.  Indeed,
equality of analytic germs would give
\[
 \sum_{\gamma\in S}t^{\gamma}\,\frac{f_{\gamma}^{(n)}(0)}{n!}=n!
 \qquad(n\ge0),
\]
and uniqueness of Hahn coefficients forces
$f_{\gamma}^{(n)}(0)=0$ for every $n$ and every $\gamma\ne0$.  Ordinary
analyticity and connectedness of $U$ then make every such $f_{\gamma}$
identically zero, so the exponent-zero coefficient function would be an
ordinary holomorphic function on a neighbourhood of $0$ with Taylor
series $\sum n!\,X^{n}$, of radius zero.  No such function exists.
\end{example}

\begin{remark}[The two statements are about different
charts]\label{f:rem-chartreconciliation}
Example~\ref{f:ex-euler} and the scale-embedding bridge
\ref{c:p4:bridge} look contradictory and are not.  The bridge
asserts that after \emph{one} monomial rescaling every formal series
over $\K$ becomes a coherent family on all of $\CC$; concretely,
\[
 E(tz)=\sum_{n\ge0}n!\,t^{n}z^{n}
\]
is a section of $\Hc(\CC)$ whose Hahn coefficient at the exponent $n$ is
the \emph{monomial} $n!\,z^{n}$, an entire function.  Example
\ref{f:ex-euler} concerns the \emph{original} chart, where the
exponent-zero coefficient function would have to carry all the growth
itself.  The rescaling changes which function is being represented.
Scale is part of the analytic data, and by
Remark~\ref{f:rem-somechart} and Theorem~\ref{f:thm-rigidity} having
some sufficiently small coherent chart carries no implication of a chart
of every larger radius.
\end{remark}

\subsection{An exact infinitesimal Lambert-type root}
\label{f:sub-lambert}

\begin{example}[Lambert-type solution, valid for every
infinitesimal]\label{f:ex-lambert}
Let $\eta\in\mm$ and consider the coherent equation
\[
 z-\eta\,\ev{(e^{z})}=0,
\]
where $e^{z}$ denotes the ordinary entire coefficient function.  The
leading coefficient function is $z$, with a simple zero at the origin
and no other zero in a small disk, so by the zero-cluster theorem
\ref{d:thm:zerocluster} the equation has exactly one root in the halo
$\mm$, counted with multiplicity.  Lagrange--B\"urmann inversion
\ref{b:lagrange} gives that root explicitly:
\begin{equation}\label{f:eq-lambert}
 z=\sum_{n\ge1}\frac{n^{\,n-1}}{n!}\,\eta^{n}
  =\eta+\eta^{2}+\tfrac32\eta^{3}+\tfrac83\eta^{4}
   +\tfrac{125}{24}\eta^{5}+\tfrac{54}{5}\eta^{6}+\cdots .
\end{equation}
The sum is Hahn summable and solves the equation exactly, by formal
substitution, for \emph{every} surcomplex infinitesimal $\eta$,
regardless of the complexity of its positive support.  The identity
\eqref{f:eq-lambert} was verified symbolically through degree $40$.  It
asserts nothing about possible infinite roots of an equation formed with
a globally chosen exponential of \ref{e:thm-twisted}.
\end{example}

\subsection{Preparation at a double root}\label{f:sub-prepexample}

\begin{example}[First preparation coefficients for
$w^{2}-t\,\ev{(e^{w})}$]\label{f:ex-prep}
Take $A(w)=w^{2}-t\,\ev{(e^{w})}$ with $t=\omega^{-1}$, so the leading
coefficient function is $w^{2}$ and $m=2$.  Put
\[
 B(w)=\frac{e^{w}-1-w}{w^{2}},
\]
with its removable value $1/2$ at $w=0$.  The first homogeneous terms of
the preparation recursion \ref{d:thm:preparation}, graded by
positive-support word degree in $t$, are
\[
 p^{[1]}=-t(1+w),\qquad q^{[1]}=-t\,B(w),
\]
and
\[
 p^{[2]}=-t^{2}\Bigl(\tfrac12+\tfrac23 w\Bigr),
\]
because $p^{[1]}q^{[1]}=t^{2}(1+w)B(w)$, whose constant and linear
ordinary Taylor coefficients are $1/2$ and $2/3$.  The prepared monic
polynomial therefore begins
\[
 P(w)=w^{2}-t(1+w)-t^{2}\Bigl(\tfrac12+\tfrac23w\Bigr)+\cdots,
\]
and its two roots in $\K$ are precisely the two roots of
$w^{2}-t\,\ev{(e^{w})}$ in $\mm$, counted with multiplicity.  The
recursion computes them without assuming an a~priori square-root
splitting of the perturbation; $77$ exact coefficients were computed and
checked.
\end{example}

\subsection{A cubic, by both constructive algorithms}\label{f:sub-cubic}

\begin{example}[Perturbed double zero, an escaping third root, and
contour moments]\label{f:ex-cubic}
Let $t=\omega^{-1}$ and
\begin{equation}\label{f:eq-cubic}
 F(Z)=Z^{2}-t+tZ^{3}\;\in\;\Hc(\CC),
\end{equation}
a coherent section with support $\{0,1\}$ and polynomial coefficient
functions.  Three quantities must be kept apart here, and the example
exists to separate them: the algebraic degree of $F$ over $\K$, which is
$3$; the ordinary zero count of the leading coefficient function
$f_{0}(Z)=Z^{2}$ on the disk $\abs Z<1/2$, which is $2$; and the number
of surcomplex zeros of $\ev F$ in the halo of that disk, which by the
zero-cluster theorem \ref{d:thm:zerocluster} is exactly $2$.

\smallskip
\emph{(a) The preparation algorithm.}  Preparation
\ref{d:thm:preparation} gives $F=UW$ with $W$ monic of degree $2$.  The
homogeneous pieces, all obtainable by polynomial division by $Z^{2}$
alone, are
\[
 \begin{gathered}
 p_1=-t,\quad q_1=tZ,\quad p_2=t^2Z,\quad q_3=-t^3,\\
 p_4=-t^4,\quad p_5=2t^5Z,\quad q_6=-2t^6,\quad p_7=-3t^7,
 \end{gathered}
\]
the unlisted earlier pieces being zero.  Assembling them,
\begin{align}
 W(Z)&=Z^{2}+\bigl(t^{2}+2t^{5}+7t^{8}+30t^{11}+\cdots\bigr)Z
      -\bigl(t+t^{4}+3t^{7}+12t^{10}+\cdots\bigr),
      \label{f:eq-cubicW}\\
 U(Z)&=1+tZ-t^{3}-2t^{6}+\cdots .\label{f:eq-cubicU}
\end{align}
Substituting $Z=sw$ with $s=t^{1/2}$ turns the root equation into
$w^{2}(1+s^{3}w)=1$, whose two finite branches give
\begin{equation}\label{f:eq-cubicroots}
 \xi_{\pm}=\pm t^{1/2}-\tfrac12t^{2}\pm\tfrac58t^{7/2}-t^{5}+\cdots .
\end{equation}
The third root of \eqref{f:eq-cubic} over $\K$ is infinite: the sum of
all three roots is $-t^{-1}$, so
\[
 \xi_{\infty}=-t^{-1}+t^{2}+2t^{5}+\cdots .
\]
A higher-degree term of positive valuation can therefore add a root
outside the finite halo without changing the zero count inside it.

\smallskip
\emph{(b) The contour-moment algorithm.}  Let $\Gamma$ be the ordinary
circle $\abs w=1/2$, positively oriented.  The exact identities are
\[
 \frac1{2\pi i}\ointH_{\Gamma}\frac{F'}{F}\,\dd w=2,
 \qquad
 \frac1{2\pi i}\ointH_{\Gamma}w\,\frac{F'}{F}\,\dd w
 =\xi_{+}+\xi_{-}=-t^{2}-2t^{5}-\cdots,
\]
the first an ordinary integer with no infinitesimal error term, as
\ref{d:eq:argument} requires.  The higher moments
$s_{k}=\frac1{2\pi i}\ointH_{\Gamma}w^{k}F'/F\,\dd w=\sum_{j}\xi_{j}^{k}$
together with Newton's identities \ref{d:thm:moments} recover every
coefficient of $W$ in \eqref{f:eq-cubicW}.  Thus the contour detects the
infinitesimal displacement of the cluster's centre even though both
roots have standard part zero, and the two algorithms produce the same
$W$ by independent routes.
\end{example}

\subsection{Scales beyond every finite power}\label{f:sub-scales}

\begin{example}[An exponent beyond all integral powers of
$t$]\label{f:ex-tomega}
Let $t=\omega^{-1}$ and $F(Z)=Z^{2}-t-t^{\omega}$.  Since
$t^{\omega}/t=t^{\omega-1}\in\mm$, the binomial family is Hahn summable
and the roots are \emph{exact}:
\begin{equation}\label{f:eq-twoscale}
 \xi_{\pm}=\pm t^{1/2}\sum_{n\ge0}\binom{1/2}{n}t^{\,n(\omega-1)},
\end{equation}
whose square is $t\bigl(1+t^{\omega-1}\bigr)=t+t^{\omega}$.  The
quantity $t^{\omega}$ is smaller than every positive ordinary integral
power of $t$.  No truncation of \eqref{f:eq-twoscale} at a finite
$t$-order captures its role, and neither does any \emph{sequence} of
such truncations: by \ref{a:thm:discrete} a set-indexed convergent net in
the fine topology is eventually equal to its limit, so there is no
limiting process to appeal to.  The support formulation includes
$t^{\omega}$ without any change to the proofs of preparation or zero
counting.  A variant with the same moral is
$F(Z)=Z^{2}-t-t^{\omega}Z$, whose roots are
\[
 \lambda_{\pm}=\frac{t^{\omega}}{2}
 \pm t^{1/2}\Bigl(1+\tfrac{t^{2\omega-1}}{4}\Bigr)^{1/2},
\]
the square root being expanded at $1$; the root count is two regardless
of the scale separation.
\end{example}

\begin{example}[Two roots in one monad at incomparable
scales]\label{f:ex-incomparable}
Let $\tau=\omega^{-1}$, $\sigma=\omega^{-\omega}$, so
$\sigma\prec\tau^{n}$ for every positive integer $n$, and let
$P(z)=z^{2}-\tau z-\sigma$.  Put $q=\sigma/\tau^{2}\in\mm$.  The
quadratic formula and the binomial expansion give
\[
 z_{+}=\tau\bigl(1+q-q^{2}+2q^{3}-5q^{4}+\cdots\bigr),
 \qquad
 z_{-}=-\tau\bigl(q-q^{2}+2q^{3}-5q^{4}+\cdots\bigr),
\]
so that $v(z_{+})=1$ and $v(z_{-})=\omega-1$.  Both roots lie in the
\emph{same} monad, the monad of $0$, yet one is smaller than every
finite power of the scale of the other.  The zero-cluster theorem
accommodates this with no Archimedean or rank-one restriction on the
value class.
\end{example}

\begin{example}[Two genuinely separated surreal scales, to second
order]\label{f:ex-twoscale}
Let
\[
 \lambda=\omega^{-1/2},\qquad s=\omega^{-\omega},\qquad
 F(z)=z^{2}-\lambda^{2}-s\,\ev{(e^{z})},
\]
so the perturbation scale $s$ is smaller than every positive ordinary
power of $\lambda$.  The leading coefficient function $z^{2}$ gives
exactly two roots in $\mm$.  Set $z=\lambda w$ and divide by
$\lambda^{2}$: the normalized leading coefficient function is $w^{2}-1$,
so one root lies in the monad of $w=1$ and one in the monad of $w=-1$.
Writing $r=\lambda$ or $r=-\lambda$,
\begin{equation}\label{f:eq-twoscaleroot}
 z_{r}=r+s\,\frac{e^{r}}{2r}
       +s^{2}\,\frac{e^{2r}(2r-1)}{8r^{3}}
       \pmod{s^{3}},
\end{equation}
where $e^{r}$ is the infinitesimal Taylor exponential.  To verify this,
put $z=r+sa+s^{2}b$ and use $r^{2}=\lambda^{2}$; the coefficients of
$s$ and $s^{2}$ in $F(z)$ are
\[
 2ra-e^{r},\qquad a^{2}+2rb-e^{r}a,
\]
which vanish exactly for $a=e^{r}/(2r)$ and
$b=e^{2r}(2r-1)/(8r^{3})$.  Higher coefficients are determined the same
way.
\end{example}

\begin{remark}[\eqref{f:eq-twoscaleroot} is an $s$-adic
statement]\label{f:rem-sadic}
The congruence in \eqref{f:eq-twoscaleroot} takes place in
$\CC((\lambda))[[s]]$, embedded in $\K$ by the ordered exponents
$n\omega+k/2$.  It says that the difference lies in the ideal generated
by $s^{3}$.  It is \emph{not} an assertion that the remainder is bounded
by a fixed real multiple of $s^{3}$, and no such Archimedean reading is
available: there is no norm on $\Hc(U)$ anywhere in this article.
\end{remark}

\subsection{Poles off the ordinary plane, and residues}
\label{f:sub-residueexamples}

\begin{example}[A pole that moves off the ordinary complex
plane]\label{f:ex-movingpole}
Let $\varepsilon\in\mm\setminus\{0\}$ and
$H(z)=1/(z-\varepsilon)$.  This function has exactly one actual pole,
at the surcomplex point $\varepsilon$, and
$\Resat{\varepsilon}\bigl(H(z)\,\dd z\bigr)=1$.  On any ordinary circle
about $0$,
\[
 H(z)=\sum_{n\ge0}\varepsilon^{n}z^{-n-1},
\]
a Hahn summable family whose coefficient functions have poles at the
\emph{ordinary} point $0$ of arbitrarily large order; only the term
$n=0$ contributes an ordinary residue there.  Hence
\[
 \ResH{0}H=1,\qquad
 \frac1{2\pi i}\intH_{\abs z=\rho}\frac{\dd z}{z-\varepsilon}=1
 \qquad(\rho>0\text{ ordinary}).
\]
But the \emph{local} residue at the exact point $0$ is
\[
 \Resat{0}\bigl(H(z)\,\dd z\bigr)=0,
\]
because on the fine neighborhood $\varepsilon\mm$ one has the expansion
$H(\zeta)=-\varepsilon^{-1}\sum_{n\ge0}(\zeta/\varepsilon)^{n}$, with no
pole at $\zeta=0$ at all.  So
$\ResH{0}H=1=\Resat{\varepsilon}H$ while $\Resat{0}H=0$: the cluster
residue at an ordinary centre and the local residue at that same point
are numerically different.  This is the simplest instance of Bridge~1,
\ref{d:thm:bridge1} --- the cluster residue equals the finite sum of
actual residues in the monad --- and it is the standing reason the three
residue symbols may never be merged.  Note also that the criterion of
Bridge~2, \ref{d:thm:bridge2}, fails here exactly as it must: the
coefficient attached to $\varepsilon^{n}$ has a pole of order $n+1$ at
$0$, so pole orders are not uniformly bounded in the Hahn exponent, and
$\ResH{0}$ is under no obligation to agree with $\Resat{0}$.
\end{example}

\begin{example}[Split poles with infinite residues and a finite
sum]\label{f:ex-splitpoles}
Let $t=\omega^{-1}$ and
\[
 H(z)=\frac{\ev{(\exp)}(z)}{z^{2}-t},
\]
with the standard unit circle positively oriented.  The poles in the
finite halo are $p_{\pm}=\pm\sqrt t$, with
\[
 \Resat{p_{+}}H=\frac{\ev{(\exp)}(\sqrt t)}{2\sqrt t},
 \qquad
 \Resat{p_{-}}H=-\frac{\ev{(\exp)}(-\sqrt t)}{2\sqrt t}.
\]
Each of these is infinite in modulus, of valuation $-1/2$, yet their sum
is finite:
\[
 \Resat{p_{+}}H+\Resat{p_{-}}H
 =\frac{\ev{(\sinh)}(\sqrt t)}{\sqrt t}
 =\sum_{n\ge0}\frac{t^{n}}{(2n+1)!}.
\]
Therefore
\begin{equation}\label{f:eq-splitresidue}
 \ointH_{\abs z=1}\frac{e^{z}}{z^{2}-t}\,\dd z
 =2\pi i\sum_{n\ge0}\frac{t^{n}}{(2n+1)!},
\end{equation}
where in the integrand $e^{z}$ denotes the ordinary coefficient
function.  Expanding instead $(z^{2}-t)^{-1}=\sum_{n\ge0}t^{n}z^{-2n-2}$
on the contour and taking ordinary residues coefficientwise gives the
same answer.  The example therefore compares two genuinely different
organizations of one exact result: split poles at half-integral scales,
summed as $\Resat{\cdot}$, and an integral-power Hahn expansion at the
ordinary centre, summed as $\ResH{0}$.  Their agreement is Bridge~1,
not a coincidence, and the cancellation of the two infinite summands is
exact rather than asymptotic.
\end{example}

\subsection{A Cauchy formula at an infinite centre}
\label{f:sub-infinitecentre}

\begin{example}[Infinite centre, hyper-small radius, surcomplex
coefficient]\label{f:ex-infcentre}
Choose $a=\omega+i\omega^{2}$ and $s=\omega^{-\omega}$, and let
\[
 F(a+su)=\ev{(\sin)}(u)+t^{\sqrt2}\,\ev{(e^{u})},
 \qquad u\in\halo{B(0,2)},
\]
a coherent datum in the chart of Definition~\ref{f:def-chart} with
support $\{0,\sqrt2\}$.  For every infinitesimal $\eta$, the affine
circle $\Gamma=a+s\{u:\abs u=1\}$ satisfies
\[
 F(a+s\eta)=\frac1{2\pi i}\intH_{\Gamma}
   \frac{F(z)}{z-(a+s\eta)}\,\dd z,
\]
the right-hand side being the coefficientwise integral
\eqref{f:eq-chartintegral} in the stated chart.  The example combines an
infinite centre, a radius smaller than every finite power of
$\omega^{-1}$, and a genuinely surcomplex coefficient $t^{\sqrt2}$,
none of which the theory treats as exceptional.  A companion at an
ordinary centre with a displaced evaluation point is
\[
 \frac1{2\pi i}\intH_{\abs z=1}
 \frac{e^{z}+t/(2-z)}{z-a'}\,\dd z
 =e^{1/3}\sum_{n\ge0}\frac{\varepsilon^{n}}{n!}
  +\frac{t}{5/3-\varepsilon},
 \qquad a'=\tfrac13+\varepsilon,\ \varepsilon\in\mm,
\]
exact because the joint support is well ordered and each exponent
receives only finitely many contributions.
\end{example}

\subsection{No growth hypothesis on the coefficient family}
\label{f:sub-nogrowth}

\begin{example}[Arbitrary factorial growth in the scale
index]\label{f:ex-nogrowth}
Let
\[
 F(Z)=Z^{2}-t+\sum_{n\ge2}(n!)^{2}t^{n}e^{nZ}.
\]
This is a member of $\Hc(\CC)$: its support is a subset of $\NN$, hence
well ordered with finite multiplicity at each exponent, and every
coefficient function is entire.  On any ordinary circle centred at zero
the leading coefficient function $Z^{2}$ is nonzero, so $\ev F$ has
exactly two zeros in the corresponding halo, counted with multiplicity,
and its preparation polynomial and contour moments exist as Hahn-summed
surcomplex coefficients.  No positive ordinary radius of convergence in
a numerical parameter replacing $t$ is required, and no boundedness of
the family $\bigl((n!)^{2}\bigr)_{n}$ is used.

Contrast this with Example~\ref{f:ex-euler}.  Growth in the \emph{scale}
index $\gamma$ is entirely unrestricted, because the summability
condition is about supports and not about magnitudes.  Growth in the
ordinary \emph{Taylor} index of a single coefficient function is exactly
what the obstruction of Example~\ref{f:ex-euler} forbids, because each
$f_{\gamma}$ must be an ordinary holomorphic function on a common
ordinary domain.
\end{example}

\subsection{The permitted and forbidden support boundary}
\label{f:sub-support}

\begin{example}[An arbitrary well-ordered family of perturbation
scales]\label{f:ex-support}
Let $E\subset\No^{>0}$ be \emph{any} well-ordered set, let
$(c_{\gamma})_{\gamma\in E}$ be arbitrary ordinary complex constants,
and let $m\ge1$.  Then
\[
 F(z)=z^{m}+\sum_{\gamma\in E}t^{\gamma}\,e^{c_{\gamma}z}
\]
is a coherent section on the finite plane, and $\ev F$ has exactly $m$
roots in $\mm$, counted with multiplicity.  No boundedness condition on
the family $(c_{\gamma})$ is needed or available.  The positivity and
well ordering of $E$, however, are essential: a decreasing sequence of
exponents tending to $0$ is not a permitted Hahn support.  Thus
$\{1/n:n\ge1\}$ is excluded, while the increasing set
$\{1-1/n:n\ge2\}$ is a perfectly legal support.  The separating
elementary example is the one recorded with
\ref{a:def:summable}: $\sum_{n\ge0}2^{-n}$ is \emph{not} Hahn summable,
because infinitely many of its terms sit at the single exponent $0$;
its ordinary complex sum must be computed first and used as a single
coefficient.
\end{example}

\subsection{What the verification script checks, and what it does not}
\label{f:sub-verify}

The nine preserved companion scripts together perform the finite
verifications cited above, in exact rational and symbolic arithmetic, with no floating
point.  It checks: the Lambert-type series \eqref{f:eq-lambert} and its
substitution into the defining equation through degree $40$; the
preparation identity of Example~\ref{f:ex-prep} in a finite bivariate
truncation, with $77$ exact coefficients; the preparation recursion for
the cubic \eqref{f:eq-cubic}, the cancellation of the product error
$F-UW$ to the requested order, the root expansions
\eqref{f:eq-cubicroots}, and the low-order contour moments; the
multiscale binomial identity \eqref{f:eq-twoscale}; the second-order
two-scale root residual \eqref{f:eq-twoscaleroot}; and the split-pole
residue cancellation \eqref{f:eq-splitresidue}.

\begin{remark}[Standing disclaimer on the finite
checks]\label{f:rem-verifyscope}
These calculations test signs, factors and coefficients in finite
truncations of displayed formulas.  They certify nothing about the
class-theoretic foundations, nothing about the summability arguments,
and nothing about any unrestricted theorem statement; in particular they
are not a machine-checked development, and no proof assistant has been
used.  The all-order and general-support arguments are the proofs in the
body, not the script.
\end{remark}

\begin{remark}[What is constructive, and in what
sense]\label{f:rem-constructive}
Preparation and reversion are constructive whenever the input supplies
the coefficient functions together with an effective presentation of the
relevant supports.  A finite computation is organized by
positive-support word degree: to reach degree $N$ one runs the
preparation recursion for $1\le n\le N$, and at each degree only finite
products and ordinary Taylor division operators occur; with finitely
many formal perturbation generators, the degree-$n$ coefficient is a
finite sum of analytic expressions indexed by binary trees.  For a
completely arbitrary surreal support the existence theorem is \emph{not}
a promise that all coefficients can be enumerated by a finite program.
Moreover, in higher rank a bound on the valuation need not leave only
finitely many terms, so an algorithm must specify a finite support
region or an explicit truncation scheme, rather than merely saying
``compute below precision $t^{\delta}$.''
\end{remark}

% =====================================================================

\section{Boundaries: what is and is not established}
\label{f:sec-boundaries}

The qualifications collected here are results, not defensive padding.
Each of them is the exact hypothesis that separates a true statement
from a false one, and in several cases a counterexample in the body
shows that dropping it produces a false theorem.  They are stated
per-theorem in the body and consolidated here so that no reader has to
reconstruct the pattern.

\subsection{Which class each classical principle lives in}
\label{f:sub-hypchart}

In the table, \emph{germ-analytic} means a Hahn summable power series at
one point with arbitrary surcomplex coefficients; \emph{coherent} means
a member of $\Hc(U)$ with the stated ordinary domain conditions;
\emph{lift} means the canonical classical lift $\ev f$ of an ordinary
holomorphic $f$, that is, a coherent section supported at the exponent
$0$; and \emph{all-scale} means Definition~\ref{f:def-allscale}.

\begin{center}\small
\begin{longtable}{@{}>{\raggedright\arraybackslash}p{0.215\textwidth}
                    >{\raggedright\arraybackslash}p{0.245\textwidth}
                    >{\raggedright\arraybackslash}p{0.44\textwidth}@{}}
\toprule
\textbf{Classical principle} & \textbf{Class in which it is true} &
\textbf{The qualification that cannot be dropped}\\
\midrule
\endfirsthead
\toprule
\textbf{Classical principle} & \textbf{Class in which it is true} &
\textbf{The qualification that cannot be dropped}\\
\midrule
\endhead
\bottomrule
\endfoot
Power-series realization and Taylor calculus &
Germ-analytic &
A sufficiently small scale; Hahn sums, not sequential convergence.  No
radius-of-convergence formula exists from ordinary coefficient
magnitudes.\\[3pt]
Cauchy--Riemann equations &
Germ-analytic &
The equations on a fine neighbourhood.  Existence of two partial
derivatives at a point does not imply complex differentiability.\\[3pt]
Inverse, implicit and ramification theorems &
Germ-analytic; coherent local inverse in one variable &
Invertible Jacobian for the germ inverse/implicit theorem; finite order
for ramification. The coherent one-variable inverse is obtained after
a scale change. No multivariable coherent implicit theorem is proved here.\\[3pt]
Open mapping and local maximum modulus &
Germ-analytic locally; coherent globally &
Valid at arbitrary surcomplex points, via the monomial normalization
radius $r=t^{\rho}$.  In global form both fail for merely
germ-analytic functions, witnessed by $L$ of
the locally constant witness of Section~\ref{b:part2}.\\[3pt]
Identity theorem &
Coherent, $U$ connected &
Fails in the germ-analytic class: $L$ vanishes on a nonempty fine-open
class and is not zero.  Accumulation is of \emph{standard parts} in
$\CC$; no set-sized fine-convergent sequence is claimed.\\[3pt]
Cauchy, Morera, primitives &
Coherent, with $\intH$ &
Integrals are coefficientwise, on ordinary contours or in transported
affine charts.  Uniqueness of a primitive up to a constant holds
\emph{inside the coherent category only}.\\[3pt]
Preparation, division, Rouch\'e &
Coherent &
Leading coefficient function holomorphic near $\overline\Omega$ and
nonzero on $\Gamma=\partial\Omega$.  Rouch\'e's strict inequality is an
ordinary complex margin, or carries an ordinary real margin $q<1$.\\[3pt]
Argument principle, exact zero counts, Bridge~1 &
Quotients of coherent families &
Nonvanishing boundary reduction; finite zero clusters.  The index is an
ordinary integer with \emph{no} infinitesimal error term.\\[3pt]
Coefficientwise residue theorem &
$\Mc(U)$: Hahn series with meromorphic coefficients &
Poles may move with the exponent and infinitely many ordinary centres
may contribute.  \emph{No} conclusion about actual surcomplex poles
transfers to this class.\\[3pt]
Liouville &
See \S\ref{f:sub-threehierarchies} &
Three inequivalent hypotheses with three different conclusions; a single
surreal bound on $\halo\CC$ is vacuous.\\[3pt]
Removable singularities &
Coherent &
Coefficientwise removal is necessary and sufficient; a scalar bound is
insufficient ($t/Z$, $t\,\ev{(e^{1/z})}$).  A punctured halo deletes a
whole \emph{monad}, not a point.\\[3pt]
Schwarz lemma and Schwarz--Pick &
Lifts $\ev f$ &
False for general coherent self-maps of $\halo{\mathbb D}$: $(1+t)z$
fixes $0$, maps $\halo{\mathbb D}$ into itself, and has
$\abs{F'(0)}=1+t>1$.  The same-monad case needs the divided
difference.\\[3pt]
Riemann mapping &
Lifts $\ev f$ &
Halos of ordinary simply connected proper domains only.  The existence
input is the \emph{ordinary} Riemann mapping theorem; no uniformization
of arbitrary surcomplex open classes is claimed.\\[3pt]
Montel / normal families &
Coherent with \emph{countable} support &
Convergence is a topology on coefficient data, not pointwise
all-radii convergence of values.  Countability is not cosmetic.\\[3pt]
Global logarithm &
Germ-analytic on $\K^{\times}$ &
Exists globally, with branch chosen by a fine-locally constant
quantity; it is \emph{not} a coherent primitive of $1/z$ on
$\halo{(\CC^{\times})}$.\\[3pt]
Classification of entire and meromorphic functions &
All-scale coherent &
Exactly the polynomials, respectively the rational functions.  This
rigid category excludes $\Exp$ and every twist $E_{\lambda}$.\\
\end{longtable}
\end{center}

\subsection{The analytic categories, and what each one buys}
\label{f:sub-categories}

\begin{center}\small
\begin{longtable}{@{}>{\raggedright\arraybackslash}p{0.21\textwidth}
                    >{\raggedright\arraybackslash}p{0.30\textwidth}
                    >{\raggedright\arraybackslash}p{0.40\textwidth}@{}}
\toprule
\textbf{Category} & \textbf{Defining structure} &
\textbf{Conclusions and restrictions}\\
\midrule
\endfirsthead
\toprule
\textbf{Category} & \textbf{Defining structure} &
\textbf{Conclusions and restrictions}\\
\midrule
\endhead
\bottomrule
\endfoot
Differentiable &
A bare $\varepsilon$--$\delta$ complex difference quotient over all
positive surreal radii; \emph{no} power-series requirement &
Uniqueness of the derivative and the sum, product and chain rules, and
the counterexamples. No regularity theorem asserting a power-series
representation from bare differentiability is proved here.\\[3pt]
Germ-analytic &
A Hahn summable power series at one point with arbitrary surcomplex
coefficients, realized near the centre; no coefficient growth condition &
The full formal Taylor calculus, including classically divergent
series; the germ algebra at any point is $\K[[X]]$ as a differential
$\K$-algebra. Local identity, openness at nonconstant germs and primitive
uniqueness hold; the corresponding unrestricted global propagation fails.\\[3pt]
Canonical lift $\ev f$ &
Taylor extension of one ordinary holomorphic $f$; a coherent section
supported at exponent $0$ &
Classical continuation; the sharp Schwarz lemma and Schwarz--Pick
theorem, which are false one level up.  Creates no infinitesimally
displaced zeros, so surcomplex coefficients are necessary for a genuine
perturbation.\\[3pt]
Coherent, $\Hc(U)$ &
One common connected ordinary domain carrying one ordinary holomorphic
coefficient function per exponent, with well-ordered set support &
Faithful evaluation on $\halo U$, identity, Cauchy and residue
calculus, preparation and division, exact zero counts, open mapping and
local inverse theorems, coherent primitives.  This is not
sheaf-coherence in the sense of complex geometry.\\[3pt]
All-scale coherent &
A coherent representative on a fixed ordinary disk after \emph{every}
positive surreal dilation, supports free to vary with the scale &
Exactly the finite-degree polynomials
(\ref{f:thm-rigidity}); exactly the rational functions in the
meromorphic version (\ref{f:thm-rationalrigidity}).  A rigidity
condition, not a proposed definition of ``entire''.\\[3pt]
Canonical global exponential and its twists &
The real surreal exponential plus a normal-form phase convention on the
purely infinite imaginary part &
A surjective germ-analytic homomorphism with
$\ker\Exp=2\pi i\,\Oz=i(\Jpi+2\pi\ZZ)$; locally coherent in finite
translated charts, but not all-scale coherent, by
Remark~\ref{f:rem-infinitebeta}.\\
\end{longtable}
\end{center}

The inclusions are
\[
 \{\ev f\}\;\subsetneq\;\Hc(U)\;\subsetneq\;
 \{\text{germ-analytic}\}\big|_{\halo U}\;\subseteq\;
 \{\text{differentiable}\}.
\]
Only the first two strict inclusions are established here: positive
support gives a section that is not a lift, and Example~\ref{f:ex-euler}
gives the obstruction to coherence. The locally constant witness $q$
of Example~\ref{b:q} gives a function on a whole halo. Strictness of
the last inclusion is not proved; $L$ is germ-analytic and cannot
witness it.  The restriction of an all-scale
Hahn-entire function to any bounded halo is coherent there, and that
containment is strict as well, since $\Exp$ restricted to a finite halo
is coherent while $\Exp$ is not all-scale coherent.

\subsection{The three hierarchies that must never be flattened}
\label{f:sub-threehierarchies}

\paragraph{Three function classes.}
The identity theorem, the global maximum principle, Liouville's theorem,
uniqueness of primitives and the open mapping theorem have coherent
versions under the stated domain hypotheses; Liouville in particular
requires bounded entire coefficient functions. Their unrestricted global
forms fail in the germ-analytic class.  The sharp Schwarz and
Schwarz--Pick inequalities hold for lifts $\ev f$ and \emph{fail} in
$\Hc(\mathbb D)$.  Results in the germ-analytic class do not descend to
$\Hc(U)$: $\sum n!\,\varepsilon^{n}$ is a legitimate germ
(Example~\ref{f:ex-euler}) and is not the restriction of any member of
$\Hc(U)$ on an ordinary neighbourhood of $0$ with those derivatives.
The one bridge that has been proved, and that must be invoked
explicitly whenever a result is moved between the two derivatives, is
that for a coherent section the coefficientwise derivative equals the
fine $\varepsilon$--$\delta$ derivative.

\paragraph{Three residues.}
$\Res$ is formal, the coefficient of $X^{-1}$ in $\K((X))$ with a finite
principal part, at any surcomplex centre.  $\ResH a$ is the cluster
residue at an \emph{ordinary} centre $a$, the Hahn sum of the ordinary
residues of the coefficient functions.  $\Resat b$ is the actual local
residue at a surcomplex point $b$.  They are not interchangeable:
Example~\ref{f:ex-movingpole} gives $\ResH 0=1$ and $\Resat 0=0$ for the
same function at the same point.  Bridge~1 identifies $\ResH a$ with the
finite sum of the $\Resat{\cdot}$ over the actual poles in $a+\mm$, for
quotients of coherent families only; Bridge~2 identifies $\ResH a$ with
$\Res$ at the exact point $a$ when pole orders are uniformly
bounded in the Hahn exponent. This is a sufficient condition, not a
necessary condition for numerical equality.  On the larger class $\Mc(U)$ there need
be no actual pole at all, so $\Resat{\cdot}$ has no meaning there and
the finite-sum conclusion must not be imported.

\paragraph{Three Liouville hypotheses.}
(a) A single \emph{surreal} bound on the finite halo is vacuous: every
nonzero $F\in\Hc(\CC)$ already satisfies
$\abs{\ev F}<t^{\vH(F)-1}$ there.  (b) A single \emph{ordinary real}
bound is exactly characterized and does \emph{not} give constancy:
$\abs{\ev F(z)}\le M$ on the finite plane for an ordinary real $M$ if
and only if $F=c+R$ with $c\in\CC$ and $R$ of positive support, so the
algebra of real-bounded macroscopically entire functions is precisely
$\CC\oplus\Hc_{>0}(\CC)$.  The function $F=tz$ is real-bounded by $1$ on
the whole finite plane and is not constant; it is $0+tz$, exactly as
that characterization predicts.  (c) \emph{Coefficientwise} bounds give
constancy, and coefficientwise polynomial growth of degree $d$ gives a
polynomial of degree at most $d$.  Importing (c)'s conclusion under
hypothesis (b) or (a) is false. Conversely the constant $F=\omega$
satisfies coefficientwise boundedness and has no ordinary-real bound;
(b) and the boundedness case of (c) are incomparable. Note finally that the counterexamples
$tz$ and $\varepsilon Z$ are a different obstruction from the locally
constant $L$: they are coherent, on the halo of the ordinary plane,
whereas $L$ is germ-analytic on all of $\K$ and restricts to the
coherent constant $1$ on each finite halo.

\paragraph{Two Cauchy estimates, and two logarithms.}
One must not read the warning that a single coarse surreal bound does
not control every \emph{coefficient function} --- which is true, and
witnessed by $F(w)=w(1+it)$ --- as forbidding the $\mathrm{ML}$
inequality, which bounds the surreal \emph{modulus} of a derivative and
is proved from the positivity of the coefficientwise integral.  The two
are about different quantities and both are retained; see
Remark~\ref{f:rem-whichestimate}.  Similarly, the global fine-analytic
logarithm on $\K^{\times}$ and the principal logarithm on the cut plane
$\K\setminus\No_{\le0}$ are \emph{different functions}, differing by
exactly $2\pi i$ on the lower half of the monad of every negative ordinary real;
both are legitimate branches and both are named separately.

\subsection{A practical proof checklist}\label{f:sub-checklist}

A proposed surcomplex analytic argument should identify, explicitly:
\begin{enumerate}[label=\textup{(\arabic*)},leftmargin=2.6em]
\item its coefficient field, and whether the work happens inside a
 set-sized $K_{\Gamma}$ or in all of $\K$;
\item its set support, and the well-ordering and finiteness clauses that
 make it one;
\item its summability convention, and where the finiteness clause is
 used;
\item the topology appearing in any limit, and whether that limit is a
 Hahn sum rather than a limit of partial sums;
\item the coherence mechanism connecting its local formulas, or the
 explicit absence of one;
\item for a contour argument, whether integration is ordinary,
 coefficientwise, definable, or a separately constructed surreal
 integral;
\item when counting zeros, whether multiplicities refer to ordinary
 leading coefficient functions, to actual surcomplex roots, or to
 formal polynomial factors --- the prepared setting connects all three
 exactly;
\item which of $\Res$, $\ResH a$, $\Resat b$ is meant;
\item which of the three Liouville hypotheses is being assumed;
\item which of the two logarithms is in use.
\end{enumerate}

\subsection{Assertions deliberately not made}\label{f:sub-notmade}

The following distinctions delimit the proved statements; none of the
stronger assertions they exclude is being assumed.

\begin{itemize}[leftmargin=1.4em]
\item There is no radius-of-convergence formula based on ordinary
 complex coefficient magnitudes for arbitrary surcomplex coefficients.
\item A Hahn sum is not a fine-topological limit of its ordinary partial
 sums.
\item Agreement at ordinary points does not determine a germ-analytic
 function.
\item A real-parameter contour is not a fine-continuous path: every
 continuous map $[0,1]\to\K$ is constant, so an ordinary circle is not
 a path in $\K$.
\item A bound in the ordered field does not bound every coefficient
 function.  (This does not forbid the $\mathrm{ML}$ inequality; see
 \S\ref{f:sub-threehierarchies}.)
\item The canonical single-valued exponential on all of $\K$ is
 constructed using a specified infinite-phase convention, but $E'=E$, $E(0)=1$, agreement with $\exp_{\No}$ and
 $\exp_{\CC}$, and the modulus law do not determine the infinite
 imaginary phase.
\item No unrestricted global contour integral for arbitrary class
 functions is constructed, and the coefficientwise functional $\intH$ is
 not proposed as the unique notion of surreal integration.
\item No universal analytic continuation across scales is claimed, and
 no unrestricted gluing or chart-independence theorem across surreal
 centres and scales is asserted (Remark~\ref{f:rem-nogluing}).
\item No general theory of essential singularities or Picard theorem
 is proved for arbitrary germ-analytic functions. The specific
 value-distribution result for $\Exp(1/z)$ in \ref{e:thm-picard}
 also holds for the twists $E_\lambda(1/z)$, as proved in
 Remark~\ref{e:rem-picardscope}.
\item No composition law compatible with the Berarducci~\cite{BMder}--Mantova~\cite{Mantova2026b} field
 derivation is claimed on $\No$; and the variable derivative $D_{z}$,
 which fixes every surreal coefficient, is provably incompatible with a
 field derivation normalized by $\partial\omega=1$
 (\ref{b:derivation-obstruction}).
\item Nothing asserts that a maximum of the modulus is \emph{attained}
 on a closed surcomplex disk.  The missing compactness is precisely why
 the maximum principle yields no extremal-value theorem.
\item Nothing asserts the existence of a nontrivial set-sized
 fine-convergent sequence; the discreteness theorem forbids it.
\item The discreteness, realization and rigidity theorems are theorems
 about $\K$, not about an arbitrary fixed set-sized Hahn subfield with
 only its internal radii: their proofs use positive surreal bounds that
 need not belong to such a subfield.  Nor does the fine topology induced
 on a nontrivially valued set-sized Hahn subfield coincide with that
 subfield's own valuation topology (the trivially valued case is
 discrete for both).
\item No uniformization of arbitrary surcomplex open classes is claimed;
 the Riemann mapping corollary takes the ordinary theorem as input.
\end{itemize}

\subsection{A concrete summability boundary at an essential
singularity}\label{f:sub-essential}

It is worth exhibiting exactly where the theory stops, rather than
merely saying that it does.

The canonical lift $\ev{(e^{1/z})}$ is coherent on
$\halo{(\CC^{\times})}$, but that domain excludes the entire monad of
$0$.  At a nonzero infinitesimal $z$ the formal expression
\[
 \sum_{n\ge0}\frac{z^{-n}}{n!}
\]
is \emph{not} Hahn summable: its term valuations $-n\,v(z)$ form a
strictly decreasing sequence, so their union is not well ordered.  Thus
the ordinary Laurent series does not itself define any extension into
that monad.

One may instead choose a global exponential $E_{\alpha}$ of
\ref{e:thm-twisted} and form $E_{\alpha}(1/z)$.  This is a germ-analytic
extension on all of $\K^{\times}$, but its values depend on the phase
convention: at $z=-i/\omega$ one has $1/z=i\omega$, and the values for
$\alpha=0$ and $\alpha=\pi$ differ, since $E_{0}(i\omega)=1$ while
$E_{\pi}(i\omega)=-1$.  Restriction to ordinary nonzero complex points
and the ordinary Laurent series do not remove this ambiguity.

Accordingly, what has been proved is a finite-order meromorphic theory
and a coherent contour theory, not a universal classification of
essential singularities on all surreal scales.  A Casorati--Weierstrass
or Picard statement must specify its domain, its topology and its
extension operator before it is even a well-posed target.  The
value-distribution theorem \ref{e:thm-picard} is an exact statement
about a specified exponential and, by Remark~\ref{e:rem-picardscope},
its explicit twists. It is not a general Picard theorem.

\subsection{Research directions, with their prerequisites}
\label{f:sub-directions}

These are proposals for further development, not theorems proved here,
and not claims that their broad subjects are unstudied.

\paragraph{Non-affine coherent atlases.}
The affine chart construction of \S\ref{f:sub-charts} is completely
explicit.  The next task is to specify classes of nonlinear chart
transitions preserving local Hahn support control, and to prove that
their coefficientwise contour functionals agree on admissible overlaps.
The overlap conditions must be formulated \emph{before} asserting any
theory of surcomplex Riemann surfaces: arbitrary topological gluing
permits the locally constant counterexamples and is too weak.  The
difficult point is not local differentiation, which already works, but
compatibility between different summation descriptions over a common
ordinary base.  Theorem~\ref{f:thm-rigidity} shows the direction is
forced: a global transcendental theory must weaken all-scale coherence,
because retaining it admits only polynomials even when the support
is allowed to vary with the scale.

\paragraph{Relative preparation in several variables.}
The local realization, Cauchy--Riemann, inverse and implicit arguments
already work in finitely many variables. The coherent inverse theorem
proved here is one-dimensional; a multivariable coherent
implicit-function theorem requires an additional statement and proof
over a specified common ordinary parameter domain
(Remark~\ref{b:coherent-implicit-elsewhere}).  A relative preparation theorem,
with one distinguished polynomial variable and coherent parameter
domains, requires a coefficientwise division operator that preserves the
chosen parameter domain, together with one common parameter domain for
all coefficient functions.  The one-variable Hermite division used here
is not by itself a proof of that generalization; the positive-support
recursion identifies exactly where those inputs would enter.

\paragraph{Broader meromorphic and cohomological theories.}
The residue theory covers quotients of coherent families with a
controlled boundary reduction, and $\Mc(U)$ with moving poles.  Allowing
arbitrary meromorphic coefficient families, essential singularities or
infinitely many poles requires explicit locally finite pole and support
conditions.  Cohomological statements would require an admissible-cover
or coefficient-domain formalism; the ordinary connectedness of the fine
topology supplies nothing, since $\K$ is totally disconnected.

\paragraph{Essential singularities and resummation.}
A natural problem is to determine which resurgent complex functions
admit sectorial surcomplex extensions with compatible residue and
monodromy data.  Allowing two-sided formal expansions requires new
summability hypotheses: a support harmless in one variable may cease to
be well ordered after evaluation, as \S\ref{f:sub-essential} shows.
Equality of formal Taylor coefficients is not sufficient asymptotic
data.

\paragraph{Normal families and approximation.}
These should be formulated in a coefficient topology, or in an
explicitly fixed Hahn workspace.  In the full fine topology every
set-indexed convergent net is eventually equal to its limit, so an
ordinary sequential Montel theorem would be the wrong statement; the
coefficientwise Montel theorem \ref{e:thm-montel} is the correctly
shaped one, and its countable-support hypothesis is not cosmetic.
Quantitative control of coefficient growth and of support changes is the
appropriate starting point.

\paragraph{Effective support presentations.}
For a finitely generated positive support, the preparation recursion and
the contour moments give implementable exact algorithms up to a
specified finite support cutoff.  Termination at a requested cut is a
separate algorithmic question, and in higher rank a valuation bound need
not leave finitely many terms.  General surreal coefficients are not
finite input data, and the unrestricted construction should not be
called an executable algorithm without first specifying an effective
subfield and a support language.

\subsection{Status of the arguments, and relation to existing work}
\label{f:sub-status}

The following are used as established inputs and are cited rather than
claimed: the surreal cut property and the normal-form theorem; real
closedness of $\No$ and the algebraic closedness of $\K$ together with
the Hahn-field results that supply it; Neumann's support lemma; Alling~\cite{Alling}'s
foundations of analysis over surreal number fields and the restricted
analytic structure of van den Dries and Ehrlich~\cite{vdDE}; the existence, group
law and kernel of the canonical surcomplex exponential, due to Ehrlich
and Kaplan~\cite[Theorem~11.1]{EK}; and ordinary complex analysis on the ordinary domains, which
enters only through the coefficient functions and the canonical lifts,
never through an unstated transfer principle to all of $\K$.

The coefficientwise contour functional is not asserted to coincide with
every integration operator in the surreal integration literature.  Where
two constructions agree coefficientwise on a common Hahn expansion,
their agreement can be checked directly; outside that overlap an
additional comparison theorem would be needed, and none is proved here.
Likewise, work on transseries as germs of surreal functions, and on
Taylor expansions over generalized power series, addresses composition
and derivation questions that should not be conflated with the fact that
a formal series in an independent variable has a sufficiently small
surreal realization.

This article is a mathematical research exposition with proofs.  It has
not been externally refereed, and neither it nor its foundational inputs
have been machine-checked in a proof assistant; the finite checks of
\S\ref{f:sub-verify} have been executed and are not a substitute for
that.  No exhaustive novelty or priority determination has been made for
any formulation, and none is asserted.
\clearpage
\section*{References}
\addcontentsline{toc}{section}{References}

This list is the union of the reference lists of the nine source
manuscripts: 94 bibliography entries naming 44 distinct works.  It is
reproduced in full so that the merged document does not lose the
literature those manuscripts consulted, and the comment on each entry
records which of the nine cited it.
Only some are cited in the text above; the others were consulted by one or
more of the sources for background. The historical provenance record
is \texttt{MERGE\_NOTES.md}; it is not a current certification of every
bibliographic or mathematical comparison made by those sources.

\begin{thebibliography}{99}

%% Merged from the bibliographies of the nine source manuscripts:
% 94 bibitem entries, 44 distinct works.  The supplied expository
%% manuscript is not listed here: it is Part I of this document.

\bibitem{Gonshor}
H. Gonshor, \emph{An Introduction to the Theory of Surreal Numbers}, London Mathematical Society Lecture Note Series 110, Cambridge University Press, 1986. \href{https://doi.org/10.1017/CBO9780511629143}{doi:10.1017/CBO9780511629143}.
% cited by source manuscripts 01, 03, 04, 05, 06, 08, 09

\bibitem{Alling}
N. L. Alling, \emph{Foundations of Analysis over Surreal Number Fields}, North-Holland Mathematics Studies 141, North-Holland, 1987. \href{https://shop.elsevier.com/books/foundations-of-analysis-over-surreal-number-fields/alling/978-0-444-70226-5}{Publisher's description and bibliographic record}.
% cited by source manuscripts 01, 04, 05, 06, 08, 09

\bibitem{Neumann}
B. H. Neumann, ``On ordered division rings'', \emph{Transactions of the American Mathematical Society} \textbf{66} (1949), 202--252. \href{https://doi.org/10.1090/S0002-9947-1949-0032593-5}{doi:10.1090/S0002-9947-1949-0032593-5}.
% cited by source manuscripts 01, 03, 05, 06, 07, 08

\bibitem{MS}
J. M\"uller and A. Strohmaier, The theory of Hahn-meromorphic functions, a holomorphic Fredholm theorem, and its applications, \emph{Analysis \& PDE} \textbf{7} (2014), no.~3, 745--770. \href{https://doi.org/10.2140/apde.2014.7.745}{doi:10.2140/apde.2014.7.745}. \href{https://arxiv.org/abs/1205.0236v3}{arXiv:1205.0236v3}.
% cited by source manuscripts 03, 05, 06, 08, 09

\bibitem{RSS}
S. Rubinstein-Salzedo and A. Swaminathan, ``Analysis on surreal numbers,'' \emph{Journal of Logic and Analysis} \textbf{6}(5) (2014), 1--39. Preprint and revised version: \url{https://arxiv.org/abs/1307.7392}.
% cited by source manuscripts 01, 05, 06, 08, 09

\bibitem{BMgermsc}
A. Berarducci and V. Mantova, \emph{Transseries as germs of surreal functions}, Transactions of the American Mathematical Society \textbf{371} (2019), 3549--3592. \href{https://doi.org/10.1090/tran/7428}{doi:10.1090/tran/7428}; \arxiv{1703.01995}. In particular, Section 3 on summation and substitution.
% cited by source manuscripts 04, 05, 06, 08

\bibitem{Conway}
J. H. Conway, \emph{On Numbers and Games}, 2nd ed., A K Peters, 2001. Foundational source for surreal arithmetic, normal forms, and surcomplex numbers.
% cited by source manuscripts 01, 03, 05, 09

\bibitem{Ahlforsb}
L. V. Ahlfors, \emph{Complex Analysis}, 3rd ed., McGraw--Hill, 1979. Reference for the ordinary Cauchy, identity, Morera, residue, Rouch\'e, maximum-modulus, Liouville, and Schwarz theorems used on coefficients.
% cited by source manuscripts 04, 05, 09

\bibitem{BMder}
A. Berarducci and V. Mantova, ``Surreal numbers, derivations and transseries'', \emph{Journal of the European Mathematical Society} \textbf{20} (2018), 339--390. \href{https://doi.org/10.4171/JEMS/769}{doi:10.4171/JEMS/769}; \href{https://arxiv.org/abs/1503.00315}{arXiv:1503.00315}.
% cited by source manuscripts 05, 06, 08

\bibitem{CEb}
O. Costin and P. Ehrlich, ``Integration on the surreals'', \emph{Advances in Mathematics} \textbf{452} (2024), article 109823. \href{https://arxiv.org/abs/2208.14331}{arXiv:2208.14331v5}.
% cited by source manuscripts 05, 08, 09

\bibitem{Mantova2026b}
V. Mantova, ``Monotonicity and a Taylor approximation theorem for transseries'', January 2026, \href{https://arxiv.org/abs/2601.07747}{arXiv:2601.07747}.
% cited by source manuscripts 04, 05, 08

\bibitem{vdDE}
L. van den Dries and P. Ehrlich, Fields of surreal numbers and exponentiation, \emph{Fundamenta Mathematicae} \textbf{167} (2001), 173--188. In particular, Sections 1--2. \url{https://www.impan.pl/shop/en/publication/transaction/download/product/89226}
% cited by source manuscripts 01, 03, 09
\bibitem{Ahlfors}
Lars V. Ahlfors, \emph{Complex Analysis}, 3rd ed., McGraw--Hill, 1979. Classical coefficient-level Cauchy, Schwarz--Pick, harmonic, and residue theorems.
% cited by source manuscripts 02, 07

\bibitem{Allingb}
Norman L. Alling, \emph{Foundations of Analysis over Surreal Number Fields}, North-Holland Mathematics Studies, vol.~141, North-Holland, 1987. See also ``Fundamentals of analysis over surreal number fields,'' \emph{Rocky Mountain Journal of Mathematics} \textbf{19} (1989), 565--573, \href{https://doi.org/10.1216/RMJ-1989-19-3-565}{doi:10.1216/RMJ-1989-19-3-565}.
% cited by source manuscripts 02, 07

\bibitem{BMderiv}
A. Berarducci and V. Mantova, \emph{Surreal numbers, derivations and transseries}, \arxiv{1503.00315}. In particular, the normal-form and exponential preliminaries in Sections 2--3.
% cited by source manuscripts 03, 04

\bibitem{BMderivations}
Alessandro Berarducci and Vincenzo Mantova, ``Surreal numbers, derivations and transseries,'' \emph{Journal of the European Mathematical Society} \textbf{20} (2018), 339--390. \href{https://ems.press/content/serial-article-files/32275}{Publisher full text}.
% cited by source manuscripts 02, 07

\bibitem{BMgerms}
Alessandro Berarducci and Vincenzo Mantova, ``Transseries as germs of surreal functions,'' \emph{Transactions of the American Mathematical Society} \textbf{371} (2019), 3549--3592. \href{https://arxiv.org/abs/1703.01995}{arXiv:1703.01995}; in particular Section~3 on surreal analytic functions.
% cited by source manuscripts 02, 07

\bibitem{CE}
Ovidiu Costin and Philip Ehrlich, ``Integration on the surreals,'' \emph{Advances in Mathematics} \textbf{452} (2024), article 109823. \href{https://arxiv.org/abs/2208.14331}{arXiv:2208.14331}; \href{https://www.sciencedirect.com/science/article/pii/S0001870824003384}{publisher record}.
% cited by source manuscripts 02, 07

\bibitem{Conwayc}
J. H. Conway, \emph{On Numbers and Games}, Academic Press, 1976; second edition, A K Peters, 2001.
% cited by source manuscripts 04, 08

\bibitem{EK}
P. Ehrlich and E. Kaplan, Surreal ordered exponential fields, \emph{The Journal of Symbolic Logic} \textbf{86} (2021), 1066--1115. \href{https://doi.org/10.1017/jsl.2021.59}{doi:10.1017/jsl.2021.59}. Author preprint: \href{https://arxiv.org/abs/2002.07739}{arXiv:2002.07739v3}, especially Section 11. See also the reference-correction erratum, \href{https://doi.org/10.1017/jsl.2022.5}{doi:10.1017/jsl.2022.5}.
% cited by source manuscripts 01, 09

\bibitem{Gonshorb}
Harry Gonshor, \emph{An Introduction to the Theory of Surreal Numbers}, London Mathematical Society Lecture Note Series \textbf{110}, Cambridge University Press, 1986.
% cited by source manuscripts 02, 07

\bibitem{Orloff}
J. Orloff, \emph{Complex Variables with Applications}, MIT course 18.04, Spring 2018, lecture notes. Background for ordinary Cauchy theory, residues, the argument principle, and related classical results. \href{https://ocw.mit.edu/courses/18-04-complex-variables-with-applications-spring-2018/pages/lecture-notes/}{MIT OpenCourseWare lecture notes}.
% cited by source manuscripts 01, 06

\bibitem{AllingShort}
N. L. Alling, Fundamentals of analysis over surreal number fields, \emph{Rocky Mountain Journal of Mathematics} \textbf{19} (1989), 565--573. \href{https://doi.org/10.1216/RMJ-1989-19-3-565}{doi:10.1216/RMJ-1989-19-3-565}.
% cited by source manuscripts 01

\bibitem{BMgermsb}
A. Berarducci and V. Mantova, Transseries as germs of surreal functions, arXiv:1703.01995. In particular, Section 3 on summability, analyticity, and composition. \url{https://arxiv.org/abs/1703.01995}
% cited by source manuscripts 03
\bibitem{BagMant}
V. Bagayoko and V. Mantova, \emph{Taylor expansions over generalised power series}, \arxiv{2509.08473}, v1, September 10, 2025.
% cited by source manuscripts 04

\bibitem{BagayokoMantova}
Vincent Bagayoko and Vincenzo Mantova, ``Taylor expansions over generalised power series,'' preprint, 2025, \href{https://arxiv.org/abs/2509.08473}{arXiv:2509.08473}.
% cited by source manuscripts 02

\bibitem{ConwayComplex}
J. B. Conway, \emph{Functions of One Complex Variable I}, second edition, Graduate Texts in Mathematics \textbf{11}, Springer, 1978. \href{https://doi.org/10.1007/978-1-4612-6313-5}{doi:10.1007/978-1-4612-6313-5}.
% cited by source manuscripts 06

\bibitem{Conwayb}
John H. Conway, \emph{On Numbers and Games}, second edition, A K Peters, 2001; first edition, 1976.
% cited by source manuscripts 02

\bibitem{Conwayd}
John H. Conway, \emph{On Numbers and Games}, Academic Press, 1976; 2nd ed., A K Peters, 2001.
% cited by source manuscripts 07

\bibitem{CostinEhrlich}
O. Costin and P. Ehrlich, \emph{Integration on the surreals}, \arxiv{2208.14331}, initially submitted August 30, 2022.
% cited by source manuscripts 04

\bibitem{EKb}
Philip Ehrlich and Elliot Kaplan, ``Surreal ordered exponential fields,'' \emph{The Journal of Symbolic Logic} \textbf{86} (2021), 1066--1115, \href{https://doi.org/10.1017/jsl.2021.59}{doi:10.1017/jsl.2021.59}. \href{https://arxiv.org/abs/2002.07739}{arXiv:2002.07739}; Section~11 treats surcomplex exponentiation.
% cited by source manuscripts 02

\bibitem{Mantova2026}
Vincenzo Mantova, ``Monotonicity and a Taylor approximation theorem for transseries,'' preprint, 2026, \href{https://arxiv.org/abs/2601.07747}{arXiv:2601.07747}, version~2, May~9, 2026.
% cited by source manuscripts 02

\bibitem{MantovaMatusinski}
V. Mantova and M. Matusinski, Surreal numbers with derivation, Hardy fields and transseries: a survey, \href{https://arxiv.org/abs/1608.03413}{arXiv:1608.03413}, 2016. Background on normal forms, summability, and the distinction between function differentiation and derivations of surreal fields.
% cited by source manuscripts 01

\bibitem{Marker}
D. Marker, \emph{Model Theory of Valued Fields}, lecture notes, University of Illinois at Chicago, 2018, Sections 1--2, especially Lemma 2.36. Available from the author's course page: \url{https://homepages.math.uic.edu/~marker/math512-f18/valued_fields_1-2.pdf}.
% cited by source manuscripts 09

\bibitem{McMullen}
C. T. McMullen, \emph{Advanced Complex Analysis}, Harvard University course notes, author-hosted version consulted 21 September 2026. \url{https://people.math.harvard.edu/~ctm/papers/home/text/class/harvard/213a/course/course.pdf}
% cited by source manuscripts 03
\bibitem{Neumannb}
Bernhard H. Neumann, ``On ordered division rings,'' \emph{Transactions of the American Mathematical Society} \textbf{66} (1949), 202--252, \href{https://doi.org/10.1090/S0002-9947-1949-0032593-5} {doi:10.1090/S0002-9947-1949-0032593-5}.
% cited by source manuscripts 02

\bibitem{PS}
Y. Peterzil and S. Starchenko, Expansions of algebraically closed fields in o-minimal structures, \emph{Selecta Mathematica, New Series} \textbf{7} (2001), 409--445. \href{https://doi.org/10.1007/PL00001405}{doi:10.1007/PL00001405}. See also the authors' \href{https://math.haifa.ac.il/kobi/Correction.pdf}{correction}.
% cited by source manuscripts 01

\bibitem{PSsurvey}
Y. Peterzil and S. Starchenko, Complex analytic geometry in a nonstandard setting, in \emph{Model Theory with Applications to Algebra and Analysis}, Cambridge University Press, 2008, 117--166. \href{https://math.haifa.ac.il/kobi/Newton.pdf}{Author manuscript}.
% cited by source manuscripts 01

\bibitem{Poonen}
B. Poonen, Maximally complete fields, \emph{L'Enseignement Math\'ematique} (2) \textbf{39} (1993), 87--106. \url{https://math.mit.edu/~poonen/papers/amsval.pdf}
% cited by source manuscripts 03
\bibitem{PoonenUnits}
B. Poonen, \emph{Units in Hahn--Mal'cev--Neumann rings}, manuscript, 14 January 2026. \url{https://math.mit.edu/~poonen/papers/malcev.pdf}
% cited by source manuscripts 03
\bibitem{Poonenb}
Bjorn Poonen, ``Maximally complete fields,'' \emph{L'Enseignement Math\'ematique} (2) \textbf{39} (1993), 87--106. \href{https://math.mit.edu/~poonen/papers/amsval.pdf}{Author's full text}.
% cited by source manuscripts 07

\bibitem{vdDEb}
Lou van den Dries and Philip Ehrlich, ``Fields of surreal numbers and exponentiation,'' \emph{Fundamenta Mathematicae} \textbf{167} (2001), 173--188. \href{https://doi.org/10.4064/fm167-2-3}{doi:10.4064/fm167-2-3}. Erratum, \emph{Fundamenta Mathematicae} \textbf{168} (2001), 295--297. The present article uses normal forms, restricted analytic evaluation, and the exponential group structure, not the ordinal estimates corrected in the erratum.
% cited by source manuscripts 07

\bibitem{vdDEerr}
L. van den Dries and P. Ehrlich, Erratum to ``Fields of surreal numbers and exponentiation'', \emph{Fundamenta Mathematicae} \textbf{168} (2001), 295--297.
% cited by source manuscripts 09

\bibitem{vdH}
J. van der Hoeven, ``Operators on generalized power series'', \emph{Illinois Journal of Mathematics} \textbf{45} (2001), 1161--1190. \href{https://doi.org/10.1215/ijm/1258138061}{doi:10.1215/ijm/1258138061}. Author version: \url{https://www.texmacs.org/joris/noeth/noeth.html}.
% cited by source manuscripts 05

\end{thebibliography}
\end{document}
